\quellseite{2}
\begin{titelseite}
\werktitel{Der Evangelische Geistliche.}\\[1.5em]
\werkuntertitel{Zweites Bändchen.}
\end{titelseite}

\quellseite{3}

\quellseite{4}
\begin{titelseite}
\werktitel{Der evangelische Geistliche.}\\[1.5em]
Dem nun folgenden Geschlechte evangelischer Geistlichen dargebracht\\[1.5em]
von Wilhelm Löhe,\\
weil. luth. Pfarrer zu Neuendettelsau in Franken.\\[1.5em]
\werkuntertitel{Zweites Bändchen.}\\[1em]
Ansichten aus den verschiedenen Arbeitsgebieten des geistlichen Amtes.\\[1.5em]
Dritte mit der zweiten gleichlautende Auflage.\\[1.5em]
Gütersloh.\\
Druck und Verlag von C. Bertelsmann.\\
1876.
\end{titelseite}

\quellseite{5}

\quellseite{6}
\begin{titelseite}
Den Zöglingen\\
der beiden engverbundenen Pflanzschulen,\\
des\\
evangelisch-lutherischen Predigerseminars\\
Wartburg\\
in St. Sebald am Quell, Iowa,\\
und\\
der evangelisch-lutherischen Missionsschule\\
zu\\
Neuendettelsau\\
in Mittelfranken.
\end{titelseite}

\quellseite{7}

\quellseite{8}
Der erste Theil meines sogenannten „Evangelischen Geistlichen“ ist aus Dictaten entstanden, welche ich den Schülern der hiesigen Missionsanstalt zum Behuf ihrer Ausbildung vor einer Reihe von Jahren gegeben, und zum Zweck der Veröffentlichung hernach umgeschrieben hatte. Ein zweiter und dritter Theil sollte auf gleichem Wege entstehen. Die Dictate, deren Inhalt den zweiten Theil bilden sollte, liegen seit Jahren bereit, die Zeit aber, sie zum Zwecke der Veröffentlichung etwas eingehender umzuschreiben, wollte bis jetzt nicht kommen, so daß ich entweder den Gedanken, aus ihnen einen zweiten Theil der schon genannten Schrift zu bilden, aufgeben oder es wagen mußte, Euch, meine werthen jüngeren Brüder, und denen, welche sonst noch für etwas so Unvollkommenes ein Interesse fassen können und wollen, meine Dictate vom Jahre 1853 in ziemlich unveränderter Gestalt zu geben. Ich wagte das letztere. Es ist mir längst bekannt, daß ich ein Buch von eigentlich wissenschaftlicher Gestalt nicht schreiben kann. Auch wenn ich die Dictate umgeschrieben hätte, würden sie doch, wenn auch vielleicht in etwas geringerem Maße, die aphoristische Gestalt behalten haben, welche Ihr an ihnen bereits kennt. Dabei ist es ja heutzutage mit der Veröffentlichung von schriftlichen Arbeiten nicht mehr so, daß man es bloß bei dem Bewußtsein, etwas recht Vollendetes zu geben, wagen könnte und dürfte. Das Drucken ist jetzt nicht viel mehr, als das Schreiben oder Reden, und wer sich von herein schon ergiebt, nichts Besonderes in die Welt zu bringen, sondern nur anzuregen und andere zu Besserem zu veranlassen, der mag sich am Ende mit der Druckerschwärze befassen, fast wie mit der Tinte des Schreibers. So geht denn der zweite
\quellseite{9}
Theil dieses „Evangelischen Geistlichen“ in der armen Gestalt hinaus, die er hat, und es ist mir angenehm, wenn ich Euch, meine Lieben, und Eures Gleichen, damit einen, sei es auch kleinen und geringen Dienst erwiesen habe.

Als Kern dieses Bändchens gilt der mittlere Theil. Wie Ihr aber sehet, hat dieser zweite Theil im Ganzen eine ähnliche Gestalt, wie der erste. Der erste verbirgt sich hinter dem berühmten Gedicht Valentin Andrea’s über das gute Leben eines rechtschaffenen Dieners Gottes, wie hinter einem Schilde; ebenso der zweite hinter dem Sendschreiben des guten Eberlin von Günzburg. Es wird Euch mit diesem Sendschreiben gehen, wie mit vielen andern Schriften derselben Zeit. Man findet sich nicht bloß durch die Trefflichkeit der Sprache, sondern auch durch das edle Gold des Inhalts angesprochen und gefördert, ohne daß man geradezu alles und alles unterschreiben möchte. Ich denke aber, Ihr, deren Beruf es mit sich bringt, so vielfach neu zu pflügen und auf Neubruch zu säen, werdet euch dem Eberlin besonders verwandt finden, der im Grunde auch nichts anders gibt, als eine Anleitung, wie man in solchen Fällen säen und arbeiten müsse. Es ist auch nicht bloß mein Einfall, dies alte Sendschreiben wieder zu veröffentlichen; es wurde bereits mehrere Male und zuletzt, so viel ich weiß, von August Hermann Francke in seinem Collegium pastorale an’s Licht gezogen. Seht einmal zu, lieben Brüder, ob nicht Eberlin wirklich ein Mann des rechten Maßes ist, und ob nicht also wir, deren ganzes Streben in allen Dingen dahin geht, unbekehrt durch die Meinung des jetzigen „menschlichen Tages,“ das rechte Maß zu suchen, ganz recht thun, den alten Eberlin wenigstens in effigie in unsern Orden aufzunehmen. Hat er es nicht allewege getroffen, so thut das zur Sache nichts: kein Mensch des rechten Maßes will es allewege getroffen haben; und ich gestehe, daß mir Leute, die nur Recht haben, verdächtiger sind als die Pharisäer,

\quellseite{10}
welche die Mißbilligung Christi getroffen hat. Es muß an einem Menschen noch etwas zu bessern überbleiben, so lang er hienieden wallt, und wir gehen alle, so heilig wir seien, nur über Mängel hin der Vollendung zu. — Als Anhang des ersten Theils meines Büchleins erschien ein Wort von der Pfarrfrau. Anhang dieses zweiten Theils ist ein Dictat für Diaconissen. Die Diaconissin, wenigstens die antike, darf gar wohl neben dem Pfarrer erscheinen, und sollte jetzt wie sonst die rechte Hand eines Haushalters über die göttlichen Geheimnisse unter den Frauen sein. Was nun aber diesen Anhang selbst betrifft, so hat er gar keine Absicht, als die große Bedeutung und Wichtigkeit der psychischen Zustände für die Seelsorge am Krankenbett hervorzuheben, und die Nothwendigkeit für alle, die Seelsorge üben, sich mit denselben bekannt zu machen. Wenn man durch den Dienst dieses Büchleins hie und da aufmerksamer wird auf diesen Punkt, und es ein wenig faßt, wie schwer die Seelsorge am Krankenbette ist, und wie gar nicht man für dieselbige reif und geschickt ist, wenn man vom Universitätsstudium frisch in’s Amt hinein tritt, so habe ich meinen Zweck schon erreicht. Man kann mir allerdings sagen, ob ich denn einer Jungfrau, die als Diaconissin dient, mehr Weisheit zutraue, als dem Jüngling, der von der Schule her in’s Amt tritt? Darauf aber antworte ich ganz einfach: „Nein. Ich alter Pfarrer bin zwar der Thor nicht, daß ich die Seelsorge bloß für mich und meines Gleichen pachten wollte; ich bin meinen Kirchenvorstehern und Pfarrkindern böse, wenn sie mir nicht seelsorgen helfen, und kann drum auch wohl die Diaconissin in meiner Gemeinde leiden, wenn sie sich mit der edlen Seelsorge befaßt. Aber eben deshalb, weil ich das wohl dulden und haben will und sehe, daß sie es eben so wenig von selbst versteht, als der Jüngling der in’s Amt tritt, unterrichte ich die Dienerinnen Jesu, und habe geglaubt, ihnen auch in dem Dictate,
\quellseite{11}
von welchem ich zunächst rede, etwas nützliches an die Hand zu geben.“ Ich habe es übrigens auch nicht von mir selbst gelernt, und habe im Texte selber meinen Hauptgewährsmann Vering genannt, der mir übrigens selten für die anzuwendenden psychischen Mittel, dagegen aber recht oft bei der Darstellung des gegenseitigen Einflusses Leibes und der Seele gedient hat. Ich meine für Euch, und wer sonst lernen mag, ist es weit besser, daß ich einen medicinischen Gewährsmann hatte, als wenn ich alles aus meiner eignen Erfahrung genommen hätte.

So nehmt nun Ihr, meine lieben Brüder, dies kleine Büchlein hin und benutzt es so gut Ihr könnet und wollet. Steigt auf unsre Schultern und hebt Euch auf denen zu höherer Weisheit. Bleibt Euren Lehrern und Vorfahren getreu, aber nicht so, daß Ihr nicht auf unsrem Wege weiter gehen wollet. Behaltet, was Ihr habet; aber wenn Euch der heilige Geist zum alten Segen neuen gibt, oder bei Euch anfrischt, was wir vergessen haben; so bildet Euch nicht ein, daß wir in der Ewigkeit Euch darüber zürnen werden. Wir haben Euch oft gesagt, Gott sei der Creatur gegenüber ein unermessliches Meer der Erkenntniß, und es gebe einen ewigen Fortschritt wie in der Gotteserkenntniß, so in der Erkenntniß unsres sich hier auf Erden abschließenden Sündenlebens. Wer aber von einem ewigen Fortschritt redet, den kann es nicht verdrießen, wenn die Nachkommen in der Zeit auch ihre Fortschritte machen. Demgemäß verhalten wir uns ja auch in unsrer Zeit gegenüber den Vätern. Vor 20, 30 Jahren wollte man in der protestantischen Kirche von der Beute des sechzehnten Jahrhunderts nichts wissen, man verachtete sie. Da thaten wir, was wir konnten, um das Heerlager Christi zu den alten Symbolen zu rufen, — wir nämlich nicht allein, aber unter anderen auch wir. Wir thun auch noch jetzt so und rufen namentlich in diesen unsren Tagen nicht mehr bloß
\quellseite{12}
zu den alten Symbolen, sondern auch zu einem denselben gemäßen kirchlichen Verhalten. Wie wenig Erfolg wir im letzteren Stücke haben, wisset Ihr; es zeigt sich eben, daß eine pure Meinungsänderung noch keine rechte Buße und Lebensänderung ist. Mehr Erfolg hat der Ruf zum Worte des alten Bekenntnisses gehabt. Da nun aber viele, mit dem Umtausch der Meinung zufrieden, ihr Tagewerk abschließen und auf der alten Bahn keinen neuen Fortschritt gestatten wollen, auch wo er von Gottes Wort und Führung erheischt und gefordert wird; so beharren wir mit eben demselbigen Ernste und Nachdruck auf dem „Vorwärts“ wie auf dem „Rückwärts.“ Dem sogenannten Lutheraner, d. i. dem wahren katholischen Christen, gehört die gesammte Vergangenheit vor und nach Luther; ihm muß auch die Zukunft gehören. Alles ist sein, was wahr und schriftmäßig ist, wann, wo und wie es gesagt wird, und die norma normata des sechzehnten Jahrhunderts ist ihm nicht in dem Sinne der norma normans congruent, daß diese in jener erschöpft und es Gott selbst nicht mehr erlaubt wäre, seiner Kirche noch etwas zu geben, was man im Normaljahr 1580 entweder nicht hatte oder nicht beachtete. — Es versteht sich von selber, daß man beim Vorwärtsschreiten Vorsicht bedarf und den Leitfaden der alten Zeit in der Hand behalten muß. Auch kann man ja freilich fehlen und irren, und muß daher auch Warnung, ja Vorwurf und Geschrei zur Linken und zur Rechten beachten, anhalten im Lesen und Forschen und Beten. Der Herr aber wird es dennoch den Aufrichtigen gelingen lassen und die Richtung segnen, welcher wir angehören, in welcher sich, wie am Tage Marien Heimsuchung das alte und neue Testament auf den Bergen von Juda, so die vergangene alte Zeit und eine, will's Gott, bessere Zukunft findet und faßt!

Der Geburtstag des Herrn ist vor der Thür. Freuet Euch in dem Herrn allewege, und abermal sage ich Euch:

\quellseite{13}
\begin{quote}
Freuet Euch! Eure Lindigkeit laßt kund werden allen Menschen. Der Herr ist nahe. Alle Eure Sorge lasset im Gebet und Flehen mit Danksagung vor Gott kund werden. Und der Friede Gottes, welcher höher ist, als alle Vernunft, bewahre Euch Herzen und Sinne in Christo Jesu! Amen.
\end{quote}

\begin{unterschrift}
N. D. 17. Dez. St. Wunnibalds Abend. 1857.\\[0.8em]
W. L.
\end{unterschrift}

\chapter{Vorwort zur zweiten Auflage}

So gering kam mir dieser 2. Theil meines Evangelischen Geistlichen vor, als ich ihn zum Behuf einer zweiten Auflage durchlas, daß ich ihn beinahe gar nicht mehr hätte erscheinen lassen. Nach viel Schwanken ergab ich mich, eine neue Auflage zu bearbeiten. Es geschah mit einer Art von Resignation. — Haben sich Leute gefunden, welche die 1. Auflage zu benutzen wußten, die schlechter als diese war; so konnte ich diese, die meines Bedünkens besser als jene ist, am Ende auch vom Stapel laufen lassen.

Die Eingangsschrift Eberlins von Günzburg ist geblieben; ich werde dafür keinen Tadel zu fürchten haben. Was ich in der Vorrede zur 1. Auflage von dieser Schrift gesagt habe, ist meine Ansicht auch noch jetzt. Dagegen ist der Anhang der 1. Auflage von den psychischen Einwirkungen leiblicher Krankheiten weggeblieben, aus Gründen, welche § 47 p. 243 ff. dargelegt sind. Doch ist ein Stück davon (§ 63 p. 314 ff.) auch in diese Auflage aufgenommen worden. An die Stelle der psychischen Einwirkungen ist in dieser Auflage etwas
\quellseite{14}
anderes getreten: Dr. Hochstetteri de recta concionandi textumque sacrum cum exponendi tum applicandi ratione Commentariolus. S. p. 315—348. Vielleicht hat der Anhang der 1. Auflage bei aller Unvollkommenheit doch einige Frucht getragen, — und vielleicht trägt auch der Anhang der 2. Auflage gute Frucht. Er scheint mir bei aller Kürze doch Freunde und Anerkennung zu verdienen.

Die Mitte des Ganzen, der eigentliche 2. Theil meines Evangelischen Geistlichen, findet sich p. 71—324. Obgleich reicher geworden und sorgfältiger gearbeitet, trägt er doch seine alte aphoristische Natur. Ich kann nicht mehr als das. Ein Schelm gibt mehr und besseres als er hat. — Ich denke mir eben keine jungen Candidaten von hoher wissenschaftlicher Bildung als mein Publicum: was sollten diese mit meinem armen Büchlein thun? Ich denke mir Leute, die von der Schule mit dem Gefühle der Armuth in’s Amt treten und froh sind, von einem alten Pfarrer dies oder jenes zu hören, was ihnen bei beginnender Praxis Irrfahrt ersparen kann. Möge ich mich nicht umsonst bemüht haben, und der Herr die Schrift mit Seinem Segen begleiten! Sie gehe hin und leiste, was sie durch Gott leisten kann.

Mehr zu sagen weiß und habe ich nicht.

Friede sei dem Leser! Amen.

\begin{unterschrift}
N. D. 24. IV. 1866.\\[0.8em]
W. L.
\end{unterschrift}

\quellseite{18}
\chapter{Einleitung}
\begin{titelseite}
\werkuntertitel{Sendschreiben}\\[1em]
\werktitel{Johann Eberlins,}\\[1em]
Der zuerst zu Ulm das Evangelium geprediget,\\[0.8em]
an\\[0.8em]
\werkuntertitel{Johann Jacob,}\\
Pfarrern zu Leipheim im Ulmer-Gebiet.\\[1.5em]
Nach der Edition D. Georg Cölestini abgedruckt.

\begin{quote}
Wie sich ein Diener Gottes Worts in all seinem Thun halten soll, und sonderlich gegen denen, welchen das Evangelion zuvor nicht geprediget ist, daß sie sich nicht ärgern.
\end{quote}

\begin{center}
Zu Wittenberg erstlich im Jahre 1525\\[0.8em]
mit D. Luthers Rath ausgegangen, aber jetzund wiederum treulich nachgedrucket.\\[1.5em]
Anno MDLXXIII.\\[1.5em]
Der evangelische Geistliche. II. 3te Aufl.
\end{center}
\end{titelseite}

\quellseite{19}
Den Wohlgebohrnen, Edlen, Gestrengen und Ehrenfesten beyden Ständen, den Herren und Ritterschaft des Erzherzogthums Oesterreich Deputirten in Religions-Sachen, seinen Gn. und Günstigen Herren, wünschet Georg. Colestinus, D. und Thum-Probst,
\quellseite{20}
Gottes Gnade, Fried, Segen und Wohlfahrt zur ewigen Seligkeit, mit Erbietung seiner treuen Dienste, bevor. Wohlgebohrne, Edle, Gnädige, Großgünstige Herren, daß Ew. Gn. und G. sich so emsig und fleißig um treue Diener göttliches Wortes bewerben, und an Ihnen nichts an Mühe oder Unkosten erwinden lassen, damit Sie derselben etliche, so das liebe Wort nach der Propheten und Apostel Schriften, auch Augspurgischer Confession, rein und lauter, und wie Petrus vermahnet, mit Furcht und Sanftmuth, ohne Bitterkeit und Lästerung predigen und fortsetzen helfen, an sich bringen mögen: an diesem thun E. G. und G. GOtt zu Gefallen und der armen betrübten Kirchen zur Erbauung, auch Ihnen selbst zur Seligkeit. Danke auch GOtt für diese Treue und Wohlthat, und bitte mit allen frommen lieben Christen, Er wolle E. G. und G. in solchem Christlichen Werck und Vorhaben gnädiglich stärken und erhalten. Amen.

Und nachdem auch Ew. Gn. und G. als Deputirte in Religions-Sachen, mich unwürdigen zu einem Diener des Worts gen Wien gnädig und günstig berufen, sich auch meines Raths, Christlicher Hülffe und Beiwohnung getrösten, hierum auch meinen Gnädigsten Herrn, den Churfürsten zu Brandenburg, um Erlaubniß ersuchen; ich aber meines tragenden Amts und anderer Ursach halben so gar eilend und bald zu E.
\quellseite{21}
Gnad und G. nicht kommen können: Als habe ich gleichwol zu der Sachen Anfang und einem Schul-Recht dies Büchlein, wie sich ein Prediger bei denen halten soll, welchen zuvor-Gottes Wort nicht geprediget, der Meinung wieder in Druck geben und E. G. und G. zuschreiben wollen, daß Ew. Gn. und G. hiermit einen Grund und Wahrzeichen hätten, wie und welcher Gestalt ich bei Derselben das Wort zu pflanzen, fortzusetzen und zu befördern vermeinete, auch unterdeß E. G. und G. Prädicanten, so es nicht besser wüßten, eine Anleitung hätten, sich in Gottes Sachen sanftmüthig und still zu erzeigen. Gefällt es nun E. G. und G. wohl, wie ich hoffe, daß es allen Gottseeligen, Friedliebenden gefallen und behagen soll, so freue ich michs im Herrn. Es mögen die unruhigen, sichern Klüglinge und Meister davon lallen und waschen, was sie wollen, so sei uns gnug Gottes Ehre und seines Worts Fortsetzung, der wird am Jüngsten Tage wol ans Licht bringen, was und worauf ein jeder gebauet. Dem treuen, ewigen Gott thue E. Gn. und G. ich in seinen Schutz, Schirm und Segen befehlen, der erhalte Sie alle sammt und sonderlich zur ewigen Seligkeit, Amen. Geben im Churfürstl. Brandenburgischen Hoflager Cölln an der Spree.

Anno 1573. den 8ten Tag Septembris, auf welchem vor 1500 Jahren die Stadt Jerusalem erstiegen und erobert von den Römischen Kriegsleuten Titi Vespasiani, verbrannt, verheeret und geschleift ist, also, daß, nach der Weissagung des Herrn Christi, kein Stein auf dem andern, und wie Josephus schreibet, kein Juden-Haar allda geblieben und ge%
\quellseite{22}
laßen worden, uns zur Lehr und Warnung, daß wir der Gnaden-Zeit, darin GOttes Wort verkündiget und geprediget wird, wahrnehmen, auf daß wir nicht gleicher Weise gestraft und vertilget werden, wie der HErr Luc. 13. dräuet: So ihr euch nicht beßert, werdet ihr alle auch also umkommen.

\begin{center}
2 Timoth. 2.
\end{center}

\begin{quote}
„Ein Knecht des HErrn soll nicht zänkisch sein, sondern freundlich gegen jedermann, lehrhaftig, der die Bösen tragen kann mit Sanftmuth, und strafe die Widerspenstigen, ob ihnen GOtt dermaleins Buße gäbe, die Wahrheit zu erkennen.“
\end{quote}

\quellseite{23}
\section[Dem Würdigen Herrn Johann Jacob]{Dem Würdigen Herrn Johann Jacob, Pfarrherrn der Christlichen Gemeine zu Leipheim an der Donau, meinem günstigen Herrn und lieben Vettern}

Ihr wißet wohl, lieber Herr und Bruder in Christo, wie gar schön und meisterlich der heilige Apostel St. Paulus in den zween Episteln an Timotheum, und in der an Titum einen Bischof oder Pfarrherrn abmahlet, daß es ohne Noth ist, etwas anders dazu zu thun.

Aber dieweil solche gemeine Lehren und Vermahnungen viel nützlicher und heilsamer sind, wenn sie nach Gelegenheit der Umstände, Zeit, Orter, \&c. auf sonderliche Personen und Zufälle angestellet und accomodieret werden; und ich nicht allein aus Christlicher Liebe, sondern auch aus leiblicher Blutsverwandniß euch zu schreiben bedacht gewesen: So habe ich, für einen oder zween Briefe, dieses Büchlein voll freundlicher und Christlicher Vermahnungen, wie wir unser heiliges Predig-Amt gottseelig, fruchtbarlich und unärgerlich führen sollen so viel GOtt der Allmächtige Gnade geben, euch öffentlich durch den Druck zuschicken wollen. GOtt gebe, daß es bei euch und andern, denen ich auch gern hiemit dienen wollte, viele Frucht schaffe, Amen.

Ehe ich aber ferner etwas schreibe, will ich euch erstlich erinnern und vermahnen, daß ihrs dafür halten sollt, daß ich euch nichts zuschreiben will, es habe denn guten Grund in der H. Schrift, oder sei mit bewährlichen vernünftigen Ursachen und Beweisungen angezeigt, wie ihr hören werdet.

\quellseite{24}
\randnotiz{Gott danken um Erkenntniß der Wahrheit.}Anfänglich und vor allen Dingen sollt ihr nimmer vergeßen der großen Gnaden und Gaben GOttes, so er euch gegeben hat, durch welche ihr gewaltiglich gerissen seid aus dem Schlund und Rachen des höllischen Löwen, der euch so tief versenket hat in Lust, Ehre und Gunst dieser Welt, da ihr ohne Geiz und große Anreizung zu vielen andern Lastern, auch zu Verfolgung und Verachtung des Evangelions, nicht wol sein konntet. Aber gelobet sei GOtt, unser Vater, der alle Bande des Teufels gewaltiglich zerbrochen, und alle seine verführerische Arglist umgestoßen hat, und euch Erkenntniß seiner Wahrheit gegeben, ja, nach dem Spruch: So ist euch gegeben zu thun, daß ihr nicht allein an Christum gläubet, sondern auch um Ihn leidet. Ihr seid in Bann gethan, und GOtt hat doch den Bann geordnet zu eurer und eurer Schäflein Heil und Seligkeit. Ihr stehet doch in großer Gefahr eures Lebens alle Stunde; dennoch gibt euch GOtt Gnade, sein Wort beständiglich ohn alle Scheu zu predigen mit großer Lust und Begierde der Zuhörer, so daß auch die umliegenden Völker, dem Wort Gottes ferne nachzureisen beweget werden, daß mich deucht, der Spruch des Propheten Sachariä wird jetzt erfüllet, da er sagt, daß zu den Zeiten des Evangelions die Heiden ergreifen werden den Saum an dem Kleide eines Christen, und werden sagen: Wir wollen mit dir gehen, denn wir haben gehöret, daß GOtt bei euch ist. Ja wahrlich ist GOtt bei den Leipheimern, da sein Wort geprediget wird. Freuet euch, ihr frommen Christen daselbst, und leset das sechzigste Capitel im Propheten Jesaia von euch gesagt sein; leset dasselbige, und lobet GOtt, bittet Ihn, daß Er seine Verheißungen wolle erfüllen. Diese Gnade und Gabe GOttes sollt ihr täglich betrachten und hoch wägen, ihm fleißiglich darum danken, und bitten um Befestigung und Mehrung solcher Gnade. Hütet euch, daß ihr nicht undankbar erfunden werdet, daß GOtt nicht seine Klage über euch führe, wie ge%
\quellseite{25}
schrieben steht Jesaia 1. Amos 2. Matth. 12. und an andern Orten mehr.

Zum andern. Nichts soll so fürnehmlich und förderlicher bei euch sein, als ein emsiges, ernstliches, herzliches und zuversichtiges GEBET für euch und für eure Zuhörer, für oberste und unterste, für alle Freunde und Feinde. Und wisset, daß im Alten und Neuen Testament GOtt uns so oft und freundlich vermahnet, im Glauben zu beten, denn ohne das Gebet schaffet ihr nichts wider den Teufel. Derhalben hat auch St. Paul das Gebet unter den Christlichen Harnisch gezählet. Ihr sollt auch eure Zuhörer vermahnen zu dem Gebet, wie auch Paulus einem Bischof oder Pfarrherrn dasselbige vorschreibet 1 Timoth. 2., da er also spricht: So vermahne ich nun, daß man vor allen Dingen zuerst thue Bitte, Gebet, Fürbitte und Danksagung für alle Menschen, für die Könige und für alle Obrigkeit, auf daß wir ein geruhig und stilles Leben führen mögen in aller Gottseligkeit und Redlichkeit, denn solches ist gut, dazu auch angenehm vor GOtt, unserm Heiland, welcher will, daß alle Menschen genesen und zur Erkenntniß der Wahrheit kommen.

Wenn das Volck versammelt ist, sollt ihr erzählen, wie in so großer Fährlichkeit wir sind des Leibes, der Seelen, des Guts, der Ehre, und wie der Teufel, die Welt und das Fleisch einen unüberwindlichen Streit und Kampf wider uns ohn Unterlaß führen, welchen keine menschliche Weißheit noch Kraft möge widerstehen. Vom Teufel spricht GOtt selbst im Hiob\footnote{Hiob. 41.}, daß auf Erden ihm niemand zu gleichen sei. Ja Christus nennet ihn\footnote{Johann. 14.} einen Fürsten, und St. Paul\footnote{2 Corinth. 4.} einen GOtt dieser Welt. Von der Welt spricht St. Johannes\footnote{1 Johann. 5.}, daß allein der Glaube an JEsum Christ die Welt überwinden möge. Von dem Fleisch beklagt sich auch der heilige Paulus\footnote{Röm. 7.}, und sagt, daß dawider nichts helfe,
\quellseite{26}
denn die Gnade GOttes. Solche große Fährlichkeit dieser dreien Feinde halber sollt ihr aus Sprüchen der heiligen Schrift vor eurem Volke lesen, und aus täglicher Erfahrung das Volk damit vermahnen, zu bitten GOtt um Hülfe und Beistand wider solche heftigen Feinde.

\subsection{Was von GOtt zu bitten}

Auch sollt ihr ihnen erzählen, wie wir in so großem Schaden stecken, die wir nicht wissen, noch finden mögen, die Schrift zeige sie uns denn, so blind und toll sind wir. Sollt also das Volk derhalben vermahnen zum Gebet um Erlösung von solcher Fährlichkeit und Schaden. Auch schickt uns GOtt täglich so viel Unglücks über den Hals, daß wir Ursache genug haben um Erlösung zu bitten. Item, man soll das Volk vermahnen, auch für andere zu beten, und sonderlich für die Verfolger des Worts GOttes, daß sie ablassen GOttes Wort zu lästern, auf daß sie Leben und Heil erlangen. Item, man soll auch dem Volke allda vortragen die Wohlthat und Güte, so uns GOtt beweiset an Seel, Leib, Ehre, Gut, Freunden und dergleichen, wie ein jeglicher oft bei ihm selbst erfahren hat. Und daß wir auch GOtt danken um seine göttliche Verheißungen, die er uns aus lauter Gnad und Barmherzigkeit gethan hat und zugesaget.

Diese Stücke sollen dem Volk also vorgetragen werden, daß man sie in das Herz ziehe, daß also unser Herz dadurch zur Erkenntniß eigenen Unglücks beweget werde, Hülfe von GOtt zu begehren, ihm ganz und gar vertrauen, danksagen und loben in allen Creaturen, wie die Psalmen thun, der 104. 148. und dergleichen mehr. „Denn wenn die heilige „Schrift gelesen, oder geprediget, oder gehöret, nicht gezogen wird in ein Beten und Danksagen, so ist alle Mühe und Arbeit verloren.“

\subsection{Studieren in der heiligen Schrift}

Zum dritten. Alle Tage sollt ihr etliche Stunden, wenns euch wohl gefällt, in der Bibel, das ist, in der heiligen Schrift
\quellseite{27}
lesen. Aber sehet zu, daß ihr dem Lesen Recht thuet. Erstlichen sollt ihr das geringste Wort, in der heiligen Schrift verfaßet, höher denn aller Welt Gut achten. Denn Gottes Wort ist eine göttliche Kraft\footnote{Röm. 1.}, welches unsere Seele selig machet\footnote{Joh. 1.}, ja es ist Gott selbst\footnote{D. i. es ist Gott selbst, der seine Kraft in dem Wort und durchs Wort erweiset.}, das Wort im Fleisch und in der Stimme ist gleich\footnote{D. i. wie sich Gott in dem Fleisch gewordenen Wort offenbaret, so offenbaret er sich gleichermaßen in dem gepredigten und geschriebenen Wort, welches von Christo zeuget.}. Darum sollt ihr nicht allein gedenken an den Buchstaben, oder an die Stimme über Gottes Wort gezogen, sondern an die göttliche Kraft, die darinnen verborgen ist. Derhalben so habt großen Fleiß, daß ihr nicht mit wenigerm Ernst die heilige Schrift leset, denn wenn ihr das Sacrament des Leichnams und Bluts Christi empfahet. Thut ihr anders, so wird euch Gott verschwinden vor euren Augen, ihr werdet Gift für Arznei aus der Bibel lesen, wie denn Gott dräuet im Propheten Jesaia am 6. und am 28., im Ezechiel am 14., im Matthäo am 13., item zu den Römern am 9. und 10. Capiteln.

Hütet euch vor Verkleinerung Gottes Worts, es sei mit Gedanken oder mit Worten. Wahrlich, wahrlich, Gott läßt es ohne große Verblendung nicht so hingehen, er straft mit großem Ernst solche frevele, freudige Mißhandlung seines Worts, denn keine Frucht kommt davon, kein Geist ist in dem Menschen, der nicht ob Gottes Wort erschrickt, wie Jesaias sagt am 66. Capitel. Auf wen werde ich sehen, spricht der Herr, denn auf den Armen und auf einen demüthigen zerbrochenen Geist, der da erschrickt über meinen Worten? Darum sollt ihr nicht nachfolgen dem freudigen, losen, rohen, falschen Christen, „die jetzt leider zu unsern „Zeiten ein Tisch-Mährchen, ein unnütz Geschwätz und eine spitz%
\quellseite{28}
findige Disputation, die zu Stolz und Hader dienet, aus dem „Worte GOttes machen,“ derhalben sie auch GOtt plaget, daß sie aller Laster voll sind, mehr denn kein Papist\footnote{Röm. 1.}, welcher Früchte werden bald heraus brechen.

Zum andern. Was ihr in der heiligen Schrift leset, das nehmt an für euch selbst. Es gehet euch am meisten an, was in der Schrift geschrieben stehet, es rühret euren Sinn, eure Gebehrde und eure Weise. Euer Leib und Seele regieret nach GOttes Wort. Euren Rath, Trost und Hülfe suchet in der heiligen Schrift. Wenn ihr euch denn selbst also habt wohl unterweiset, so habt ihr schon einen Kunstbrunnen in euer Herz und Gedanken, ja in eure Glieder gegraben, der nimmermehr austrocknen, noch versiegen mag, wie Christus saget Johann. 4. und 7. Davon ihr denn alle Welt möget lehren und tränken, ohne ein mühseliges und arbeitsames Suchen und Lesen zu der Zeit, wenn man predigen, oder Rath geben muß. Denn so man allein aus Büchern, und nicht aus Erfahrung, oder aus dem Herzen lehret, folget keiner oder gar kleiner, Nutzen heraus.

Zum dritten. Machet euch ein kurz- und kleines Registerlein, darein ihr verzeichnen möget die fürnehmsten Puncte und Artikel eurer Lection, daß also die vorige Arbeit die nachfolgende gering mache, und möget allweg zur Hand haben, was euch zu eurer anliegenden Noth dienet.

Zum vierten. Ehe denn ihr anfahet zu lesen, so bittet GOtt, daß er euch wolle führen und leiten in der Schrift, und eröffnen, was euch dienet, und verhalten, was euch nicht dienet; daß er euch wolle behüten vor dem, daß mehr Frage, denn Besserung bringe. Fahret gemeinlich für und für\footnote{D. i. fahret immerdar fort.} im Lesen; habt ein Aufsehen
\quellseite{29}
im Herzen auf GOttes Hand; wie er euch führet, also folget hernach: eilet nicht. Wenn er euch forttreibet, so gehet fort; wenn er still stehet, so stehet auch still, wie geschrieben ist im vierten Buch Mose am 9. Capitel.

\section{Wie die Schrift zu lesen}

Zum fünften. So ihr überdrößig wollt werden zu lesen, so höret auf, daß nicht dieser Ueberdrieß euch zuletzt einen Greuel und Ekel mache über dem Himmel-Brod, zum Schaden eurer Seelen, wie da stehet im vierten Buch Mose am 11. Cap. und zum Römern am 12.\footnote{Soll wohl 1. Cor. 10 heißen.} Wollet auch nicht zu viel wißen, daß euch nicht das Manna zum Schaden diene, wie im andern Buch Mose stehet am 16. Werdet ihr nicht klug sein im lesen, hören oder reden von GOttes Wort, so sollet ihr unsäglichen Schaden erleiden müssen. Die Corinther wollten auch zu viel spielen mit GOttes Wort, drum warnet sie Paulus vor Schaden.

Zum Lesen in der heiligen Schrift vermahnet uns Christus im Johanne\footnote{Joh. 5}, da er zu den Juden spricht: Suchet in der Schrift, denn ihr habt das Leben drinnen, und sie ists, die von mir zeuget. Desgleichen auch Paulus in der ersten Epistel an Timotheum am 4. und in der andern am 3. Auch da er die Eigenschaft eines Bischofs erzählet, spricht er unter andern, er soll lehrhaftig sein. Und zu Tito am 1. sagt er, er soll halten ob dem gewissen Wort der Lehre, auf daß er mächtig sei zu ermahnen durch die heilsame Lehre, und zu strafen die Widersprecher.

\randnotiz{Concordantien.}Wenn ihr leset in der Bibel, so suchet auch die Oerter, welche neben aus an dem Rande verzeichnet sind, wie man es nennet, die Concordantien, und besehet in denselbigen, was zuvor, und was hernach stehet in dem Capitel, wie es zusammen stimme.

Also thut auch in andern Büchlein und Lehrern. Merket auf den Titel des Buchs, darnach besehet, wie alle nachfolgende
\quellseite{30}
Reden sich auf den Titel richten, wohl, oder übel, und überleset ihr alle citation oder allegation aus der Bibel im Original, wie jetzt gesagt ist. Thut ihr das ein halb Jahr, ihr werdet großen Nutzen befinden.

\randnotiz{Ordnung im Neuen Testament zu lesen.}Die Ordnung aber, die ihr im lesen halten sollet in der Bibel, sei diese. An Sanct Matthes Evangelion fahret an, und leset den Text schlecht hinweg, ohne Comment, ein mal oder drei, daß ihr wohl darinnen kundig werdet. Darnach das Evangelion Sanct Marcus, welches in einer Kürze Sanct Matthes Evangelion verfasset, und vergeßet ja nicht die Concordantien zu suchen. Nach diesen leset die Epistel Sanct Paulus an Titon mit großem Fleiß ein mal oder drei, darauf die zwei an Timotheum, als der vorigen Auslegung. Darnach leset die Epistel an die Colosser auch mit großer Acht, auf dieselbige alsdenn die Epistel an die Epheser, welche ist eine weitere Auslegung der an die Colosser. Dann so leset die an die Galater; dazu die an die Römer, als eine Glosse der an die Galater. Zuletzt leset das Evangelion Johannis, als einen Johannes-Segen nach der Mahlzeit der heiligen Schrift, dergleichen auch seine Episteln, die nichts anders sind denn väterliche Vermahnungen zu denen Dingen, so in seinem Evangelio sind beschrieben. So findet ihr auch in Sanct Peters erster Epistel eine Summa aller Dinge, so ihr in den vorigen Büchern gelesen habt, mit Lehren vom Glauben, von Zucht und Liebe, also gemenget und gezieret, daß sie eurem Volk zu predigen soll eine Form und Anfang sein. Die Geschichten der Aposteln reizen den Leser selbst zu fleißigem Lesen und Erkenntniß des Fortgangs und Zunehmens der Evangelischen Lehre; darnach wir uns auch sollen richten zu unsern Zeiten.

\randnotiz{Ordnung im Alten Testament zu lesen.}Am fünften Buch Mose, Deuteronomion genannt, fahret an im Alten Testament zu lesen, darein wird sich der ganze Mose schicken, durch fleißiges
\quellseite{31}
\randnotiz{Concordantien.}übersehen und nachsuchen der Concordantien, so an dem Rande verzeichnet sind. Niemand mag glauben ohne Erfahrung, wie großen Nutzen das bringet, so man die Concordantien und Allegaten hin und her in der Bibel mitzubesiehet. Denn also wird der Leser bekannt in der heiligen Schrift ohne Mühe und Arbeit, findet auch zu Zeiten etwas, das ihm anders wohin dienet, darauf er nie gedacht hatte.

\randnotiz{Historien-Bücher.}Die vier ersten Bücher Mose sind so viel desto lustiger zu lesen, wie viel sie lustiger Historien in sich haben, die den Leser sehr erlustigen. Also thun auch die andern Bücher, in welchen die Historien stehen, als: Josua, das Buch der Richter, Ruth, Samuel, das Buch der Könige, die Chronica, Esra, Nehemia, Esther, und dergleichen. Aber in allen Historien merket darauf, wie GOtt die Gläubigen in seiner Hand trägt, und die Ungläubigen schwerlich strafet.

\randnotiz{Hiob.}Im Buch Hiob findet man eine hübsche Disputation davon: Ob GOtt nicht allein die Sünder und Gottlosen, sondern auch die Frommen und Gottesfürchtigen mit Unglück überfalle? dawider doch die Freunde des Hiob gar stark handeln und disputiren. Hiob aber hat Recht in seinem Fürnehmen, daß GOtt auch den Frommen zu Zeiten Unglück zuschicket, wie denn auch ihm geschehen sei. In dieser Disputation und Handlung findet man gar schöne Dinge von GOttes Gewalt und Weisheit, wie die in den Creaturen erscheinen und erkannt werden.

\randnotiz{Psalter.}Der Psalter soll euer Gesangbüchlein sein, aus welchem ihr Bußliedlein singen sollt eurem lieben Gotte.\footnote{D. i. in herzlicher Liebe zu GOtt dieselbe singen sollt (Ef. 6, 2, 4. 5).} Also, verzeichnet euch die Summa und Inhalt der Psalmen, daß ihr allweg zur Hand haben möget, was euch noth ist. Zu solchen Liedlein vermahnt uns St. Paul\footnote{Coloss. 3.} in der Epistel
\quellseite{32}
an die Colosser, da er spricht: Lehret und vermahnet euch selbst mit Psalmen und Lobsängen und geistlichen Liedern in der Gnade, und singet dem HErrn in euren Herzen. Und alles, was ihr thut mit Worten, oder mit Werken, das thut alles in dem Namen des HErrn JEsu, und danket GOtt dem Vater durch ihn. Dergleichen auch in der Epistel an die Epheser sagt St. Paul:\footnote{Ephes. 5.} Werdet nicht unverständig, sondern verständig, was da sei des HErrn Wille, und saufet euch nicht voll Weins, daraus ein unordig Wesen folget; sondern werdet voll Geistes, und redet unter einander von Psalmen und Lobsängen und geistlichen Liedern. Singet und spielet dem HErrn in euren Herzen, und saget Dank allezeit für jedermann GOtt dem Vater in dem Namen unsers HErrn JEsu Christi; und seid unter einander unterthan in der Furcht GOttes. Wenn ihr ihm also thut, wird euch der Geist Christi solche Andacht dazu geben, daß eure größte Freude darinnen sein wird. Glaubet mir, kein Poet, kein Buhlen-Lied wird euch je so wohl gefallen haben als eben die Psalmen, wenn ihr diese Andacht dazu überkommet. \randnotiz{Salomons Bücher.}Die Sprüche Salomons laßt euch sein für eine Kinder-Schule in der Christenheit. Der Prediger Salomon sei euch ein Spiegel des Welt-Laufs, und wie keine Freude nicht ist, wo nicht ist göttliche Erkenntniß, und wie allein dem gottseligen Menschen wohl sein mag auf Erden. Das Hohe Lied Salomons soll euch der Geist Christi lehren verstehen zu seiner Zeit.

\randnotiz{Propheten.}Unter den Propheten weiß ich euch keine Regel, noch kein Ziel zu stecken, wie ihr euch haltet im Lesen. Denn etliche fahren all den kleinen Propheten an, am Hosea, und richten sich nach dem Hosea in alle nachfolgende kleine Propheten. Etliche gebrauchen des Propheten Amos, als zu einem Ziel und Register, wie denn des Propheten Jesaia unter den großen Propheten. Meines Bedünkens aber, wenn einer den
\quellseite{33}
Jesaiam einmal oder drei fleißig durchlesen hätte, und wäre etwas ein wenig drinnen erfahren, der sollte sich darnach aus allen andern Propheten wohl richten, und sie darauf wissen zu ziehen.

\subsection{Annotationes}

Man hat auch gute Anweisung in etliche Bücher, als: D. Mart. Luthers Auslegung über die Epistel St. Paullus an die Galater, in etliche Psalmen, in die Episteln St. Peters, in das fünfte Buch Mose, und dergleichen Büchlein mehr. Item Philipp Melanchthons Anweisung in die heilige Schrift, in Matthäum, in Johannem, in St. Paulus Episteln an die Römer und Corinth(er). Item Johann Bugenhagens Commentarien in den Psalter, seine Auslegung in die kurzen Episteln St. Pauli, ins Deuteronomion, in Samuelem, und ander Ding mehr. Wenn ihr abgemeldte Auslegungen gelesen habt, könnt ihr euch wohl behelfen in der Bibel fortzufahren, daß euch forthin nicht noth sein wird viel Comments zu lesen. Aber für aller Lection leset ja mit großem Fleiß Locus communes Theologicos, das ist, die gemeine Anweisung in die heilige Schrift Philipp Melanchthons, und suchet alle Sprüche, von ihm citirt und allegirt, aus der Bibel, daß ihr einen Grund der Lehre fassen möget, welches euch großen Nutzen und Frommen bringen wird.

\subsection{Hinderniß des Zunehmens im Lesen}

Ich muß hie auch sagen von Hinderniß des Zunehmens im Lesen. Daß etliche viel lesen, und kleinen, oder keinen Nutzen davon haben, ist dies die Ursach: Sie lesen und studieren nicht darum, daß sie selbst wissen wollen und gelehrt werden in denen Dingen, die sie lesen; sondern lesen allein für andere Leute, dadurch sie ihre Ehre, Ruhm und Gut suchen mögen, oder thuns aus lauterm Fürwitz. Denn wo das nicht in ihnen wäre, würden sie das Lesen wohl lassen anstehen. Das verstehet also: Der rechte Verstand und das rechte Wissen in der heiligen Schrift liegt daran, daß mans so befinde im Herzen, wie mans lieset,
\quellseite{34}
und daß der Leser also gesinnet werde, wie die Lection anzeigt. Denn, die heilige Schrift wird nicht verstanden durch hohe Beschaulichkeit der Vernunft, sondern durch ein heftiges Befinden und Fühlen des Herzens, daß also unser Herz und Wandel der Lection gleich gestalt und ähnlich werde: also, daß alles andere Wissen in der Bibel, ob einer auch alle Sprüche so wohl wüßte, als das Vater Unser, und viel davon könnte schwatzen, dazu alle Auslegungen der Väter und Doctoren auswendig wüßte, so ist es nicht mehr, denn ob er die Worte allein lesen könnte ohne allen Verstand, gegen dem Fühlen und Befinden der Dinge, die einer lieset. Wo das Befinden nicht ist bei dem Wissen, so folget nichts denn unnütz Geschwätz und Disputiren, vergebene Fragen und Leichtfertigkeit der Gemüther, dawider Paulus viel schreibet in den Episteln an Timotheon und Titon.

Dieweil denn aber wenig nach dem Befinden streben, so sie in der heiligen Schrift lesen, so folget, daß auch ihrer wenig die Bibel recht verstehen, werden also lose Schwätzer draus zum Aergerniß der Einfältigen. Denn wie St. Paul sagt zu den Corinthern: Das Reich GOttes stehet nicht in Worten, sondern in Kraft: also thun ihrer auch viele, die in menschlichen Künsten studieren, wo sie nicht mehr das gründliche Wissen suchen, denn Gut und Ehre zu erjagen und den Fürwitz zu büßen, so werden sie wohl ungelehrt bleiben, ob sie gleich Magistri und Doctores werden. Ich hatte mir vorgenommen und auch angefangen euch zu vermahnen, den Pfarr-Stand recht anzusehen, so bin ich zu ferne neben aus gefahren in diesen vorigen Artikeln und Punkten, aber doch nicht ohne Noth und sonderlichen Nutzen, als ich hoffe. Nun aber wollen wir die andern Stücke nach einander sehen.

Zum vierten. So ihr nun alle Tage also reformieret und informieret, abbrechet und aufbauet an euren Sinnen und Sitten,
\quellseite{35}
(wie auch dem St. Paul in den Episteln an Tim.\footnote{1 Timot. 4. 5. 6.} sein lehret, sonderlich da er spricht: Fleuch die Lüste der Jugend,\footnote{2 Timot.} und euch selbst ganz, wie ein Kind, unter die Sporen nehmet, zu lehren, zu ziehen innen und außen nach der Regel der Schrift: so werdet ihr euch selbst erst lernen kennen, viel in euch selbst erfahren, davon euch alle Welt nicht sagen möchte. Derhalben folget dem Geist, der euch wohl treiben wird zu seiner Zeit.

Aber in andern Uebungen eures Gemüths, durch Studieren in menschlichen Künsten, Historien, Rechten, natürlichen Dingen und dergleichen, oder in Arbeit des Leibs, sollt ihr eine Ordnung und Regiment haben. Denn ohne vernünftige Ordnung schaffet ihr gar nichts. Habt auch eine ordentliche, zeitliche, mäßige Zeit im Schlafen, und ohne große Noth verändert sie nicht, daß ihr nicht vom Schlaf kommet. Also thut auch mit Speis und Trank. Werdet ihr euren Leib nicht in eine leidliche, füglich[e] Ordnung richten, im Schlafen, Essen und Trinken, so habt ihr nimmer Ruhe, noch Freude, werdet auch nichts ausrichten. Davon lehret St. Paul an den Timotheon\footnote{1 Timot. 5.} in seiner ersten Epistel.

\randnotiz{Hausgesinde.}Eures Hauses und Hausgesindes Regiment soll wohl geordnet sein, denn so spricht Paulus zu Timotheon\footnote{1 Tim. 3.}, daß ein Bischof unter andern so soll geschickt sein, daß er seinem eigenen Haus wohl fürstehe, und gehorsame Kinder habe, mit aller Redlichkeit, und sagt: So aber jemand seinem eigenen Hause nicht weiß fürzustehen, wie wird er die Gemeine Gottes versorgen? Auch euer Studieren soll eine Ordnung, Maß und Ziel haben, wie ihr in andern Büchlein gelehret werdet.

\randnotiz{Arbeiten ist ehrlich.}Dazu fleißet euch auch irgend eine Arbeit zu lernen zur Uebung des Leibes, und zu verzehren nützlich die Zeit, darinnen ihr zu andern Uebungen unlustig seid. Arbeiten ist ehrlich, so daß man auch etliche Für%
\quellseite{36}
sten findet, die das Drechsler-Handwerk und dergleichen lustige Arbeit mehr gelernt haben, die Zeit dadurch zu kürzen und zu vertreiben. Denn ohne Arbeit des Leibes mögt ihr nicht lang ohne Schaden eurer Sinne, eures Dings und Geschäfts warten, ihr müßt zuletzt auslaufen, hin und herspazieren für die lange Weile, welches ohne euren merklichen Schaden nicht wohl kann geschehen. Auch wird es euch nützlich sein im Hause zu eurem Gebrauch. Ihr vertreibet auch hiermit viel melancholischer Phantasei. Es ist auch besserlich euren Nächsten und Zuhörern. Es macht gar böses Geblüt, wenn die Bauern in ihrer Arbeit gedenken an der Pfaffen faul Leben, sonderlich zu unsern Zeiten. Wie fleißig der große Prediger Paulus gearbeitet habe mit seinen Händen, leset ihr in den Geschichten der Apostel\footnote{Apg. Gesch. 20.} und in den Episteln an die Thessalonicher\footnote{1 Thess. 2. 2 Thess. 3.}.

Zum fünften. Zwei Dinge sind, welchen alle andere Arbeit des Studierens weichen sollen: \randnotiz{Andacht.}Andacht und eures Bruders Noth. Als: Wenn euch ein guter Gedanke und Andacht zufället von GOtt, oder von eurer Seelen Heil und Seligkeit, item ein guter Verstand eines Spruchs irgend in der heiligen Schrift: so wartet seiner aus mit \randnotiz{Göttliche Einsprechen.}allem Fleiß, und wisset, daß euch GOtt väterlich heimsuchet, und jetzt GOttes Zeit sei, mit euch zu handeln. Alda sollt ihr denn, als ein gehorsamer Diener, alle Arbeit lassen anstehen, wo sie euch daran hindert. Ist es aber um die Arbeit also gethan, daß dabei solch Aufnehmen auf GOttes Einsprechen\footnote{Dadurch verstehet er keine unmittelbare Eingebung, sondern daß man drauf merken solle, wenn der Heilige Geist durchs Wort an das Herz anklopfet und dasselbe zur Andacht erwecket.} mag bestehen, so vollführet die Arbeit mit für und für. Hütet euch ja, daß ihr kein solches Heimsuchen GOttes
\quellseite{37}
vergebens lasset fürüber gehen, welches wahrlich nicht ohne großen merklichen Schaden geschehen mag. Denn unser Schicken zu andächtigem Lesen oder Beten hilft allein nichts, sondern Gott kommt selbst durch das gelesene oder gehörte Wort, und klopfet an; wir sollen auf ihn fleißig warten, und dies Warten ist die beste Schickung zu der Andacht.

Also haben die heiligen Patriarchen und Propheten oft gehört Gott mit ihnen reden ins Herz, aber wir wissen nichts mehr davon, darum sein wir auch so gottlos.

\randnotiz{Unterscheid der Einsprechungen.}Dazu wird euch denn nun die heilige Schrift noth sein, daß ihr merkt, ob solches Einsprechen der Schrift gemäß sei, oder nicht, daß ihr wisset, heilsame und betrügliche Einsprechung zu unterscheiden und zu erkennen, bei der Regel des Glaubens in der Schrift gelehret, wie St. Paul sagt zu den Römern\footnote{Röm. 12.}: Hat jemand Weissagung, so sei sie dem Glauben ähnlich. Dergleichen handelt er in der Epistel an die Corinthier\footnote{1 Cor. 2.}, und St. Johannes\footnote{1 Joh. 4.} in seiner ersten Epistel am vierten Capitel. Davon ließ viel gutes Dings im Buch, Collationes Patrum genannt, darinnen ein guter Theil des Christenthums stehet, daß man wisse Unterschied zwischen rechter und falscher Andacht, zwischen der Wahrheit und Betrug der Einsprechungen.

\randnotiz{Des Nächsten Noth helfen.}Das andere Stück, dem alle selbst fürgenommene Uebungen weichen sollen, ist Nothdurft eures Nächsten, eures Weibes, eurer Kinder, eures Hausgesinds, eurer Freunden und Feinden, einheimischen und fremden; und das also: Eurem Weib und Kindern sollt ihr förderlich dienen, sie trösten, ihnen helfen und rathen, und dergleichen Werk der Liebe erzeigen in Gesundheit und Krankheit. Davon soll euch keine selbsterdachte Uebung hindern, und es soll euch keine Scham davon ziehen. Denn was ihr Weib und Kindern thut, das habt ihr Gott selbst gethan. Unsere Eltern habens auch gethan. Darnach sollt ihr
\quellseite{38}
hülflich und tröstlich sein eurem Hausgesind, euren Nachbarn und andern, die es bedürfen, ihnen zu-willfahren nach Vermögen. Denn ein Mensch ist dem andern zu gut geschaffen. Derhalben, welcher seinem Nächsten nicht hilft, noch räth, so viel er vermag, der ist weder Christ, noch Mensch. Es soll kein Mensch in eurem Pfarrspiel sein, der nicht Dienst, Rath, Hülff und Trost von euch empfahe.

\subsection[Zum sechsten]{Zum sechsten. Sonderlichen Fleiß sollt ihr ankehren, die Betrübten zu trösten. Dazu gehöret große Weisheit und Fürsichtigkeit}

\randnotiz{Betrübte zu trösten.}Mein lieber Herr und Vetter, ich bitte euch gar freundlich, ihr wollet in Besten aufnehmen, daß ich euch solch Ding zuschreibe, und euch hierinnen lehre und unterweise. Ich thue es wahrlich euch und vielen zu gute. Ich habs an mir und andern mehr erfahren, was für Fährlichkeit in diesem Stück liegt. Darum sollt ihr erstlich diese Tugend an euch haben, daß ihr verschweigen könnt alles, was ihr hört von den Leuten, die euch ihre Noth klagen und fürtragen.

\randnotiz{Heimlichkeit verschweigen.}Heimlichkeit verschweigen ist eine große Kunst, die wenig Leute können. Laßt euch weder Weib, noch Freunde so lieb sein, daß ihr vor ihnen davon etwas redet; trauet niemand, denn der Teufel fleißet sich auf allen Orten, Unrath und Verleumdung auszurichten, und zu treiben auf die Prediger und Lehrer des Evangelions. Denn so man sollte erfahren, daß ihr einem andern seine Heimlichkeit offenbaretet, würde er, oder andere Leute sagen, ihr wäret ein Schalk und Schmäher; daraus denn dem Evangelio Schmäh und Schande entstünde. Auch sollt ihr keinem Menschen seine Klage oder fürgetragene Noth verlachen. Einem jeglichen dünket sein Anliegen das größte sein. Denn so ihr einmal offenbaret eines andern Heimlichkeit, oder seine Klage und Noth für Spott hieltet, erschrecket ihr ihn und viel andere mehr, daß sie euch hinfürder nicht mehr dürften um Rath fragen und ihr
\quellseite{39}
Gewissen euch eröffnen; welches denn zu großem merklichen Schaden der betrübten Gewissen reichen würde. Ein Pfarrer oder Prediger soll die letzte Zuflucht sein der Betrübten und Geängsteten auf Erden, sonderlich wo es die Gewissen belanget, also, ob einer keinen Rath noch Hülfe bei euch finde, daß er dennoch einen freundlichen Gruß, ein herzliches Mitleiden und einen verschwiegenen Mund bei euch gewarten möge.

\subsection{Zum siebenden}

\randnotiz{Rathen in schweren großen Händeln.}Es ist bißher seltsam gewesen, daß man den gemeinen läglichen Pfaffen, wie man sie genennet hat, viel größer, schwerer Händel fürgeleget hätte, sonderlich aber habens weise und die fürnehmsten Leute nicht gethan. Man ist gangen irgend zu großen Doctoren, zu den berühmten Beichtvätern und Predigern in die Klöster, darum auch euer einer nicht so viel Erfahrung haben mag, noch viel Ursach gehabt, großen, schweren und tapfern Händeln nachzutrachten, nachzudenken, lesen oder fragen, damit man redlichen, ehrbarlichen, getreuen, guten Rath geben möchte. Ich weiß, was ich sage. So aber ihr und unsers gleichen anfahet das Evangelion zu predigen, und viel Volks zulauft, und glaubt eurer Lehre, und ein Mißfallen hat an den vorigen Rathleuten, wird man anfahen in großen Sachen Rath bei euch zu suchen. Darum ists noth, daß ihr die Leute freundlich empfahet, ernstlich und fleißig zuhöret, die Händel wohl begreifet, euch auch züchtig in Geberden und Worten haltet, niemand nachredet, niemand scheltet, niemand lästert, niemand verachtet, mit jedermann redet, als sichs ziemet. Ihr werdet nicht jedermann kennen, es werden viel Täuscher und falsche Brüder unter guter Gestalt zu euch kommen, Frauen und Männer, und wo sie euch unzüchtig\footnote{D. i. unbescheiden, und unverständig.} erfänden, würden sie euch ausschreien für jedermann, zu Spott dem heiligen Evangelio.

\quellseite{40}
Nicht ohne Ursach hat St. Paul so fleißig gelehret an vielen Orten, wie ein Christ in Worten sittig und züchtig sein soll, und sonderlich zu Timotheo\footnote{1 Tim. 5.} spricht er: Der jungen Witwen entschlage dich, denn weil sie geil worden sind wider Christo, so wollen sie freien. Und bald hernach: Halt dich selber keusch. Für allen Dingen sollt ihr euch lassen finden getreu, redlich und verschwiegen, darnach weislich rathen, denn der Heilige Geist gibt auch Fürsichtigkeit und Rath in menschlichen Handlungen, wie im Propheten Jesaia\footnote{Jesaia 11.} stehet. Wisset ihr nicht dazu zu rathen nützlich, so schweiget lieber stille, daß ihr niemand verführet, und daß nicht eure Unfürsichtigkeit getadelt werde. So ihr Gottes Geist und Weisheit, Verstand und Klugheit bittet, wird euch Gott nicht lassen, wie Christus im Luca\footnote{Jac. 1. Luc. 11.} sagt.

Hierzu werden euch nicht übel dienen die Historien der heiligen Schrift, die Lehre Salomonis, St. Paulus und andere dergleichen Bücher mehr. Item wird euch auch nicht schädlich sein, daß ihr zu Zeiten leset die heidnischen Historien, die Bücher Senecae, Ciceronis und anderer, item Kaiserliche, Römische, Stadt- und Land-Recht, welches alles noth sein wird, in zeitlichen, äußerlichen menschlichen Sachen zu rathen. Aber die Sachen, so das Gewissen betreffen, soll man handeln mit Gottes Wort. Und eigne Erfahrung in viel Leiden und Aengsten, Trübsalen und Widerwärtigkeit hilft viel dazu. Darum mögt ihr auch wohl lesen das Buch von dem Leben der Altväter, Collationes Patrum, die Episteln Hieronymi, Augustini und anderer dergleichen mehr, etliche Büchlein Gersonis, Keisersbergs, Tauleri, und was des Dings mehr ist, in welchen man findet große Erfahrung und treue Warnung, zu trösten die ängstigen
\quellseite{41}
Herzen und Gewissen, die von dem Teufel ohne Unterlaß werden angefochten.

Es wird auch fürkommen mancherley Leiden, Angst, Noth, Bekümmerniß, Widerwärtigkeit, Trübsal und Anfechtung, dazu euch vonnöthen sein wird göttliche Weisheit, ob ihr gleich allen Fleiß ankehret.

\randnotiz{Böse, lose Leute.}Davon wissen rohe, lose Leute nichts zu sagen, wischen auf, schelten Pfaffen und Mönche, verwerfen alle alte Ordnungen, schwatzen vom Evangelio ohne Erfahrung, ohne Befinden und Fühlen geistlicher Dinge; wissen so viel von gemeinem menschlichen Leben, als eine Kuh vom Mittag; sind ohne Zucht, ohne Ehr, und zu nichts nütz, denn Schand, Laster, Unruhe und Unglück anzurichten, ja Land und „Leute schmähen, oder hindern in andern Sachen Gottes Wort durch ihr rohes, loses Leben.“ Aber ihr, mein lieber Vetter und Herr, sollt Christo und euch die Schande nicht anthun, daß ihr solchen Gesellen folgen wollet. Dies will ich euch auch vermahnen, welches ich doch schier vergessen hätte, daß, wenn ihr einem Handel nicht klug genug seid, wollet andere getreue, weise Leute hierinnen Rath fragen, oder die Menschen daselbst hin mit gutem Abschied weisen.

\subsection{Zum achten}

Lernet von Paulo in seinen Episteln\footnote{1 Thess. 2.}, und auch in den Geschichten der Aposteln\footnote{2 Cor. 1. 2 Cor. 12. Röm. 15. Gal. 6. Ap. Gesch. 20.}, wie herzlich und mütterlich er mit Bitten, Weinen und allem Mitleiden einen jeglichen gelehret hat, ermahnet und getröstet Tag und Nacht. Leset dieselbigen Capitel, werdet ihr gar hübsche Sprüche drinnen finden. Auch sagt St. Paul in der andern Epistel an Timotheon\footnote{2 Tim. 2.}: Ein Knecht des Herrn soll nicht zänklich sein, sondern väterlich gegen jedermann,
\quellseite{42}
lehrhaftig, der die Bösen tragen kann, der mit Sanftmuth strafe die Widerspenstigen. Desgleichen in der ersten Epistel\footnote{1 Tim. 3.} spricht er unter andern, daß ein Bischof \randnotiz{Bischof.}sein soll fleißig, nüchtern, züchtig, sittig, gastfrei, lehrhaftig, nicht weinsüchtig, nicht bissig, nicht schändliches Gewinns gierig, sondern gelinde, nicht haberig, nicht geizig. Und in der Epistel an Titon\footnote{Tit. 1.} sagt er fast dieselbigen Worte: Ein Bischof soll untadelig sein, als ein Haushalter GOttes, nicht hoch von ihm selbst halten, nicht zornig, nicht weinsüchtig, nicht bissig, nicht schändliches Gewinns gierig, sondern gastfrei, gütig, züchtig, gerecht, heilig, keusch, und halte ob dem gewissen Wort der Lehre, auf daß er mächtig sei zu ermahnen durch die heilsame Lehre, und zu strafen die Widersprecher. Mein lieber Herr und Vetter, bittet auch GOtt oft und fleißig für angefochtene Menschen, denn groß ist die Qual und die Angst der Betrübten.

\randnotiz{Treulich den Gewissen rathen.}Zum neunten. Habt große Sorge dazu, daß ihr niemand verkürzet oder verführet mit eurem Rathen, daß ihr nicht darnach ein ewiges Reuen müsset tragen euer Lebenlang, und sonderlich in euer letzten Noth deshalben in Angst fallen, wenn ihr meinet, GOtt habe vergessen. Darum so thut alle Dinge mit gutem Gewissen und wohlbedachtem Gemüth, dazu mit möglichem Fleiß. Sehet GOtt in die Hände, Er wird euch helfen. Welches vorhin ist ein ängstlich, mühsam Beichten gewesen, das soll jetzt unter den Christen gelehret werden in ein zuversichtigs getreues Rathfragen und Rathgeben, aus brüderlicher Liebe für Christo. Denn er sagt also im Matthäo: Wo zween unter euch eins werden auf Erden, warum es ist, daß sie bitten wollen, das soll ihnen widerfahren von meinem Vater im Himmel. Denn wo zween oder drei versammelt sind in meinem
\quellseite{43}
Namen, da bin ich mitten unter ihnen. O wie ein theurer Mensch ist das, bei dem jedermann Rath, Trost, Mitleiden, Treue, einen verschwiegenen Mund und ein christlich, lieblich Herz findet in seinen Nöthen und Anfechtungen! Von einem solchen Manne sagt Salomon in den Sprüchen:\footnote{Spr. Sal. 10.} Der Mund des Gerechten ist ein lebendiger Brunn. Item: Der Mund des Gerechten bringet Weisheit. Dergleichen Sprüche viel findet man in der heiligen Schrift. Wenn ihr die Worte Sanct Pauls wohl wäget in den Epifteln an Timotheon und Titon, findet ihr Grund und Ursache der Dinge aller.

\subsection{Zum zehenden, von dem Predig-Amt}
\randnotiz{Predig-Amt.}%

Erstlichen und vor allen Dingen bittet GOTT, daß er sein Wort lege auf eure Zunge, und dieselbige regiere nach seinem Willen, zu Heil und Seligkeit der Zuhörer. Denn euer Anschlag und Fürnehmen gilt hier nichts. GOTT regieret die Zunge, wie Salomon sagt in den Sprüchen:\footnote{Spr. Sal. 16.} Der Mensch setzt ihm wohl für im Herzen, aber vom HErrn kommt das Antworten der Zungen. Bittet GOTT, daß er euch nicht laße irren, euch und den Zuhörern zu Schaden, wie Esekiel\footnote{Esek. 14.} sagt im vierzehenden, und bittet ihn, daß er euch gebe also zu halten im Herzen, wie ihr lehret mit dem Munde. Denn wo nicht Herz und Mund gleich ist, sündiget man ohn Unterlaß, ob man auch nichts denn von GOTT und Christo redete. Was nicht aus dem Glauben hergehet, das ist Sünde, wie Sanct Paul zu den Römern\footnote{Röm. 14.} sagt.

\randnotiz{Fleiß haben im lesen und reden.}Zum eilften. Versuchet GOTT nicht, thut alles, wozu euch GOTT hat geruffen und Kraft geben, das ist, leset mit Fleiß, schreibts zusammen, haltets gegen die Schrift, bedenckts vorhin wohl, obs den Zuhörern diene, oder
\quellseite{44}
nicht, ob nicht mehr Schaden, denn Nutz, daraus komme: so werdet ihr oft finden, wenn ihr eine Predigt vorhin zwei- oder dreimal bedacht habt, wie viel euch darin mißfallen wird, und wie nöthig sei ein fleißiges Wahrnehmen der Materien und der Wörter. Man redet bald ein Wort, das zu großem Schaden dienet, welches man nimmermehr kann herwieder bringen. So läßt GOtt nicht schimpfen mit seinem Worte. Das Wörtlein und diese gemeine Rede, da man spricht: Ich hätts nicht gemeinet, daß es also ergehen sollte, ist auch für den Menschen spöttlich, wie viel mehr für GOtt und sonderlich wenn es GOttes Sache und der Seelen Seligkeit antrifft. Das wisset fürwahr, daß kein Wörtlein dem Gerichte GOttes entrinnen wird, Stein und Holz wird Zeugniß von eurer Predigt geben, denn sie eben geschaffen sind und erhalten werden durch das göttliche Wort, welches ihr prediget.

\randnotiz{unnütze Schwätzer und unbescheidene Prediger.}Zum zwölften. Ich habe jetzt gesagt, daß ihr ohne fleißiges, andächtiges, fürgehendes Gebet, nicht sollt auf den Predigtstuhl gehen. Denn wahrlich, der Teufel feiert nicht, er wird euch hindern an nützlichen Worten, oder wird euch treiben etwas schädliches zu reden ohn euren Dank, wo er nur mag, davon euch groß Leid erwachsen wird: Als jetzt viele unnütze Schwätzer unbescheiden auf der Kanzel schwätzen, deshalb sie in große Pein und Leiden fallen, sagen denn: ja ich leide um der Wahrheit willen. Nicht also, du leidest um deiner Narrheit und Unbescheidenheit willen. Wenn du in sterbender Noth liegst, wird dir dein eigen Gewissen absagen, wenn du dich erkennen wirst, und sprechen: Ists doch noch nicht Zeit, nicht statt, nicht noth gewesen, daß ich dies, oder jenes sagete; so hab ich GOttes Zorn dadurch verdienet, und leide das nicht um GOttes willen, sondern von meiner Narrheit und Thorheit wegen, und habe dadurch un%
\quellseite{45}
zeitlich gereizet die Feinde, wider Gottes Wort zu handeln und sich daran zu vergreifen, habe ihnen mit meiner Unbescheidenheit guten Schein der Verfolgung gemacht, und die Zuhörer Gottes Worts beraubet, daß sie hätten mögen im Friede lernen, hätte meine Thorheit sie nicht gehindert. Wie willt du dann bestehen, so dir dein eigen Gewissen in der Wahrheit absaget? Du wirst wahrlich schreien: O ihr Berge, o ihr Hügel, fallet auf mich, und bedecket mich für dem Zorn Gottes, welchen ich mit meinem predigen verschuldet habe! Darum, mein lieber Herr und Vetter, wappnet euch fürhin mit Gebet und anderm christlichen Harnisch, davon Sanct Paul schreibet in der Epistel an die Epheser\footnote{Ephes. 6.}. Denn so ihr auftretet zu predigen, so stehet ihr an die Spitzen, zu fechten und zu kämpfen, nicht wider Fleisch und Blut, sondern wider die Fürsten und Gewaltigen, mit den Welt-Regenten der Finsternis in dieser Welt, mit den Geistern der Bosheit unter dem Himmel. Fehlet ihr derselbigen, sie werden wahrlich eurer nicht fehlen.

\subsection{Zum dreizehenden}
\randnotiz{Demüthiglich predigen.}%

Gottes Wort sollt ihr predigen mit ernstlichem demüthigen Herzen und Gebehrden. Hütet euch für trotzigem stolzen Sinn und Gebehrden, oder ihr fallet dem Teufel in die Stricke. Nicht mit stolzem Trotz, sondern mit Demuth und Sanftmuth soll man Gottes Wort lehren, das macht den Teufel matt und müde. Der demüthige, sanftmüthige Christus will demüthiglich und sanftmüthiglich geprediget werden. Thut ihr anders, so sündiget ihr für Gott, und werdet zu Spott vor dem Teufel, auch vor der Welt. Denn so man den demüthigen, sanften, freundlichen Christum prediget, soll man ja demüthige, sanfte Worte und Gebehrden gebrauchen, welches auch das gemeine Urtheil der Menschen anzeiget. Euer innerlich
\quellseite{46}
Auge soll mehr über sich zu GOtt sehen in eurer Predigt, (als ob ihr alle Worte von ihm höret und für ihm redet,) denn daß ihr auf irgend etwas anders merket. Dies Pünctlein behaltet und übet es, den Nutz und Frommen werdet ihr bald greifen.

\subsection{Zum vierzehenden}

\randnotiz{Rhetorica ist nutz.}Die Rhetors und Kunst-Redner haben etliches Ding geschrieben, welches euch nütz sein mag zu fügligem, förmlichen Fürhalten eurer Lehre, welche Regeln und Weisen die Rhetorica zeiget, auch wie ihr anderer Lehrer Schrift baß verstehen möget. Dies sollt ihr nicht verachten, denn der Heilige Geist schüttet es mit seinem Trichter ein wunderbarlich, so man ein Ding wohl natürlich haben mag. Derhalben unterlaßet nicht zu lesen, was hierinnen Cicero, Quintilianus, Erasmus, Philippus Melanchthon und andere mehr schreiben, oder geschrieben haben.

\subsection{Zum fünfzehenden}

\randnotiz{Die Zuhörer nicht mit Predigten zu überschütten.}Ueberschüttet eure Zuhörer nicht mit zu vielen und langen Predigten, und das auch nicht mit zu Unzeiten, wie die thun, die täglich wollen eine lange Predigt machen, als hätte man sonst nichts zu schaffen, denn ihrem Predigen zuzuhören, wissen nicht, daß die Zuhörer nicht mit des Predigers Ohren, sondern mit ihren eigenen Ohren hören. Die Feiertage sollen in Ruhe des Leibes und in Christlicher Handlung, beide der Lehre und der Sacramente, zugebracht werden. An den andern Tagen soll man kurz predigen, und ausgelesene Pünctlein, welche ein einfältiger Christ leichtlich und wohl behalten möge, und denselbigen Tag über in dem Gedächtniß umwälzen und seine Seele damit speisen.

Ich kenne einen Evangelisten oder Prediger, da man ihn fragte, obs auch nütz wäre, daß man alle Tage in einer, auch sehr kleinen Stadt zwo oder drei Predigten sagete, auch vormittags, antwortet er: Es wäre gut, und zog herzu den Spruch
\quellseite{47}
Sanct Pauls in der Epistel an die Colosser\footnote{Coloss. 3.}, da er also sagt: Laßt das Wort Gottes in euch reichlich wohnen. Das heißt, meine ich, die Schrift mit den Haaren ziehen auf unser Gutdünken und Wohlgefallen. Dieser Prediger ist noch ein junger Müller, hat nicht viel Säcke gebunden, darum ist er noch frisch und lustig, läßt ihm wohl sein, wenn man ihm zuhört. Er ist vorhin eines solchen Zulaufens und Ruhmes ungewohnt gewesen. Gott wolle ihm geben Gnade, daß ihm nicht geschehe, wie dem Narcisso, von dem die Poeten fabuliren. Ihm wäre gut, daß er läse Gersonem über den Evangelisten Marcum, da fände er treue Warnung, die einem neuen und jungen Prediger fast nütz und noth sein zu wissen.

\randnotiz{Ein Esel wollte ein Löwe geachtet sein.}Ich muß sagen, was ich gedachte. Da der Cardinal Raymundus einst gen Heilbrunn kam, war das Volk solcher Herrlichkeit ungewohnt, und ihm doch fast angenehm. Da das des Cardinals Hofleute merkten, bracht ein jeglicher Stallknecht herfür, und wollte so fast prangen, als der Cardinal selbst. Also thun ich und meines gleichen auch. So GOtt der Allmächtige erhöhet hat Doctor Martin Luther, Philipp Melanchthon, und andere dergleichen mehr, in großer Kunst, in heilsamer Christlicher Lehre, und in einem ehrbaren Leben auch für der Welt, also, daß die ganze Christenheit ein \randnotiz{Unverschämte stolze Prediger.}Auge auf sie habe: so komme ich und mancher Fischer daher, und nehmen uns an, wir sind Lutherisch oder Evangelisch, wie man sagt, und wollen nicht weniger gehalten sein, denn sie gehalten werden, und ohne Kunst, ohne Zucht, ohne geistliche Erfahrung und ohne Christliche Bescheidenheit wollen wir schelten die Pfaffheit Mönchheit, und alte Gebräuche (welche eines theils möchten vielleicht dem Christenthum unschädlich sein) mit Füßen treten, und sagen
\quellseite{48}
(doch daneben so wenig und kalt vom Reich Christi, daß man wohl merkt, daß wir inwendig Affen sind und Esel, von außen mit Menschen- oder Löwen-Häuten bekleidet, und nichts weniger sind denn Lutherisch oder Evangelisch. Ja lieber Gesell, man höret dir wohl zu, und du meinest, man soll dich für den Luther halten; aber die Zuhörer gedenken: Vox quidem est vox Jacob, manus autem manus sunt Esau, das ist, die Stimme ist Jakobs Stimme, aber die Hände sind Esaus Hände\footnote{1 Buch Mose 27.}. Item: Dat sine mente sonum, das ist, es ist ein Getön ohne Verstand. Item: Sermo est, non Virtus, das ist, Worte seien’s ohne Kraft. Wirst du vorhin von GOtt, mit abgemeldeten Lehrern, in die Hölle der Anfechtung geführet, und wieder heraus gen Himmel, daß die Welt, auch der Teufel befinde GOttes Kraft in dir, denn magst du gebärden wie Luther und andere. Wir loben abgemeldte nicht allein darum, daß sie den Pabst aus unsern Herzen vertrieben, sondern, daß sie uns das Reich Christi gelehret, und dasjenige mit Worten und guten Exempeln fürgebildet haben, und auch sonst weise, gelehrte Leute sind. Wir aber, die weder gaben, noch Eier legen können, und uns baß auf weintrinken, denn auf geistliche Anfechtung verstehen, wissen auch unsers Nächsten Anliegen, so er uns beklagt, weder zu schweigen, noch zu rathen, ja spotten seiner oft dazu, wollen doch gar herrlich solchen gelahrten, trefflichen Leuten gleich sein und ihnen gleich gehalten werden; man siehet wohl, daß wir Esel sind, und nicht Löwen. Aber mein lieber Herr und Vetter, GOtt der Allmächtige hat euch von Natur eine Bescheidenheit geben, GOttes Geist wird sie euch wohl mehren. Derhalben so hütet euch für solchem ärgerlichen prangen, bleibt in der Maaß, die euch GOtt hat dargemessen, wie St. Paul anzeiget zu den Römern\footnote{Röm. 12.} und Corin%
\quellseite{49}
thern\footnote{1 Cor. 12.}, und prediget ja den demüthigen Christum demüthiglich, so werdet ihr bei GOtt und bei den Menschen Ehre und Lob überkommen.

\subsection{Zum sechzehenden}

\randnotiz{Anfang Evangelischer Lehre.}Ihr habt bereits einen Anfang gemacht, wäre wohl nicht noth für euch zu lehren, wie bescheiden der Anfang Evangelischer Lehre sein soll. Aber andern guten Freunden zu gefallen, welchen dies Büchlein zu Nutz, so viel GOtt Gnade gibt, kommen wird, will ich ein wenig etwas davon sagen. Ein Prediger muß ihm selbst fürsetzen die zwei Reiche, eines des Teufels, das andere Christi. Der Teufel besitzt das Herz durch ein gottlos Wesen\footnote{1 Buch Mose 6, 8.}, mit Blindheit, eigner Liebe, Verstockung und andern Lastern mehr. Denn alle Sinne und Gedanken des Menschen sind zum Bösen gerichtet von Jugend auf. Derhalben er in GOttes Zorn ist, sein Gewissen zappelt alleweg für GOtt, mag auch kein Vertrauen zu GOtt haben, und keine Liebe zum Nächsten, suchet also stets eigen Lob, Ehre, Nutz und das Seine. Er fähet wohl viel seltsames Wesens an, ja mancherlei Weise GOtt zu dienen, aber kein herzliches Vertrauen setzet er auf GOttes Barmherzigkeit. Viel Menschen sind so blind und verwegen, daß sie sich nimmer nicht unterstehen mit GOtt zu handeln. Fähet aber ein solcher Mensch an mit GOtt zu handeln, so folget er nicht GOttes Wort, sondern seinem eigenen Dünkel, will mit eigenen erlesenen Werken GOttes Gnade und Gaben kaufen, das ist, er getrauet ihm selbst mehr, denn GOTT.

\randnotiz{Teufels Reich.}Also sündiget ein solcher Mensch in all seinem Thun und Lassen, wie gut, wie klug auch sein Ding für der Vernunft scheinet. Der Mensch läßt sich davon auch nicht abweisen, als hätte er unrecht. Summa Summarum, Fleisch und Blut will unverachtet sein in seinem Urtheil und Fürnehmen. Und das ist
\quellseite{50}
eigentlich des Teufels Reich, Gott nicht erkennen, ihn weder fürchten, noch lieben in der Wahrheit. Daraus denn folget, daß auch der Mensch keine Creatur recht erkennen kann, noch keiner wohl gebrauchen.

\subsubsection{Reich Christi}

Wider dies Reich des Teufels ist das Reich Christi, welches da stehet in rechter Erkenntniß Gottes, in Gottesfurcht und Liebe, in Erkenntniß Gottes Willens und seiner Gebote, in Erkenntniß eigener Sünde und Gottes Barmherzigkeit. Ein solch Herz befindet, wie schwer ihm der verschuldete Zorn Gottes zu tragen ist, wie unsägliche Pein dem Sünder für Augen stehet, wie ihm selbst der Mensch so gar weder rathen, noch helfen mag, sondern alles Vermögen, Rath und Hülfe muß herkommen aus Gottes Barmherzigkeit, ohne unsern freien Willen und Vernunft, ohne unsere undienstliche gute Werke.

\randnotiz{Alles haben wir von Gott durch Christum.}Kürzlich: Ein solch Herz erkennt im Grunde, wie böß es selbst ist, und wie gut Gott gegen sich\footnote{D. i. gegen ein solch Herz, oder solchen Menschen.} ist; wie Gott, aus Gnaden, ohne unser Verdienst, seinen Willen uns durch Christum gelehret, und Christum für uns geopfert hat, dazu in Christo uns alle seine Güte beweiset und versprochen, also, daß alle, die solches Vertrauen zu Gott durch Christum haben, sollen hinfüro liebe Kinder sein. Gott gibt uns auch durch Christum seinen Heiligen Geist, der uns in unser Herz drucke solche Liebe und Erkenntniß Gottes.

Ein solch Herz fähet an Gotte zu vertrauen im Leben und im Tode, setzet alles Heil auf Gottes Hulde, verachtet sich selbst, thut jedermann Guts, angesehen die große unverdiente Liebe Gottes gegen sich. So nun Gott diesen Menschen also gnädiglich ansiehet, so gefallen ihm auch alle Werke desselbigen Menschen wohl. Gott regieret diesen Menschen, und ob er etwan sündiget, so
\quellseite{51}
hilft ihm GOtt wieder auf. Ein solch Mensch erkennet, daß GOtt kein Werck das gefällt, denn ihm also durch Christum vertrauen. Er hat auch in allen Dingen ein Aufsehen auf GOttes Wort, als auf eine Lucern in der Nacht dieser Welt, wie der Psalm sagt: Dein Wort ist meines Fußes Leuchte, und ein Licht auf meinem Pfad. Desgleichen sagt auch St. Peter\footnote{2 Petr. 1.} in seiner andern Epistel am ersten Capitel.

\randnotiz{Wie ein Prediger die zwei Reiche predigen soll.}Diese zwei Reiche nehme ihm ein Prediger für, und gedenke, wie er auch seine Zuhörer solches lehre erkennen, und sahe an vom Reich Christi zu predigen, daß es in Furcht und Liebe GOttes stehe, und theile die Rede nach Geschicklichkeit der Zuhörer. Sind die Zuhörer frech und muthig, so erzähle er GOttes Zorn über unsere Sünde. Damit erschrecke er der Gewissen Hochmüthigkeit seliglich, daß sie darnach erfreuet werden in GOttes Barmherzigkeit gnädiglich. Wenn aber die Gewissen wohl erschreckt sind mit GOttes Zorn, so predige er denn von GOttes Gnad, und ziehe also mit dem Wort den Sünder in das Reich Christi, darin nichts denn Gnade regieret. Sind aber die Gewissen der Zuhörer erschrocken, so sahe der Prediger an zu sagen von GOttes Barmherzigkeit in Christo, also, daß es alles durch Christum muß zugehen. Denn dadurch wird man seliglich erschreckt, so man siehet den großen Zorn GOttes im Leiden Christi, den kein Mensch versühnen mochte, Christus mußte es mit seinem Blut thun, und zahlen für alle Sünde der ganzen Welt.

\randnotiz{Das Leiden Christi.}Auch findet man einen gnädigen Trost im Leiden Christi, daß GOtt aus lauter Gnade seinen Sohn für uns in Tod gibt. Und das ist recht GOttes Reich gepredigt, also, daß man für allen Dingen soll GOttes
\quellseite{52}
Zorn und Gnade uns verkündigen, und unsere Sünden uns zu erkennen geben.

Zum siebenzehenden. Hie ist Fürsichtigkeit vonnöthen, daß man nicht einen Splitter ziehe aus den Augen der Zuhörer, und sich selbst oder andere hindere, einen großen Balken zu gewinnen. Also fahen etliche an, und reizen das Volk wider die \randnotiz{Schwärmerei.}Pfaffen und Mönche, sagen, ihr Wesen sei böß und gottlos, ihre Lehre falsch, ihre Beiwohnung sei schädlich, daß gewöhnliche Fasten, Beichten, Meßhören, Sacrament empfahen, Beten, Kirchengang, Feiertag gelten nichts zu der Seligkeit, die Werke thuens nicht, der Glaube macht allein selig. Denn fallen die Zuhörer drauf, nehmens an, nicht den Glauben an Christum, sondern den Wahn und Gefallen über diese Rede, daß man alles Gewöhnliche ablege, und sehen zu, wie ein Spiel sich machen wölle. Denn ohne alle Gottesfurcht, ohne alles Gewissen und Bescheidenheit fahren sie zu, und kehren alle Dinge um, wo sie nur mögen, werden Frevler und verwegen zu und in allen Dingen, und freuen sich, daß sie überkommen haben einen Deckel ihres Muthwillens, daß sie mit Glimpf mögen zerbrechen den Zaun aller Zucht und Ehrbarkeit, darinnen sie vorhin ungerne beschlossen waren.

Darneben stehet dann der Schwärmer auf der Kanzel, und stärket das gemein Pöbel in ihrem Fürnehmen, als thun sie GOtt einen Dienst daran, sind vorhin Pfaffen und Mönchen nicht hold, und ist jedermann der Zucht und Ehrbarkeit feind, darum führet man für und für\footnote{D. i. darum führet man immer weiter fort in allem Muthwillen.}. Weise Leute sehen zu, und dieweil sie sich bedenken, so nimmt der Muthwill überhand, daß man denn nicht mehr wehren kann. Darüber lachet der Teufel, und hat aber ein gewonnen Spiel. Denn wie vorhin ihm gedienet ward in fleißiger
\quellseite{53}
\randnotiz{Ceremonien.}Haltung der Ceremonien von wegen des gottlosen Wesens im Herzen: also wird ihm jetzt gedienet im unsinnigen Abreißen der Ceremonien, so es ohne Gottesfurcht und ohne Gewissen geschieht, und eben der meiste Theil unter solchen Predigern und Zuhörern Gott weniger erkennen denn kein Papist, in dem, daß sie Hurer sind, Trunkenbolde, Gotteslästerer, Afterreder, Geizige und dergleichen Laster mehr an ihnen haben.

Wo man denn findet ehrbare, züchtige, gewissenhaftige Menschen, welche ob diesem Frevel erschrecken, nicht sich darein wollen geben, derselbigen spotten die Muthwilligen, nennen sie Heuchler, Gleisner, und wider sie stellet man alle Predigt, und wird ein solcher Jammer in der Welt, ehe denn man anfähet, das Reich Christi zu predigen, daß man gleich schier gezwungen wird, die Lehre zu verbieten, als sei sie schädlich, aufrührisch und dergleichen. Dann haben wirs wohl ausgericht. So man soll anfahen dem Herzen mit Gottes Wort sein eigen gottlos Wesen zu beweisen, in Gottesfurcht und Liebe fürhalten, so haben wir mit unsrer teuflischen Weise gemacht, daß man weder Christum, noch Teufel\footnote{D. i. weder vom Reich Christi, noch von dem Reich des Teufels, wie im vorhergehenden gezeigt ist.} predigen darf, und das ist die Ursach des großen Unwillens in der Welt über dem Evangelio. Man vermeinet, daß Evangelion sei das schwärmig Wesen, wie es solche Leute' fürgeben. Nein, nein, lieben Freunde, das heißt nicht evangelisch geprediget, sondern dem Teufel ein Ei braten.

\randnotiz{Evangelische Predigt.}Zum achtzehenden. Ein evangelischer Prediger weiset den Menschen von seinem mißtrauigen, eigensüchtigen Herzen zu gutem Vertrauen in Gott durch Christum, zu wohlthätigem Gemüth gegen den Nächsten. Das Evangelion schrecket ab
\quellseite{54}
daß stolze, frevle Herz von eignem Wohlgefallen, Vernunft, Stärke, Weisheit, Witz, Anschlägen und dergleichen, und treibt zu GOttes Gewalt, Furcht, Weisheit, dieselben herzlich und wahrlich zu bedenken, und durch Furcht in Liebe zu kommen. Wenn man nun lange, viel und wohl gegründet hat diesen Haupt-Punct und Artikel im Herzen und Sinne der Zuhörer, darnach soll man gemächlich und sittig anzeigen, in welchem Thun oder Lassen wir wider GOttes Furcht und Liebe handeln, ja die Zuhörer mögens aus ihnen selbst merken. Alsdenn findet man äußerlichen Mißbräuchen allezeit wohl Rath, sie werden selbst täglich abfallen, so viel mehr, wie viel GOttes Furcht und Liebe zunimmt im Herzen, denn mag man viel Dings lassen bleiben als gering und nicht schädlich, wo das Licht göttlicher Furcht und Liebe im Herzen leuchtet, bis daß unsere Herzen und unsere Nächsten stärker im Glauben werden. Etliche Dinge mag man wohl von außen bleiben lassen, so das Herz von inwendig neu worden ist durch das göttliche Licht.

\randnotiz{Wie man äußerliche Mißbräuche soll angreifen.}Und doch jene GOtt keinen Gefallen thun, ob sie gleich alle Bilder, Tempel, Klöster, Fastel-Tage, Beichten, Meßhalten und dergleichen äußerliches Dings mehr abthäten, denn sie an keine Gottesfurcht gedenken. Ihrer keiner ist mit Furcht und Urtheil GOttes je getroffen worden, hat sein Herz nie für dem Richtstuhl Christi, des Abbrechens halben, examiniert, ob es am letzten Ende bestehen möge für GOtt, oder nicht, laufen also hin wie unsinnige, rasende Menschen.

\randnotiz{Ceremonienstürmer.}Zum neunzehenden. Ich wollte so ungern rathen und helfen zu gemeinem irrseligen Gottesdienste im Volk ohne ernstliches Treiben des Worts GOttes, als ungerne ich wollte einem morden helfen; ich besorgte, der da hulfe, thäte das nicht aus Glauben, und darum wäre hie dem Teufel so wohl gedienet, als dort.

\quellseite{55}
Und wie jene sich verlassen auf das Aufrichten gemeldter Ceremonien, und vermeinen GOtt einen Dienst daran zu thun, ohne alle Gottseligkeit, und gefallen ihnen selbst wohl: also viel mehr verführen diese im Abbrechen, verlassen sich darauf, sie haben die Abgötterei helfen zerstören, jetzt seien sie Christen, und kommen nimmermehr zu rechter Gottseligkeit, davon oben gesagt ist. Ja welcher ihnen dawider redete, der müßte mehr Leides von ihnen erdulden, denn von den Papisten. Ich wollte lieber predigen in einer papistischen Stadt, da nie kein solcher Schwärmer gewesen wäre, denn in einer solchen Stadt, da das Volk so zerfallen, frevel und muthwillig worden ist. Aber wir werden unsers Frevels halben Rechnung geben müssen ohne allen Zweifel. Ich danke meinem GOtt, daß er mich geführet hat zu dem frommen Herrn Phil. Melanchthon, der solchen Frevel in mir gestraft hat, und mich treulich gelehret diese Bescheidenheit, wie ich jetzt geschrieben habe. Und welche Doctor Martin Luthers Bücher und Lehre fleißig lesen, oder hören, die werden auch nichts anders finden. Ob es aber jemand anders verstehet und seinen falschen Verstand verfechten will, dem ist zu fürchten, daß GOttes Fluch sei schon über ihn kommen, daß ihm GOttes Wort zu großer Verblendung diene, wie der Prophet Jesaias\footnote{Jes. 6.} sagt, und Christus im Evangelio\footnote{Luc. 8.} dasselbige wiederholet.

\randnotiz{Wie mit den Schwärmern zu handeln.}Zum zwanzigsten. Darum lieben Herren, die ihr solche Schwärmer zu Predigern habt in euern Landen und Städten, thut in Zeiten dazu, ehe denn euer Volk muthwillig und los werde. Das Evangelion Christi lehret Geduld, Gehorsam, Zucht, Ehrbarkeit, als man in der Epistel St. Paulus an die Römer 12. 13. 14. lieset, und anderswo mehr
\quellseite{56}
befindet. Alles, was nicht das herzliche Vertrauen in GOtt durch Christum hindert, das trägt und thut ein Christ. Lieben Herren, wollt ihr aus gemeinem Rath eine Ordnung eurer Gemeine ändern, oder gar abthun, wohl und gut, euer Glaube euer Richter. Wollet ihr das nicht thun, aber gut. Doch sollen sie eure Zuhörer nicht rotten, und dies oder jenes wollen verändern. Ein jeglicher reformiere sich selbst nach GOttes Wort und Gabe, und vermahne seinen Nächsten auch dazu. Will es helfen, wohl und gut; wo nicht, so befehle man die Sache GOtt im andächtigen Gebet, der wirds alles wohl machen. Und thut man also, so wird GOtt mit uns sein, und wird das antichristisch Reich zerstören. Darum so nehmet Prediger an, die euch Christlichen Glauben und Zucht lehren aus Grund der heiligen Schrift, der Schwärmer gehet mäßig.

\randnotiz{Ceremonien ändern.}Zum ein und zwanzigsten. Nicht sage ich das darum, daß man nicht soll einen ernstlichen Widerstand thun denjenigen, so mit falscher Lehre wollen die Wahrheit GOttes aus den Herzen reißen; sondern lies St. Pauli Epistel recht an Timotheum, da er also spricht unter andern: Ein Knecht des HErrn soll nicht zänkisch sein, sondern väterlich gegen jedermann, lehrhaftig, der die Bösen tragen kann, der mit Sanftmuth strafe die Widerspenstigen, ob ihnen GOtt dermaleinst Buße gebe, die Wahrheit zu erkennen, und wieder nüchtern zu werden von des Teufels Strick, die von ihm gefangen sind nach seinem Willen.

\randnotiz{Wie ein Diener des Herrn soll geschickt sein.}Merckt, ein Knecht des HErrn soll nicht zänkisch sein, nicht alles wollen verantworten, nicht jedermann anfallen, nicht mit dem Kopf hindurch, und mit dem Schwerdt alles wollen verfechten. Nein, nein, sondern er soll sein väterlich gegen jedermann, und lehrhaftig, auch wissen die Bösen zu tragen, und mit Sanftmuth strafen die Widerspenstigen. Man soll in Sanftmüthigkeit
\quellseite{57}
ernst sein. Vom Ernst redet St. Paul in den Episteln an Titon\footnote{Tit. 1.} und Timotheon\footnote{2 Tim. 4.}; von Sanftmuth aber in der andern Epistel an die Corinthier\footnote{2 Cor. 2.}, und in den jetzt erzählten Worten zu Timotheo\footnote{2 Tim. 2.}. Drum halt dich also: Bis im Ernst sanftmüthig, und in Sanftmuth ernstlich, und wie St. Paul sagt, handle mit aller Langmüthigkeit und Lehre, daß man nicht Scheltwort, sondern Lehre von dir empfange. Hippelbuben und alte Weiber können auch wohl schelten, aber nicht wohl lehren.

\randnotiz{Schelten der Propheten und Apostel.}Sprichst du denn: Wer kann Ernst und sanften Muth also vermischen und bei sich zusammen haben? Antwort: Kannst du das nicht, so laß dein Predigen anstehen, und fahre dafür zu Acker, grabe und arbeite. Christus siehet wohl eine Weile zu, aber zuletzt wird er die Mißhandlung seines Wortes am höchsten strafen. Ja, sprichst du denn, die Propheten, Christus selbst und die Aposteln, auch der Luther, sind scharf, und schelten übel wider ihre Widersacher. Antwort: Du willst auf Einen Tag alle Scheltworte ausschütten, die du in allen ihren Büchern findest, daran sie vielleicht viele Jahre geredet haben. Auch haben sie es gethan mit Fug und Bequemlichkeit der Zeit, und wenn sie die ernstliche Lindigkeit des Geistes am meisten befunden haben. Also auch, da Christus die Juden hart straft im Garten, dennoch war sein Herz so mild so lind und sanft, daß er dem Malchus\footnote{Luc. 22.} das Ohr wieder ansetzte, ja er starb für sie. Wenn wir auch also geschickt wären, würde uns das Schelten wohl anstehen, würde auch mehr Frucht bringen bei den Zuhörern; aber wir sind Schwärmer: schelten, verachten, nachreden und dergleichen, ist uns allen von Natur angeboren. Das üben wir jetzt unter der
\quellseite{58}
Gestalt des Evangelions wider die Papisten, und ist doch kein Geist, sondern eitel Fleisch in uns.

Lieber Herr Better, ich schreib darum euch und andern so viel davon, denn ich bin auch „zu viel geneigt auf gähe, scharfe Handlung, und habe viel Schiffe damit verführt, bin aus eigenem Schaden witzig worden. Wenn man das Spiel übersehen hat, so untersteht man sich den Misrath zu beschirmen mit Glimpf und Unglimpf, niemand will Unrecht haben: aber Gott läßt seiner nicht spotten.“ Betrüge sich selbst einer nicht, ist mein getreuer Rath.

\subsection{Zum zwei und zwanzigsten}

Wir sollen von den Propheten und Aposteln lernen, wie fein und gemächlich man fahren soll mit Predigen. Ist nicht der Artikel von der Gottheit Christi fast nöthig zur Seligkeit? Dennoch saget St. Peter anfänglich nichts davon zu Jerusalem, wie in den Geschichten der Aposteln\footnote{Ap. Gesch. 2. 10. 13. 14. 17.} stehet. Auch saget er dem Cornelio erstlich nichts davon. Item: St. Paul schweiget stille zu Antiochia in Pisidien von der Gottheit Christi; allein obenhin rühret ers mit einem kleinen kurzen Sprüchlein. Desgleichen thät er auch zu Lystris und zu Athen, wie in den Geschichten der Aposteln allenthalben geschrieben ist. Christus verbot auch seinen Jüngern, sie sollten vor seinem Tode niemand sagen, daß er Christus wäre\footnote{Matth. 16.}. Den Artikel von der Jungfrauschaft Mariä hat das ganze Neue Testament nie öffentlich fürgehalten, denn allein Matthäus und Lucas. Vom Sacrament des Altars schweiget Johannes und Petrus in ihren Schriften, und Paulus redet allein in Einer Epistel davon, und haben doch auch alle das Evangelion beschrieben. Also haben auch die Propheten viel Artikel gar dunkel fürgetragen.

\quellseite{59}
\randnotiz{Rechte Weise und Ordnung zu lehren.}Warum das alles? Darum, man muß gemacht fahren. Anfänglich soll man hören Gottes Allmächtigkeit, wie er allein alles gut schaffet und thut, und wie auch alle unsere Habe, Thun und Lassen aus Gottes Gewalt und Ordnung herkomme und für sich gehe. Bis daß man diesen Artikel wohl in die Zuhörer treibe, darf man viel Zeit und Fleiß dazu, denn dadurch wird das Herz getrieben von eigener Vermessenheit, und lernet das Auge auf Gott kehren, siehet auch, wie viel Gutes er uns täglich thut.

\randnotiz{Christum predigen.}Danach soll man die Gebote Gottes erzählen und rechten Grund derselbigen fürlegen. Das bedarf auch viel Zeit und Arbeit. Wenn denn nun das Herz seine Sünde dadurch lernet erkennen, und sich vor Gott ansähet zu fürchten, alsdann soll man anheben Christum zu lehren, wie er uns zu gut von Gott gesandt sei, für unsere Sünden zu sterben, und wie er aller Welt sei zu einem Richter gesetzt, werde aber aller verschonen, die an Ihn gläuben. Desgleichen auch soll man sagen, daß wie er sei von Tod und Hölle gewaltiglich erstanden, also werden unsere Körper auch wieder aufstehen, und daß unsere Seele frei sei vom Tod, Hölle und Teufel. An diesem Stück hat man lange zu predigen, und soll das oft treiben, bis daß der Zuhörer eine Liebe und Vertrauen zu Christo gewinne. Alsdenn so ists gut von andern Artikeln mehr zu reden.

\randnotiz{Wie man fahren soll mit denen, die Christum erst lernen erkennen.}Alle Bücher des Neuen Testaments sind zu denen geschrieben worden, die jetzt bereits an Christum gläubeten. Derhalben es nicht wohl daraus zu lernen ist die Form und Weise anfänglich zu predigen denen, die Christum noch nicht recht erkennen. Aber in den Geschichten der Aposteln allein findet man solche Form zu predigen den Glauben. Auch aus der Form der Apostel-Predigt für dem Tode Christi mag man lernen anfangen zu
\quellseite{60}
predigen. Denn wir müssen ja so glimpflich mit denen fahren, die Christum noch nicht recht erkannt haben, ob sie gleich unter dem Pabstthum gewesen sind, schier als mit den Ungläubigen, wie die Apostel erstlich haben thun müssen mit den Heiden. Und wenn mans nach aller Weise und Form vorgetragen hat, wie uns obgemeldte Exempel der Apostel lehren, und man jetzt Christum anfähet zu erkennen, noch soll man nicht eilen mit den Ceremonien, sie zu geringern oder abzubrechen, also wenig es Petrus und Jacobus zu Jerusalem thäten. Davon lies die Geschichte der Apostel\footnote{Ap. Gesch. 15.}, auch St. Paulus Episteln an die Römer\footnote{Röm. 14.} und Corinther\footnote{1 Cor. 8 und 10.}. Das soll man aber thun, täglich predigen, wie unser Heil und Seligkeit nicht in diesen Ceremonien und äußerlichen Satzungen, sondern am Glauben an Christum liege: So werden diese Ceremonien je länger je mehr abfallen, doch in einem mehr, denn in dem andern, nach eines jeglichen Glauben, ungezwungen, ungedrungen, mit Fried und Liebe aus dem Glauben, aus gutem Gewissen mit Gottesfurcht. Denn wird des Teufels Reich zerstöret mit Gottes Wort, nicht mit unserm Frevel.

Zum drei und zwanzigsten. \randnotiz{Unterscheid zwischen uns und den Heiden.}Beileibe wolle keiner sagen: Ja ich soll so fast die Ceremonien verjagen, als Paulus gethan hat. Lieber Gesell, du verstehest Paulum nicht recht. Die Heiden hatten vorhin der Juden Gesetz nicht, und lernten jetzt Christum mit Willen. Da verhütete Paulus, daß man den Heiden nicht sollte auflegen das fremde Gesetz der Juden, denn der Glaube an Christum wäre ihnen genug zur Seligkeit, wie St. Paul schreibet und ausleget in den Episteln an die Römer\footnote{Röm. 3.} und

\quellseite{61}
Galater\footnote{Gal. 2.}. So hatten die Heiden auch Landrecht genug zu weltlichem Regiment, darum bedurften sie nicht des Gesetzes Mose. Auch hielt Paulus viel heidnischer Gebräuche und Gewohnheiten, wenn er bei den Heiden war, daß er möchte die Heiden zum Glauben Christi bringen, wie da stehet in der ersten Epistel an die Corinther, da St. Paul also saget: Wiewohl ich frei bin von jedermann, hab ich doch mich selbst jedermann zum Knecht gemacht, auf daß ich ihrer viel gewinne: den Juden bin ich worden als ein Jude, auf daß ich die Juden gewinne. Denen, die unter dem Gesetz sind, bin ich worden als unter dem Gesetz, auf daß ich die, so unter dem Gesetz sind, gewinne. Denen, die ohne Gesetz sind, bin ich als ohne Gesetz worden, (so ich doch nicht ohne Gottes Gesetz bin, sondern bin in dem Gesetz Christi,) auf daß ich die, so ohne Gesetz sind, gewinne. Den Schwachen bin worden als ein Schwacher, auf daß ich die Schwachen gewinne. Ich bin jedermann allerlei worden, auf daß ich allerdings ja etliche selig mache. Er erlaubet auch andern Christen zur Wirthschaft der Heiden zu gehen, und nicht achten, ob sie, doch unwissend, auch vom Götzenopfer essen, wie er in der jetzt gemeldten Epistel\footnote{1 Cor. 10.} anzeiget, da er also spricht: So aber jemand von den Ungläubigen euch ladet, und ihr wollt hingehen, so esset alles, was euch vorgetragen wird, und forschet nichts, auf daß ihr der Gewissen verschonet.

\subsubsection[Aber mit den Juden]{Aber mit den Juden handelt Paulus anders}

Er hielt mit ihnen äußerlich das Gesetz, wie in den Geschichten der Aposteln\footnote{Ap. Gesch. 16. 21.}, und wie jetzt aus der ersten Epistel an die Corinther\footnote{1 Cor. 9.} gehöret ist, auf daß er andere zum Glauben zöge. Und da sie schon an Christum glaubten, und nicht mochten das Gesetz so bald gar lassen, gehet er, man sollte ihrer nicht derhalben spotten, noch sie verwerffen oder urtheilen, sondern tragen und dulden, bis daß sie täglich durch
\quellseite{62}
Erkenntniß Christlicher Freiheit vom Gefängniß des Gesetzes erlöset würden, wie in der Epistel zum Römern\footnote{Röm. 14.} stehet. Und die Aposteln haben viele Jahr auch zu Jerusalem ein Mitleiden mit den Juden gehabt, und Paulus mit den Römischen Juden.

\randnotiz{Pabsts Gesetz.}Wir aber sind bisher behaftet gewesen, von unsern Eltern her, mit des Pabstes Gesetzen, und mögens so schwerlich laßen, als die Juden ihr Gesetz. Ob wir denn nun nicht so bald mögen gar davon abstehen, ists kein Wunder. Darum soll man ein Mitleiden mit uns haben, bis daß wir wachsen im Glauben so fast, daß wir die Gesetze gar mögen vom Herzen bringen aus dem Glauben an Christum, und nicht aus Frevel. Denn ohne Aufrichtung der Gottseligkeit des Pabstes Gesetz halten, oder brechen, gilt gleich. So wir aber uns die Gottseligkeit laßen gefallen, oder noch nicht so fast darinnen gewachsen sind, daß wir die Gesetze möchten verachten aus dem Glauben: so ists beßer, wir wandeln noch im Gesetz, doch mit stetem unaufhörlichen Beten zu GOtt, er wolle unsere Schwachheit stärken. Und daneben soll ein Prediger „bescheidenlich, freundlich GOttes Wort treiben, damit auch anzeigen, wie unser Heil am Glauben an Christum liege, und an keinem Gesetze,“ und daß keiner den andern soll urtheilen, verachten oder ärgern, er halte das Gesetz, oder laße das Gesetz; sondern einer soll den andern tragen, bis daß GOtt hilft, damit der freventliche muthwillige Haufe, der allein aus Schwärmerei die Gesetze läßt, und darneben weder der Ehre, noch der Zucht achtet, „nicht verspottet, noch betrübe ander ehrbare, gewissenhaftige, doch schwachgläubige, aber gutherzige Menschen, welche warten wollen, bis sie lernen, nicht aus Frevel, sondern aus Glauben, die Gesetze und menschliche Satzungen verachten.“

\quellseite{63}
\randnotiz{Pabstes Gesetz beschweret unser Gewissen.}Hätten wir des Papstes Gesetz nicht so hart auf unserm Gewissen liegen, wir wolltens nicht lassen einbrechen, als wenig Paulus Mose Gesetz auf die Heiden wollte fallen lassen. Da wir aber, leider, im Pabstthum geboren sind und erzogen, wollen wir gerne mit Fleiß Christum lernen, bis wir aus Glauben die Gesetze mögen verlassen. Drum, lieben Herren Prediger, thut uns gemach und hübschlich, wartet unserer Blödigkeit, wir sind nicht so stark, als ihr, wir mögen noch nicht unsere Ceremonien so frei und muthig lassen, als ihr. Lobet ihr euren GOtt, daß er euch erlöset hat, und lehret uns freundlich, auch wartet unser geduldiglich, bis unsere Gewissen auch von GOtt gestärket werden. Es sei denn, daß unsere Gewissen stärker werden, so mögen wir nicht davon lassen, denn allein aus Frevel, welches uns zu großem Schaden und Beschwernis des Herzens reichen möchte, und bittet neben uns GOtt um Stärke unsers Glaubens.

\randnotiz{Prediger hören hie zu!}„Werfet uns den Luther nicht für, wir halten noch nicht einen jeglichen unter euch für den Luther, euer eines Theils können Luthers Schriften noch nicht recht lesen, viel weniger verstehen ihn etliche recht.“ Der geringste Theil ist nicht so weise, so klug, so gelehrt, so erfahren, so fromm, so bewährt, als Luther ist. Wir wissen auch, was Rüben, oder Birn sein. Dazu hat Luther kräftiglich mehr denn drei oder vier Jahr zu Wittenberg öffentlich täglich disputirt, gepredigt, gerathen, geboten, mit Hülfe und Beistand des bescheidenen und gelehrten Mannes Philipps Melanchthon, und dergleichen anderer mehr, dennoch wollte Luther „nichts lassen verändern in öffentlichen Gebräuchen, er besorgte, es geschehe mehr aus Thorheit, denn aus Glauben,“ darüber der Teufel pfleget zu lachen. Erst innerhalb einem Jahr hat er so gemächlich eins nach dem andern angefangen abzustellen nach seinem Vermögen; „und etliche aus euch
\quellseite{64}
\randnotiz{Schwärmer.}ungelehrte, unerfahrne Neulinge überfallen uns in Städten und Dörfern mit neuen Lehren, wollen in Einem Monden alle Ceremonien bei uns abtreiben; uns mehr leichtfertig, denn gottseelig zu machen, hänget das Pöbel-Volk an euch, das weder GOtt erkennet, noch Vernunft hat;“ und so wir das aus Blödigkeit nicht thun mögen, sind wir euch zu Spott und zum Ziel im Tempel und auf dem Markt, alle Bößlein an uns zu schießen.

Ich muß noch eins sagen. Lieben Freunde, man möchte auch zu unser etlichen also sagen: Wir müssen eurer so wohl in Geduld warten, als ihr unser. Wir haben noch \randnotiz{Geduld zu haben mit dem Schwachen.}nicht gelernt Fleisch fressen am Freitag, die Beicht und Meß verlassen, unser Gebetlein zu den Heiligen hinlegen, und dergleichen; aber wir wollen möglichen Fleiß ankehren solches zu lernen. Darum seid nicht ungeduldig mit uns, denn euer eines Theils auch noch nicht gelernt haben abzustellen Hurerei, Füllerei, Gotteslästerung, ohne Noth schwören, schelten, nachreden, und dergleichen viel mehr. Ihr habt auch noch nicht gelernt, was zu einem friedlichen, züchtigen, bürgerlichen Leben dienet, ihr übel's ja nicht, so habt ihr noch keine Erfahrung auch in menschlichen und weltlichen Sachen, wisset uns wenig zu trösten und zu lehren von dem Reich Christi. Dazu seid ihr auch noch schwach, unsere Heimlichkeit, die wir euch aus Herzen klagen, zu verschweigen. Eure Liebe und Mitleiden mit uns erzeiget sich wenig, ihr wollet nicht unterwiesen sein, noch vermahnet werden, als ob ihr alle Dinge wisset und schnureben treffen möchtet. An welchen allen Stücken, die doch zum Theil fast ärgerlich, tragen wir mit euch Geduld, darum es auch wiederum billig wäre, mit uns in unserer Blödigkeit Geduld zu haben. Wahrlich, man findet ehrbare Leute,
\quellseite{65}
die in der Wahrheit mit mir und mit meines gleichen unerfahrenen Predigern also mögen reden, und noch wohl härter.

Dies alles, mein lieber Herr, hab ich so viel länger und lieber zu euch geschrieben, wollet es im Besten von mir annehmen, denn ich wohl weiß, wie übel euch solches Schwärmen gefällt. Doch muß ich noch eins hinzu setzen. Ob jemand wollte freventlich verharren im Unglauben, und andere an guter Lehre hindern, denen sollen wir mit Bescheidenheit widerstehen, so viel uns GOtt Gnade gibt. Hilfts nicht an ihnen, so sollen wir uns von ihnen abziehen, und der Gutwilligkeit mit Fleiß pflegen, wie Paulus thät zu Ephesus, als in den Geschichten der Apostel\footnote{Ap. Gesch. 19.} stehet, und wie Christus mit den Pharisiäern handelt im Matthäo am funfzehenden Capitel.

\subsection{Zum vier und zwanzigsten}

Wenn ihr merken möget, daß ein Mensch gefasset hat GOttes Wort, sollt ihr euch nicht fast bekümmern um die Veränderung seines äußerlichen Wesens und Wandels halben, als etliche thun, \randnotiz{Aeußerliche Dinge.}die mit allem Fleiß rathen und treiben, die Pfaffen von ihren Pfründen zu lassen, die Mönche und Nonnen aus den Klöstern zu gehen, und dergleichen äußerliches Dings mehr zu thun, als ob kein Heil wäre, wo man nicht auch dies alles nicht allein innerlich, sondern auch äußerlich von sich abwürfe; so doch nicht viel an äußerlichen Dingen gelegen ist, wenn nur GOttes Wort ins Herz geführet wird. Denn alsbald, wo das ist, da ruhet es nicht, mit der Zeit wird des Dings viel von \randnotiz{Pabstthum.}sich selbst abfallen. Auch hat man das größte erstritten in der Welt an dem Pabstthum, man höret täglich mehr und mehr, worin das Reich Christi stehet, und man verachtet nun, GOtt lob, des Antichrists Reich. Man ärgert sich auch nicht viel an alten Kranken, untauglichen, unbehelflichen Klosterleuten und Pfaffen, welche sonst nicht mögen
\quellseite{66}
aufkommen, ob sie im Kloster-Stand oder bei den Pfründen bleiben mit Erkenntniß eigener Blödigkeit, und mit Begierde der Stärke GOttes und seiner Erlösung.

\randnotiz{Schutzrede Eberlins.}Daß ich aber abgetreten bin vor zweien Jahren vom Kloster-Stand, hat andere Ursach gehabt. Denn obgleich Luther nicht wäre aufgestanden, dennoch war meine Sache also gestaltet, daß ich in allen billigen Rechten durch Vernunft allein wäre von dem Orden ledig gesprochen worden. Es hätte auch der Pabst wohl mit mir dispensiret, wenn ich Geld gehabt. Dazu wurde ich auch im Herzen ledig durch evangelische Lehre, mir durch Lutherische Schrift angezeigt. Da ich in dem Gewissen frei war, und Ursach hatte vor der Welt genugtsam, den Orden zu verlassen, und jetzt nun der Gebrauch war, daß viele Ordens-Leute ohne Aergerniß austraten, da zog ich auch davon ohne päbstliche Dispensation und Erlaubniß. Denn Rom war mir zu fern, und hatte auch kein Geld. So zweifelt mir nicht am Pabst, hätte ich Geld geben, er hätte mit mir dispensiret und mir vom Orden erlaubet. Dieweil es mir vor den Menschen recht war, mit Geld eine Dispensation zu erwerben, war mirs auch recht, ohne Geld abzutreten vor der Welt, denn im Herzen bedurft ich keines Erlaubnisses vom Pabst, GOtt hatte mich schon erlöst lange vorhin.

\randnotiz{Aus dem Kloster gehen.}Wenn aber einer wohl mag im Kloster bleiben, und erkennet Christum, aber seine Sachen wollen sich noch nicht schicken zum Ausgang, wollt ich keinem drum GOttes Huld absprechen. St. Paul spricht: Weder Beschneidung, noch Vorhaut, gilt für GOtt, sondern eine neue Creatur.\footnote{Gal. 6.} Denn, wollen alle Ceremonien verjagen, ist auch eine Ceremonie, und wollen ohn all Gesetz leben, ist auch ein Gesetz, und eben so sie fliehen die Gesetze, machen sie Gesetze.

\quellseite{67}
Auch ist es zu besorgen, daß ein Mensch wohl sollte in größere Fährlichkeit kommen, wenn er seine Pfründe verließe, denn er vorhin war, so derselbige Mensch unerfahren, übelkundig und schwach wäre, und hätte nicht gewohnt eines andern Lebens, und dürfte sich also selbst in Bettel, in ungewohnte Arbeit, Mühe und Sorge, da sollte man wohl sehen, wie sich der Teufel regen würde. Nicht ohne Ursach bittet der König Salomon\footnote{Spr. Sal. 30.} GOtt den HErrn, er soll ihm weder Reichthum, noch Armuth geben, daß er nicht dort zu stolz werde, und hie zu verwegen und verzweifelt. Denn so spricht er in den Sprüchen: „Zweierlei bitte ich von dir, die wollest du mir nicht weigern, ehe denn ich sterbe. Eitelkeit und Lügen laß ferne von mir sein, Armuth und Reichthum gib mir nicht, laß mich aber mein bescheiden Theil Speise dahin nehmen. Ich möchte sonst, wo ich zu satt würde, verleugnen und sagen: Wer ist der HERR? Oder, wo ich zu arm würde, möcht ich stehlen und mich an dem Namen meines GOttes vergreifen.“

So spricht nun einer: Ja man muß GOtt vertrauen. Antwort: So springe mit gleichen Füßen in die Hölle, ists GOttes Wille, so wirft dich die Hölle wieder heraus. Soll man sich für den Teufel nicht hüten in deinem Sinne, so will ich nicht mit dir disputiren. Wenn dich GOtt will vom Kloster oder Pfaffen-Stand ziehen, wird er dir wohl solche Mittel anrichten, daß du davon kommest und weißt selbst nicht, wie. Bist du ein Christ, so erkennest du, daß deine Nahrung von deinem himmlischen Vater kommt. Magst du aber sie nicht haben durch bequeme Mittel, so bekenne deinen Schaden vor den Menschen, und klags deinem GOtt\footnote{Psalm 27.}, bitte ihn um Hülf und Rath, getraue ihm, er wirds wohl machen, und harre auf ihn, so schadet dir weder Kloster, noch Pfaffen-Freunde.

Ich will euch schreiben einen Rath eines gelehrten, christlichen, erleuchteten Mannes, welcher mich deucht ganz nutz sein den
\quellseite{68}
Predigern, nemlich, daß man anfänglich nicht zu viel auf einmal soll umstoßen, als Meß, Vigilien, Jahrtage, Beichte und dergleichen, sondern glimpflich anfangen, also sagen, man solle um keines Gewinns willen, noch um Ehre, noch aus Gewohnheit Meß lesen, sondern allein um Gottes willen, für sich und andere zu bitten. Wenn nun einer etwan lange freundlich und ernstlich davon gelehret hat, so fahre er fort und sage, man soll auch keine andere Ceremonien, es seien Tagzeiten oder Vigilien, zum Pomp und Gepräng gebrauchen, dazu keinen Nutzen, noch Ehre damit suchen, sondern allein das Gebet an ihm selbst ansehen. Nach etlichen Tagen dann lehre man die Leute, wozu sie sollen der Beichte und des Sacraments gebrauchen.

\section{Ceremonien}

\subsection{Beicht}

\randnotiz{Absolution.}In der Beicht komme der Sünder von keinem Menschen getrieben, sondern von eigener Noth, seine Sünde nach Vermögen, ohne sonderliche Angst auf die Zahl, mit ganzer Treue dem Priester vor Gott zu erzählen, Trost der Absolution vom Priester zu empfahen, und Rath zu hören, wie sich der arme betrübte Mensch forthin soll halten. Denn Christus hat gesagt im Matthäo:\footnote{Matth. 18.} „Wahrlich, ich sage euch, was ihr auf Erden binden werdet, soll auch im Himmel gebunden sein, und was ihr auf Erden lösen werdet, soll auch im Himmel los sein.“ Weiter, sage ich euch, spricht Christus: „Wo zween unter euch eins werden auf Erden, worum es ist, daß sie bitten wollen, das soll ihnen widerfahren von meinem Vater im Himmel. Denn wo zween oder drei versammelt sind in meinem Namen, da bin ich mitten unter ihnen.“ Auf diese Zusage Christi komme der sündige, trostlose Mensch zum Priester, als zu einem Christen, Gottes Wort, Trost und Rath durch ihn zu hören. Wenn ihr diese Meinung oft und wohl gelehret habt, und Gottes Geist wirkt auch im Herzen der Zuhörer: werden
\quellseite{69}
sie bald verstehen mögen, daß mehr Fleiß soll gehabt werden, wie man auf Gottes Wort, durch den Priester geredt, vertrauen soll, und darauf sich trösten, denn auf eigene Beichte.

\randnotiz{Des Priesters Gewalt.}Darum so saget eine Weile davon, darnach so erzählet, wie ein Priester diese Gewalt habe, nicht darum, daß er vom Bischof geweihet ist; sondern darum, daß er vom „Heiligen Geist“, in der Taufe, mit der Gnade Gottes gesalbet ist, wie auch alle Christen. Und so mag man merken, daß alle Christen ingemein Gewalt haben die Sünde zu vergeben(?), lehren und zu trösten, alle gleich. Wiewohl aber solche Gewalt von Gott allen Christen ingemein geben ist(?), dennoch soll sich niemand öffentlich annehmen, diese Gewalt zu üben, er sei denn von der Gemeinde dazu erwählet; und die \randnotiz{Priester.}also erwählet sind, hat man Priester genennet.“ Eine solche Erwählung wird bedeutet durch des Bischofs Weihe, wenn sie auch am besten ist, und nicht mehr. Also kommt man täglich \randnotiz{Bischofsweihe.}für und für dazu, daß man wohl mag erkennen, daß Bischofsweihe wohl mag erlassen, oder gehalten werden, wie das einer Gemeine gefällt, und daß man kein Gelübde soll nehmen von den erwähleten Priestern, denn daß sie mit Gottes Hülfe wollen und sollen getreulich Gottes Wort predigen, bekennen, und jedermann nach Vermögen Gute thun und dienen. Ein solch Gelübde oder Bekenntniß that Timotheus, als ich den Text verstehe, in der ersten Epistel, da Sanct Paul\footnote{1 Tim. 6.} unter andern Worten also sagt: „Du Mensch Gottes, fleuch solches, jage aber nach der Gerechtigkeit, dem Glauben, der Liebe, der Geduld, der Sanftmuth. Kämpfe einen guten Kampf des Glaubens, ergreif das ewige Leben, dazu du auch berufen bist, und bekannt haft ein gut Bekenntniß vor vielen Zeugen. Siehe, an diesen Stücken hat man lang zu predigen mit Fug und Bescheidenheit, daß man ohne schädliches Aergerniß verstehen mag, was die Weihe
\quellseite{70}
sei, darnach auch, wann und wie sie nöth, oder nicht nöth sei. Sollte man aber anfänglich die Beichte so ganz verwerfen, ehe denn gnug und wohl davon gelehret und geprediget wäre; solltest du wohl alle gute Herzen von deiner Lehre abschrecken. Auch fähret kein ehrbar, vernünftig Mensch so unfüglich im Anfang mit der Lehre, es thuns allein unschamhaftige Leute, welche gern wollten gesehen werden, daß sie etwas seltsames fürbrächten.“

\randnotiz{Sacrament des Leichnams und Bluts Christi.}Also thue man auch mit dem Sacrament des Leichnams und Bluts Christi. Erstlichen sage man, daß dies Sacrament sey ein Siegel Gottes Worts, also, wenn ein trostlos Mensch höret Gottes Trost, und ob er noch schwachgläubig ist, so empfähe er das Sacrament des Altars zu mehrer Befestigung seines Herzens, in Vertrauen auf Gott wider alles Leiden und Anfechtung. Auch zeige sich ein Mensch selbst allen andern zu Trost, daß er auch Christum bekenne mit dieser Empfahung, also, daß andere auch gezogen und gereizet werden zum Glauben an Christum, und sich wolle lassen finden für jedermann zu möglichem Dienst und Wohltat, zum Exempel der Zucht und Ehrbarkeit, aller Welt.

\randnotiz{Beider Gestalt.}Fället nicht gleich drauf, man solle unter beider Gestalt das Sacrament empfahen, und der Papst sei ein Tyrann, daß er uns eine Gestalt genommen hat, und was dergleichen mehr ist; sondern fahret immer für und für in der Lehre vom Sacrament, lehret den Glauben auf die Worte Christi, damit dies Sacrament eingesetzt ist worden. Also, daß man allweg mehr Fleiß ankehre, das Volk zu weisen auf den Glauben an Christum, denn auf das äußerliche Zeichen. Wenns nun Zeit ist, so sage man ein wenig von beider Gestalt des Sacraments, aber man soll keine ohne Glauben empfahen. Saget auch, daß der Mensch sei ungezwungen, weder von Gott, noch von Menschen,
\quellseite{71}
zu diesem Sacrament, sondern es soll geschehen aus einem freien Herzen und Andacht, zu diesem Ende und Nutz, der droben erzählet ist. Vermahnet ja das Volk, daß sie nicht leichtlich, noch verächtlich unterlassen, die Beicht und das Sacrament zu empfahen. Denn der Teufel hat eben so großen Fleiß, die Leute aus Leichtfertigkeit, oder Hinlässigkeit, von der Beicht und vom Sacrament zu ziehen, als fast er vorhin sich geflissen hat, ohne Glauben mit Gewalt und großen Haufen hinzu zu treiben. Durch diese Weise kommt man dann dazu, daß der Mißbrauch der Sacramente abgestellt werde.

\randnotiz{Sacrament.}Und sonderlich wenn man das fast treibet, was ein Sacrament sei, nemlich daß es dies sei, wenn GOtt ein äußerlich Zeichen einsetzt zu seiner Verheißung, bei welchem Zeichen der Mensch vergewissert wird auf GOttes Huld und Gnade.

\randnotiz{Zwei Sacramente.}So kann man darnach fein anzeigen und beweisen, daß die Ehe, Weihe, Oelung, Firmung, dieweil sie nicht für Zeichen sein von Christo aufgesetzt, auch nicht Sacrament Christi sollen genennet werden; sondern allein die Tauffe und das Sacrament des Altars sind Zeichen und Sacramente Christi. Also führet ihr das Volk von dem Irrthum, den sie haben von den Sacramenten, und auch zu Empfahung unter beider Gestalt, ohne alle Aufruhr und Aergerniß. Ach! wenn „man Prediger hätte, die mehr der Seelen Heil, denn eigen Nutz und Ehre suchten, würde man GOttes Wort ohne solch Poltern und Unruhe wohl predigen in unserm Landen.“ Denn Doctor Martin Luther und etliche getreue seine Helfer haben das größte und schwereste abgehauen, das am Wege lag, und haben das heilig Grab göttlicher Schrift geöffnet, daß fast an allen Orten der Greuel des Papstthums erkannt ist, und daß jedermann einen Gefallen an der Bibel hat, darum uns nicht nöth wäre,
\quellseite{72}
auf ein neues aufzublasen und Narren jagen ohne Nutz und ohne Noth\footnote{1 Cor. 1.}. Wir möchten wohl auf das Fundament, vom Luther gelegt, bauen Gold, Silber, Edelgesteine, gute tröstliche Lehre der Gewissen, ehrbare, züchtige Sitten, und dergleichen. Aber ich habs an mir selbst erfahren, daß der Teufel treibet zu „solchem Poltern, auf daß er guten Glimpf habe, unsere Lehre gar abzutreiben, oder an bessern Dingen uns zu hindern.“

Wenn nun solches alles von euch geschehen, darnach möget ihr anzeigen aus den Worten Christi, vorhin so oft gepredigt, daß die \randnotiz{Messe.}Messe kein Opfer möge für uns sein, so thue der Priester auch nicht mehr, wenn er das Sacrament empfähet, (auch im Glauben) denn ein anderer gläubiger Laie thut. Darum soll man nicht halten, daß der Priester opfere das \randnotiz{Sacrament ist kein Opfer.}Sacrament für die Lebendigen und für die Todten, und alle diese Meinung der Messe halben soll, als unchristlich, abgestellt werden. Christus ist unser Pfaff, der opfert für uns, und kein Mensch. Der Priester sei allein ein Diener der Gemeine, und nehme das Sacrament für sich, und gebe andern auch davon. Diese Lehre wird das Volk auch wohl mögen fassen, wenn ihr vorhin gelehret habt dasjenige, davon droben gesagt ist, denn so werden die Votiven und Opfer-Messen alle von sich selbst abfallen.

Ihr sollt oft und viel sagen, daß „ein Christ unsicher und ungewiß sei, wenn er etwas annimmt oder glaubt, dazu ihm keine Schrift vorgetragen wird.“ Denn wenn dieser Beschluß wohl gegründet ist im Volk, so mögt ihr viel Mißbräuche damit in ihrem Herzen abthun, nemlich, dieweil die Schrift kein \randnotiz{Fegefeuer.}Fegefeuer anzeiget, noch Weise und Mittel den Seelen zu helfen, und man ist auch nicht sicher, noch
\quellseite{73}
\randnotiz{Vigilien, Jahrtage, Seelmessen.}gewiß, daß man predige und glaube, daß ein Fegfeuer sey; so fallen alle Jahrtage, Vigilien und Seelmessen dahin ohne Mühe und ohne Arbeit.

\randnotiz{Heiligen Dienst.}Und dieweil man keine Lehre, noch Exempel hat in der Schrift, die abgestorbenen Heiligen anzurufen um Hülfe und Fürbitte, so stehet man auch billig ab von so viel Wallfahrten, Gebetlein, Gelübden und andern Narrenwercks mehr, die man den Heiligen gethan hat, sonderlich dieweil solche Zuversicht zu den Heiligen oft schädlich ist dem Glauben an Christum, bei welchem wir alle Dinge gewarten sollen.

\randnotiz{Kloster-Gelübde.}Also auch, dieweil GOtt an keinem Ort gebeut, Kloster-Gelübde zu halten oder zu geloben, und auch nicht sicher ist, solch Gelübd zu thun, oder zu halten, denn sie sind wider GOtt, in dem, daß man die Gewissen damit bindet und regiert, so doch die Gewissen sollen frei sein von allen Gesetzen, und allein durch GOttes Wort regiert werden: So folget, wenn ein Pfaff oder Kloster-Mensch nicht mag Keuschheit halten, mag es und soll ehelich werden, viel Fährlichkeit zu vermeiden, wie St. Paul erzählet in der ersten Epistel an die Corinthier\footnote{1 Cor. 7.} im siebenden Capitel.

\randnotiz{Pabstthum.}Nach diesem allen möget ihr sein zeugen, sintemal Christus allein das Haupt ist der Kirchen, wie Paulus sagt zu den Ephesern\footnote{Eph. 1.} und Colossern\footnote{Col. 1.}, so soll weder St. Peter, noch der Pabst, für das Haupt Christlicher Gemeine gehalten werden. Und dieweil Christus selbst bei uns ist bis ans Ende der Welt, wie er uns verheißet und zusaget nach seiner Auferstehung im Matthäo\footnote{Matth. 28.}, so bedarf er keines Statthalters. Darum
\quellseite{74}
\randnotiz{Christus ein Haupt der Kirchen.}so soll sich weder St. Peter, noch der Pabst, Christi Statthalter nennen. Auch so sich Petrus und Paulus nicht anders, denn als Diener der Gemeine nennen, und dafür wollen gehalten werden, wie sie selbst in ihren Episteln und Sendbriefen schreiben, und sonderlich St. Paul zu den Corinthiern\footnote{1 Cor. 4. 2 Cor. 4. Col. 1. 1 Petr. 5.}, da er also spricht: Dafür halte uns jedermann, nemlich für Christus Diener und Haushalter über GOttes Geheimnisse: So sollen nun unsere Päbste und Bischöfe, wenn sie am frömmsten sind, auch nicht anders sein und erscheinen, denn als Diener der Gemeine, und nicht als Herren, sondern als Helfer unseres Glaubens\footnote{2 Cor. 9. 1 Petr. 5.}.

\randnotiz{Concilia.}Auch sintemal die Concilien nichts anders sind denn Versammlung der Christen, alle Christen aber sind durch das Wort GOttes geboren und Christen worden, folget, daß GOttes Wort über alle Christen ist, sie sind in oder außerhalb der Concilien. Darum so soll kein Concilium etwas ordnen, zu binden oder zu trösten die Gewissen, wenns nicht in der heiligen Schrift ist gegründet; wird es aber ohne das geordnet, soll es nichts gelten. Jetzt weiß man sich schon aus allen Concilien, Doctoren und Schulen zu richten, und aus ihnen sich zu zerren und abzureißen mit GOttes Wort. Sehet ihr nun, wie friedlich, beßerlich, tröstlich möchte man GOttes Wort lehren, wenn man die Sache mit Fuge anfinge. Nicht ohne Ursach schreibt Paulus zu den Colossern\footnote{Col. 4.}, man solle GOtt für ihn bitten, daß er möge GOttes Wort predigen und also reden, wie er soll reden. Die Erfahrung lehret, wie viel Weisheit dazu gehört, daß man rede GOttes Wort, wie man es reden soll. Auch lehret Paulus Timotheon\footnote{2 Tim. 2.}, daß er GOttes Wort recht scheide.

\quellseite{75}
Zum fünf und zwanzigsten. Für allen Dingen sollen wir Christen aufsehen, daß wir uns nicht selbst Hindernis und Verfolgung auf den Hals laden, so wir wohl möchten \randnotiz{Friedlich leben.}friedlich leben in Gottesdienst und Wort, in aller Ehrbarkeit, wie St. Paul zu Timotheo\footnote{1 Tim. 2.} sagt, da er also spricht: So ermahne ich nun, daß man für allen Dingen zuerst thue Bitte, Gebet, Fürbitte und Danksagung für alle Menschen, für die Könige und für alle Obrigkeit, auf daß wir ein geruhlich, stilles Leben führen mögen in aller Gottseligkeit und Redlichkeit. Dann aber laden wir uns Hindernis und Verfolgung auf den Hals, wenn wir anfahen ohne Schrift und ohne Vernunft zu murmeln wider gemeine Gebräuche und Gewohnheiten, daran allein liegt Beschwerung des Säckels, des Gelds, des Leibs und der Ehre, doch ohne Schaden der Seele, wo sie allein an Christum gläubet, als den „Zehenden geben, \randnotiz{Zehenden geben, Zins reichen, opfern \&c.}Zinse reichen, vier Opfer halten, Frohndienste leisten, schadet niemand am Gewissen, und soll niemand weniger dadurch murmeln,“ denn eben die Christen, dieweil ein Christ höhern Trost hie hat von Gott, und eines größern gewärtig ist, denn alles Irdische sein mag. Es sollt doch ein Christ verschonen dem Worte Gottes, das da verhindert wird von den Ungläubigen, Pfaffen, Prälaten, Fürsten und Herren, so sie hören, daß wir gestracks wider ihren Nutz predigen. Denn so fahen sie an, uns mit dem Wort zu verjagen, sonst ließen sie uns etwan wohl bleiben, wenn wir ihnen ihre Zinse und Frohndienste gäben und reicheten. Und das ist an vielen Orten die Ursach unserer Verfolgung. Ein Christ sollte viel lieber sterben, ehe denn er wollte solche Unruhe erwecken. Ist nicht der Teufel drinnen, daß niemand unwilliger und ungehorsamer erfunden wird, denn eben etliche und viel, die
\quellseite{76}
sich evangelisch oder lutherisch nennen! Hat doch weder Christus, noch Luther solches gelehret. Verstehen sie aber diese Lehre also auf ihre Weise, so hat sie wahrlich ihre eigene Bosheit verblendet, und nicht diese Lehre.

\randnotiz{Obrigkeit zu ehren.}Nicht ohne Ursach vermahnet St. Paul so fleißig zu den Römern\footnote{Röm. 13. Col. 3. 1 Tim. 6. Eph. 5. Eph. 6.} und an viel andern Orten mehr, daß man die Obrigkeit soll in Ehren halten, ihr gehorsam sein, pflichtige Dienste leisten. Die Knechte, Unterthanen oder eigenen Leute sollen alle Treue beweisen und gehorsam sein ihren Herren, und die Frauen ihren Männern, die Kinder den Eltern, daß nicht die Lehre Christi gelästert werde, als sei sie eine Ursach solches Ungehorsams. Auch gebeut St. Paulus seinem Jünger Timotheo\footnote{1 Tim. 2.}, daß er soll vermahnen, wie jetzt neulich gesagt, daß man Gott treulich bitten soll für alle Obrigkeit, daß es ihnen wohl gehe, auf daß wir unter ihrem Regiment mögen ein geruhlich und stilles Leben führen in aller Gottseligkeit und Ehrbarkeit. Dazu vermahnet uns auch St. Peter\footnote{1 Petr. 2.} gar ernstlich. Und Summa Summarum, wie Christus im Matthäo\footnote{Matth. 7.} sagt, alles, was ihr wollet, daß euch die Leute thun sollen, das thut ihnen auch ihr. Desgleichen spricht er auch an einem andern Ort\footnote{Matth. 22.}: Gebt dem Kaiser, was des Kaisers ist. Item zu Petro\footnote{Matth. 17.} sprach er: Daß wir sie nicht ärgern, so laß uns den Zoll geben.

Wenn uns Gott gleich gnädige, fromme Fürsten gäbe, so verderben wirs „mit unserer muthwilligen Weise, daß sich Fürsten und Herren, Edlen und Städte oft unseres Frevels schämen müssen.“ Solche muthwillige Leute sind ärger, denn die Papisten. Darnach so muß es alles das Evangelion über sich nehmen, und sagen denn dieselbigen armen Leute: Ich leide viele Verfolgung um des Evangelions wil%
\quellseite{77}
Ien. Item, man will an diesem oder jenem Orte das Evangelion nicht predigen lassen, da hast du es denn fein getroffen. „Du leidest nicht um der Wahrheit und um des Evangelions willen, sondern von wegen deines Frevels und Unbescheidenheit, die du unverschämt führest und dem unverständigen Volke unweislich fürträgst. Wenn die Leute merkten, daß wir wären, wie wir sein sollen, nach Laut der Episteln St. Paulus\footnote{2 Cor. 6. Phil. 4. Röm. 12.} an vielen Orten, man würde uns wohl lassen predigen; aber wir sind so viel verleumdet worden durch unser schwärmen, daß man uns nicht mehr so bald getrauen will.“ Und das ist unsere Schuld. Der allmächtige Gott wolle sein Wort nicht abziehen, noch gute Prediger verhindern von unserer etlichen Schwärmerei wegen.

\randnotiz{närrische Fragen.}Zum sechs und zwanzigsten. Für allen Dingen gedenket, wie Paulus so fleißig, so oft und so ernstlich Timotheon\footnote{Tim. 1 und 6. 2 Tim. 2.} und Titon\footnote{Tit. 3.} vermahnet, daß sie sich sollen entschlagen der närrischen, hoffärtigen zänckischen Fragen, aus der Schrift abgeschäumt, welche allein zu Leichtfertigkeit des Gemüths dienen. Auch soll sich ein Prediger entziehen von solchen zänckischen Menschen. Der Teufel treibt jetzt viel Leute dazu, daß sie nichts thun, denn nur fragen, fragen, und wollen doch nicht geweiset werden, sondern mehr gesehen sein für andern; das ist Teufels-Dreck, der stinkt zumal übel. Und wenn die evangelischen Prediger auch beginnen anzufahen mit solchen Fragen umzugehen, werden sie ärger denn kein Sophist. Gott behüte uns dafür. Auch hindern solche Fragen gute Herzen, machen sie unfriedlich, ungewiß und unsicher. Darum nach der Lehre St. Paulus soll man solche Fragen lassen fallen und Gott befehlen, sich mehr befleißen fromm zu werden aus der Schrift, denn gelehrt. 1 Tim. 5.

\quellseite{78}
\subsection{Zum sieben und zwanzigsten}

\randnotiz{Frevele Weise zu reden von dem Christenthum.}Es sind etliche unter uns so freudig, so frevel und so unhöflicher Gebärdes, wenn wir vom Christenthum reden, daß sie meinen, man muß also wild sein; so doch St. Paul\footnote{Col. 4.} allenthalben lehret und vermahnet, ein Christ soll sanft, freundlich und lieblich sein in der Rede, und sonderlich ein Prediger, wie er es anzeigt in den Eigenschaften eines Bischofs in der ersten Epistel an Timotheon\footnote{1 Tim. 3.}. Auch wenn wir werden ersucht um Grund unsers Glaubens und unserer Hoffnung, sollen wir antworten mit Sanftmüthigkeit und Furcht. Ich bin dabei gewesen, daß ein evangelischer Prediger, in einer fast großen Stadt, so trotziglich, freudig und zänkisch vom Christenthum, auch vor vielen am Tisch, redete, daß ich mich selbst schämen mußte, also, daß einer darnach sagte: Ich meine, daß die evangelische Lehre erfordere eine solche Weise zu reden, denn alle, die ich höre davon reden, stellen sich also. Ich antwortete ihm: Nein, sondern Paulus lehret uns, „freundlich und bescheidenlich von Gottes Wort zu reden.“ Aber solche trotzige Weise ist unsrer Thorheit schuld. Mein lieber Herr Better, gebraucht eurer angebornen Freundlichkeit zu dem Worte Gottes. Es gehet gar wohl zu Ehren, „wenn man von dem sanftmüthigen, demüthigen Christo sanftmüthiglich und freundlich redet.“

\subsection{Zum acht und zwanzigsten}

Jetztgemeldte Warnung, lieber Herr Better, sollt ihr nicht verachten, sondern ihren Grund aus der heiligen Schrift baß lernen, und stellet euch nichts mehr für, denn daß ihr allen Menschen ein Spiegel seid eines ehrbaren, frommen, redlichen Lebens, wie Paulus seinen Titon\footnote{Tit. 2} vermahnet, da er also spricht: Allenthalben aber stelle dich selbst zum Fürbilde guter Werke, und halte dich in der Lehre unschädlich und redlich, und das Wort heilsam und untadelich, auf daß der Widermärtige sich schäme, und nichts habe, daß er von
\quellseite{79}
uns möge Böses sagen. Hierzu leset das Büchlein, das man nennet Pastorale St. Gregorii, und Bernardum de consideratione ad Eugenium. In einem züchtigen frommen Leben durch den Glauben werdet ihr viel Schrift lernen verstehen, auch viel Weisheit erlangen, andern Leuten auch zu rathen.

\randnotiz{Eheweib haben.}Zum neun und zwanzigsten. Ein Eheweib sollt ihr haben, und dürft ihr sie nicht öffentlich haben, so findet ihr davon einen Rath in dem Büchlein Doctor M. Luthers an den Teutschen Orden geschrieben (?). Euer Weib, Kind und Hausgesind soll regieret werden nach der Regel, die St. Paul fürschreibet in der ersten Epistel an Timotheon\footnote{1 Tim. 3.}, und das soll geschehen zur Besserung der ganzen Gemeine.

\randnotiz{Vater und Mutter ehren.}Zum dreißigsten. Und daß ich meiner lieben Verwandten, eurer Eltern, nicht vergesse, sollt ihr auch euer Weib, Kinder und Hausgesind dahin halten, daß sie dieselbigen eure alten und frommen Eltern in Ehren haben, daß also die guten Leute Freude und Förderniß an euch in ihrem Alter empfahen. Darum wird euch Gott seinen Segen geben zeitlich und ewiglich, nach der Verheißung, so er bei diesem Gebet geben hat, wie im andern Buch Mosi\footnote{2 Buch Mos. 20.} geschrieben stehet. Und also sollen alle Kinder an euch lernen Vater und Mutter zu ehren.

\randnotiz{Christlich Leben hat Verfolgung.}Zum ein und dreißigsten. Ich will euch das auch nicht unerinnert lassen, daß alle, die da wollen gottselig leben in Christo Jesu, wie Paulus sagt\footnote{2 Tim. 3.}, müssen Verfolgung leiden. Darum, alsbald so ihr anfahet also christlich zu predigen, wird groß Unglück, Angst, Noth, Anfechtung, Widerwärtigkeit, Verfolgung auf euch fallen, hie dräuen, dort schelten, anderswo heimliche Nachstellung, und was der Anstoß werden mehr sein. „Derhalben so seid gewarnet, leset euch zusammen hübsche Sprüche aus den Worten
\quellseite{80}
Christi und seiner heiligen Aposteln, desgleichen auch aus den Propheten, zu stärken euer Herz in solchem Unfall, und erschrecket nicht, lasset darum nicht alsobald ab,“ euch ist gnug, daß GOtt eurem Herzen Zeugniß gibt, daß solche Handlung ihm gefalle. Auch sollt ihr eure Hoffnung auf keinen Menschen setzen, hütet euch dafür, daß ihr keinem Herrn, keiner Gemein, keiner Stadt getrauet, oder Gott wird euch lassen zu Schanden werden. Getrauet allein GOtt, der wird euch durch und durch mit Fried und Ehren helfen.

\randnotiz{Kinderzucht.}Zum zwei und dreißigsten. Zuletzt ist dies mein getreuer Rath, daß ihr die jungen Kinder in der Wochen einmal oder drei zusammen berufet, und ihnen von Zucht, von Ehrbarkeit und von den Geboten Gottes deutlich, klärlich, kindlich saget, denn sie mögen sonst eure Predigt in dem Tempel nicht wohl verstehen, sie sind zu blöde. Glaubet das sicherlich, daß die größte Besserung der Christenheit liege an fleißiger Warte und Unterweisung der Kinder.

\randnotiz{De tradendis parvulis ad Christum.}Davon leset ein hübsches Büchlein Gersonis, wie man die Jugend auf Christum leiten soll, und lasset ja die Lehr-Schulen nicht abgehen bei euch in dem Städtlein, zu welchem euch reizen soll das Büchlein, so Doctor Martin Luther geschrieben hat, wie man soll Schulen anrichten.

\section{Beschluß}

Also habt ihr, mein lieber Herr und Vetter, eine lange unordentliche Schrift von mir ungelehrten und unordentlichen Menschen, darinnen mein einfältiger Rath verfasset ist von etlichen Artikeln, die euch, mir und andern vielen, dienen werden. Ich verhoffe, dies Büchlein soll vielen Ursach geben zu größerem Verstand und fernern Nachtrachten, welche auch für mich GOtt bitten werden
\quellseite{81}
um Gnad und Hülff, mein Leben zu bessern, denn ich täglich seufze, wie ich möge im Christenthum frömmiglich leben, aber es gehet leider noch wenig von statten. Meine fleißige Bitte ist, ihr wollet sammt andern diese meine Schrift im Besten annehmen, denn ich das herzlich und gut gemeinet habe. Und so ihr etwas hieraus gebessert werdet, wollets GOtt dem allmächtigen, von dem alles herkommt, zuschreiben und wieder dahin tragen, da es ist herkommen. Bittet GOtt für mich, und grüßet mir eure Eltern, meine Blutsverwandten, und alle andere Brüder in Christo Jesu, unserm HErrn. Geben zu Wittenberg auf den Grünen Donnerstag, Anno 1524.

Die Gnade Gottes sei mit uns allen! Amen.

\begin{unterschrift}
C. R. Vetter\\
Johann Eberlein.
\end{unterschrift}

\quellseite{82}
\begin{titelseite}
\werkuntertitel{Der}\\[1ex]
\werktitel{evangelische Geistliche}\\[2ex]
\werkuntertitel{Zweites Bändchen.}
\end{titelseite}

\quellseite{83}
\quellseite{84}
\chapter{Homiletisches}
\section{1. Hirten und Lehrer}

Unter den Gaben, welche der Herr Jesus Christus bei seiner Auffahrt von seinem Vater empfangen und seiner Kirche geschenkt hat, werden ausdrücklich genannt: die Hirten und Lehrer, ποιμένας καὶ διδασκάλους. Eph. 4, 11. Durch zwei Worte ist ein und dasselbige Amt angedeutet: die Hirten sind auch die Lehrer. (So erkannte es auch Augustinus Ep. 59 ad Paulin. Hieron. ad Ocean. Ep. 33 Balduin. p. 113.) Daher ist es auch ganz richtig, wenn man den Charakter und die höchste aller Hirtentugenden in der Lehrhaftigkeit findet, wie Theophylact, der zum ersten Kapitel an Titus sagt: „Est doctrina virtus et character episcopi;“ — wie Hilarius, der lib. 8 de Trin. schreibt: „Summa omnium virtutum episcopalium est scientia et doctrina.“

Aus der Bezeichnung eines und desselbigen Amtes mit den Ausdrücken „Hirten und Lehrer“ erkennt man nicht bloß den Zweck des Amtes, die Gemeinden zum ewigen Leben zu leiten, wie Hirten ihre Schafe zu der grünen Aue führen; sondern auch das Hauptmittel, welches den Hirten zu diesem ihren Zwecke von Gott geschenkt ist, nemlich das göttliche Wort, das aus dem Munde treuer Hirten als heilige Lehre kommt
\quellseite{85}
und die Herzen ihrer Gemeinden selig machen muß. In der That hat der Geistliche kein anderes Amtsmittel als das göttliche Wort: alles andere, was man sich sonst noch als Mittel des heiligen Amtes denken mag, wird doch nur durch Verbindung mit dem Worte befruchtet und hilfreich. Selbst im Sacramente ist es nicht sowohl das Element, als vielmehr das Wort, worauf es hauptsächlich ankommt: „Wasser thuts freilich nicht,“ Essen und Trinken thuts freilich nicht, sondern das „Wort, das mit, und bei den Elementen ist,“ das aus den Elementen Träger himmlischer Güter macht, das beides, Element und Himmelsgut zusammenschließt, die Seelen zum Empfange vorbereitet, und sie für den Genuß der geheimen Wirkungen und Segnungen des Sacramentes offen und geschickt macht. Das Wort ists gar: es ist das Wort des Herrn, durch welches die Welt geschaffen ist, und durch welches sie auch wieder erneut und geheiligt wird, und die Knechte des Herrn vollbringen daher auch alles und alles in ihrem heiligen Amte durchs Wort.

Man könnte nach diesem die Frage aufwerfen, ob nicht Amtserfahrung, Amtsgeschick, Amtsweisheit, Amtseifer und Benützung der äußeren Verhältnisse, in welchen ein Geistlicher lebt, doch auch Mittel zum Zwecke seien, ob nicht auch sie Amtsmittel genannt werden können? Die Antwort ist jedoch nicht schwer, wenn man nur richtig unterscheidet. Daß ein Geistlicher seine Verhältnisse z. B. als Vorstand des Armen-, Schul- und Stiftungswesens benützen könne, um sich in menschlicher Weise auch für das Wort
\quellseite{86}
seiner Predigt Aufmerksamkeit und Gehör zu schaffen, — daß Erfahrung, Geschick und Weisheit manches Hinderniß beseitigen können, welches sich außerdem dem Worte entgegenstellen würde, daß sie psychologisch pädagogisch dem Worte Bahn machen und Herzen in menschlicher Weise zum Empfang göttlicher Segnungen vorbereiten können, — daß ein heiliger Eifer geeignet ist, dem Worte Gelegenheiten zu erspähen und die sich darbietende Gelegenheit schnell zu ergreifen und zu benützen: wer könne und wollte das leugnen? Man würde ja damit jede Amtsanweisung überflüssig und unnütz nennen, und die ganze Pastoraltheologie, welche doch alles Preises werth ist, in den Staub treten. Ja, man würde vergessen, daß der HErr, der heilige Geist, der einem jeglichen Seines, sein bescheiden Theil von seinen Gaben mittheilt, manches von dem Erwähnten grade als seine Gabe und Kraft bezeichnet. (1. Cor. 12, 4 ff. Διαίρεσις χαρισμάτων εἰσί, τὸ δὲ αὐτὸ πνεῦμα. λόγος σοφίας, λόγος γνώσεως etc. v. 8. — πάντα δὲ ταῦτα ἐνεργεῖ τὸ ἓν καὶ τὸ αὐτὸ πνεῦμα, διαιροῦν ἰδίᾳ ἑκάστῳ καθώς βούλεται.) Des Geistes Gaben aber darf man nicht verachten, sondern um sie eifern soll man. Ζηλοῦτε δὲ τὰ χαρίσματα τὰ κρείττονα. 1. Cor. 12, 31. Aber so hoch man auch alles das anschlage, und von so großer Wirkung das Amtsbenehmen und die Gaben eines Pfarrers sein können, so bleibt es doch wahr, daß alles nur heilsam wirkt, wenn und weil es von dem göttlichen Worte durchdrungen ist; das Benehmen des Geistlichen gibt weder noch mehrt es dem Worte die Kraft, sondern es ist und bleibt doch immer nur
\quellseite{87}
das Wort, welches die selige Wirkung hervorbringt. Man kann das treffliche Amtsbenehmen eines Geistlichen und die Gaben der silbernen Schaale vergleichen, auf welcher die goldenen Früchte des göttlichen Wortes liegen. Die Schaale ist der Frucht würdig, und empfiehlt sie in menschlicher Weise den Menschen, aber Erquickung, Stärkung und Leben geben und mehren können am Ende doch nur die Früchte des Lebensbaumes selbst, die gütigen Kräfte des Wortes.

Daher kommt es auch, daß Geschick, Eifer und ein kluges Benehmen, losgetrennt von dem göttlichen Worte, zwar nicht unwirksam sein müßen, wohl aber verderblich wirken können. Ja man kann sagen, daß jede geringe Abweichung vom göttlichen Wort im Amtsbenehmen eines Geistlichen ihre schädlichen Folgen hat. Laß z. B. einen Geistlichen in seinem Amtsbenehmen, bei dem guten Satz, den er vertheidigt oder vorträgt, nur das rechte Maß verfehlen, das er haben sollte, laß ihn zu eifrig und zu feurig sein, so ist schon das verderblich. Gleichwie der Wagenlenker mit der einen Hand seine Thiere antreibt, mit der andern aber sie im Zaume hält; so muß der menschliche Wille des Seelsorgers von dem göttlichen Worte beides gehemmt und getrieben sein, wenn das heilige Amt und seine Werke, und nicht das Gegentheil geschehen soll.

Diese Unterscheidung zwischen Amtsmittel, Amtsgabe und Amtsbenehmen zu machen, thut man gewiß sehr wohl. Gar leicht überschätzt man das Amtsbenehmen im Vergleich mit den Amtsmitteln, und es dürfte das z. B. ein bedeutender Unterschied zwischen den kirchlichen Pastoraltheologen und denen eines
\quellseite{88}
Gottfried Arnold und Baxter \&c. sein, daß jene mehr auf die Mittel, diese mehr auf’s Benehmen sehen, jene auch in Betreff des Maßes, welches ein Geistlicher für seinen Fleiß und Eifer anzuwenden hat, ungesund und fast möchte man sagen pelagianisch urtheilen, während die kirchlichen Theologen das Maß der heiligen Mitte zwischen menschlicher Extreme einhalten. Wie nöthig jene Unterscheidung zwischen Mittel und Weise oder Benehmen sei, sollten namentlich diejenigen Geistlichen recht beherzigen, die, sei es durch Temperament, oder Führung, zu einer methodistischen Amtsführung gekommen sind. Methodistische Lehrer setzen, oft ohne es zu wissen oder zu wollen, gerne zu große Hoffnung auf die eigene menschliche Zuthat zum Wort, auf das Amtsbenehmen; sie verlieren die höhere Zuversicht auf das Wort der Wahrheit und geraten dafür in eine hastige, unruhige Vielthuerei und Menschenfischerei hinein, während doch die Ruhe und stille Größe jener Zuversicht die Herzen der Menschen am Ende mehr anzieht, als ein maßloses Wirken. Es muß ja auch die genannte Ruhe keineswegs in Trägheit enden, und die Ueberzeugung, daß am Ende alles am Wort und an der Gnade Gottes liege, muß keineswegs das Bewußtsein ertödten, daß man für Art und Weise seiner Amtsführung eine große Verantwortung habe, daß man durch weises Thun dem Licht zum Leuchter, der Gabe zum dankbaren Empfang, durch verkehrtes Beginnen aber dem Licht zum Schäffel, der Gabe zur Verachtung helfen könne. Dies Bewußtsein und jene Ueberzeugung zusammengenommen, geben dem seelsorgenden Manne sein gerechtes Maß, und damit wie viel!

\quellseite{89}
Ein Seelsorger, der alle Hoffnung auf seine Arbeit setzte, und den himmlischen Mitarbeiter vergaße, wäre, um bei einem schon gebrauchten Gleichniß zu bleiben, einem Fuhrmann ähnlich, der seine Pferde nur immer antreibt und eben deshalb zu Schaden, nicht aber zum Ziele kommt. Ein Seelsorger des gerechten Maßes hingegen gleicht einem Fuhrmann, der, während er seine Pferde mit der Geißel antreibt, zugleich mit dem Zaume sie lenkt und eben dadurch schnell und sicher zu seinem Ziele kommt.

Aehnlich ist es mit dem guten Beispiele eines Geistlichen. Wehe dem, der es verabsäumt, und der Heerde Christi das göttliche Wort ohne Empfehlung durch’s Beispiel, vielleicht gar unter dem Hinderniß des Beispiels, nemlich eines bösen Beispiels, predigt. Böses Beispiel ist wahrlich nicht bloß ein menschliches, sondern ein satanisches Hinderniß des Wortes. Die Menge ruft mit tausend Zungen: Vita clerici Evangelium populi (Apstg. 5, 13.); sie hat ein arges Auge, sie läßt kaum ein reines und heiliges Beispiel unbekrittelt und unbetitelt, sie hascht nach jedem Flecken eines Predigers und wäre es auch nur, um Aergerniß zu nehmen, das eigene böse Gewissen zu stillen, die eigene Sünde zu entschuldigen. Es kann daher einem Diener des Evangeliums nicht tief genug eingeprägt werden, wie verantwortungsvoll es sei, dem Worte, das er predigt, durch das eigene Beispiel zu widersprechen. Aber bei aller Werthschätzung des Beispiels, bei allem Preise des heiligen Lebens eines Seelsorgers, bleibt doch Leben und Beispiel nur ein menschliches Mittel, welches dem göttlichen Gna-
\quellseite{90}
Gnadenmittel des Wortes nicht gleich zu stellen ist. Wie es vom Menschen kommt, wirkt es auch nur menschlich, vielleicht hie und da einmal gewaltig und überwältigend, aber doch immer nur menschlich und nur Menschliches. Manches Beispiel trieft vom göttlichen Worte, aus dem es auch selber seinen Ursprung genommen hat, es ist dem Leibe vergleichbar, von dem nach dem Worte des Herrn Wasser des Lebens fließt; es wird verehrt, angestaunt und nachgeahmt. Aber wird es zur Buße, zum Glauben, zur Heiligung führen; wird es wahrhaft geistlich wirken, wenn es nicht durchs Wort erklärt wird, wenn nicht gewissermaßen das Wort und seine Kraft von ihm losgelöst und in der Gestalt der Lehre, der Strafe, der Besserung und Züchtigung den Ohren und Herzen nahe gebracht wird? Ist ein Christenleben, so wohl es den Zeitgenossen gefalle, so große Anerkennung es finde, ohne Wort auch nur verständlich? Wird ohne Wort die Quelle der wahren Nachfolge entdeckt? Gewiß nicht! Eine solche Einsicht in dies Verhältnis zwischen Wort und Beispiel hatten die Alten, daß Hieronymus ad Oceanum ep. 33 schreibt: Innocens et absque sermone conversatio, quantum exemplo prodest, tantum silentio nocet. Fast übertrieben klingt dies Wort, aber es ist jedenfalls aus der richtigen Ueberzeugung und Erfahrung entsprungen, daß Wort und Lehre weit höher als das Beispiel zu schätzen sind. Jene können ohne dieses wirken, dieses aber wird und wirkt nicht ohne die Kraft von jenem, nämlich nicht zum Guten. Zum Bösen hilft freilich böses Beispiel mächtig vorwärts, weil es in der Verderbniß des Menschenherzens einen weichen und
\quellseite{91}
kräftigen Boden findet. — Mit alle dem dürfte denn doch Einsicht und Erfahrung christlicher Lehrer zusammenstimmen und so der Satz fest stehen, daß auch das Beispiel im eigentlichen Sinne kein Amtsmittel des Geistlichen genannt werden kann. Mögen viele heilige Beispiele in der christlichen Kirche blühen, aber ohne Ueberhebung, bescheidenlich, ohne Hoffnung und Anmaßung solcher Wirkungen, wie sie allein vom Wort und Sacrament ausgehen.\footnote{Gleichwie der Geistliche außer dem Wort kein Mittel hat, geistlich und wahrhaft umwandelnd auf die Seele zu wirken; so gibt es überhaupt und in der ganzen Welt keines außer dem Wort. Losgerissen vom Worte und ohne Träger desselben zu sein wirkt keine Creatur zur Seligkeit. Wie vielen Segen hofft mancher für die Umänderung der Gemüther, z. B. vom Wechsel der Zeit, von Ortsveränderungen, von Lebensschicksalen, vom Gebrauch natürlicher Dinge, der Betrachtung der Natur, von anderem Umgang und von dergleichen mehr, da doch alles das nicht vermögend ist, eine wahre Aenderung des Herzensgrundes hervorzubringen. Wohl sagt Jesaias 28, 19.: „Allein die Anfechtung lehret auf das Wort merken“; aber sie lehrt nur aufs Wort merken, das zuvor schon da ist, und weist ebendeswegen aufs Wort, weil sie selber zu arm ist, die gewünschte Aenderung im Menschenherzen hervorzubringen. Den alten Menschen in andern Gestalten umzuwandeln, dazu taugen auch andere Mittel; aber aus dem alten einen neuen Menschen zu machen, das vermag allein Gottes Wort und seine schöpferische Kraft. Daher ist auch die Erfahrung des Wortes die wünschenswertheste und höchste Erfahrung, die segensreichste für Zeit und Ewigkeit, während alle weltlichen und menschlichen Erfahrungen, abgesehen davon, daß sie unzuverlässig sind, doch keine wahrhaft bessernde Kraft besitzen.}

\quellseite{92}
„Dein Wort ist eine rechte Lehre“, betet der heilige Sänger, und welcher Geistliche sollte nicht hoch erfreut sein, dies Wort und diese rechte Lehre so nahe zu haben, in ihm sein gewaltiges Amtsmittel zu finden? Besitzesfroh wollen wir aber dem Ausspruch des heiligen Sängers, ehe wir weiter gehen, noch ein wenig in’s Auge sehen. „Gottes Wort ist eine rechte Lehre“: dieser Psalmensvers stellt wohl Gottes Wort eines Theils in eine und dieselbe Reihe mit andern Lehren; aber er hebt es auch wieder von allen andern heraus, indem er eben sagt „es ist eine rechte Lehre.“ Andere Lehren gleichen ihm nicht, nicht dem Inhalte nach (und das leugnet niemand!) aber auch nicht der Kraft nach, mit welcher es sich an die Menschenherzen wendet. Zwar wendet auch das Wort des Herrn sich an die natürlichen Kräfte des Menschen, an sein natürliches Erkenntnißvermögen; es setzt auch die Kräfte des natürlichen Willens voraus, welcher dasjenige zu bewerkstelligen sucht, was die Vernunft als recht erkannte. Es geht also in den Menschen auf demselbigen Wege ein wie ein anderer menschlicher Gedanke. Auch wirkt es im Innern des Menschen durch die Macht seiner Gründe eine menschliche Ueberzeugung wie ein menschlicher Gedanke. Dennoch aber ist es kein bloßer Menschengedanke und wirkt nicht bloß menschlich, auch nicht bloß eine menschliche Ueberzeugung, sondern es bringt eine Kraft mit sich, welche den Widerstand der natürlichen Vernunft, oder vielmehr Unvernunft, niederlegen und dem Worte einen Nachdruck verleihen kann, vermöge dessen es eine göttliche jeden Widerspruch und Zweifel überwältigende Ueberzeugung, den Glauben, zu
\quellseite{93}
wirken vermag; denke nur an die Kraft der Absolution. Das meint die Kirche, wenn sie sagt, es werde durchs Wort und Evangelium der Geist gegeben, das Wort sei Organ und Träger des heiligen Geistes. Eben dasselbe meint auch die heilige Schrift, wenn sie von einer Erleuchtung des heiligen Geistes durch das Wort, von einem Schmecken des Wortes Gottes und seiner Kräfte redet Ebr. 6, 4 ff., wenn sie vom Worte sagt, daß es lebendig und kräftig, daß es schärfer sei, denn ein zweischneidig Schwert, daß es durchdringe, bis daß es Seele und Geist, auch Mark und Bein scheide, daß es ein Richter der Gedanken und Sinne des Herzens sei, ja daß keine Creatur vor ihm unsichtbar, hingegen alles bloß und entdeckt sei vor seinen Augen Ebr. 4, 12 f. Daraus erkennen wir eine mehr als natürliche, eine übernatürliche und unbegreifliche Art des Wortes Gottes und seiner Wirkung. Wir sehen, das Wort ist menschlich und göttlich zugleich: menschlich, denn es redet in menschlicher Sprache zu menschlichem Verständnis; göttlich, denn es wirkt mit göttlicher Kraft eine himmlische Erkenntnis zum ewigen Leben. Welch ein Amtsmittel ist also dem Geistlichen in die Hand gegeben, wenn ihm Gottes Wort vertraut ist! — Zwar hat man gesagt, das Wort, welches ein Geistlicher bei seinen Amtsverrichtungen ausspricht, sei nicht Gottes Wort im eigentlichen Sinne. Der Inhalt könne wohl dem göttlichen Worte entsprechen, aber der Vortrag, die Fassung dieses Wortes sei ein menschliches Werk; nur die Apostel hätten Gottes Geheimnisse in Worten geredet, die der heilige Geist lehrte. Man mag
\quellseite{94}
in einem gewissen Sinne Recht haben. Aber hört denn Gottes Wort durch eine richtige menschliche Auffassung auf, Gottes Wort zu sein? Es zieht menschliche Gestalt an, aber bleibt deshalb Gottes Kraft nicht bei ihm? Lies einmal aufmerksam die Pastoralbriefe Pauli und überzeuge dich, welche göttliche Kraft und welche göttliche Segnungen dem menschlich aufgefaßten Worte beigelegt werden.

Wir erinnern an ein schlagendes Gleichniß der Väter.\footnote{Cf. V. 2. Löscher’s vollständigen Timotheus Verinus. I. p. 276.} Das Sonnenlicht fällt in einen Spiegel und dieser wirft es zurück: ist nun bloß das Sonnenlicht, was in den Spiegel hinein fällt, oder auch das, was widerstrahlt? Offenbar auch das aus dem Spiegel zurückstrahlende Licht. Ist dies wahr, so ist es auch richtig, daß das aus dem Munde eines Predigers hervorbringende Wort, so lange es dem Worte der Schrift, aus welcher er schöpfte, getreu bleibt, im Grunde nichts anderes ist, als ein und dasselbige Gotteswort, von dem die Schrift sagt, daß es einen hellen Schein in unsre Herzen gebe, auf daß durch uns die Erleuchtung anderer entstünde. Es ist der Wille Gottes, daß sein heiliges Wort im Herzen seiner Knechte die Art und Weise dieser seiner Knechte annehme, auf daß es aus Glauben gepredigt werde und in göttlich-menschlicher Kraft einen Glauben wirke, der selber göttlich und menschlich sei. Es geht hier aus Glauben in Glauben: das Wort wird lebendig gepredigt und wirkt auch Leben. Darum erinnert auch der heilige Apostel nicht bloß an Apostel, wenn er vermahnt: „Gedenket an eure Lehrer, die euch das Wort Got-
\quellseite{95}
tes gesagt haben.“ Ebr. 13, 7. Nicht bloß die Apostel, sondern auch andere dem Wort getreue Prediger können sich getrösten, indem sie predigen, das große Amtsmittel zu besitzen, von dem wir eben reden.\footnote{Es kann gar keine Frage sein, sondern es ist offenbare gewisse Wahrheit, daß auch ein einfacher Vater, eine fromme Mutter, ein junges Kind Gottes Wort mit reichem Segen reden und predigen können. Wer sollte die Kraft des göttlichen Wortes vom Amte abhängig machen wollen wider das Zeugnis der Jahrhunderte und die tägliche Erfahrung? Die göttliche Kraft, die sich mit der ewigen Wahrheit verbindet, wohnt ihr inne und wirkt aus ihr, sie werde gepredigt, von wem es auch sei. Nichtsdestoweniger ist und bleibt es aber doch auch wahr, daß der Herr seinen Knechten gewisse Verheißungen zur Führung des heiligen Amts gegeben hat. Nicht bloß den Aposteln wird eine mitwirkende Gnade (ἡ χάρις θεοῦ, ἡ σὺν ἐμοί. 1. Cor. 15, 10. — 3, 10. Gal. 2, 9. Eph. 3, 2. 7. 8.), nicht bloß Evangelisten ein χάρισμα (1. Tim. 4, 14. 2. Tim. 1, 6.) zugeschrieben, sondern auch einem einfachen Lehrer, dem Apollos, schreibt St. Paulus nicht minder wie sich selber die διακονία πνεύματος zu (2. Cor. 3, 8), also doch wohl auch das πνεῦμα, den Geist, ohne welchen das Amt nicht segensreich werden kann. Auch ruft er 1. Cor. 3, 5 ff.: „Wer ist nun Paulus? Wer ist Apollos? Diener sind sie, durch welche ihr seid gläubig worden, und dasselbige wie der Herr einem jeglichen gegeben hat etc.“ Und V. 9 sagt er von Apollos eben so wie von sich selbst: „Θεοῦ ἐσμεν συνεργοί, wir sind Gottes Mitarbeiter.“ Ist aber Apollos ein Diener und Mitarbeiter Gottes, so ist auch Gott, sein Herr, der durch ihn arbeitet, sein Mitarbeiter; Apollos und alle, die ihm gleich, alle Lehrer des Evangeliums, sind demnach nicht alleine, wenn sie des Amtes walten, sondern es ist mit ihnen die Gnade und der Geist des Herrn, der allen seinen Knechten besonderen Beistand zusagte, da er sprach: „Siehe, Ich bin bei euch alle Tage bis an der Welt Ende!“ Ihnen allen wird also mit Recht das Amt des Geistes zugeeignet und von allen kann man deswegen mit Wilhelm Lyser sagen: „Vocatio ministerium constituit, ita ut hominibus per illam potestas super humana concedatur, et homines in gratia dispensanda cooperarii Dei reddantur — —, ita ut ultra humanam sortem operatione Dei per ipsam vocationem eleventur. — — Minister non tantum significative divinum quod indicat, sed Deus revera effective per illa ipsa organa divina opera perficit etc., ita ut ἐνέργεια organi propter accedentem causae principalis conjunctissimam cooperationem et elevationem concurrat ad commune ἀποτέλεσμα producendum, quod vero divinus effectus est etc., ut ab ipsis effective divinus effectus producatur.“ (S. V. 2. Löschers vollständ. Timotheus Verinus I. p. 290.) Wenn deshalb überhaupt nicht zu leugnen ist, daß der evangelische Geistliche als vornehmstes Mittel zur Ausrichtung seines Amtes das göttliche Wort besitze; so muß ihm dies auch zugestanden werden, daß er sich beim treuen Gebrauche des Wortes besonderer göttlicher Verheißungen zu erfreuen habe, und zwar müssen diese Verheißungen, ganz abgesehen von dem Gnadenstande des Lehrers, seiner Amtsführung in solange zugeschrieben und belassen werden, als er sich in derselben von dem Worte Gottes nicht entfernt. Im umgekehrten Falle würde ja Gottes Amt von dem menschlichen Verhalten, Göttliches von Menschlichem abhängig gemacht. Mag immerhin zuweilen die Mühe eines Predigers nicht viel größer sein, als die eines Buchdruckers, der die ganze Schrift unverfälscht setzen und abdrucken kann, ohne deshalb einen seligen Gebrauch für sich selbst davon zu machen; so ist doch die Mühe und der Fleiß in dieser Sache gar nicht das punctum saliens, sondern die wortgetreue Amtsführung. Ueberdies liefert ja die Erfahrung Beispiele genug, daß offenbar und unwiderstreitbar der reinen Lehre selbst unsittlicher Lehrer von dem Herrn eine segensvolle Wirkung beigelegt ward. Greifen wir also nicht an, was unwiderstreitbar ist, und freuen wir uns vielmehr, daß es also ist. Die Amtsmittel müssen ja auch desto höher und herrlicher erscheinen und desto größeres Vertrauen erwecken, je weniger sie vom Menschen abhängen und je mehr hingegen der Mensch vom Wort abhängt.}

\quellseite{96}
\section{2. Lehrgabe}

Wer ein so herrliches Amt und so gewaltige Amtsmittel besitzen will, wie Gottes Wort und Sacrament, der soll auch billig Fähigkeit und Gabe haben, Amt und Mittel zu gebrauchen. Keiner soll ohne diese das Amt bekommen. Da-
\quellseite{97}
her verlangt auch Sanct Paulus 1. Tim. 3, 2. von allen Bischöfen, daß sie διδακτικοί oder lehrhaft seien. Wort und Lehre ist Amtsmittel, Lehrhaftigkeit also nothwendige Amtsgabe. Demgemäß erinnert der heilige Apostel seinen Timotheus, 2. Timoth. 2, 2., daß er nur solchen das Lehramt übergebe, welche tüchtig seien, andere zu lehren. Ja nach Tit. 1, 9—11. sollen die Aeltesten oder Bischöfe mächtig sein, zu ermahnen durch die heilsame Lehre, zu strafen die Widerspenstigen und den frechen, unnützen Schwätzern und Verführern das Maul zu stopfen.

Man wird allerdings die Lehrhaftigkeit nicht zu den von der Mittheilung des heiligen Geistes abhängigen, besonderen geistlichen Gaben rechnen wollen; sie ist eine natürliche Gabe, ohne welche keiner in die Wahl für das geistliche Amt kommen soll, eben so wenig, als einer zu irgend einem anderen irdischen Beruf zugelassen werden soll, der die dazu nöthigen natürlichen Gaben und Eigenschaften nicht besitzt. Damit aber daß die Lehrgabe in die Reihe der natürlichen Gaben einge-
\quellseite{98}
ordnet wird, wird sie nicht vermehrt: die Gaben des Schöpfers sind Gottes Gaben, wie die des Erlösers und des heiligen Geistes. Ja die natürlichen Gaben und Fähigkeiten werden oft durch den heiligen Geist nicht bloß an und für sich selber gereinigt, gehoben und verstärkt, sondern auch Träger und Organe der eigentlich geistlichen Gaben. So ist 1. Kor. 12, 8 ff. unter den eigentlich geistlichen Gaben die Gabe zu reden von der Weisheit aufgeführt, desgleichen die andere Gabe zu reden von der Erkenntniß. So wenig die Weisheit und die Erkenntniß, von welchen hier die Rede ist, natürliche Gaben genannt werden dürfen, und so gewiß sie in das Gebiet des dritten Artikels gehören, eben so wenig ist die Gabe, von dieser Weisheit und Erkenntniß zu reden, eine natürliche Gabe, und eben so gewiß ist es eine besondere Wirkung und ein Geschenk des Geistes Jesu, von Weisheit und Erkenntniß des Neuen Testamentes reden zu können. Wer sieht nun nicht, daß sich eigentlich geistliche Gaben hier an natürliche geistige Gaben anschließen und sich zu diesen verhalten etwa, wie die Krone eines Königs zum König selbst. Was ist die Krone ohne den König, und eben so könnte man sagen: was ist die geistliche Gabe ohne die ihr entsprechende natürlich geistige Gabe, die für jene voraus gesetzt wird? Das alles ist bloß zur Erhebung und zum Preise der natürlichen Gabe der Lehrhaftigkeit gesagt, ohne deren Vorhandensein das große Amtsmittel des göttlichen Wortes weder zum Gebrauche, noch zu seiner Wirkung und zu seinem Segen kommen kann.

Diese Lehrhaftigkeit ist nun aber freilich nicht eins und
\quellseite{99}
dasselbe mit der Beredtsamkeit oder Rednergabe, mit welcher man sie unter uns Protestanten so gern verwechselt. So schätzbar die Gabe der Beredtsamkeit ist, und so verwandt mit der Lehrgabe, so kann man doch, denken wir, die Lehrgabe im hohen Grade besitzen, während die rhetorische Gabe mangeln kann: es könnte ein Diener Jesu vermöge seiner Lehrhaftigkeit die größte Wirkung auf eine Gemeinde haben, ohne daß er jemals als Redner sein Glück machen könnte. Bei Festhaltung dieses Unterschiedes und bei der aus dem göttlichen Worte unzweifeligen Gewißheit (1. Kor. 12, 31.), daß auch unter Gottes geistigen und geistlichen Gaben ein Unterschied ist, daß es größere und geringere, vorzüglichere und unwichtigere gibt, dürfte wohl die Lehrhaftigkeit über die Beredtsamkeit zu erheben und nach Würden zu pflegen und auszubilden sein, wie doch jede Gabe durch Pflege gehoben und gestärkt, durch Vernachlässigung aber geringer und unbedeutender werden kann.

Freilich ist aber auch ein Maß der geistlichen und geistigen Begabung vorhanden, welches selbst durch Pflege und Bildung nicht überschritten werden kann: über die vom Herrn gesetzten Schranken kommt niemand hinaus. Diese Wahrnehmung veranlaßt die nachfolgenden Erfahrungssätze: Wozu haben wir in der heiligen Schrift so vielfache Belehrung über die geistlichen Gaben, wenn nicht dazu, daß wir die Belehrung auf uns anwenden sollen? So wahr der Herr, der heilige Geist, lebt, so gewiß wirkt Er in seiner Kirche, und so gewiß die Kirche selbst noch besteht, so gewiß hat sie auch noch Gaben des heiligen Geistes, wahrhaftige geistliche Gaben, für
\quellseite{100}
deren Dasein wir freilich kein Auge haben, und nach denen wir nicht forschen, die uns aber, mag auch ihr Maß durch unsere Schuld geringer sein, als früher, mit erfreuendem Glanze in die Augen leuchten würden, wenn wir nur die Augen öffnen und die Gaben suchen möchten. Wer aber unter uns weiß etwas von seiner geistlichen Gabe, oder doch: wer weiß etwas Gewisses und Sicheres darüber? — Ist es doch gerade so mit der leichter erkennbaren natürlich geistigen Gabe. Wie viele wissen, wofür sie nur ganz im allgemeinen begabt sind? Jede Berufswahl ist durch die Gabe bedingt; weil man aber die Gabe nicht erkennt, geräth so oft die Wahl des Berufes so schlecht. Ist dann einmal ein Beruf ergriffen, so ist die Entdeckung, daß die nöthige Gabe nicht vorhanden sei, wie man sich denken kann, so niederschlagend, daß man ihr möglichst ausweicht und das Mißlingen des Berufes lieber aus allen anderen Ursachen herleitet, als aus der wahren, dem Mangel an Begabung. Ein Pfarrer, der keine Lehrhaftigkeit hat, gibt das zuletzt zu. Freunde und Verwandte, vor allen das eigene Weib, reden es ihm aus, wenn er etwa auf die rechte Spur gekommen ist, und die Erfahrung bezeugt es, wie oft eine ganze Gemeinde die Ueberzeugung hat, daß ihr Pfarrer die eigentliche Gabe der Lehrhaftigkeit nicht besitzt, während nur er selbst und die Seinen davon weder wissen, noch wissen wollen. Sündenerkenntniß ist leichter als Gabenerkenntniß, und Armensünderdemuth eine weit geringere, weit leichtere Tugend, als die Demuth, die sich darein ergiebt, eine Gabe nicht zu haben, die man eigentlich nothwendig brauchte, aber
\quellseite{101}
auch nur ein geringeres Maß dieser Gabe zu besitzen. Wie schwer ist es dem Menschen, Mängel seiner Begabung und Anlagen zu erkennen! Und doch läge in der richtigen Selbsterkenntniß oft sogar Anfang, Keim und Trieb einer Besserung. Aber auch das wird versäumt, die letzte Hoffnung auf Besserung und größeren Erfolg der Amtsführung dadurch weggenommen, daß man die Selbsterkenntniß scheut und die Demuthswege nicht mag, auf denen man werden kann, was möglich ist, und Schaden verhüten, der aus der eitlen Ueberhebung und aus der Affectation von Gaben und Gabenmaßen kommt, die nun einmal nicht da sind. Es läßt sich kein unglücklicherer Mensch denken, als einer, der lehren soll, und nicht kann; keine größere und schmerzliche Armuth als diese, die auch in der That dadurch nicht geringer wird, daß man sich, wie viele Pfarrer, einbildet, man sei nicht so arm, sondern reich. Man wird Mühe haben, diese Einbildung fest zu halten, wenn man alt wird, und nun je länger, je mehr die Gemeinde des armen Redners müde wird. Da gibt's alsdann ein Sterben, das man sich hätte ersparen können, wenn man einen angemesseneren Beruf erwählt hätte.

\section{3. Verschiedene Formen der Lehre}

Wenn man das Wort Lehre genau nach dem Wortsinne faßt, so bezeichnet es nichts anders als Mittheilung von Erkenntniß, im Sinne, von dem hier die Rede ist, Einführung in die Erkenntniß des Heils. So wird dann jede Art und
\quellseite{102}
Weise der Mittheilung zu diesem Zwecke Lehre. Ob der Zweck durch einen fortlaufenden und zusammenhängenden Vortrag oder in der Form des Gesprächs erreicht wird, das ist alsdann der Hauptsache nach gleichgültig, und man kann daher eigentlich nicht sagen, daß der zusammenhängende Vortrag auf das Gebiet der Homilie gehöre, während die dialogische Form die Form der Katechese sei. Die Katechesen des heiligen Cyrillus von Jerusalem, die von Alters her bei der Kirche in hohen Ehren stehen, sind zusammenhängende Vorträge, während die bei unserem Volke noch so sehr beliebten Homilien von Spangenberg ähnlich wie das in römisch-katholischer Kirche weit verbreitete, keineswegs auf das Gebiet der Katechese gehörende „Unterrichtsbuch von Goffine“ in Frage und Antwort predigen. Nicht die Form, wohl aber Inhalt und Absicht vertheilen die menschlichen Vorträge auf die Gebiete der Katechese und der Homilie. Wird bloß Erkenntniß mitgetheilt, und ist keine andere Absicht vorhanden, als diese, so dient der Vortrag der Lehre im eigentlichen Sinne. — Es gibt nun aber noch andere Absichten der göttlichen und menschlichen Rede, als die der Belehrung, wie wir das II. Tim. 3, 16. deutlich ausgesprochen finden, wo es nach Luthers Uebersetzung heißt: „Alle Schrift, von Gott eingegeben, ist nütze zur Lehre, zur Strafe, zur Besserung, zur Züchtigung in der Gerechtigkeit, daß ein Mensch Gottes sei vollkommen, zu allem guten Werk geschickt.“ Wie nun die von Gott eingegebene Schrift, so hat auch die menschliche Rede oder Schrift der Diener Gottes nicht bloß Einen Zweck, den der Belehrung%
\quellseite{103}
, sondern einen mehrfachen, wie das II. Tim. 3, 16. in den Ausdrücken des Grundtextes so vollkommen und von der Kirche anerkannt, vierfach theilend und zusammen faßend dargelegt wird. Sowie nun ein Vortrag nicht ausschließlich den Lehrzweck verfolgt, sondern die gütigen Kräfte des göttlichen Wortes in ihrer Mannigfaltigkeit auf die Herzen der Hörer anzuwenden bestrebt ist, wie das bei den Vorträgen der alten und neuen Lehrer innerhalb der Kirche je und je der Fall gewesen ist, so hört er auf, katechetisch, d. i. im besonderen Sinne lehrhaftig zu sein, und tritt, wenn man so sagen will, in den Dienst des christlich gemeindlichen Gesammt-Zweckes, wird homiletisch. Bei alle dem aber läßt es sich nicht leugnen, daß die dialogische oder Gesprächs-Form ganz besonders zur Katechese und Einführung in die Lehre und Erkenntniß des Heils paßt; man darf nur an kein bloßes Abfragen und Einlernen von Antworten denken, sondern an das lebendige Wechselverhältniß eines Lehrers und Schülers, die beide, wie Christus unter den Lehrern im Tempel, fragen und antworten und auf diese Weise am vollkommensten das Ziel der Erkenntniß erreichen. Kein acroamatischer Vortrag kann alle Hindernisse und Bedürfnisse des Schülers berücksichtigen; was diese Art des Vortrags unerreicht läßt, wird durch die aufrichtige Frage des wißbegierigen Schülers und durch die Antwort des Lehrers, also durch das eingehende katechetische Gespräch erreicht, welches am Ende die für Lehrer und Schüler erwünschteste Form der Belehrung sein, und für alle Katechesen als Ziel und Muster stehen wird. Wie nun aber die dialogische Form
\quellseite{104}
auf dem Gebiete der Katechese, so ist die des zusammenhängenden Vortrags in Mitte der unterrichteten und gläubigen Gemeinde die natürlichste und von sich selbst sich ergebende. Die verschiedenen Bedürfnisse der Gemeinde sprechen sich in den Worten Strafe, Besserung und Züchtigung aus, und werden in der öffentlichen Versammlung sicherlich am allerbesten in der zusammenhängenden Rede des Bischofs oder Pfarrers befriedigt. Daher wird man wohl sagen dürfen: der lehrhafte Diener Christi erreicht seinen Zweck eben so wohl durch den zusammenhängenden Vortrag, als durch die dialogische Form; will er die Gemeinde, oder überhaupt viele zugleich lehren, so wird der zusammenhängende Vortrag herrschen; gilt es die Belehrung Einzelner und Ungeschickterer, so wird die Gesprächsform vorwalten, und weil wir so selten in einem anderen Fall sind, als in letzterem, so begreift es sich, weshalb man bei uns die katechetische Form ganz einfach als Gesprächsform nimmt, die Katechese vom Gespräch, oder nicht besser, aber doch der Wirklichkeit entsprechender zu reden, von Frage und Antwort nicht trennen kann. Frage und Antwort geben allerdings noch kein Gespräch, aber sie sind das, was bei unseren Katechesen von dem meist nicht zu erreichenden Gespräche übrig bleibt.

Zur Zeit, da die Reformatoren auftraten, war der Erkenntnißzustand unseres Volkes ein so tief gesunkener, daß man vor allen Dingen reine Lehre verbreiten mußte. Lehrhaftigkeit wurde nicht bloß das erste, sondern schier das einzige Erforderniß, schier die einzige Tugend eines Lehrers. Alles
\quellseite{105}
Singen und Sagen, ja wahrhaftig auch das Singen, mußte in den Dienst der Lehre, der reinen Lehre treten. Da konnte es schon nicht mehr anders sein, es mußte auch die Liturgie nicht bloß lehrhaft werden, was sie in einem gewissen Maße, wie jedes geistige Produkt der Kirche sein muß, und gar nicht anders kann, sondern es wurde ihr die Verbreitung der reinen Lehre zum eigentlichen Zweck gegeben. Allerdings sind die Liturgien der alten und neuen Kirche unter die Quellen der Erkenntnis, für diejenigen zu rechnen, die begierig sind, die Verschiedenheiten der kirchlichen Lehre bei den mancherlei Kirchenparteien kennen zu lernen. Auch kann man nicht leugnen, daß die Jugend insgemein reiche Schätze der Erkenntnis aus der Liturgie entnimmt, an der sie heran wächst. Aber wird man deswegen sagen dürfen: die Liturgie sei nichts anderes, als ein Lehrmittel der Kirche, ein Mittel in die kirchliche Lehre einzuleiten, oder eine besondere Form der heilsamen Lehre? Gewißlich nicht! Gewiß betet die Gemeinde nicht deshalb Gott an, damit die Unwissenden aus ihrer Anbetung Gott kennen lernen. Wenn die Gemeinde ihre Herzen zu Gott erhebt, so ist die Anbetung nicht bloß eine ihrer Absichten, etwa neben der Lehre, sondern ganz einfach ihre Absicht. Wer Gott anruft, um die Unwissenden zu lehren, daß und wie man Gott anrufen solle, würde sich in Acht nehmen müssen, daß er nicht in die Gesellschaft der Heuchler käme, die Gott anrufen, um gesehen und gehört zu werden. Die Unwissenden und Kinder, die das kirchliche Leben der Gemeinde mitmachen, lernen gewiß recht viel; aber unversehens, wie man eben aus allem Leben
\quellseite{106}
lernt; ganz etwas anderes ist es aber, wenn man fragt, ob die Liturgie unter anderem auch lehrhaft sei, oder ob man fragt, ob die Kirche bei der Liturgie die Absicht habe, zu belehren, und die Liturgie eine Form der Lehre sei. So wie man nicht lehrt, um ein Beispiel zu geben, obgleich man es nicht leugnen kann, daß alles Leben nothwendig zum Beispiel werden muß; so hält man auch nirgends Liturgie, um zu lehren, so viel Licht und Lehre auch von der Liturgie ausgehen mag. Die Absicht der Liturgie ist entweder die, im Namen Gottes mit der Gemeinde, oder im Namen der Gemeinde mit Gott zu handeln. Soferne die Lehre ein Handeln ist zwischen Gott oder seinen Stellvertretern mit den Menschenkindern, kann sie freilich auch einen Platz im liturgischen Leben haben, und hat ihn auch; aber Lehre ist weder die einzige, noch die vornehmste Handlungsweise Gottes und seiner Stellvertreter unter den Menschen, ganz abgesehen davon, daß die gesammte, sacrificielle Abtheilung der Liturgie ihrer Natur nach mit der Lehre gar nichts zu schaffen haben kann. So gewiß wir es daher im Folgenden mit der Einführung der Lehre und Lehrgabe in die amtliche Wirksamkeit zu thun haben, so gewiß haben wir es nicht mit der Liturgie zu thun, wohl aber mit Homiletischem und Katechetischem.

\quellseite{107}
\section[4. Formale Vorbildung]{4. Formale Vorbildung für das Amt der Lehrer}

Die Lehre, die Lehrgabe, die Form der Lehre haben wir berührt, — aber wir dürfen nicht vergessen, daß es fürs Amt der Lehre und des Wortes auch vorbereitende Studien gibt. Manch großer Amtsträger hat freilich von diesen Studien gar keine Erinnerung, er hat sie nicht gemacht. Die große Gabe und die Macht des Lebens hat ihn mitten in die Praxis hineingeführt, ohne großes Lehrgeld und ohne Schaden. Allein für uns gewöhnliche Leute und unsre gewöhnlichen Kinder sind die Vorbereitungsstudien, die wir jetzt meinen, gar nicht zu verachten, nicht zu vernachlässigen. — Als ich jung war, hatte ich das Glück, im Nürnberger Lyceum, welches sich an den Schluß der Gymnasialzeit anlehnte, Logik, Dialektik und Rhetorik zu hören, die letzteren jeden Falls im Sinne und in der Weise der früheren Zeit. Es war mir, wie wenn erst sie meiner empfangenen formalen Bildung Werth, Klarheit und Abschluß gäben. Später hörten die Lyceen wieder auf, und andere Gymnasiasten mußten ohne diesen Abschluß auf die Universität gehen, wo jene formalen Studien, nach Weise der neuen Zeit, für formale Bildung gerade nichts auszutragen pflegten. Es hat mir geschienen, wie wenn den schriftlichen Arbeiten und Predigten vieler Candidaten, so viel classische Muster sie gehört und gelesen haben mochten, doch nichts mehr als formale Bildung fehlte. Desto mehr freute mich die Be%
\quellseite{108}
kanntschaft, die ich, erst als Mann und Lehrer (s. unten), wie zufällig mit Melanchthon’s Lehrbüchern machte. Meine Lycealzeit lebte wieder auf. Als ich später nicht mehr lehrte, sondern an meine Stelle ein für dergleichen Dinge fähigerer und gründlicherer Mann trat, fand auch er an den genannten Büchern nicht weniger Geschmack als ich, und ihm (ich meine den Inspector der hiesigen Missionsschule, Herrn Fr. Bauer) wird es vielleicht noch gegeben werden, die edlen Schriften für unsre studierende Jugend aufs neue zugänglicher und recht genießbar zu machen. Mögen einstweilen die nachfolgenden Capitel der vorigen Auflage dieses Buches dazu dienen, einige Lust und Begier zu erwecken. — An diesen Studien kann der Jüngling von eingehendem Fleiß eine treue Hand finden, die seine Gabe in die Praxis zu leiten vermag, daß sie zur Fertigkeit, zur Geschicklichkeit, zur Kunst und Tugend werde. Möchte recht bald die Handleitung durch Veröffentlichung der Schriften Melanchthon’s nach Bedürfniß unsrer Tage möglich und dann reichlich angenommen und gebraucht werden!

\section{5. Dialektik}

Alles menschliche Denken bewegt sich innerhalb Begriff und Urtheil. Definieren, dividieren und schließen, — vergleichen und unterscheiden — Falsches vom Rechten, das ist’s, was alle, die denken und reden, immerfort thun und thun müssen, auch wenn sie es selbst gar nicht bemerken sollten. Es ist allen Menschen von gesunden Sinnen von Gott, dem Herrn, als
\quellseite{109}
Mitgabe in diese Welt geschenkt. Was aber an sich Gabe ist, wird zur Fertigkeit, Kunst und Tugend durch Ausbildung. Diese Ausbildung aber gibt dem willigen Schüler für seine innere Verstandesthätigkeit die Dialektik. Est enim dialectica ars sive via, recte, ordine et perspicue docendi, quod fit recte definiendo, dividendo, argumenta vera connectendo et male cohaerentia sive falsa retexendo et refutando. (Melanchthon.) Es ist weder eines jeden Menschen Sache die philosophischen Lehrbücher der gegenwärtigen Zeit zu studieren, noch fällt es dem Schreiber dieses ein, das Studium derselben ohne weiteres allen zu empfehlen, zumal er sich selbst kaum jemals mit demselben abzugeben vermochte. Die meisten jungen Geistlichen werden am Ende wie er Ursache haben, sich nur ganz mittelmäßige oder geringe Gaben für die philosophischen Wissenschaften zuzuschreiben; das Studium derselben führt schon deshalb, wie Claudius meint, die meisten vom Ziele ab, nach dem sie streben, oder raubt ihnen doch viel Zeit und Kraft, ohne daß sie dadurch tüchtiger für einen Beruf oder gar für das heilige Amt werden; sie bleiben in der Mühsal hängen, ohne daß sich die Mühe lohnte. Dagegen aber gibt es aus früheren Zeiten noch einfache, philosophischen Systemen weniger unterworfene Lehrbücher der formalen Wissenschaften, welche auch ein für Abstraction nur mittelmäßig und gering begabter Kopf, leicht verstehen und mit großem Nutzen zur Ausbildung der vorhandenen natürlichen Gabe gebrauchen kann. Hieher gehören eben die schon erwähnten Lehrbücher Melanchthon’s, aus denen man in der That erst recht begrei%
\quellseite{110}
fen lernt, wie er zu dem hohen Ehrentitel eines praeceptor Germaniae gekommen ist. Von ihm besitzen wir besonders eine lateinische Schrift unter dem Titel: „Erotemata Dialectices continentia integram artem, ita scripta, ut juventuti utiliter proponi possint\footnote{z. B. Witebergae. Anno 1581.}.“ Neben einer edlen Simplicität findet sich in diesem Buche eine wunderbare Vereinigung sogenannter classischer und theologischer Gelehrsamkeit; Christus herrscht in dieser Dialektik, obwohl er nicht auf dem Lehrstuhl sitzt, und auch Magister Philipp nicht unmittelbar von ihm redet. Dazu ist die Wahl der Beispiele eine überaus schöne und treffende; und die Sprache des guten Lehrers erleichtert das Lernen auch dem, der wenig Fertigkeit erlangt oder übrig hat, lateinische Bücher zu lesen. Der Verfasser dieser Blätter hat erst als Lehrer von Missionszöglingen, die keine Gymnasialbildung empfangen hatten, Melanchthon’s ebengenannte Schrift genau kennen gelernt. Missionsschüler von der eben bezeichneten Art bedürfen einer formalen Bildung, ohne daß doch der Zweck, zu dem sie lernen, das Alter, in dem sie gewöhnlich stehen, und die Schulbildung, welche sie mitbringen, es möglich oder räthlich macht, sie die Gymnasialstudien nachholen zu lassen. Da fand ich denn Melanchthon’s Erotemata sehr geschickt, das dringende Bedürfniß zu stillen, und es hat mich zuweilen bedünken wollen, als könnte ein fleißiges und eingehendes Studium derselben in manchem Betracht mehr formale Bildung geben, als das Lesen der Classiker in den Schu%
\quellseite{111}
len ohne das nebenhergehende Studium einer einfachen Dialektik den meisten Jünglingen gibt. Es sei daher hier auch gewagt, jungen Geistlichen, die den Mangel ihrer Erziehung inne werden und es an sich selber spüren, daß ihnen das Geschick abgeht, recte, ordine et perspicue docendi, das Studium des obigen Buches oder eines ähnlichen anzurathen. Wer den Rath befolgt und hernach den Rathgeber schilt, soll Recht haben. So wenig fürchtet sich aber der Rathgeber noch zur Stunde, daß er sogar noch mehr wagt, nemlich den bereits im Amte Stehenden anzurathen, daß sie die Erotemata oder etwas dieser Art studieren. Er verspricht ihnen, daß sie nach gemachtem Studium die Sehnen ihrer Gabe gestählt, und das Geschick, die Gabe anzuwenden, gemehrt und gestärkt fühlen werden.

\section{6. Rhetorik}

Wenn die Dialektik richtig reden lehrt, so ist es das Amt der Rhetorik, zweckmäßig reden zu lehren. Schon hieraus kann man vermuthen, wie nöthig und nützlich es sein wird, daß der, welcher seine vorhandene Lehrgabe ausbilden will, die Rhetorik studiere. Ist die Dialektik eine Vorbereitung für alles Lehren, so ist die Rhetorik insonderheit eine Vorbereitung für diejenige Art zu lehren und zu reden, welche Gegenstand der Homiletik ist. Der heilige Geist bedient sich bei allen seinen Offenbarungen einer Sprache, die, gleich der Sprache der menschlichen Vernunft, den natürlichen Gesetzen des Denkens, also den Regeln der Dialektik entspricht, — und nicht
\quellseite{112}
minder sind die Vorträge der heiligen Propheten und Apostel so eingerichtet, daß sie den allgemeinen Regeln der Rhetorik nicht widerstreiten. Und ob auch die Weise der heiligen Redner Gottes von der anderer, bloß menschlicher Redner sehr verschieden ist, so wird doch auch diese specifische Verschiedenheit dem am klarsten, welcher die Erzeugnisse himmlischer Redegabe mit den Regeln der nur menschlichen Redekunst vergleichen und das Gemeinsame finden kann. Der heilige Geist reinigt und verklärt die natürlichen Gaben, aber er tilgt sie nicht aus. Neugeburt, nicht neue Schöpfung ist sein Werk im Neuen Testamente. So nimmt er auch die menschliche Rede so wie sie ist, tilgt ihre Gesetze, die Er selbst gestiftet, nicht aus, sondern lehrt sie untadelig und des himmlischen Lichtes würdig gebrauchen. Deshalb schäme sich keiner, die Rhetorik der Alten zu studieren, zumal wenn sie in einem so heiligen Sinne und mit solcher Lehrergabe vorgelegt wird, wie es der Praeceptor Germaniae Ph. Melanchthon thut.\footnote{Daß deshalb nicht alles, was man bei Melanchthon findet, einfach angenommen werden kann, versteht sich von selbst. Auch seine Bücher sind menschlich. Wer dürfte z. B. seine Definitionen (S. Definitiones multarum appellationum, quarum in ecclesia usus est, traditae a Philippo Melanth. Torgae et Witebergae. anno 1552 et 1553.) nachschreiben? Es geht hier wie bei der Casuistik der nachreformatorischen Lehrer. Wie empfehlenswerth ist die Casuistik! Aber darf man sich denn wirklich auf alle alten casuistischen Entscheidungen verlassen? Nur der Unkundige könnte ja antworten.} S. dessen Schrift: Elementorum Rhetorices libri duo, recens re%
\quellseite{113}
cogniti ab auctore Philippo Melanthone. His adjectae sunt Epistolae contrariae Pici et Hermolai Barbari, una cum dispositione Philip. Melanth. Accessit demum Index et rerum et verborum locupletissimus. Cum gratia et privilegio. Witebergae. Clemens Schleich excudebat. 1582.

\section{7. Die Reden der Alten}

Die Alten, nemlich die Griechen und Römer, hielten keine Reden zum Zweck religiöser Unterweisung, oder besser: ihre Rhetorik gibt uns keinen Einblick in die religiöse Unterweisung, welche bei ihnen etwa gebräuchlich war. Sie redeten aber vor der Volksversammlung und vor den Richtern. Ihre Reden waren Staatsreden, gerichtliche Reden, auch wohl Lobreden auf Helden, Staatsmänner u. dgl. Die Rhetorik hat deswegen bei ihnen tria genera causarum: demonstrativum, deliberativum, judiciale, (das darlegende, erwägende, gerichtliche). Melanchthon setzt ein genus didascalicum (ein lehrhaftes) hinzu, mehr weil es seinem Zwecke gemäß war, als weil er es bei den Alten vertreten fand\footnote{„Vulgo tria numerant genera causarum: Demonstrativum, quo continetur laus et vituperatio. Deliberativum, quod versatur in suadendo et dissuadendo. Judiciale, quod tractat controversias forenses. Ego addendum censeo διδασκαλικόν genus, quod, etsi ad Dialecticam pertinet, tamen, ubi negociorum genera censentur, non est praetermittendum, praesertim cum hoc tempore vel maximum usum in ecclesiis habeat, ubi non tantum suasoriae conciones habendae sunt, sed multo saepius homines dialecticorum more de dogmatibus religionis docendi sunt, ut ea perfecte cognoscere possint.“ p. 15 f. — „Demonstrativum genus continet laudationem et vituperationem, sed multum interest, utrum personae tractentur, an facta, aut res. Cum laudamus personam, ordine narramus historiam et rerum seriem in dicendo sequimur. Sunt igitur loci personarum: Patria, Sexus, Natales, Ingenium, Educatio, Disciplina, Doctrina, Res gestae, Praemia rerum gestarum, vitae Exitus, Opinio post mortem.“ p. 61. — „Cum factum aut rem certam laudamus, Locos mutuamus ex genere deliberativo: honestum, utile, facile etc.“ p. 63. 17. p. 57. „Deliberativum versatur in suadendo et dissuadendo, adhortando et dehortando, petendo et precando, consolando et similibus negotiis, ubi finis est non cognitio, sed praeter cognitionem actio aliqua.“ p. 56. — Judiciale. p. 28. f. S. Anm. * zu § 8. p. 106.}. — So manches uns nun beim
\quellseite{114}
Lesen der rhetorischen Unterweisungen der Alten fremd bleiben wird, namentlich für unsern Zweck; so werden wir doch auch sehr viel für uns abstrahieren können. Es werden aber nicht bloß die genera causarum, christlich umgewandelt, bleiben, sondern auch die Eintheilung der Rede. Man denke sich eine christliche Rede — wir sprechen aber zunächst nicht von der Predigt oder Homilie, sondern von der eigentlichen Rede: welche Theile wird man ihr zuweisen müssen? Es wird eben doch bei exordium, narratio, propositio, confirmatio, confutatio, peroratio bleiben, es gibt keine andern Theile der Rede, — und alles, was man bei Melanchthon I, c. von p. 32\footnote{Tria in exordio efficienda sunt, ut reddamus auditores benevolos, attentos et dociles. p. 32. Narratio est facti expositio, quae statim exordio subjicitur. Hanc sequi debet propositio, quae contineat summam rei. Propositio nunquam omitti potest. p. 35. — Confirmatio difficillima pars est oratorii operis, haec enim persuadere judici debet, ut propositioni assentiatur. p. 35. — Confutatio est dissolutio argumentorum, quae objiciuntur. p. 38. — Peroratio est conclusio orationis, in qua repetitur propositio principalis. Partim enim repetitione constat propositionis et potissimorum argumentorum, partim affectibus. p. 40.} an über diese einzelnen Theile der Rede lernt, wird
\quellseite{115}
fruchtbringend sein auch für die geistliche Rede. Hat sich doch, wie wir bald sehen werden, die Predigt der späteren christlichen Zeit bis auf den heutigen Tag eng an die Rede der Alten und ihre Einrichtung angeschlossen und anschließen müssen. Abermals ein Beweis, wie gut und nöthig es nun umgekehrt sein wird, ehe man zur Homiletik schreitet, die Rhetorik zu studieren. Wer Melanchthons Dialektik und Rhetorik kennt und mit dem Inhalt vertraut ist, wird sehr gerüstet zur Homiletik kommen, und wer Beispiele alter Reden, wie sie Melanchthon zum Theil im Anhang seiner Rhetorik ausführlich gibt, zum Theil citiert, mit dem Gelernten vergleicht, wird zum besonderen Studium der christlichen Rede ein reiches Maß ersprießlicher Vorbildung mitbringen.

\section{8. Die erste Predigtweise}

Die heiligen Apostel predigten, d. i. verkündigten der Welt unter Beweisung des Geistes und der Kraft die große Thatsache der Erlösung durch unsern Herrn Jesum Christum. Dazu waren sie von dem Herrn gesandt; schon
\quellseite{116}
die erste Predigt, die sie irgend wo halten konnten, mußte deshalb an den Lebenslauf oder die Geschichte Jesu Christi oder des Reiches Gottes anknüpfen. Evangelium und Geschichte wurde darum innigst verwandt und fast identisch, wie schon aus dem Vortrag der vier Evangelisten ersichtlich ist. Wer nun die apostolische Predigt — oder eines von den vier Evangelien — unter ein altes genus causarum subsumieren wollte, der würde sie ohne Zweifel zum genus demonstrativum rechnen müssen. Und würde man nun p. 61 bis 65 in Melanchthon’s Rhetorik nachsehen und bedenken, so würde man auch bald erkennen, wie einfach so apostolische Predigt wie Evangelium nach den dort angegebenen Regeln geordnet ist. Nicht, daß die heiligen Schreiber oder Redner Regeln vor Augen gehabt hätten, sondern daß sich bei der Vorlegung eines Lebenslaufes einer Geschichte und großer Thatsachen nicht anders verfahren läßt, als jene Regeln sagen. Die Regeln sind eben nicht willkürlich, sondern der Natur und Erfahrung entnommen. Ihnen folgt auch der ungeschulte Redner, wenn auch unbewußt.

So sehr nun aber die Evangelien und die ersten Predigten der Apostel den Charakter des genus demonstrativum an sich tragen müssen, so offenbar ist es doch, daß auch Apostel durch die Umstände gedrungen werden konnten, zu einem andern genus abzubiegen, sei es zum deliberativum oder judiciale. Man prüfe z. B. die erste Pfingstpredigt Petri, die Rede des heiligen Stephanus, die Pauli zu Athen \&c. So wie nun das geschieht, wird man auch im Ganzen dieselbigen
\quellseite{117}
Grundsätze beachtet finden, welche bei Melanchthon p. 26 ff.\footnote{Disputationes ecclesiasticae magna ex parte similitudinem quandam habent forensium certaminum. Interpretantur enim leges, dissolvunt ἀντινομίας, videlicet sententias, quae in speciem pugnare videntur; explicant ambigua; interdum de jure, interdum de facto disputant; quaerunt factorum consilia. Ideo hoc genus in his nostris moribus etiam magnum habet usum. p. 26.} vom judiciale oder p. 56 ff. vom deliberativum zu lesen sind, natürlich ohne künstliche Methode und Manier, so wie es Aposteln und den lebendigen Verhältnissen geziemte, aus deren Boden die apostolischen Reden entsprossen sind.

Wer da will, der studiere die heilige Rhetorik und Dialektik der großen Apostel aus ihren Reden und vergleiche dieselben mit dem, was er aus Melanchthon lernte. Es wird die alte Rede verklärt wieder finden, nur daß eben, was der Geist derselben eingeschaffen hat, desto unvertilgbarer und ehrwürdiger erscheint, wenn es im Glanze der edelsten Brauchbarkeit für die Kirche Gottes daher schreitet.

\section{9. Predigt und Text}

Die Predigt oder Verkündigung des Evangeliums ist ihrer Natur nach ohne Text. Was fragt der Heide nach dem Text eines Buches, von dem er nichts weiß oder hält? Denn sonst wäre er kein Heide. Ihm kommt die Predigt ohne Anschluß an frühere Offenbarungen. Sie wirkt auf ihn mit ihrer
\quellseite{118}
angestammten Kraft, welche sich nicht verleugnen kann, auch wo das Wort einfach, beweislos, in der Gestalt einer göttlichen Botschaft und Offenbarung oder eines prophetischen, apostolischen Zeugnisses kommt.

Die Heiden-Missionspredigt ist eine textlose. Auch die Predigt an die ungläubigen Juden bedarf keines Textes. Wenn gleich die großen Thatsachen, welche ihnen verkündigt werden, Erfüllungen alttestamentlicher Weißagungen sind, so nehmen doch die heiligen Apostel insgemein nicht einen Text aus dem Alten Testamente\footnote{Vergl. jedoch die Rede Jesu zu Nazareth Luc. 4, 16 ff.}, sondern sie gebrauchen die alttestamentlichen Texte nur zu Beweisen für ihr Evangelium, da es auch an die Juden in Gestalt und Kraft einer göttlichen Offenbarung kam. Auch die Juden-Missionspredigt ist daher meist textlos. Einzelne Texte treten überhaupt als Grundlage der Predigt erst dann ein, wenn die Predigt selbst nicht mehr bloße Berufung, nicht mehr bloß grundlegend ist, sondern wenn es gilt, die bereits gewordene Gemeinde tiefer in das heilende und heiligende Wasser göttlicher, im Glauben angenommener Rede einzuführen. Ja die christliche Rede selbst kann ohne Text gedacht werden, eben weil sie bei ihrer Richtung auf bereits gewordene Gemeinden eigentlich nicht mehr Predigt ist, sondern entweder mehr die Natur der Schriftauslegung oder die Gestalt der alten Rede annimmt. Es kann also Predigten mit und ohne Text geben und zwar die letzteren nicht bloß bei den Missionaren, sondern auch innerhalb
\quellseite{119}
der Gemeinde. (Z. B. die 50 Homilien des Macarius, welche Ermahnungen an Asceten enthalten.)

\section{10. Die Homilie}

Von den Predigten der Apostel und Missionare verschieden ist die sogenannte Homilie, welche sich eng an die Perikope anschließt. Die Kirche las in ihren Versammlungen von Anfang an ausgewählte Texte aus dem Worte Gottes, an welche sich die freie Rede des Bischofs oder Presbyters anschloß. Der Anschluß konnte aber ein sehr verschiedener sein, und das Bedürfniß der Gemeinde wurde natürlicher Weise ein Schlüssel der Invention für die verschiedenartigsten Ansprachen. Je länger je mehr schmiegte sich jedoch die freie Rede dem Texte an, wurde eine Emanation des göttlichen Wortes und eine menschliche Parallele göttlicher Gedanken.\footnote{Siehe Beispiele dieser Art in D. Haymonis, episcopi Halberstatensis († 853), Homiliae in evangelia dominicalia totius anni et de sanctis praecipuis, nunc recens excusae et a multis erroribus restitutae. Antverpiae, in aedibus Joannis Steelsii. Anno 1559. Cum privilegio. Z. B. die Homilie auf Quinquagesima. „Haymo, Angelsachse, Beda’s Blutsfreund, Herzensfreund des Rhabanus Maurus von Kindesbeinen an, sein Cellengenoße in Fulda, Abt zu Corvei und Fürstfeld, Bischof von Halberstadt 13 Jahre.“} Der Zweck wurde Texterkenntniß, Textanwendung, und diesem Zwecke gemäß war es auch, wenn die Form der alten Rede, gleich viel ob sie der Redner kannte oder nicht, vor der Einfalt des Beginnens zurücktrat. Da ward kein exordium,
\quellseite{120}
keine narratio, propositio etc. Der Text herrschte und die einfältige Vertiefung ins Wort Gottes genügte der im Worte lebenden Gemeinde vollkommen. Das gab dann Homilien, d. i. einfache Unterhaltungen über das Wort, einfache Speisung mit demselben.

Da wir noch gegenwärtig wie in der ersten Zeit das göttliche Wort in der Gemeinde lesen und Auslegung des Wortes das große Hauptgeschäft eines jeden Hirten jetzt wie damals bleibt; so kann es auch bei uns wie in den ersten Zeiten Homilien der einfachsten Art geben, bei denen sich das besondere Bedürfniß der Gemeinde dem Texte völlig unterordnet, nur gelegentlich seine Befriedigung findet, und Texterkenntniß Hauptabsicht ist.

\section{11. Ausbildung der Homilie}

Dr. Andreas Adam Hochstetter, weiland Professor in Tübingen, gab im Jahre 1701 eine kleine Homiletik heraus, deren Titel ist: De recta concionandi textumque sacrum tum exponendi tum applicandi ratione commentariolus. Die kleine treffliche Schrift\footnote{Im Anhang zu diesen Blättern zu Nutz und Frommen derer mitgetheilt, welche die rechte Weise zu predigen sich aneignen möchten.} wurde vergessen, aber von andern wieder hervorgezogen und 1767 zum drittenmal von dem Tübinger Professor Christ. Friedrich Sartorius herausgegeben. In diesem lateinisch geschriebenen, zur Einleitung in die Homiletik
\quellseite{121}
wohl dienenden Büchlein wird vor allem auf Texterkenntnis gedrungen. § 4 heißt es: Textus sacri natura ipsa ex his quinque dignoscitur: Investigandum primum objectum, progrediendumque deinceps ad genus et speciem, inquirendum in argumenta, ipsumque denique thema ex istis judicandum. p. 6. Die angegebenen einzelnen Theile werden in den weiteren §§ 5—18 abgehandelt. Beim Lesen kann man oft auf die Gedanken kommen, der Autor rede nicht vom Text, sondern von der Predigt, so ganz ist jeder Text selbst als ein Ganzes, als eine Rede mit Haupt und Gliedern behandelt. Obwohl nun diese Auffassung nicht auf alle Texte passen wird, so liegt doch etwas völlig Richtiges zu Grunde. Auch der heilige Geist hat nach den Regeln der von Ihm geschaffenen Grammatik, Dialektik, Rhetorik geredet: — die Kirche hat nur Stücke, die ein rundes Ganzes bilden, zu Lectionen ausgewählt: man wird darum diese Lectionen am richtigsten auffassen, wenn man die Gliederung, den Zusammenhang, die Einrichtung aufsucht, — wenn man wie bei einer menschlichen Rede nach Thema, Partition, Confirmation, Confutation, etwa gar nach Exordium und einer Peroratio forscht und der Gemeinde das wohl erkannte Ganze in der wohl erkannten Gliederung zum Verständnis vorlegt. Dadurch wird sich ebensowohl Klarheit als Deutlichkeit, Uebersicht und Einsicht mehren, — und wenn der Homilet so sein Werk gethan hat, wird die nach dem Wort hungrige Seele satt und froh sein. Der Totaleindruck des Ganzen wird durch die Kenntniß der schönen Gliederung verstärkt sein, so wie man
\quellseite{122}
ein Bild desto mehr liebt und desto höher schätzt, wenn der Künstler die einzelnen Theile bedeutend und erklärend durchgesprochen hat.

Auch das ist Homilie. Je klarer und deutlicher der Homilet seinen Text erkannt hat, desto einfältiger wird er ihn vorlegen können, und die Kunst wird deshalb dem himmlischen Sinne des Ganzen trefflich dienen und zu Statten kommen. Je vollendeter ein solcher heiliger Exeget in seinem Dienste ist, in desto ungesuchterer Selbstverläugnung wird er verfahren. — Es wird am Ende scheinen, als hätte er nichts gethan, da er doch so viel geleistet hat, und ihm es zu danken ist, daß der Text in seinem Glanze sonnenklar und kräftig in die Seele dringt.

Hierbei kann sich als ein zweiter Theil eine Application anschließen, in welcher man unter Beibehaltung der eigenen Disposition des Textes das Licht des göttlichen Wortes auf die Verhältnisse treffen und fallen läßt und also Lehre, Besserung, Strafe und Züchtigung wirkt. Wer diese Weise versteht, hat eine große und heilige Gabe, welcher kaum eine andere Predigergabe gleichzustellen sein dürfte.\footnote{Gegenüber dieser Hochstetter’schen Anweisung zu predigen, welche so viel Kunst wie Einfalt zuläßt, steht warnend der ganze Wust der protestantischen Predigtweisen. Balduin zählt 7, Rebhahn 25, Carpzov 100 Methoden, Löscher reducirt die Zahl auf 25 etc. etc. Da gibt es eine jenaische, helmstädtische, eine gekünstelte, geschnörkelte Leipziger Methode mit 3 Exordien, — eine niederländische, englische, französische Methode, eine methodus Loescheriana, Balduina, ja eine methodus classica (von der Einführung der Anwendung durch das Wörtchen „daß“ — daher. dassica) \&c. \&c. Man könnte denken, aller Ernst der Predigt sei vor Kleinigkeiten verschwunden, — und doch ist’s nicht so. So viel kann man dem geduldigen Geist und seinem heiligen Wort und Amte zumuthen!}

\quellseite{123}
\section{12. Die christliche Rede}

Die Homilie in ihren beiden Formen ist das Gegentheil der christlichen Rede, des textlosen Vortrags über heilige Gegenstände in Mitten der Gemeinde. Bei der Homilie hat der Redner keinen eigenen Gedanken, keine andere Absicht als die, seinem Texte zu dienen, ihn als ein Licht auf den Schäffel zu stellen. Für den Homileten paßt so das Wort des Täufers: „Christus muß wachsen, ich muß abnehmen.“ Ganz anders ist es bei der Rede. Der Redner hat sich unabhängig von einem besondern Texte, höchstens zufälliger Weise von einem Texte ausgehend oder angeregt, ein christliches Thema ausgewählt, worüber er nun, aus was für Anlaß oder Nothdurft es geschehe, sprechen will. Sein Thema ist entweder eine quaestio simplex oder conjuncta. Im ersteren Fall geht er vom Thema in die Theile durch Definieren und Dividieren. Der alte Satz: Quis, quid, ubi, quibus auxiliis, cur, quomodo, quando? wird die Partition regeln, und was Melanchthon vom genus didascalicum causarum p. 20—25 sagt\footnote{methodus perfacilis est, si quis mediocrem exercitationem adhibuerit. p. 20.}, wird völlig anwendbar sein. Im zweiten Falle wird die conclusio irgend eines Syllogismus zum
\quellseite{124}
Hauptsatz werden; Major und Minor werden in der Confutatio wie in der Confirmatio regieren; am Ende wird alles auf’s Thema zurückgehen, und das Thema wird die ganze Rede regieren. — Ein christlicher Redner wird nur christliche Themata abhandeln. Alle seine Definitionen, Divisionen, Syllogismen, Enthymemen etc. werden mit der heiligen Schrift zusammenstimmen. Wort und Geist der Schrift wird den gesammten Inhalt seiner Rede durchdringen, aber die Form wird rein die Form der alten Rede sein, die hier in verklärter Schönheit und als eine heilige Magd unseres Herrn Jesu Christi erscheinen wird. Nicht im Dienste eines Textes, aber im Dienste des Wortes und der Wahrheit selbst steht diese heilige Magd, am Ende doch wohl eine ebenbürtige Schwester der Homilie. Sie gleicht dem bekehrten Abendländer, dem Griechen oder Römer, der unter Beibehaltung seiner Lebensformen dem einigen Christus aller Gläubigen dient, während die Homilie dem heiligen Morgenländer gleicht, der nichts eigenes beizubehalten scheint, keinen Fleiß anwendet, seine Individualität zu wahren, ganz in das Betrachten des Göttlichen aufgeht und eben damit, ohne es zu wissen, im schönen Glanze seiner Individualität erscheint. Dort ist das Menschliche durch Göttliches verklärt: hier ist das Göttliche zum Menschlichen herabgestiegen und hat es allzumal bewältigt.

\quellseite{125}
\section[13. - Vereinigung der Homilie und der Rede]{13. - Vereinigung der Homilie und der Rede}

Die Homilie und die Rede, so verschieden an Form, haben nichts desto weniger eine Vereinigung gefunden. Wenn man nemlich die besonders bei unseren lutherischen Vätern gebräuchliche Form und Gliederung der Predigt betrachtet; so findet man sie mit der Form und Gliederung der alten Rede völlig übereinstimmend, nur daß an die Stelle der alten narratio die enarratio getreten ist, d. i. Homilie, die Text-Aus- und Darlegung. Mag man nun das Maß der einzelnen Theile der Rede verändert, daß der enarratio im Vergleich zur alten narratio verlängert, anderes gekürzt oder zusammengefaßt haben: die Gliederung ist doch eine und dieselbe. Da ist ein Exordium, welches zur Textauslegung (enarratio) überleitet, von dieser ein Uebergang zur propositio, welche in Form der confirmatio und confutatio ausgeführt wird; da ist eine peroratio, welche nicht selten mit starker Kraft auf die Affekte wirkt; kurz was eigentlich zu einer Rede vom alten Styl fehlen soll, ist nicht zu ersehen. Auch scheint die Natur der Rede durch den engen Anschluß an einen Text und die enarratio gar nicht wesentlich verändert.

So eigen diese Form der neueren Christenpredigt zu sein scheint, so ist sie doch gründlicher, in Schrift und Wahrheit gründender als jene andere neuere Form, bei welcher die Predigt eigentlich nichts als eine Rede ist, nur unter Vorausstellung eines als Motto figurierenden Textes, von welchem unbeirrt man dahin geht, es wäre denn, daß man im Verlauf
\quellseite{126}
sich hie und da einmal auf ihn, wie auch auf andere nicht textmäßig vorausgestellte Bibelstellen, besänne und bezöge. Es muß auch keineswegs durch die oben angegebene homiletisch-oratorische Einrichtung schwerfällig werden, zumal wenn die enarratio ein erster Haupttheil ist, Alles darauffolgende aber in bündiger Zusammenfassung, wie ein zusammengehöriger zweiter Theil, wie eine Applicatio erscheint.

Die homiletische Weise zu predigen nennt man auch die analytische, die andere oratorische dagegen die synthetische. Herrscht die klar erkannte und vorgelegte Disposition des Textes in der Predigt, so daß diese ein gegliederter, zu einem schönen Ganzen vereinigter Wiederhall des Textes ist, so ist das die analytisch-synthetische Weise zu predigen; ist hingegen der Inhalt des Textes unter einen besondern Gesichtspunkt, unter den eines erwählten Themas gestellt, und einer selbstgemachten Partition an= und eingeschmiegt, so redet man von synthetisch-analytischer Predigtweise. Es kann je nach Zeit und Umständen eine jede stattfinden, bald diese, bald jene den Vorzug verdienen. — Bei dem fortlaufenden Lesen und Erklären der heiligen Schrift, wie es in den sogenannten Bibelstunden zu sein pflegt, wo man das Lesen für Hauptsache erkennt, der Dienst des Homileten nur so viel nöthig der Schwachheit des Hörers nachhilft, wird die einfachste Art der homiletischen Schriftauslegung, die Auslegung der heiligen Schrift um ihrer selbst willen, an Ort und Stelle sein. Bei den Lectionen des Hauptgottesdienstes wird die Homilie der zweiten Art, d. i. die analytisch%
\quellseite{127}
synthetische oder auch die synthetisch-analytische Weise die geziemendste sein. Bei den Casualpredigten, also bei Tauf-, Hochzeit-, Leichenpredigten und anderen dieser Art wird sich die synthetische Weise aufdringen.\footnote{Überhaupt muthet die synthetische Predigtweise dem hörenden Publicum am wenigsten zu und die lutherischen Kirchenordnungen, welche sie besonders empfehlen, haben bei dem geistigen Standpunkt unserer meisten Gemeinden, auch der gegenwärtigen, ganz recht. Das wörtlich ausgeprägte Thema und die namhaft vorgelegten Theile unterstützen und erleichtern die Aufmerksamkeit, sowie auch der Prediger durch diese Predigtweise gezwungen ist, einfach, klar und deutlich zu disponieren. Es muß deshalb die Predigt doch nicht immer ein opus tripartitum sein.} Die reine textlose Rede aber wird in vielen Casualfällen (Tauf-, Trau-, Leichen-, Gedächtnisreden, Danksagungsreden, Anreden an verehrte Personen) allezeit Ruhm und Ehre behalten.

\section{14. Wahl der Predigtweise}

Die Predigtweise hängt theils, wie schon gesagt, von den Umständen, theils aber auch von der Begabung ab. Wer aber für jede Weise begabt ist, sich in jeder ausbilden will und kann, oder seine Gabe zu erforschen Versuche anstellt, dem könnte man vielleicht für seine Fest- und Sonntagspredigten den folgenden Rathschlag geben. Den ersten Jahrgang deiner Predigten laß einfach homiletisch sein, lege bei dir und deinen Hörern den Grund deiner amtlichen Redethätigkeit damit, daß du ganz am Texte bleibst und seine Einzelheiten klar machst. Befleißige dich der Deutlichkeit. Im ersten Jahre kennst du die Gemeinde, und sie dich noch nicht. Da ist es nicht
\quellseite{128}
bloß natürlich, sondern auch gut, wenn du weniger auf die Gemeinde, als auf den Text siehst und ihn auslegst, als hörte ihn die Gemeinde zum erstenmal, als wäre sein Inhalt eitel Neuigkeit. — Im zweiten Jahrgang befleißige dich, deinen Text nicht bloß deutlich, sondern klar zu machen. Da behandle ihn als ein wohlangelegtes gegliedertes Ganze und zeige seine Theile, aber im Zusammenhang, seinen Hauptgedanken und wie derselbe in allen Theilen lebt. Damit gibst du der Gemeinde den Text zu eigen und zwar unter dem Eindruck, welchen die Form, die Zusammenfassung zum Ganzen, die Unterordnung aller Einzelheiten unter den Hauptgedanken gibt. Hiermit wird deine Predigt, auch wenn sie keine eigene Applicatio hätte, von selbst praktischer und anmaßender. — Im dritten Jahrgang stelle deine Texte unter besondere Gesichtspunkte: zeige das Ganze und seine Theile von der dogmatischen, von der ethischen Seite oder nimm διδασκαλία, ἔλεγχος, ἐπανόρθωσις, παιδεία, je nach Art und Inhalt des Textes, heraus \&c.; dabei kannst du immer das Ganze und seine Theile vollständig berücksichtigen, und auch die indeß gemachte Erfahrung, die Kenntniß der Gemeinde, welche du gewonnen hast, bekommt Gelegenheit, sich geltend zu machen und in Wirksamkeit zu treten. — Nach diesem dritten Jahrescurs wirst du in so vielen Cursen, als du willst — denn du wirst nicht in Verlegenheit kommen — einzelne Gedanken und Themata aus dem Texte nehmen können, nach Bedarf der Gemeinde, und das seelsorgerliche, eigentlich praktische Predigen wird in die Macht kommen und an den vorhergegangenen Jahrgängen gute Vorbereitung haben.
\quellseite{129}
— Endlich wirst du zum Anfang zurückkehren und echt praktisch ganze Texte behandeln können, auch analytisch-synthetische Predigten vom höheren Chor halten lernen. Denn das ist doch gewiß, daß keine Kunst höher und herrlicher ist, als wenn man in ganzen Texten die eigenen Zustände wie im Spiegelbilde und im Lichte eines reichen Textes die Gegenwart in ihrer wahren Gestalt und in ihrem Werthe zeigen kann.

\section{15. Schriftstudium}

Eine Predigt, sagt man, soll sein eine emanatio scripturae sacrae. Man will damit das Verhältnis der Predigt zum Texte andeuten. Im Grunde deutet man aber mit denselben Worten auch das richtige Verhältnis einer jeden, selbst der textlosen Rede zum göttlichen Worte an. Qui loquitur, spricht Petrus 1. ep. 4, 11. loquatur eloquia Dei. Wie könnte es auch anders sein! Ein Prediger soll im Namen Gottes reden. Kann er auch in Gottes Namen etwas reden, was nicht Gottes Wort ist? Muß er nicht die Zuversicht haben, daß er ebensowohl durch das, was er sagt, als durch den göttlichen Auftrag zu reden, den vollen Anspruch machen könne, als ein Bote Gottes angesehen zu werden? Wenn einer redet, was weder er, noch die Gemeinde für Gottes Wort erkennen kann, was ohne Gottes Wort gesagt ist, verdient er keinen andern Namen als den eines frevlen Schwätzers, und Dehling wendet mit Fug und Recht auf ihn
\quellseite{130}
den Spruch des heiligen Hieronymus an: Garrulitas non habet fidem.

Ist aber das der Fall, so ist einem Prediger nichts mehr zu wünschen, als daß er, mit den Alten zu reden\footnote{Doctrina ex ipsis visceribus textus biblici petenda est, quod ut fiat, textus ordine disponatur et recte explicetur. Nam doctrinae, quae ex textu fluunt, sunt optimae, ideo beatus Lutherus Theologos textuales maxime laudare solebat. Nam de hujusmodi doctrinis, quas textus suppeditat, certi esse possumus, quod eas intenderit Spiritus Sanctus, ideo et suavius afficiunt, et facilius arripiunt animumque docentis corroborant, quod non figmenta hominum, sed Dei verbum protulerit. . Balduini brevis institutio ministrorum verbi. 1622. p. 148 s.}, ein guter textualis sei. Um aber ein guter Textualis zu werden, muß man nicht bloß seine Predigttexte studieren. Man muß auch das, aber es muß dieses besondere Textstudium auf Grund eines täglichen, fortgehenden Studiums des gesamten göttlichen Wortes überhaupt geschehen. Uebersicht schafft Einsicht, und das Einzelne wird nur aus seinem Zusammenhang mit dem Ganzen richtig erkannt. Wer einzelne Texte studieren will, ohne unablässig auf die Erkenntniß der heiligen Schrift im Ganzen und Allgemeinen los zu gehen, der kann unmöglich das Einzelne und die Theile aus dem Ganzen recht verstehen lernen. — Ist es nun daraus gewiß, daß man die Schrift als Ganzes und ihren Inhalt als ein Ganzes studieren müsse, so wird auch nicht zunächst das sogenannte statarische, sondern das cursorische Lesen der heiligen Schrift empfohlen werden müssen. Denn Einsicht in Inhalt und Zusammenhang des Ganzen der heiligen Schrift
\quellseite{131}
kommt aus der cursorischen Lectüre des göttlichen Wortes. Das cursorische Lesen der heiligen Schrift kann aber freilich auch auf eine verkehrte Weise angefangen werden. Daher vermeide man den Fehler. Das cursorische Lesen muß mit dem Studium der biblischen Geschichte und mit demjenigen zusammen gehen, was man die Einleitungswissenschaft nennt. Geschichte des Reiches Gottes, Geschichte des Volkes Gottes, Geschichte der außertheokratischen Völker im A. T., Geschichte Jesu und seiner Apostel, Chronologie, heilige Geographie, Alterthümer, Geschichte der Bibel und des Bibeltextes u. dgl. schaffen eine große Fähigkeit, den Inhalt der heiligen Schrift zu verstehen, und die Einleitung in die einzelnen Bücher der heiligen Schrift gibt mehr Boden für einen richtigen Fortschritt in der Erkenntniß der einzelnen Theile der heiligen Schrift, als selbst die Grammatik.\footnote{Circumstantiae illuminant dicta, Hilarius v. Pictavium. S. Melanchthon’s Rhetorik p. 52 f.}

Deshalb darf man aber freilich das statarische Lesen der heiligen Schrift nicht verwerfen. Uebersicht hilft viel zur Einsicht. Aber Einsicht hilft auch viel zur Uebersicht. Und hat denn der HErr etwa bloß die Erkenntniß des Ganzen gewollt? Warum legt uns denn sein heiliger Geist so viel einzelne Sprüche und so manches Stück der heiligen Schrift vor die Seele, welches die abgerundete Natur eines Ganzen im Ganzen an sich trägt (Makrokosmos — Mikrokosmos)? Je mehr einer das große Ganze erkannt hat, desto erfreuter wird seine Seele werden, wenn er bei genauer
\quellseite{132}
Erkenntniß des Einzelnen inne wird, wie völlig dieses in jenes und zu ihm paßt, wie herrlich es sich einfügt, wie sich das Große im Kleinen, das Ganze in seinen Theilen spiegelt, und wie demnach das Ganze so durchaus vortrefflich und göttlich ist. Daher sei die cursorische Lectüre und das dazu gehörige Studium der Boden, auf welchem eine reiche Saat und Ernte der statarischen Lectüre wuchern.

Die Frucht der cursorischen Lectüre der heiligen Schrift werden auch solche, welche des Grundtextes mächtiger sind, schneller aus der Uebersetzung gewinnen, sei es, daß sie allein die lutherische Uebersetzung, oder auch eine andere von minderem sprachlichen und kirchlichen Werthe, aber vielleicht von knapperer Treue gebrauchen.\footnote{Z. B. die deutsche Elberfelder Uebersetzung („Neue Uebersetzung des zweiten Theils der heiligen Schrift, genannt Neues Testament. Aus dem Urtext übersetzt von einigen Christen. Selbstverlag der Herausgeber. Elberf. 1855.“ „Zweite revidierte Auflage. Elberf. und Barmen. In Commission bei Langenwiesche.“ Ohne Jahr (1864.).), — oder die autorisirte englische Uebersetzung, wer nemlich englisch versteht.} Auch die statarische Lectüre kann in einem gewissen Maße ganz nach einer treuen Uebersetzung vorgenommen werden, und es ist möglich, daß auch ein mit den Grundsprachen nicht Vertrauter zu einer tiefen und genaueren Kenntniß und Erkenntniß des göttlichen Wortes aus einer getreuen Uebersetzung gelangen kann. Wäre dies nicht der Fall, so wäre nicht bloß unmöglich, was doch in Nordamerika schon oft in der Wirklichkeit sich zeigte, ohne Kenntniß der Grundsprachen ein Pfarrer zu sein, sondern es wäre
\quellseite{133}
überhaupt mit dem protestantischen Grundsatze der Schriftlesung in der Uebersetzung nichts.\footnote{Doch darf ein der Grundsprachen Unkundiger sich namentlich in der Auslegung und Anwendung einzelner Sprüche der Bescheidenheit befleißigen, da so viele falsche Deutung und Anwendung bloß auf dem Mangel der Uebersetzung beruht. Man denke z. B. an den gewöhnlichen Mißbrauch des Wortes: „Was nicht aus dem Glauben kommt, ist Sünde.“}

Dabei bleibt jedoch das Studium der Schrift im Grundtext in höchsten Ehren. Wahrlich es ist ein gewaltiger Unterschied zwischen jeder Uebersetzung und dem Grundtexte. Man verliert außerordentlich viel durch jede Uebersetzung. Wer daher kann, der lerne und treibe griechisch und hebräisch\footnote{Der Leser vergesse nicht, daß das ganze Büchlein ursprünglich ein Dictat für Neuendettelsauer Missionsschüler ist.} bis ins Grab, und das nicht bloß stümperhaft, sondern insoweit, daß er von den Commentarien unabhängig wird, und sie frei benützen kann. Dasjenige Maß von Kenntniß des Grundtextes, welches die meisten unserer Pfarrer haben, befähigt sie kaum, tiefer in das göttliche Wort hineinzusteigen, als sie es bei der bloßen Uebersetzung könnten. Es ist zwar etwas Wahres, was Balduin beim Lob der Lutherischen Uebersetzung sagt, daß sie commentarii instar sei. Er drückt damit die Tugend der lutherischen Uebersetzung aus, aber zugleich auch ihre Schwachheit, weil ja jeder Commentar nur eine menschliche, oft falsche oder schiefe Auffassung des Grundtextes gibt, so daß man zwar berathen wird, aber ob allezeit und in allen Fällen recht? Was hilft dann ein Bischen Kenntniß der Sprachen des Grundtextes? Ein wenig Kenntniß kann erst an der Uebersetzung z. B. Lu%
\quellseite{134}
thers recht zweifelhaft machen, während genauere Sprachkenntniß zuweilen die Uebersetzung wieder zu Ehren bringen könnte. Luther wollte ja, wie er sagt, den heiligen Geist deutsch reden lassen; da scheint das Deutsch oft weit vom Texte zu liegen, während es doch treffend ist. Rechte Sprachkenntniß versöhnt zuweilen mit der Uebersetzung. Wo das aber auch nicht ist, ist sie desto mehr zu rühmen, weil sie dann den Text uncommentiert und eben damit richtiger und segensvoller gibt.

Was insbesondere die Erkenntniß und das Studium einzelner Texte betrifft, so wird, was man aus Hochstetter vom analysis exegetica, dialectica und rhetorica gelernt hat, ganz in seinem Werthe bleiben. Es setzt aber nicht bloß ein fortgesetztes Studium der Dialektik und Rhetorik, sondern auch der Sprachen voraus. Ohne letzteres bei bloßem Gebrauch der Uebersetzung wird auch Dialektik und Rhetorik von der Uebersetzung abhängig werden, was oft nicht bloß in Verlegenheit lassen, sondern auch zu Irrthum führen kann.

Für das Textstudium von großem Werth, es mag nun dasselbe nach dem Grundtexte oder der bloßen Uebersetzung vorgenommen werden, ist das Studium der Parallelen. Scriptura scripturam illustrat; aus den hellen Stellen fällt Licht auf die dunkeln. Eben deshalb ist einem Pfarrer nichts zuträglicher für seine Predigtvorbereitung, als eine gute Concordanz, und zwar keine mehr als eine solche, welche neben dem Hauptvortheil der Vollständigkeit und leichter Auffindung hauptsächlich auch den einer lichten Zusammenfassung des Gleichartigen gewährt. Eine solche Concordanz
\quellseite{135}
(z. B. die Büchnerische) bringt zugleich den Vortheil der bei unseren Vätern hochberühmten loci communes (S. § 16. p. 121). Sollte eine Concordanz beide Seiten nicht vereinen, so muß man sich zweierlei Bücher anschaffen, nemlich eine Concordanz der ersten und eine der zweiten Art, neben Büchner die vortreffliche von Landisch\footnote{Concordantiae Bibliorum germanico-ebraico graecae. Deutsche, ebräische und griechische Concordanzbibel etc. durch M. Friedr. Landisch, neu revidiert \&c. von M. Christ. Reineccio. Leipzig und Frankfurt 1718. 2 Folianten. Das Werk verdient ein eigenes zum Aufschlagen bequemes Stehtischchen in der Studierstube deutscher Pastoren.}, welche zugleich genau anzeigt, welches griechische oder ebräische Wort für das deutsche der Uebersetzung an jedem Ort gebraucht ist. Es gibt ja der Ausdrücke gar mancherlei, die im Deutschen mit Einem Worte ausgedrückt sind, so daß die feine Verschiedenheit des Grundtextes verloren geht.

Wie viel zu einer lebensvollen Auffassung der heiligen Schrift neben der Meditatio auch die Oratio und Tentatio thue, ist nicht nöthig hervorzuheben, weil alle darin übereinstimmen. Vergl. Balduin’s herrliche Worte p. 123—135. Eine gesegnete Amtsführung des Predigers erfordert ein Leben im Worte Gottes und ein Einführen desselben in Gebet und Anfechtung. Wer bloß amtsmäßig leben und studieren würde, ohne es zu eignem Bedürfniß zu thun, würde in seinen Predigten einen hohlen Ton geben. Durchaus gilt das Wort: Willst du der Gnade Posaune sein, räum dich der Gnade erst selber ein. — Ein Prediger muß die Schrift täglich aus Bedürfnis
\quellseite{136}
der eigenen Seele lesen. Man kann treulich und amtsmäßig täglich das Wort lesen, ohne daß man selbst gefördert und fromm wird. Nichts ist daher einem Prediger mehr zu wünschen, als der Geist der Meditation über die Schrift und das Gebet darum. Die erste Anwendung der Schrift muß auf die eigene Seele geschehen, dann erst auf die Seelen der Gemeinde. So gibt's dann ein Predigen aus der Tiefe der Seele, ein herzliches Reden, das zu Herzen geht, während bloß amtmäßiges Schriftstudium kalt läßt, kalte Predigten macht und den Hörer nicht recht faßt, kein menschliches Vertrauen, kein Ohr für das Gottes-Wort eröffnet, das er predigt. Es gibt keine emanatio scripturae ohne Schriftstudium für das eigene Herz, und wenn die Schrift „keinen hellen Schein ins eigene Herz gibt“, kann auch keiner herauskommen. 2 Cor. 4, 6. Ps. 1. zeigt den rechten Prediger. Ueberhaupt wäre es sehr zu wünschen, daß Prediger die meditatio nicht bloß als „Studium“ auslegten, sondern als geistliche Uebung, und daß wir einen Löscher\footnote{D. Valent. Ernst Löscher's Edle Andachtsfrüchte oder 68 auserlesene Oerter der heiligen Schrift, so von der Andacht handeln, zur Ermunterung des Geistes in so vielen Predigten nach XXV (S. p. 105 dieses Bändchens) unterschiedenen Methodis ausgeführt, darinnen die Theologia mystica orthodoxa in 6 Theilen vorgetragen wird, nebst einem Vorbericht von der Theologia mystica, und 2 Anhängen: I. vom Nutzen der Theologiae mysticae wider die Päbstler; II. von den Grenzen der Andacht und des Enthusiasmus. 3. Auflage. Leipzig 1741.} hätten, welcher Bücher wie z. B. Garcias exercitatorium protestantisch umwandeln und den Pfarrern zu Nutz und Frommen ihrer Seelen
\quellseite{137}
und damit ihrer Gemeinden zugänglich und lieb zu machen müßte.\footnote{Exercitatorium spirituale cum directorio horarum canonicarum. Auctore R. P. Garcia Cisnerio, et Formula orationis et meditationis. Auctore R. P. Ludovico Barbo. Editio nova. Ratisbonae. Sumtus fecit G. J. Manz. 1836.}

\section{16. Von der Invention}

Entweder redet ein geistlicher Redner mit oder ohne Text. Im ersteren Fall ist ihm der Text entweder gegeben, wie bei den Predigten de tempore, oder er hat ihn zu wählen, wie z. B. bei allen Casualpredigten.

Bei gegebenem Texte hat man entweder analytisch oder synthetisch zu predigen. Im ersteren Fall kann der Prediger nicht in Zweifel sein, was zu predigen sei, er legt seinen Text vor, und die Invention fällt hier mit der Analysis zusammen, von welcher wir bereits gesprochen haben. Soll er aber synthetisch predigen, so muß er finden, welchen Hauptgedanken er zu nehmen habe. Bei Casualpredigten ist ein Casus da, der den Gedanken Richtung und Maß geben kann, aber es muß doch erst für den Casus der bestimmt geformte Gedanke und dazu ein Text gesucht und gefunden werden. Bei der textlosen Rede ist wenigstens der Hauptgedanke zu formuliren. Diese gesammte Thätigkeit gehört ins Bereich der Invention, so wie die Auffindung der richtigen Partitionen und Divisionen der gefundenen Hauptgedanken.

\quellseite{138}
Man sollte freilich denken, daß es einem Prediger nie an Stoff fehlen könnte.\footnote{Sir. 24. Insbesondere v. 38 f. „Er ist nie gewesen, der es — das Buch des Bundes — ausgelernt hätte, und wird nimmermehr werden, der es ausgründen möchte. Denn sein Sinn ist reicher, weder kein Meer, und sein Wort tiefer, denn kein Abgrund.“} Die Schrift ist ein unerschöpfliches Meer von Gedanken — und die Gelegenheit, der Casus, die Kenntniß der gemeindlichen Zustände und der Personen \&c. sollten es zu einer ganz leichten Sache machen, aus der Schrift das Rechte und Nöthige zu finden. Wie von selbst sollte man meinen, müsse sich der Stoff zu jeder Rede und Predigt ergeben, ja aufdrängen. Allein so beziehungsreich Leben und Erfahrung auf die heilige Schrift, und so sehr diese ein wallendes Meer von Gedanken ist, so ist es doch eine unleugbare Thatsache, daß der Mensch träge ist, und daß auch ein „Gottesmensch“ (Mann Gottes, 1 Tim. 6, 11.), ein Prediger, durch die Gewohnheit der Beziehungen leicht taub und stumpf wird, sie zu finden. Auch ist es überhaupt gewiß, daß kein Beruf ein solches Maß von productiver Thätigkeit erfordert als der des Predigers, während doch gerade Productionskraft eine seltene Gabe ist, und diese, wie jede andere Gabe, auch wenn sie da ist, nicht schnell zu Geschick und Tugend erzogen werden kann. Es ist daher derjenige Theil der Dialektik und Rhetorik, welcher von der Invention handelt, so wie die alte Lehre von den locis communibus und die alte Topik\footnote{z. B. Topica theologica a clarissimo viro gravissimoque theologo D. Andr. Hyperio, sacrarum litterarum in inclyta schola Marpurgensi professore celeberrimo. Witebergae. Excudebat Petrus Seitz. 1565. Melanchthon sagt am Anfang des vierten Theils seiner Dialektik: Dialectices vetus divisio haec est: Alia pars dialectices est judicatrix, alia inventrix. Dicitur autem judicatrix illa pars, quae discernit voces, — judicat, quae recte connectantur, quae non recte jungantur; — discernit et propositiones et judicat, quando recte cohaereant membra in syllogismis et ceteris formis argumentorum. Altera pars nominatur inventrix, quae monet, quomodo res investigandae sint, aut proposito rerum cumulo, docet eligere ea, quae praesentem materiam illustrant. Haec ars vocatur τοπική, i. e. doctrina locorum, qui sunt velut indices rerum vel investigandarum, vel eligendarum, ut medicus proposito signo externo, ut velocitate vel inaequalitate pulsus arteriae, sequitur locum ab effectu, scit signi hujus quaerendam esse causam, hoc est fontem pulsus, scilicet cor, quaerit deinde, unde cor incensum sit. — Saepe autem locis utimur non tam ad investigationem, quam ad electionem rerum, quarum cumulus propositus est, ut in ecclesia docenti non sunt res inveniendae, non enim gignimus doctrinam; sed cum pars aliqua coelestium concionum proposita est, eligit res praecipuas prudens interpres, et loci monstrant, quo ordine sint explicandae, quaerenda definitio, faciendae partitiones, eruendae causae, judicandi effectus. — Später: Locus dialecticus est sedes argumenti, seu index, monstrans, ex quo fonte sumendum sit argumentum, quo confirmanda est propositio, de qua dubitas. cf. Hyperii oben angeführte Topik. p. 3. — Vergl., wenn du willst und Gelegenheit hast zur weiteren Instruction: Themata concionum sive Sylva locorum in singula evangelia etc. von Andreas Pancratius von Hof. Frankfurt a/M. 1597.} nicht
\quellseite{139}
zu verachten, sondern hoch zu achten, zu lernen, zu hegen und zu pflegen. Recht ernstlich darf auf den vierten Theil der Melanchthon’schen Dialektik, sowie auf dessen erstes Buch der Rhetorik verwiesen werden.

Zur Invention leitet Balduin vortrefflich auch dadurch an, daß er aus St. Paulo 2 Tim. 3, 16. zeigt, wie man
\quellseite{140}
aus der heiligen Schrift und ihren Texten heraus finden könne und solle, was theils zur διδασκαλία, theils zum ἔλεγχος, theils zur ἐπανόρθωσις, theils zur παιδεία dienen kann. Es läßt sich durchaus nicht leugnen, daß jene Stelle des heiligen Paulus, welcher die oft gebrauchte Division alles Inhalts der heiligen Schrift und Lehre entnommen ist, von der größten Wichtigkeit sei, sowie daß das Studium desjenigen, was die Väter über den elenchus und seine verschiedenen Arten geschrieben haben, sehr anzurathen ist und zur Invention sehr tüchtig macht.\footnote{Vgl. Balduin p. 164 f. VI. Deyling p. 328 f. „Der wahre Lehrelenchus schriftmäßig betrachtet von Bernhard Walther Marperger. Dresden 1727.“}

\section{17. Länge und Kürze der Predigten}

Länge und Kürze der Predigten sind relativ. Zwar nimmt man gewöhnlich nach Luther’s und anderer Urtheil eine halbstündige Dauer der Predigt für das rechte Maß, und man hat oft nachgesagt, es müsse ein gewaltiger Prediger sein, den man länger vertragen solle. Aber eben mit diesem Ausspruch deutet man doch auf einen maßgebenden Umstand, nämlich auf die Begabung des Predigers hin, ohne selbst ableugnen zu wollen, daß es solcher maßgebenden Umstände auch andere mehr geben könnte. Die Erfahrung lehrt z. B., daß hie und da ein solcher Hunger nach Predigt ist, vermöge dessen auch ein mittelmäßiger, ja ein geringer Prediger nicht bloß vertragen, sondern mit Lust gehört wird. Auch kann einmal eine Materie
\quellseite{141}
von der Art sein, daß sie, ohne zu ermüden, viel länger abgehandelt werden kann, als eine andere. Dauert doch manche politische Rede, man denke an das englische Parlament, stundenlang bei gespannter Aufmerksamkeit der Zuhörerschaft. Warum sollte, so hoch man auch den Ekel des menschlichen Herzens gegen das Heilige und Gute anschlage (cf. Luther’s gr. Kat. über die ἀηδία) unter gewissen Umständen nicht auch ein kirchlicher oder sonst ein heiliger Stoff die Herzen mit größerem Interesse ergreifen und länger festhalten können?

Es wird übrigens allerdings die Frage und Antwort über Länge und Kürze der Predigt anders lauten, je nachdem man die Frage vom Standpunkt des Lehrers oder des Hörers stellt. Der Hörer wird von Materie, Gabe, Zeit und Umständen sein Urtheil über die Dauer einer Predigt mit Recht abhängig machen. Der Lehrer selbst hingegen wird wohl thun, wenn er die Dauer seiner Predigten so bemißt, als wäre seine Begabung geringer und deshalb ein Hinderniß für lange Predigten. Man könnte freilich ganz einfach sprechen: „Was ich zu sagen habe, das sage ich, und wenn ich fertig bin, so höre ich auf.“ Und diese Regel wird jedenfalls Geschwätz verhüten, wenn sie gewissenhaft befolgt wird. Allein diese Regel ist doch nicht die einzige, die es gibt, auch nicht die herrschende. Ein Prediger muß bei seinen Vorträgen so nothwendig auf den Ekel des menschlichen Ohres Rücksicht nehmen und so ernstlich nach Aufmerksamkeit seiner Hörer streben, daß er lieber zweimal unvollständig von Einer Sache reden muß, wenn er sich die Aufmerksamkeit sichern kann, als einmal vollständig,
\quellseite{142}
wenn er fürchten muß, die Hörer zu ermüden. Er muß Herr des Inhalts sein, und das rechte Maß des Inhalts gibt alsdann auch das rechte Maß der Zeit. Kann man bei dieser Maßbestimmung nicht anders als entweder zu lang oder zu kurz reden, so ist ein zu kurzes dem zu langen vorzuziehen. Der Zuhörer wird immer mehr gefördert sein, wenn er am Ende des Vortrags noch länger hören möchte, als wenn er schon längst aufs Amen gewartet hat, während es immer und immer verzog zu kommen. Eberlin von Günzburg hat 1525 an den Pfarrer Joh. Jakob zu Leipheim ein Sendschreiben „Wie sich ein Diener Gottes Worts in all seinem Thun halten soll“ verabfaßt\footnote{Die treffliche Schrift ist als Einleitung zu diesem Bändchen vollständig mitgetheilt.}, worin er sagt: „die Prediger, welche zu lang reden, wüßten nicht, daß die Zuhörer nicht mit des Predigers Ohren hören, sondern mit ihren eigenen.“ Er vergleicht auch solche Prediger ganz richtig mit jungen\footnote{Es geht aber auch bei vielen Alten gerade so. Neminem libentius audiunt, quam se ipsos, et ut simia in purpura sua, ita ipsi in suo docendi dono sibi solis placent. Balduin p. 130.} Müllern, die noch nicht viele Säcke gebunden haben und sie deshalb zu groß binden. Ueberhaupt ist Eberlin ein Mann des rechten Maßes, dessen Capitel X. bis XVIII. besonders zu empfehlen sind.

Vergessen dürfen wir endlich dieses Orts nicht, daß eine zu lange Predigt auch die Harmonie und den Zweck des ganzen Gottesdienstes stört. Man kommt ja auch nicht
\quellseite{143}
bloß deshalb in die Kirche, daß man höre; man will ja auch beten, singen, zum Sacramente gehen. Man kann also nicht die ganze gottesdienstliche Zeit zum bloßen Hören anwenden.

\section[18. Von dem Einfluß der persönlichen Autorität]{18. Von dem Einfluß der persönlichen Autorität auf den Ton und Sprache des Predigers}

Die Autorität des Predigers ist bekanntlich eine göttliche, nemlich des Amtes, das an Christi Statt vermahnt, redet und bittet. Man könnte deshalb behaupten, es hätten alle Prediger einerlei Autorität und Ansehen. Allein das ist nun einmal nicht zu leugnen, daß der Mangel an persönlichem Ansehen, das auf natürlicher Gabe, Character, Erfahrung und Tugend etc. etc. beruht, der Autorität des Amtes viele Hindernisse in den Weg legt, und umgekehrt, daß persönliche Autorität die amtliche schnell zur Anerkennung bringt. Es ist zwar gewiß, daß sich die persönliche Autorität da, wo sie ist, von der amtlichen nur erhöhe, aber kaum unterschieden wird, mit derselben verschmilzt, und deshalb auch im Vortrag manches, was Fleisch ist, für Geist gehalten wird. Allein das läßt sich in dieser Welt nicht ändern, und man muß am Ende doch wünschen, daß kein Mann das Amt und seine Autorität suche, der keine natürliche Autorität besitzt. Wirkt auch persönliches Ansehen für sich allein nicht immer geistlich, so ist doch auch die Wirksamkeit einer bloß menschlichen Autorität, zumal wenn der, welcher sie hat, dem göttlichen Worte treu ist, nicht zu
\quellseite{144}
gering anzuschlagen, da sie, wie alles menschlich Pädagogische den Segen des vierten Gebotes hat und zu Gott führt. Wollte man aber auch den Segen der menschlichen Autorität noch so gering ansehen, so hängt doch gerade von ihr Ton und Sprache des Redners ab. Ein gering oder mittelmäßig begabter Prediger kann und darf nicht reden wie ein hochbegabter, ein junger nicht wie der alte, der Neuling nicht wie der bewährte, der Mann von beflecktem Gewissen nicht wie der von reinem Bewußtsein \&c. \&c. Balduin sagt in Bezug auf das Alter und dessen Einfluß: Multa profert magna cum autoritate vir gravis et dicendi peritus, ita ut ad cor auditorum loquatur, quae risu exciperentur, si dicerentur a juniore. (p. 120.) Daher ist bei unseren Vätern, nicht bloß bei Balduin, die Bescheidenheit des Predigers sehr in Ehren. Bescheidenheit ist für den Prediger eine hohe und unerläßliche Tugend. Wer sich und sein Ansehen durchweg um ein gutes Theil niedriger schätzt, als ihm sein alter Mensch ansagt, wird wohl thun, Gewissensbelastung der Eitelkeit vermeiden und seine Zuhörer anreizen, zu ihm zu sprechen: Freund, rücke hinauf. Jedermann ist geneigt, dem Bescheidenen zu vertrauen, während die Rede des Unbescheidenen Mißtrauen erweckt und Gottes Segen verscheucht. Wer sich selbst erhöht, wird erniedrigt werden. Wer daher gerne auf eine große Höhe des Vertrauens kommen will, der erkenne vor allem, daß er von Natur und durch sein Leben kein Vertrauen verdient, und nehme demgemäß seinen Standpunkt ein, und zwar nicht bloß für die jugendlichen Jahre, sondern fürs ganze Leben. Der wird ihm
\quellseite{145}
selbst am nützesten sein, sonderlich wenn Vertrauen und Ansehen wächst; denn es erscheint oft ein Lebensgang nach außen als große Bewährung, und das Vertrauen steigert sich im Fortgang, während das eigene Auge, das ins Innere schaut, je länger je mehr gerechte Ursache zur tiefsten Demüthigung findet. Der Herr erhält die Seinen klein, erhält sie aber damit in der Wahrheit und im Streben und macht sie auf diesem Wege immer treuer und vertrauenswürdiger. Es geht ihnen nach dem Spruch: Wenn du mich demüthigst, machst du mich groß. — Merken mögen sich dies besonders die jungen Geistlichen, die schnell zu einer großen Bewunderung der eigenen Gabe und Leistung kommen. Fast ist kein eitleres Geschlecht als junge Pfarrer. Alle Erfolge sind ihre Früchte. Alles Gute, was sie sehen, setzen sie in Verbindung mit ihrer Wirksamkeit. Andere erkennen, wie gering ihre Wirkung ist; nur sie selbst leben in Eigenlob. Sie wissen bei jungen Jahren alles, und zwar nicht bloß aus der Schule, sondern — wie lächerlich — gar aus Erfahrung.

\section[19. Vom Extemporiren]{19. Vom Extemporiren und der nothwendigen Ausarbeitung der Predigten}

Ueber Extemporiren und Ausarbeitung der Predigten haben sich entgegengesetzte und einander ausschließende Meinungen geltend zu machen gesucht. Die bei allem Unwerth erträglichere Uebertreibung war es immer, wenn man gar kein Sprechen ex tempore wollte gelten lassen, während die entgegengesetzte,
\quellseite{146}
gemäß welcher ein Reden aus dem Stegreif Zeichen des rechten Redners sein soll, gar keinen Werth ansprechen kann.

Vor allen Dingen ist es ein unerläßliches Erforderniß, daß einer, der reden will, wisse, was er vorbringen will, und wie er es vorzutragen habe. Dies vorausgesetzt läßt sich ganz wohl von einer Gabe reden, vermöge welcher ein Mann bei hervortretender Gelegenheit und Nöthigung, auch ohne seine Rede geschrieben und memoriert, ja ohne sie schriftlich skizzirt zu haben, vortrefflich und mit bester Wirkung sprechen kann. Die Landtags- und Parlamentsreden würden dazu Belege genug liefern können. Man kann ja aber von dergleichen ganz absehen; die Amtswirksamkeit hochbelasteter Landpfarrer, denen es zuweilen obliegt, acht Tage hintereinander täglich ein- oder zweimal öffentlich zu reden, gibt hierhergehörigen Beweis und Beleg in Ueberfluß. Es wäre nur zu wünschen, daß es noch mehr solche Männer geben möchte, als es leider der Fall ist, solche nemlich, welche die Gabe der freien Rede haben und daher auch ohne ihre Predigten zu schreiben, mit Ruhe und Kraft, mit Weisheit und Umsicht reden können. Allein nicht nur zeigt die Erfahrung, daß es dieser Art wenige gibt, sondern es ist auch die Frage, ob einer, der wirklich die nöthige Gabe hat, sie unbeschränkt und für gewöhnlich gebrauchen soll. Eine Gabe wird allerdings schön durch den Gebrauch geweckt, gehoben und gestärkt; so wird auch die Gabe der freien Rede schon durch den bloßen Gebrauch einen gewissen Grad der Ausbildung erlangen. Es kann einem Manne gelingen,
\quellseite{147}
sofern er nur selbst beständig in dem lebt, was er zu predigen hat, sofern er nur an seiner inneren Ausbildung und der Erweiterung seines Kenntnis- und Gesichtskreises unermüdlich vorwärts arbeitet. Allein wird nicht eine oftmalige, ja zur Gewohnheit gewordene genaue Vorbereitung auf jede einzelne Predigt die vorhandene Gabe noch mehr stärken, zumal ja niemand an seine Vorbereitung gebunden ist, sondern nach gründlicher Vorbereitung desto sorgloser auch frei reden kann? Wird die stille Meditation, welche so herrlich mit dem Studieren, Disponieren und Schreiben zusammengeht, nicht große Wirkung auf die Seele des Predigers haben und ihm Salbung verleihen? Wird sie nicht die Ruhe und jenes „urkräftige Behagen“ mehren, aus dem hervor die Rede wie ein milder und gewaltiger Strom dringt und „die Herzen aller Hörer zwingt?“ — Ueberdies wer darf denn in allen Fällen seiner Gabe trauen, wer, ohne irgend einmal recht bitter und vielleicht zu großer Einbuße an Vertrauen und Ansehen getäuscht zu werden, sich darauf verlassen, daß er im Augenblick der Rede das rechte Ziel und den rechten Weg finden und einhalten werde? Es ist mancher hochbegabte Mensch für sein eitles Selbstvertrauen bitter gedemütigt und hart gezüchtigt worden.\footnote{Vgl. Balduin p. 145. 146. — Concionatores timidi sunt optimi. — „Suchet in aller Einfalt Gottes Ehre, nicht Ruhm und Zufall von Menschen, und betet, daß euch Gott Verstand und Mund, und den Zuhörern ein recht rein Gehör verleihe, und lasset Gott walten. Denn das wollest mir glauben, daß Predigen nicht Menschenwerk ist. Denn ich, wiewohl ich nun ein alter und geübter Prediger bin, doch fürchte ich mich, wenn ich predigen soll, und ihr werdet gewiß diese drei Stücke erfahren. 1) Da ihr gleich die Predigt aufs allerbeste gefaßt und begriffen habt, worauf sie stehen soll, so soll es euch wohl zerrinnen und zu Wasser werden; — 2) Wenn ihr am Begriff und Concept gar verzaget, so gibt Gott Gnade, daß ihr am besten prediget, und es dem Haufen wohl gefällt, euch aber nicht gefällt; — 3) Wenn ihr nicht gefaset habet, daß es beiden, euch und den Zuhörern wird gefallen. Darum bittet Gott, und laßt’s dem befohlen sein.“ Luther. Porta’s Pastorale Lutheri. ed. Cramer. 1729. p. 113. Anmerk. ibid. „Auf einem gewissen Dorfe predigte einst ein Studiosus Theol., und da ihm der Pastor Gottes Gnade anwünschte, sagte er, er brauche es nicht, denn er habe schon mehr gepredigt und noch nie gefehlet. Es geschah aber darauf, daß er stecken blieb und von der Kanzel wieder herunter gehen mußte.“}

\quellseite{148}
Man könnte sagen: „Predigen sei doch nicht Menschenwerk, es sei etwas ganz anderes mit der Predigt, wie mit einer weltlichen Rede. Bei dieser walte der menschliche Geist allein, während dem Prediger ein höherer Geist beistehe und ihm zu seinem Werke helfe.“ Und wahrlich, dem ist so, predigen ist nicht Menschenwerk. Wie oft ist ein gründlich vorbereiteter Mann vor die Gemeinde getreten um zu predigen, dem entweder alles, was er sagen wollte, dahinschwand oder doch sich als unzulänglich erwies, so daß er verwirrt und arm vor der Gemeinde stand und mehr um ihr Erbarmen betteln mußte, als er Anspruch machen konnte, für ihren Lehrer und Prediger zu gelten. Und umgekehrt: Wie oft tritt ein Prediger bettelarm auf seinen Predigtstuhl und siehe, wenn Gesang und Orgel schweigen, hat er Reichthums die Fülle, speist die Gemeinde mit Himmelsbrod, tränkt sie mit Wein und Milch
\quellseite{149}
und mangelt ihm an keiner Gabe. Die wunderbare Erfahrung machen viele Prediger immer und immer, und welche sie machen, die leben je mehr und mehr vom Mund des Herrn und seinem gnädigen Erbarmen. Aber werden diese beiderlei Erfahrungen den Menschen in seiner Trägheit und in seinem Leichtsinn bestärken dürfen? Sind nicht beide recht eigentlich nach des Herrn Sinn und Meinung mächtige Anmahnungen zum Fleiß und zum Gebete, daß der Fleiß gesegnet werde? Der sich nicht bereitet hat, wird der sich beim Mißlingen so leicht demütigen können? Wer beim Auftritt eine Sünde der Trägheit zu bekennen hat, wird sich der so leicht der Armuth anklagen und beten können, daß ihn Christus reich mache? Wenn wir alles gethan haben, sind wir unnütze Knechte, die der Herr nütze macht. Es geschieht uns dann nach dem Wort: „Fleißige Hand macht reich“, nemlich durch des Herren Gnade. Die leichtsinnigen, trägen Frevler am Heiligen aber werden sich vielleicht nicht durch Gelingen, noch durch Mißlingen demüthigen lassen, der Geist des Herrn aber wird sie je länger je mehr durch Mißlingen und Untüchtigkeit strafen.

Nun haben wir aber bisher immer solche Männer im Auge gehabt, welche die Gabe der freien Rede haben, und es dennoch nicht räthlich gefunden, daß sie dieselbe wollten lassen, ohne sich vorher so oft und viel als möglich auf das genaueste bereitet zu haben. Denn daß ein begabter Mann, nach gründlicher Vorbereitung seine Gabe walten lasse und das Bessere, was ihm nach treuem Fleiß der Vorbereitung im Augenblick des Bedarfs gezeigt wird, sage, oder das Gefundenen
\quellseite{150}
besser sage, wer wollte dem wehren? Was soll aber mit denen werden, die keine Gabe frei zu reden, vielleicht überhaupt keine Redegabe besitzen und sich doch herausnehmen, ex tempore zu sprechen? Welch’ eine Verkennung der eignen Begabung! Welche Blindheit, die, so oft und gewöhnlich sie vorkommt, doch in der That, wie aus Unverschämtheit und Hochmuth entsprungen erscheinen muß! Und welche Mißachtung des Worts, der Gemeinde, des Amts, des Herrn und der Rechenschaft! Die Katechesen des Cyrill mögen ex tempore gehalten sein. Chrysostomus hat seine meisten Lobreden auf Märtyrer extemporiert. Was geben aber den gering Begabten solche Beispiele an? Und wie viel besser ist es, sich in diesem Stücke nach dem Beispiel eines Origenes zu achten, der auch extemporierte, aber nicht vor dem 60. Lebensjahre. Balduin sagt mit Recht p. 136 ff.: Minime laudandi sunt extemporanei concionatores, qui sibi laudi ducunt, quod quasi stantes pede in uno ad rude vulgus aliquid deblaterare et quasi e manicis, quod dicitur, concionem excutere possunt. Quae temeritas tentat Deum, fugit laborem et non auditorum utilitatem, sed privatam quandam gloriolam captat.

Also nein, man extemporiere nicht, oder besser zu reden, man extemporiere nicht, wenn man es vermeiden kann; denn es ist allerdings nicht möglich, daß ein Pfarrer, der fast alle Tage Predigten zu halten hat, sie auch alle Tage schreibe, ausarbeite, memoriere.\footnote{„Dr. Aegid. Hunnius in Method. concionandi p. 19. 6. sqq. räth, daß man sich aufs Auswendiglernen nicht legen soll, sondern eine gute Disposition machen und solche dem Gedächtnis wohl imprimieren solle.“ Bei dem angestrengten und ängstlichen Memorieren könnte es kommen, „si quis uno verbo impingat, ut in toto contextu concionis hallucinetur.“ Wer nach guter Vorbereitung auf Grund einer guten Disposition predige, sei viel unbefangener und von freierer Bewegung „Praeterea frigidior et sermo et actio videtur ejus, qui omnem vim memoriae verbis potius, quam rebus addicit; multo autem gravior, quando concionator comprehensis animo rebus ipsis liberrimum sibi inter dicendum relinquit verborum delectum.“ Porta's Pastorale Lutheris p. 112 f. Ganz richtig, aber wir haben im Texte oben auch das gute Gewissen nicht auf ein völlig und wörtlich ausgearbeitetes Concept und sicheres Memorieren, sondern auf treue Vorbereitung gegründet. Wir stellen extemporieren nicht in Gegensatz zum memorieren, sondern zur treuen Vorbereitung. Daher der Tadel des Extemporierens. Die Angst vor unüberlegtem Schwatzen empfiehlt Strenge der Vorbereitung — für junge Prediger auch wörtliches concipieren und memorieren.} Man extemporiere nicht, sondern
\quellseite{151}
halte auf treue Vorbereitung und bringe das anlangend, nie ein anderes als ein gutes Gewissen auf die Kanzel. Man sei nicht pedantisch, man hange nicht am Wort, man klammere sich nicht so an's Concept an, daß man beim Entfallen eines Wortes oder Satzes auch alle Haltung verliert und gleich dem Kinde nach mißlungener Gedächtnisaufgabe ratlos steht. Aber man sei treu im Studieren, Disponieren, Schreiben, Ausarbeiten und soviel es die Begabung erfordert und rathlich macht, auch im Memorieren. Auf dem Wege der treuesten Vorbereitung bringe man zu einer freien und anständigen Bewegung auf dem Predigtstuhl vorwärts.

\quellseite{152}
\section{20. Von der Electio}

Es muß Absicht und Wunsch eines jeden Predigers sein, seinen Zuhörern verständlich zu werden, weil eine Wirkung seiner Vorträge ohne lichte Verständlichkeit gar nicht zu denken ist. Diese Verständlichkeit, welche nicht ein und dasselbe mit dem zweifelhaften Lobe der Popularität ist, wird allerdings schon dadurch angebahnt werden, daß die Rede licht und klar angelegt und eingetheilt wird und der Gedankengang einfach und folgerichtig und eben damit auch zur Aufmerksamkeit einladend, sie fördernd und erleichternd vorwärts schreitet. Viel aber kommt doch auch auf den Bau der Sätze und auf die Wahl des Ausdrucks an.

Man kann sich in Betreff des Satzbaues entweder einfacher kurzer Sätze oder einer periodischen Redeweise bedienen. In Perioden gut schreiben, erfordert viel Gabe und Bildung. Unter Hunderten, die es wagen, gelingt es kaum einem. Im schlechten periodischen Styl aber schreiben und reden heißt so Leser wie Hörer von sich scheuchen und seine Gedanken in Gräber verbergen, aus denen man sie nicht einmal herausbohren mag. Wer liest Spener’s und Francke’s Bücher? Sie sind unzugänglich durch die schlechte periodische Schreibart ihrer Zeit. Ueberdies kann die periodische Schreib- und Sprechweise nicht bloß zum Begraben der Gedanken dienen, sondern sie begünstigt auch das Schwatzen, während der einfache, kurze Satz an und für sich dem Geschwätze widerstrebt, dagegen aber zur planen Vorlegung reifer Gedanken sich vortrefflich eignet.

\quellseite{153}
Warum liest man jetzt noch Heinrich Müller, warum ist er so erfrischend und fesselnd? Die Schreibart in kurzen Sätzen ist es\footnote{Stilus nervosus, sententiöser Styl. — Johann Laesenius. Nervös, kräftig, lebendig.}, ohne was er wenigstens nicht wäre und wirkte, was er ist und wirkt. Da man nun durch kurze Sätze so leicht verständlich werden kann, und man auch eine Vollkommenheit in dieser einfachen Weise des Vortrags gewisser erreichen kann, als im periodischen Style; so wird es namentlich denen, welche keine besondere formale Ausbildung empfangen haben, zu rathen sein, sich im einfachen Ausdruck der Gedanken in kurzen Sätzen zu üben. Allerdings darf man dabei nicht vergessen, daß der kurze Satz auch zu lakonischer Unverständlichkeit und Schwerfälligkeit ausarten kann, wie die Periode zum Geschwätz.

Man kann übrigens seine Predigt einfach angelegt haben und in kurzen Sätzen schreiben und doch den Eindruck des Gezierten und Unnatürlichen machen und eben damit abstoßend wirken, wenn man bei der Wahl der Worte lieber auf Pracht und Pomp sieht, als auf treffende und genaue Bezeichnung dessen, was man sagen will. An Demosthenes wird gerühmt, daß er seine großen Erfolge bei einer den Sachen genau anpassenden eigentlichen Ausdrucksweise erreicht habe, während an Cicero die gesuchte asiatische Pracht des Ausdrucks getadelt wird. Bei Demosthenes vergesse man den Redner, bei Cicero werde man des Mannes nicht los, auch wenn er für seine Sache einnehme. Hierin liegt eine Weisung für den geistlichen Redner, man könnte sagen, eine doppelt starke Weisung.
\quellseite{154}
Seine Gegenstände sind ohnedies nicht von dieser Welt und können im Conversationston gar nicht abgehandelt werden. Sie sind erhaben und bringen auch den einfachen Redner dahin, daß er im genus dicendi grande zu reden scheint. Eben deshalb braucht ein geistlicher Redner nicht erst nach erhabener Rede und feierlichen Ausdrücken zu suchen. Er drücke sich so eigentlich als möglich aus und vermeide Allegorien, Bilder und Redefiguren, wenn sie nicht dienen, die Sache und das Verständnis zu fördern oder sie angenehmer zu machen. Summa utilitas omnis regula: ist ein Satz, welcher auch in Anbetracht der Zier und des Schmucks der Rede gilt, und von welchem man sich auch nicht durch große Beispiele abführen lassen sollte. Veritas proponatur quam simplicissime, sagt Balduin. Er und andere reden gegen alle Singularität des Ausdrucks, und auch der lautere Geschmack eines Melanchthon wehrt sich gegen alle ξενοφωνία und καινοφωνία. Dagegen rathen die Alten, namentlich die stehende Phraseologie der Kirche festzuhalten und sich immer so auszudrücken, daß dem Hörer die Eintracht aller Prediger bemerklich werde. Gewisse Begriffe und Sätze müssen unveränderlich beibehalten, wie Marksteine der Wahrheit dürfen sie durchaus nicht verrückt werden.

Sehr oft kommt es gerade schwach begabten gedankenarmen Predigern bei, sich beständig in biblischen Worten auszudrücken; ihre Predigten sind eine Mosaik von Bibelsprüchen.\footnote{Antiochus in Kloster Saba, Mönch 629, wird als erster Vorläufer der Predigten aus Bibelsprüchen genannt. Pandectae divinae scripturae. 130 kurze Homilien. Sehr einfach.} Eine unerträgliche Weise, wenn man, ein kleiner
\quellseite{155}
Mensch, immer in der schwereren Waffenrüstung der heiligen Schrift einhergehen will, und zwar gerade zu dem Endzweck, daß man leicht gehen könne. Wer eine Predigt aus Bibelstellen zusammensetzen will, wird sicher, auch wenn er ein großer Textualis wäre, gar oft den genauen Sinn eines Spruches bei Seite setzen und mit Gottes Worten anderes sagen, als Gott damit sagt. Schier keine schwerere und verantwortlichere Sache, als ipsissimis verbis Domini die eigene Meinung zu sagen. Die Schrift ist Text, nicht Predigt, die Predigt Auslegung, nicht Textesammlung. In der Predigt kann man wohl hie und da eine andere, als die Textesstelle zur Beweisführung und Erläuterung anwenden, aber ein überfließender und überflüssiger Gebrauch der göttlichen Worte ist unziemend, und kaum ein mal wird ein Prediger so begabt sein, daß er die Bibelstellen wie Edelsteine zusammenreihen kann, so daß der Glanz der einzelnen durch den Widerschein der andern erhöht und doch nicht alteriert werde. Man stehe deswegen ab und thue, was recht ist, gebe zu Gottes Wort eine einfache, treue, menschliche Auslegung in eigenen Worten.\footnote{Scripturae sacrae testimonia non tam cumulanda et numeranda, quam apposite seligenda et illa imprimis verba, quae ad rem pertinent, urgenda et recitanda sunt. Nimia et immodica dictorum, quae parallela sunt, accumulatio et ex concordantiis repetitio ac librorum capitumque allegatio tam supervacua est quam fastidiosa. Quid enim prodest, sententias scripturae coacervare et testimoniorum veluti densum nimbum extendere in ejusmodi praesertim argumento, quod operosa non indiget probatione. Deyling. Auch einschlägig.}

\quellseite{156}
Ebenso wenig als es sich ziemt, immer mit Bibelsprüchen zu reden, ziemt es sich, die Ausdrucksweisen früherer Zeiten und Lehrer nachzuahmen. Manche von unseren Brüdern in Nordamerika z. B. nannten es Kirchenstyl, ihre Redeweise der Redeweise Luther’s ähnlich zu machen. Sie ziehen Luther’s Kleider und Saul’s Waffen an, ohne den Körper dazu zu haben. Auf den Unbefangenen kann das ebenso wenig einen guten Eindruck machen, als wenn man gegenwärtig in der Tracht des 30jährigen Krieges oder in den Allongeperücken des 18ten Jahrhunderts dahergehen wollte. Auch von den Zeiten heißt es: suum cuique. Es gehört zur männlichen Bescheidenheit, in seiner Eigentümlichkeit zu bleiben und zu reden, den Helden ihr Heldenschwert zu lassen und sich des eigenen zu bedienen. Man könnte aus amerikanischen Predigten die Beweise liefern, wie viele Lutherische Eruditäten statt Perlen man kraft dieser Lutheromanie und Lutherophonie aufgelesen hat. — Bei uns in Deutschland haben manche Scrivers Weise oder gar die von Valerius Herberger nachzuahmen gesucht. Daher eben der unerträglich manierierte Styl, den sie haben, der niemand erquickt denn sie und wer gleiches Steckenpferd mit ihnen reitet. Das große Wort, welches auch hier wieder gilt, heißt Einfalt.

\section{21. Gebrauch der Postillen}

Wenn man hört, daß Karl der Große in seinem herrlichen Homiliarium (denn herrlich ist es, und bis zur Stunde besitzt die Kirche kein zweites Werk der Art) den Predigern
\quellseite{157}
ein Buch der Vorbereitung gehalten habe, so pflegt man die arme Zeit zu bedauern, deren Prediger man in die Postille zur Schule schicken mußte. Aber wahrlich eine aufgeblasene und heuchlerische Zeit, in der wir leben! Wie wenn nicht auch jetzt noch gar viele, viele Pastoren alle Studiertage auf ihren Studierstuben säßen und ihr Ragout aus anderer Brei herausbräuten! Wie wenn nicht jetzt noch viele, viele gar nicht leben noch predigen könnten, ohne — ich sage nicht Postillen, denn wir haben keine, die diesen Namen verdienten, aber doch ohne Predigtbücher! Statt daß man so vornehm und heuchlerisch auf vorige Zeiten sähe, sollten die begabten Prediger lieber zusehen, daß sie einmal in usum communem eine Predigtsammlung, eine Postille\footnote{Postille hat den Namen von dem Uebergang vom Text zur Erklärung: „post illa verba evangelii.“ — Karl d. Gr. wollte, daß die Postille entweder in der Muttersprache (theotiscam) oder in rusticam latinam linguam übersetzt würde. — Titel des Homiliariums: „Homeliae seu mavis: sermones sive conciones ad populum praestantissimorum ecclesiae doctorum: Hieronymi, Augustini, Ambrosii, Gregorii, Origenis, Joh. Chrysostomi, Bedae, Herici aliorumque. In hunc ordinem digestæ per Alchuinum levitam, idque injungente ei Carolo Magno Rom. Imper., cui a secretis fuit. Nunc ex vetustissimis exemplaribus ab innumeris mendis vindicatae, quod facile deprehendet is, qui hanc editionem cum prioribus contulerit. Adjectis etiam nonnullis homeliis, nusquam antea impressis. Apud sanctum Ubiorum Coloniam Agrippinam in aedibus Heronis Alopesii. Ann. MDXXV. — Neuere Sammlungen von anderem Zwecke, wie Wilh. Beste’s „Bedeutendste Kanzelredner der älteren lutherischen Kirche.“ 2 Bde. Leipzig 1856. 1858, oder Nesselmann’s „Buch der Predigten und Reden aus den verschiedenen Zeiten, Ländern und Confessionen, zu einem Jahrgang geordnet“, Elbing 1862, können gerade dazu dienen, zu zeigen, wie schwer es sein würde, ein populäres Predigtbuch von der Einheit und Herrlichkeit des Carolinischen Homiliariums aus den Predigten auch nur der lutherischen Kirche herzustellen.} zusammenbrächten, welche es in ihrer
\quellseite{158}
Art verdiente, z. B. nur neben dem alten Haymo von Halberstadt zu liegen, — und die unbegabteren sollten einfach und offen die Postille gebrauchen. Summa utilitas omnis regula — heißt es hier wieder. Was kann der Gemeinde daran liegen, eigene Arbeiten ihres Pfarrers zu hören, wenn er doch nicht im Stande ist, solche Arbeiten selbst zusammenzuschreiben und zusammenzubrauen, welche sie fördern können und werth sind, vorgetragen und gehört zu werden? Ich habe alle Achtung für die Beredsamkeit manches einfachen Christen und wollte viel lieber seine regellosen Töne hören, als das geregelte leere Getöne eines und des andern sogenannten studierten, aber blutarmen Predigers. Ich habe große Achtung für die einfache zum Herzen sprechende gelegentliche Rede auch eines zum Reden gering begabten Pastors, aber gar wenig oder keine Achtung für seine abgequälten, der Armuth abgerungenen, sogenannten studierten Vorträge, mit denen er seine Bauern einlullt und statt sie zum Göttlichen empor zu heben, sie zum Gewöhnlichen und Niedrigen herunterdrückt. Ach wenn man doch den Muth hätte, weise zu sein und den Nutzen der Gemeinde zu schaffen! Ein junger Geistlicher predigte zu großem Segen und unter großer Theilnahme der Gemeinde und Umgegend, so lange er die Predigten eines Begabteren studierte, in sich reproducirte und dann auf seiner Kanzel vortrug. Als aber
\quellseite{159}
einer seine Quelle entdeckte und ihm unbescheidener Weise die Entdeckung mittheilte, fühlte sich der gute Prediger gedemüthigt, ja verrathen und anstatt einfach zu sagen: Summa utilitas omnis regula, und unter heiterem Bekenntniß seines Beginnens zum Segen seines Völkleins fortzufahren, sagte er von nun an aus seiner eigenen Feder und aus seinen Fingern seinen Vortrag, und aus war’s mit dem Segen. Sapere aude, und die elenden Vorurtheile namentlich derer, die in der Gemeinde die Gebildeteren sein wollen und deshalb nur einen beredten und begabten Prediger vertragen — die wirf dahin! So Prediger wie Gemeinden sollten einmal belehrt und überzeugt werden, daß wirklich alles am Segen der Zuhörer liege, und daß vor die Gemeinde die möglichst beste Predigt gehöre, woher sie auch genommen werde, nicht aber eine geringe, nur damit der Pastor eigene Arbeit gebe. Balduin redet gegen den Gebrauch der Postillen, aber heutzutage muß man gegen entgegengesetzte Uebel kämpfen. Ganz schön sagt er: Novitii utantur postillis instar corticum, quos tandem abjiciant postquam ipsi natare didicerunt. Er setzt hinzu: servilis enim ingenii est, duci malle quam ire. Allein was hilft’s, mit dem servile ingenium den Pfarrerhochmuth zu stacheln, wenn nun einmal die meisten ein servile ingenium haben, welches dadurch allein geadelt werden kann, daß es lieber zum Heil der Gemeinde geführt werden will, als zu ihrem Schaden auf eigenen Füßen gehen? Wahrlich, hier ist namentlich für unsere Zeiten Augustinus die bei weitem höhere Autorität, der de Doctr. 4. 26 sagt: Sunt sane quidam, qui bene pronunciare
\quellseite{160}
possunt, quid autem pronuncient, excogitare non possunt. Quod si ab aliis sumant eloquenter sapienterque conscriptum, memoriaeque commendent atque ad populum proferant, si eam personam gerunt, non improbe faciunt. Wer die Postille treulich braucht, wird überdies an ihr auch frei reden lernen, wenn Amt und Beruf, Gelegenheit und Zeit ihn dazu bringen. Und wer es mit Verstand thut, wird auf diesem Wege die freie Rede mehr lernen als verlernen.

Wenn man übrigens vom Gebrauch der Postille redet, so meint man, versteht sich, nicht, daß der Pfarrer fremde Arbeiten bloß ablesen soll. Eine Gemeinde sollte zu ihrem Segen auch lesen hören können und die sogenannten Lesezeichen, wie Lesebetstunden könnten dann ihren reichen Segen bringen. Man sollte auch die Gemeinden für’s Lesenhören ziehen und darum selbst recht lesen lernen, lesen, als ob man spräche, was jedoch bei gewissem Lesen, insonderheit bei dem liturgischen, nicht angeht. Allein die freie Rede spricht doch ganz eigen an, von Herz zu Herzen, aus einer Seele in die andere dringt die Wahrheit mit eigenthümlicher Macht. Referiert nun aber ein Pfarrer aus der Postille selbstständig, gibt er die fremde Arbeit wieder, nachdem er sie selbst aufgenommen und innerlich verarbeitet hat, strömt aus seinem Munde, was er aufgenommen, frei, herzgründlich, kräftig, so ist Ton und Wirkung anlangend zwischen einer Postillenpredigt und einer eigenen kein Unterschied, wohl aber könnte ein Pastor, der dabei ehrlich und offen handelte, selbst desto größere Ruhe und tieferes Behagen bei seinem Predigen inne werden.

\quellseite{161}
\section{22. Zu wem spricht der Prediger?}

Diese einfache und leicht zu beantwortende Frage wurde in der neueren Zeit gerade von solchen Männern, die geistlich gesinnt waren, auf eine unnatürliche und verkehrte Weise beantwortet. Man soll, sagten sie, mehr in sich und über sich zum Herrn als auf die Gemeinde sehen. Die Absicht dieser Behauptung war, den Pfarrer vor Leidenschaft, dem Pfarrer ein reines Gewissen zu bewahren und ihn beim Predigen in die heiligende Gegenwart Gottes zu stellen. Gewiß eine gute Absicht, und doch so, wie sie ausgesprochen ist, sehr mißverständlich. Beim Gebet hat man die Absicht, auf die Gemeinde zu wirken, während man doch im Namen der Gemeinde oder auch im eigenen Namen zu Gott sprechen, ihn und das Gut im Auge behalten soll, um das man bittet. So verkehrt man in guter Meinung das Gebet und gibt ihm eine schiefe Richtung und Wirkung. Und in consequenter Verkehrtheit will man nun beim Predigen, wo man doch im Namen Gottes die Gemeinde ansprechen soll, den Sinn von der Gemeinde weg in sich und über sich zu wenden. Zwei Theile des Gottesdienstes hat man damit des rechten Zieles beraubt. Die Absicht des Predigers beim Predigen ist keine andere und darf keine andere sein, als die Gemeinde durch’s Wort zu ihrem Heile zu führen, in demselben zu kräftigen und zu gründen. Das muß der Prediger mit allem Ernste wollen, mit aller Kraft vollbringen, und deshalb kann und darf es nicht anders sein, als daß er sich ernstlichst zur Gemeinde wendet, mit
\quellseite{162}
ihr in jene innere geheime Verbindung tritt, die wie ein magnetischer Rapport die Seelen faßt, aber göttlich schön und rein, und ihre Seelen mit dem Worte zu fassen, zu halten, zu leiten sucht. Dabei kann — so groß ist das Geheimniß des inwendigen Lebens — ein mächtiges Seufzen zu Gott um gnädiges Gelingen allerdings vorhanden, und die Seele des Predigers kann und soll vom Geiste des HErrn gehoben, erleuchtet und getragen, voller Andacht und himmlischer Bewegung sein. Aber bei aller Feier in Gott und aller Erhebung zu Gott geht doch die Richtung des Predigers zunächst zur Gemeinde. Es ist eine Art von cherubinischem Leben, man ist Gottes voll und gerade deshalb voll Kraft und Macht der Rede und Richtung nach außen, zur Gemeinde. Ruhe und Wirkung ist beisammen, aber die Wirkung tritt hervor und erscheint, während die Ruhe in der tiefen Stille der Seele thront.

\section{23. Pronunciation}

Beim Vortrag handelt sich’s a) um Schwäche und Stärke der Stimme, b) um Langsamkeit und Schnelligkeit, c) um Declamation, d) um den Kanzelton. Ueber alle diese Dinge ist es leichter Regeln zu geben, als sie zu halten, weil ein jeder Prediger von seiner Eigentümlichkeit gar zu abhängig ist. Niemand kann nach Regeln reden lernen. Jedermann redet, ehe er von Regeln weiß, und benützt die Regel bloß als Correctiv für seine Fehler und Mängel.

Was Stärke und Schwäche der Stimme anlangt, so mißt manchem seine Leibeskraft ein geringeres Maß zu, als
\quellseite{163}
er bedürfte. Eine zu schwache Stimme ist für einen Prediger kein Glück und sollte vornherein von Ergreifung des Predigerberufs abhalten, weil aller Beruf des Predigers mißlingt, wenn er nicht verstanden wird. Es gehört viel sonstige Redegabe, ein reicher und gediegener Predigtinhalt und große persönliche Würde dazu, wenn die Zuhörer sich in die leise und schwache Stimme eines Predigers fügen und durch ihre Stille und Aufmerksamkeit ersetzen sollen, was ihm an durchdringender Kraft fehlt. Diejenigen, welche mit einer starken Stimme begabt sind, haben, was den Aufwand an Kraft betrifft, auf die Größe des Raumes zu sehen, in welchem, und auf die Menge der Zuhörer, an welche sie zu reden haben. Ich wiederhole, nicht bloß auf den Raum, sondern auch auf die Menge der Zuhörer, weil für eine volle Kirche ein größeres Maß von Stimme nötig ist, als für eine leere, wenn man von allen verstanden werden will. Kennt man das Maß von Stimme, welches man Raumes und Publikums halber bedarf, so halte man dasselbe ein, ohne jedoch der natürlichen, von dem Gegenstand und den Umständen erforderten Abwechslung der Stärke und Schwäche der Stimme zu wehren. Wenn man lehrt, ist ein anderes Maß von Stimme zu nehmen, als wenn man warnt, oder straft oder ermuntert, oder wenn man tröstet oder klagt. Das mittlere Maß von Stimme ist eben nur das mittlere. Auf und ab steigt man nach Bedürfniß. Da der Herr befiehlt: „Verkündige meinem Volke sein Uebertreten und Israel seine Fehle“, leitet er diese Vermahnung mit den Worten ein: „Erhebe deine Stimme wie eine Posaune.“ Damit ist jedenfalls
\quellseite{164}
angedeutet, daß nicht unter allen Umständen ein und dasselbe Maß von Stimme anzuwenden sei. In keinem Fall soll ein Prediger das ihm gegebene Maß von Stimme überschreiten, nicht schreien bis zum kreischen oder gar bis zur Heiserkeit. Dadurch bringt man den gerechten Vorwurf der Maßlosigkeit auf sich und hindert sein Amt, ohne daß man sich deshalb des Worts getrösten dürfte: „Der Eifer um dein Haus hat mich gefressen.“ Noch weniger aber soll ein Prediger den Mangel an Gedanken durch großen Aufwand von Stimme zu bedecken suchen. „Viel Geschrei und wenig Wolle“ wird deshalb nicht verziehen, weil man etwa geneigt ist, viel Geschrei bei viel Wolle zu verzeihen.

Wie mit dem Aufwand der Stimme, ebenso ist es mit Langsamkeit und Schnelligkeit der Sprache. Es ist unerträglich in immer gleichem Tone, sei es nun immer langsam oder immer schnell, reden zu hören. Ein der Sache angemessener, sich beim Predigen von selbst ergebender Wechsel ist das Wünschenswerthe. Langsam anfangen, mit der Gemeinde erwärmen, mit der erwärmten, ergriffenen Gemeinde schnelleren Gedankenflugs vorwärts eilen, nach Befund der Sache die Flügel wieder langsamer schwingen, wie man ja auch nach Befund der Sache lauter und gemäßigter spricht, das empfiehlt sich von selbst und kommt auch von selbst bei dem Prediger, bei dem die Predigt ein Stück Leben ist. Insbesondere dürfte bei einer periodischen Sprache eine größere Behendigkeit des Vortrags, um nicht zu sagen Schnelligkeit, die Klarheit der Auffassung fördern, während ein langsamer Vortrag das
\quellseite{165}
Gedächtnis für eine solche Redeweise zu wenig unterstützt. Sonst und im Allgemeinen wird die Vermahnung der Alten, nicht zu schnell zu reden, ganz in der Ordnung und ganz praktisch sein.

Ist die angegebene Regel für das Maß der Stimme und der Zeit nur als Correctiv zu gebrauchen, so würde man eine Versündigung gegen die Predigt auf sich laden, wenn man den Predigern zumuthen sollte, die Regeln der Declamation in anderer Absicht und in anderem Maße anzuwenden. Es ist keine Frage, daß man eine jede Sache angemessen und möglichst zweckmäßig vortragen soll, und dazu hilft ohne Zweifel auch das Studium der Redefügung, des Tons, der Stimme und ihrer Modulation und der Gebärde. Allein es scheint doch für die geistliche Rede unpassend zu sein, nach Regeln reden zu wollen, während es allerdings ganz in der Ordnung ist, einschleichende Fehler kraft vorhandener Regeln zu fassen und abzuthun. Es gibt eine geistliche Redekunst, aber sie gleicht, wenigstens was das Aeußerliche, die Pronunciation belangt, den zehn Geboten Gottes, bei denen es mehr Verbot als Gebot gibt.

Ein besonderes Leiden der Prediger ist der sogenannte Kanzelton, ein Produkt nicht bloß des Organs, sondern auch des großen Inhalts, welcher vorgetragen wird, ja auch der Art und Weise, wie sich das Herz des Predigers gegen denselben verhält.\footnote{Sehr oft kommt er aber auch aus bloßer Nachahmung. Ein Student oder Candidat hört einen Prediger, der ihm gefällt: den macht er nach und gewöhnt sich dran, so daß er seine andre Natur wird. Er kann dann später nicht mehr anders reden, als er sich’s angewöhnt hat. So oft er predigt, hört er auf, er selbst zu sein und spielt die Rolle seines Vorbildes, zieht die andere Natur an. Trauriges Loos so vieler. — Und wie oft wird dann mit der Zeit auch der Inhalt nur Rolle und Maske, so daß der Ernst des Lebens grade da aufhört, wo er vor allem hingehörte, beim heiligen Amte.} Fast ist ein Kanzelton unvermeidlich, und
\quellseite{166}
er ist besser als der Conversationston, in welchem man eben auf der Kanzel nicht reden darf, weil ὁμιλεῖν und conversieren von einander himmelweit verschieden sind, und die ewige Weisheit und Wahrheit keine Materie für eine pure Conversation, am wenigsten aber auf der Kanzel ist. So wahr nun aber auch das ist, so ist es doch die Pflicht eines jeden Predigers, gegen ungeziemenden Kanzelton, gegen dessen Monotonie, wie gegen die oft ungeheuerlichen Abwechselungen, Höhen und Tiefen zu kämpfen und treue Freunde aufzufordern, daß sie Barmherzigkeit üben und durch offenen Tadel diese Verunzierungen des göttlichen Wortes abthun helfen. Am meisten wird jedoch dem unschönen Kanzelton gewehrt werden, wenn der Prediger sich bei jeder Predigt vergegenwärtigt, was er eigentlich zu thun hat, wenn er sich erinnert, daß er einen göttlichen Auftrag an die Gemeinde auszurichten, sie im Namen Gottes anzusprechen und mit ihr zum Heile ihrer Seele zu reden hat. Je weniger aus dem Predigen eine Gewohnheit wird, je weniger der Prediger seine Vorträge einen nach dem andern, wie man zu sagen pflegt, bloß „ablegt“; je mehr er im Bewußtsein seines Berufes und Auftrags bleibt, desto angemessener und einfacher, richtiger und ungezwungener, desto weniger manieriert wird sein Vortrag sein. Wie jener Busch im Feuer
\quellseite{167}
nicht brannte, noch verbrannte, sondern nur verklärt stand; so muß das Feuer des Amtes die natürliche Art und Weise zu reden nicht absonderlich machen, sondern verklären. Die Sprache eines Predigers muß beim amtieren ungesucht im schönsten Lichte erscheinen, dessen sie fähig ist.

\section{24. Gesticulation}

Melanchthon will über die Gesticulation nichts sagen:\footnote{Dagegen s. Bahrdt’s Rhetorik für geistliche Redner. 1778. Nur durch die Lehre von der Action ausgezeichnet. Zwei Drittheile des Buches handeln von der Action.} es hätten sich von den Regeln der Heiden abweichende Bräuche und Sitten gebildet. Vielleicht ist es auch wirklich am besten von der Gesticulation wenig zu sagen. Wenn die Alten behaupteten, es sei nicht geziemend, die Hand höher als bis zum Auge zu heben, so wird man leicht in der gegebenen Regel eine Grenzlinie des Schönen für die ruhige Rede anerkennen. Allein, wer würde sich wohl herbeilassen, für alle Fälle, für jede Gegend der Rede, für jeden Redner und für jeden Gegenstand diese Regel als maßgebend anzunehmen. Nicht alles schickt sich immer, überall und für alle. Ein jeder hat seine eigene, für ihn, wenn auch nicht für andere ziemliche Weise, die man ihm wohl auch lassen kann und muß, wenn er nur nicht über sein eigenes Maß hinüberschreitet und dadurch auffallend wird, also möglicher Weise die Aufmerksamkeit stört und erschwert und dem Zweck seiner Rede entgegen tritt. Ein
\quellseite{168}
aufrichtiger Freund thut hier, was kein System von Regeln vermag, auch wenn es wirklich ein durchgreifendes gäbe, wiewohl man nicht leugnen kann, daß treue Freunde, die es wagen, dem Freunde eine unschöne Aeußerlichkeit, also eine scheinbare Kleinigkeit zu bereden, sehr selten zu finden sind. Oft weiß alle Welt die Unart eines Predigers, nur er nicht. Er kann darüber altern und sterben und doch nicht erfahren, was er, wenn es ihm rechtzeitig gesagt worden wäre, gefaßt, gehaßt und gelassen hätte.

\section{25. Die Abkündigungen}

Es ist jedermann bekannt, daß nach der Verkündigung des göttlichen Wortes und gesprochenem „Amen“ von dem Prediger noch alles dasjenige abgekündigt wird, was zur äußerlichen kirchlichen Führung der Gemeinde gehört. So z. B. werden Geschenke und Gaben, Schulconferenzen, Armen- und Kirchenverwaltungs-Sitzungen u. dergl. von der Kanzel abgekündigt, Brautleute proclamirt, besondere Gottesdienste angesagt u. dergl. Alles das erfordert keine besondere Vorbereitung und die Mühe der Verabfassung ist klein, da man sich verordneter oder herkömmlicher Worte und Ausdrücke bedient. Es kommt nur darauf an, Alles recht genau und formgetreu zu schreiben und abzukündigen. Außer diesen bloß formalen Dingen gibt es aber noch anderes, was man unter die Reihe der Abkündigungen zu stellen pflegt, wie z. B.: es werden eintretende Festtage abgekündigt, es wird zur Barmherzigkeit in diesem oder jenem Falle aufgemuntert und eine Collecte angesagt,
\quellseite{169}
Beicht- und Abendmahlstage bestimmt, besondere Ereignisse mit dazu gehörigen Warnungen oder Ermahnungen bekannt gegeben, und was dergleichen ist. Solche Abkündigungen der edleren Gattung erregen allgemeine Aufmerksamkeit, es herrscht tiefe Stille, auch die Schläfer wachen, nichts wird überhört und man kann sicher sein, daß eine wohlgerathene Abkündigung nicht bloß von der ganzen Versammlung beachtet, sondern auch in die Häuser heimgetragen wird, wie es denn überhaupt keine Art der Veröffentlichung gibt, die in so großem und allgemeinem Ansehen steht, wie die kirchliche Abkündigung. Je länger man im Amte ist, desto mehr drängt es sich auf, und desto mehr muß man merken, daß man in den kirchlichen Abkündigungen ein sehr zu beachtendes Mittel der Einwirkung auf die Gemeinde besitzt. Man wird daher auch je länger je mehr geneigt, den Abkündigungen besonderen Fleiß zuzuwenden und sie mit großer Sorgfalt zu verabfassen. So wie man aber auf diesem Punkte angekommen ist, sieht man auch leicht ein, daß die Abkündigung ins Reich der Homiletik, und zwar unter die Casualreden gehört. Wird auch um des Ausdruckes Casualrede willen niemand auf die Meinung kommen, daß eine Abkündigung eine Rede in voller Form sein müsse, so wird andererseits doch auch jedermann leicht einsehen, daß sich in einer wohlgelungenen Abkündigung die Gestalt der Rede wiederspiegelt und daß es bei der Verabfassung fördersam ist, daran zu denken. Denken wir uns einen jungen Geistlichen, der diese unsere Sätze liest und sich vornimmt, die Erfahrung zu machen, so wird er nicht lange den Versuch machen, ohne
\quellseite{170}
daß sich auch ihm die Brauchbarkeit und der Segen der kirchlichen Abkündigung empfiehlt; er wird sie je länger je mehr pflegen und lieben. Vielleicht geht es ihm wie dem Schreiber dieser Zeilen. Da er jung war, galt ihm die Abkündigung gering. Sein Abkündigungsbüchlein war klein, und wurde bloß zur Unterstützung des Gedächtnisses gehalten; nun er aber alt geworden ist, steht er zuweilen stundenlang am Pulte, um seine Abkündigungen zu verabfassen, trägt sie mit allem Fleiße in ein stattliches Buch ein, und wenn er das Buch unter dem Arm zu seiner Kirche geht, freut er sich oft der Abkündigungen wie ein Schütze, wenn er eines auserwählten Pfeiles sicher ist. Manchmal, wenn er unter dem Predigen merkt, daß seine Rede nicht fährt, tröstet er sich, daß er noch seine Abkündigungen übrig hat. Er spart Zeit und Lust und Kraft, um sie für die Abkündigungen übrig zu haben, und weiß, daß er manche Lücke seiner Predigten durch manche Abkündigung gebüßt hat.

In Pastoral-Anweisungen findet man die Abkündigung nicht beachtet, man sucht vergeblich darnach. Um so mehr thut es dem Verfasser wohl, in Zwickenpflugs „Grundzügen einer Pastoral im Geiste Christi und seiner heiligen Kirche“ (Regensburg 1844) pag. 213 bis 221 einen Abschnitt zu finden, der ganz dasselbige ins Auge faßt, was er selbst meint. Der Abschnitt ist überschrieben: „Von besonderen Erinnerungen, welche ein Seelenhirt zu verschiedenen Zeiten im Jahr geben soll oder kann.“ Unter den besonderen Erinnerungen ist ungefähr dasselbe verstanden, was oben Ankündigungen der edleren
\quellseite{171}
Art genannt ist, und wenn Zwickenpflugs Worte für uns Protestanten nicht durchaus passend sind, so werden doch die folgenden Mittheilungen aus denselben für diejenigen von Werth sein, die überhaupt auf den Gegenstand eingehen wollen.

Zum Eingang sagt Zw.: „Die Erinnerungen sind oder sollen sein, eine jede ein Wort zu seiner Zeit; gut angebracht thun sie zuweilen mehr Wirkung als die gewöhnlichen Predigten, Ermahnungsreden und Katechismuslehren.“

Hierauf kommen einige Regeln für die Verabfassung der Erinnerungen:

\begin{quote}
„I. Der Seelsorger muß sich auf die Erinnerungen, die er zu geben hat, wohl vorbereiten, d. i. überlegen, was und wie er es zu sagen hat: es gehört dazu ein eigenes Talent.“
\end{quote}

\begin{quote}
„II. Er muß seine Erinnerungen auf gute Gründe, nicht aber auf Drohungen und Schmähungen stützen, und zwar in der Art, daß er eine Stelle oder ein Beispiel der heiligen Schrift anführt, selbige erklärt und darüber eine Nutzanwendung macht.“
\end{quote}

\begin{quote}
„III. Er soll in seinen Erinnerungen nicht weitläufig sein und sie nicht zu lang machen.“
\end{quote}

\begin{quote}
„IV. Alle Erinnerungen, die er macht, müssen ganz besonders mit dem Salze der Klugheit gewürzt, also ernst und gütig, nachdrücklich und zärtlich, aufrichtig und auf Folgsamkeit vertrauend sein.“
\end{quote}

\begin{quote}
„V. Haben Erinnerungen die gehoffte Wirkung nicht gemacht, so muß er sie wiederholen und seinen Kummer darüber
\quellseite{172}
aussprechen; haben sie aber gefruchtet, so muß er seine Zufriedenheit damit bezeugen.“
\end{quote}

Nach diesen Regeln belehrt Zw. die jungen Seelsorger, sie sollen niemals über Schäden und Uebel reden und Erinnerungen geben, die nicht offenkundig sind, wohl aber über öffentliche Aergernisse und Schäden, denen Widerstand geleistet werden muß, gerade da aber mit aller Vorsicht und mit gehöriger Vorbereitung, und nie in Gegenwart von Personen, die nicht zur Gemeinde gehören, welches letztere freilich in mancher Gemeinde kaum zu erreichen sein wird. Wenn er außerdem meint, die Erinnerung solle nie in Betreff geringer Dinge geschehen, sondern immer nur in Betreff solcher, die von großem Belang in Betreff des Heils seien, so wird man das bestreiten können. Es ist auch oft etwas Geringes werth ins Licht einer Erinnerung gestellt zu werden, weil oft das Geringe dem Samen gleicht, aus welchem Großes hervorwächst etc.

Die Zeit anlangend, warnt Zw. ganz richtig, zu oft Erinnerungen zu bringen. Ebenso richtig ist seine Mahnung für jede Erinnerung die schickliche Zeit zu wählen. Manche Erinnerungen können mehreremal im Jahre wiederholt werden, andere dürften nicht öfter als einmal kommen. Auch meint er, zuweilen könne man die Erinnerung anstatt einer Predigt oder Rede oder Katechismuslehre halten, man könne sie vor oder nach dem Gottesdienste, vor oder nach dem gewöhnlichen Vortrage halten.

Nach diesen Unterweisungen von allgemeiner Art bringt
\quellseite{173}
Zw. Anleitungen zu Erinnerungen für die Adventszeit und überhaupt die Feste des Kirchenjahres, so wie für den Anfang der Festzeiten in ganz kurzen, aber dennoch in ihrer Art instructiven Worten. Ebenso Erinnerungen für die Frühlingszeit, und zwar für die Feldarbeiter und die Hirten, für die Heu-, Ernte- und Saatzeit, für die Winterzeit, für eingetretene Unglücksfälle und Trübsale, bei Hagelschlag und schädlichem Reif, ansteckender Krankheit, einem jähen Todesfall, einer Feuersbrunst, einem öffentlichen Aergerniß, bei Zwistigkeiten in der Gemeinde, nach dem Hintritt des Seelenhirten, beim Antritt des Pfarramtes. So ganz kurz diese Anweisungen sind, so sind sie doch ganz geeignet, einem jungen Seelsorger Lust zu machen, die Erinnerungen zu pflegen. Wir wollen ein ganz kurzes Beispiel geben:

Die Erinnerung für die Feldarbeiter in der Frühlingszeit soll zum Inhalt haben: die Feldarbeiter möchten ihr Morgen- und Abendgebet nicht unterlassen, morgens und öfters unter Tags alles was sie thun „durch Erweckung einer guten Meinung heiligen.“ Beim Hinausgehen aufs Feld, während ihrer Arbeit und beim Nachhausegehen möchten sie geistliche Lieder singen, von Zeit zu Zeit an die letzten Dinge denken, keinen Eingriff in fremdes Eigenthum thun, niemandem Schaden zufügen, nicht fluchen noch schwören, das Ungemach des Wetters, die Last der Arbeit, die üble Laune des Mitarbeiters geduldig tragen zc. Ebenso werden die Hirten beim Beginn der Frühlingszeit ermahnt, sich während ihres Hütens mit einer anständigen Arbeit zu beschäftigen, nicht mit Leuten vom anderen
\quellseite{174}
Geschlechte zusammen zu hüten, zur Zeit des Gottesdienstes, wenn sie ihn nicht besuchen können, den Katechismus zu lesen und geistliche Lieder zu singen, zuweilen zum Sacramente zu gehen und zur Kirche, und dergleichen.

Es mag ein jeder von solchen Abkündigungen und Erinnerungen sich seine eigenen Gedanken machen, der Verfasser aber würde sich gestraft haben, wenn er nicht in dieser neuen Auflage seines Buches so viel gesagt hätte als er gesagt hat, weil ihm die Erfahrung eine zu tiefe Ueberzeugung gegeben hat, daß die Abkündigung oder Erinnerung der Beachtung, des Fleißes und der Sorgfalt werth sei, eben weil ihr, wie schon gesagt, so mancher Segen folgt.

\chapter{Ueber die Katechese}

\section{26. Ueber die Katechese}

Κατήχησις talem indicat institutionem, cujus echo sive sonus aures audientis personat, sagt Deyling und weist zugleich den Gebrauch des Wortes bei den Griechen nach. Im Worte läge deshalb eigentlich ein „mündlicher Unterricht“ angedeutet, eine Lehre vom Mund zum Ohr, wenn man nicht lieber sagen will vom Mund zu Mund. Denn in der letzteren Bezeichnung läge zugleich die heut zu Tage gebräuchliche Auffassung der Katechese, wonach sie ein Unterricht in der christlichen Lehre wäre, welcher gesprächsweise von Mund zu Mund gegeben wird. Das Wort kommt bei den LXX,
\quellseite{175}
und im N. T. vor: Gen. 18, 19. Exod. 12, 26. Deuter. 6, 2. Luc. 1, 4. 1 Cor. 3, 2. Hebr. 5, 12—14.; 6, 11. Act. 18, 25. Die Stellen werden sich auf die genannte Bedeutung zurückbringen lassen. Dabei war es ursprünglich nicht unterschieden, ob man den Unterricht akroamatisch oder dialogisch gab. Von beiderlei Arten des Vortrags ist der Name Katechese gebraucht, beiderlei Arten sind dem Alterthum bekannt. So wie man unter dem Namen des heiligen Basilius eine ganz in Frage und Antwort gehaltene ἔκθεσις πίστεως ἐν συντόμῳ hat, so hat man auch fortlaufende Reden des heiligen Cyrillus von Jerusalem, welche den Namen Katechesen führen, und mitten unter den Homilien des heiligen Chrysostomus findet man eine κατήχησις πρὸς τοὺς μέλλοντας φωτίζεσθαι. Augustinus de catechizandis rudibus c. 13. will, daß sich der Katechet durch Fragen überzeuge, ob er auch verstanden sei. Er vereinigt also beides, das akroamatische und das dialogische Element, er erkennt beide als katechetische Elemente an. Und das war auch in der That später der Fall, namentlich in der lutherischen Kirche. Man vergleiche z. B. die vielverbreiteten vortrefflichen Katechismuspredigten der Brandenburgischen Kirchenordnung und die zu ihrem Gebrauche gegebenen Anordnungen, und man wird sich überzeugen, daß unsere Väter von der Verbindung des akroamatischen und dialogischen Elements ausgingen. Diese Verbindung ist auch ganz naturgemäß. Der christliche Unterricht gründet sich nicht auf die notiones innatae oder auf allgemein anerkannte und zugestandene
\quellseite{176}
Sätze. Niemand kann aus solchen in der Weise des Sokrates, wie sie sich in Xenophons Memorabilien findet, die christliche Lehre entwickeln, und Sokratik im Sinne der Entwickelung aus Eingebornem und aller Welt Bekanntem kann es auf dem Felde der christlichen Lehre gar nicht geben. Bei der christlichen Katechese muß erst mitgetheilt werden, worüber man sich mit dem Schüler gesprächsweise zu verständigen und von seiner Auffassung zu überzeugen hat. Was für ein langsamer, unnützer, Zeit und Kraft fruchtlos verschwendender Weg es ist, den Inhalt christlicher Lehre gesprächsweise mitzutheilen, davon kann sich ein jeder, der es versuchen will, aus eigner Erfahrung überzeugen\footnote{Es gibt zwar articuli mixti, nicht bloß articuli puri, und bei denen kann man immer auch etwas heraussokratisieren und herauskatechisieren. Und auch bei den articulis puris läßt sich, wenn einmal der Satz gegeben ist, durch die Thätigkeit des Katecheten und die eingehende Theilnahme des Schülers oft unvermuthet viel Licht und Erkenntniß durch Frage und Antwort entwickeln.}. Die Katechese muß deshalb die beiden Elemente vereinigen. Sie kann nicht anders und eben deshalb kann und muß es bis auf den heutigen Tag katechetische Reden eben so wohl geben als katechetische Gespräche. Jene geben Inhalt und Maß des zu Lernenden, diese zeigen, wie man den Inhalt dem Hörer und Schüler nahe bringt. Und wenn deshalb die Katechetik die Wissenschaft von der Katechese ist, so muß sie sich ebensowohl mit der katechetischen Rede als mit dem katechetischen Gespräch befassen, denn sie muß uns ja Inhalt, Maß und Form des Unterrichts vorlegen.

Man könnte freilich sagen, daß heut zu Tage die kate%
\quellseite{177}
chetische Rede wenig gebraucht werde; allein das liegt doch am Ende nur daran, daß man überhaupt nicht mehr so viel von Mund zu Munde lehrt, daß man aus allerlei Gründen von verschiedenem Werth dem Schüler anstatt der mündlichen Rede einen gedruckten Unterricht (d. i. Katechismus, was eben Unterricht und metonymisch Unterrichtsbuch bedeutet) in die Hand gibt, der Inhalt und Maß bestimmt. Dies ist ein Surrogat der mündlichen Belehrung, welche dem katechetischen Gespräch vorangehen soll. Es ist so, aber es sollte nicht sein, weil die mündliche Belehrung, wenn sie rechter Art ist, durch Schrift und Druck nicht ersetzt werden kann. Daher sollte der Lehrer dennoch, und obwohl die Schüler den Katechismus gedruckt vor sich haben, jedes Pensum erst in lebendiger Rede vortragen, dann lesen lassen, dann fragen. Hierdurch entstände eine Vereinigung dreier verschiedener Unterrichtsmomente, welche desto sicherer zum Ziele führen würden.

— Man kann gegenüber der Behauptung von der Nothwendigkeit des akroamatischen Unterrichts bei der Katechese, auch nicht auf die dialogische Form der meisten Katechismen hinweisen, denn kein Katechismus hat die Form des katechetischen Dialogs und keiner soll sie haben, sondern die Form des Katechismus ist rechtmäßiger Weise die des Examens, der Rechenschaftslegung, des Bekenntnisses, und es läßt sich auch diese Form, die allerdings dem Dialoge verwandt ist, beseitigen, wie es am Tage ist, da es Katechismen genug gibt, welche in fortlaufenden Paragraphen ohne alle Frage verabfaßt sind. Desto weniger hebt aber der gedruckte Katechismus
\quellseite{178}
die zusammenhängende Belehrung, das mündliche, einleitende und vorbereitende Wort des Katecheten auf. — Mag aber auch ein Katechismus in Frage und Antwort, oder in Paragraphen abgefaßt sein, immer muß er vor allem denen, so Lehrern wie Lernenden das rechte katechetische Maß zu überliefern.

\section{27. Inhalt der Katechese}

Die Katechese hat Belehrung zur Absicht, woraus sich von selbst ergibt, daß alles, worüber eine Belehrung gegeben werden kann, sich auch zum Inhalt einer Katechese eignet. Das weite Feld der katechetischen Möglichkeit, welches dadurch entsteht, findet aber doch auch wieder seine Grenzen. Man muß bedenken, daß sich ein christlicher Katechet nicht mit allem befaßt, worüber man möglicher Weise eine Katechese halten könnte, sondern bloß mit dem, worüber die Kirche zu katechisieren pflegt, was im Verlauf und nach der Erfahrung so vieler Jahrhunderte vorzugsweise für katechetischen Inhalt erkannt wird. Man könnte z. B. Katechesen über Erfahrungen des inwendigen, verborgenen Christenlebens halten, und zuweilen wird sich auch der Katechet mit Segen auf das Gebiet der geistlichen Erfahrung begeben können. Es gibt ja gewisse Erfahrungen, welche jeder Christ macht und machen muß. Aber man wird deshalb im allgemeinen die Behauptung aufstellen: die Katechese habe ihren Stoff besonders auf den Feldern und in den oft wunderlich verschlungenen Gängen der innerlichen
\quellseite{179}
Erfahrung zu holen? Sie hat es mit dem, was gewiß ist, was in Gottes Wort klar und einfach vorliegt, mit dem, was zur Seligkeit nöthig ist, zu thun. Die besonderen Tiefen und Fernen und Höhen der christlichen Erkenntniß und Erfahrung sind schon deshalb von der Katechese auszuschließen, weil die, welche katechetisch unterrichtet werden, bei uns dem jugendlichen Alter angehören, bei welchem erst Grund zu legen ist. Es ist ὁ τῆς ἀρχῆς τοῦ Χριστοῦ λόγος, οὐχ ἡ τελειότης, das θεμέλιον ist es, wovon Hebr. 6, 1—2. steht, was der Katechet vorzutragen und einzuprägen hat\footnote{Die merkwürdige Stelle enthält, wie man sagen könnte, die Hauptstücke des apostolischen Katechismus, deren Vergleich mit dem kleinen lutherischen Katechismus so interessant ist und eine gewisse Verwandtschaft des letzteren mit jenem nachweist. Ἀφέντες τὸν τῆς ἀρχῆς τοῦ Χριστοῦ λόγον, ἐπὶ τὴν τελειότητα φερώμεθα, μὴ πάλιν θεμέλιον καταβαλλόμενοι μετανοίας ἀπὸ νεκρῶν ἔργων καὶ πίστεως ἐπὶ τὸν θεόν, βαπτισμῶν διδαχῆς, ἐπιθέσεώς τε χειρῶν, ἀναστάσεώς τε νεκρῶν καὶ κρίματος αἰωνίου.}. Und hiegegen streitet keineswegs der Unterschied, den man zuweilen zwischen der Katechese der Kinderlehre und der Christenlehre gemacht hat. Denn die Christenlehren nehmen zwar Rücksicht auf die anwesenden erwachsenen Christen, aber auch ihr Hauptaugenmerk ist und bleibt eben doch der Kinderhaufen. Und wenn auch nicht, wie selten ists, daß man mit den Erwachsenen weiter gehen kann als mit den Kindern, und wie gewöhnlich ist es im Gegentheil, daß die Erwachsenen auf eine geringere Erkenntnißstufe zurückgesunken sind, als die der Kinder ist, so daß sie der einfachsten Repetition bedürfen.

\quellseite{180}
Jener Alte, welcher auf eine Frage des Katecheten die Antwort gab: „Da ich so alt war, wie diese Kinder, wußte ich alles, jetzt weiß ich nichts mehr“ — steht nicht vereinzelt, er könnte Legion heißen. Daher wird der katechetische Inhalt so ziemlich sich gleich bleiben unter allen Umständen.

Es wird also biblische Geschichte und Katechismus Hauptinhalt der Katechese genannt werden müssen, und selbst der nach gesunder Ansicht von dem allgemeinen Unterricht zu sondernde Confirmandenunterricht unterscheidet sich, was den Inhalt anlangt, bei den meisten Geistlichen wenig von dem übrigen katechetischen Unterricht und findet seinen characteristischen Unterschied mehr in der besondern Absicht und Beziehung. So sollte es, wie wir unten hören werden, nicht sein; aber leider ist es so.

Man könnte es bedenklich finden, die heilige Geschichte mit in den Bereich der Katechese hineinzuziehen; allein die christliche Lehre ist dem größten Theile nach reine Abstraction von der Geschichte, von den großen Heilsthatsachen Gottes, oder geradezu Geschichte. Es kann niemand ohne Geschichte die christliche Lehre lernen, und wenn eines von beiden, entweder die Lehre oder die Geschichte wegbleiben müßte; so würde man an der Lehre weniger verlieren als an der Geschichte, welche den Inhalt der Lehre bildet und selbst lehrhaft ist.

\quellseite{181}
\section{28. Maß des katechetischen Inhalts}

Schon darin, daß sich gegenwärtig die Katechese nicht mit der christlichen Erkenntniß im allgemeinen, sondern mit der zum Heile nöthigen, der Jugend darzubietenden Erkenntniß befaßt, liegt etwas Maßgebendes. Aber in diesem Sinne das richtige Maß zu treffen, ist nicht Kleines, sondern Probe großer katechetischer Weisheit. Es liegt in der That weniger an der Methode und Form der Katechese, als am rechten, nicht überfließenden, sondern genügenden Maße. Wenn man gleich weiß, daß man biblische Geschichte und Katechismus zu lehren habe, weiß man damit so gut, wie noch nichts, wenn man das Maß der Lehre nicht kennt. Man weiß nicht, was man geben soll, wenn man nicht allezeit weiß, wie viel man aus dem Reichthum des Ganzen zu geben hat. Die männliche Tugend der σωφροσύνη besteht hauptsächlich in Bewußtsein und Einhaltung des rechten Maßes, wie bei allen Berufsarten, so namentlich und ganz vorzüglich bei der des Katecheten.

Wenn nun vom katechetischen Maße die Rede ist, so begreift man darunter zweierlei, nemlich 1) daß Maß des gesammten Unterrichts und 2) das Pensum der einzelnen Unterrichtsstunde. Man könnte beides auch in die Frage zusammenfassen: was soll im Ganzen geleistet werden, und wie viel muß man in jeder einzelnen Unterrichtsstunde leisten, wenn man mit Sicherheit das ganze Maß erfüllen will. Im Allgemeinen, wie in Betreff der
\quellseite{182}
einzelnen Pensa wird man nicht das möglichst große, auch nicht das möglichst geringe, sondern ein mittleres Maß festzustellen haben. Doch wird sich dies mittlere Maß selbst wieder mehr zum geringen, als zum großen neigen müssen, weil zwar die mittleren Köpfe die Mehrzahl in der Welt zu sein pflegen, der Lehrer aber, welcher das Maß für die Schüler zu bestimmen hat, sich oft verführen läßt, das mittlere Maß für kleiner zu nehmen, als es ist, während die Mehrzahl der mittleren Köpfe durch allerlei zufällige Umstände heruntergedrückt und in der Entfaltung der vorhandenen Gabe gehemmt zu werden pflegen. Das mittlere Maß muß ein kleines Maß sein. Soll ein Fehler gewagt werden, so sei es lieber der des zu geringen Maßes, als der des zu großen. Ein geringes Maß, zu dem man sicher gelangt, ist höher zu schätzen, senkt sich tiefer ein und wird fruchtbarer, als ein größeres, bei dessen Erstreben man unsicher wird, in Aufregung und Zagen geräth\footnote{Als ich den ersten Theil meines Hausbuches ausarbeitete (anno 1845), glaubte ich das Maß der in den „Fragen und Antworten zu den sechs Hauptstücken“ gegebenen Erkenntniß recht gering genommen zu haben. Seitdem, also nach 20 Jahren, find ich’s nicht mehr gering, sondern ziemlich groß und bin mit geringerer Leistung zufrieden. Man kann in unsern Schulen, wie sie sind, wenigstens in den Landschulen, nicht genug herunter steigen. — Ich habe aber auch je länger, je mehr die Erfahrung gemacht, daß von der christlichen Erkenntniß schon ein geringer Theil alle Weisheit der Welt übersteigt. Es darf nur der heilige Geist das wenige, das ein Kind gelernt hat, segnen, so schlägt es aus, erleuchtet, kräftigt und heiligt zum Erstaunen derer, die Zeugen davon sein dürfen. Was ist im Christenthum klein?}.

\quellseite{183}
Was die besondere Anwendung dieser Grundsätze anlangt, so findet man in der That für den Katechismusunterricht einen weit betretenen Weg als für den in der biblischen Geschichte. Der kleine lutherische Katechismus und seine in der Vorrede gelehrte, ganz auf Einhaltung eines rechten Maßes gegründete Methode empfiehlt sich dem Katecheten desto mehr, je mehr die Erfahrung wächst. Die Vorrede zeigt die Pensa im Großen, d. i. die Curse, wie man vom Einprägen des Wortlautes zur katechetischen Verständigung und zur Zurückführung des katechetischen Wortes auf die heilige Schrift fortschreiten soll. Selbst die einzelnen Pensa könnte man durch die Eintheilung in Hauptstücke und durch die Unterabtheilungen dieser angedeutet finden, wiewohl doch für den ausführlichen, verständigenden Unterricht und für die Schriftbegründung weitere Sorge zu tragen wäre. Das Nürnbergische Kinderlehrbüchlein, welches der Form nach gewiß nicht mehr anwendbar ist, könnte dennoch für das Maß der einzelnen Pensa eine gewisse Anleitung geben. Wenigstens ist das Maß der einzelnen Pensa in dem Büchlein klein genug. — In den Spruchbüchern ist insgemein das Maß zu klein und zu ungenau bestimmt, und doch sollte in ihnen das geistliche Brot rein und klein vorgeschnitten sein; bis ins Einzelste sollte der Fingerweis vorliegen. Dahin könnte man es auch bringen, wenn man statt sich zu sehr zu besinnen, welches das rechte mittlere Maß sei, ein geringeres Maß feststellen möchte, da man ja, wenn dies erreicht ist, leicht mehr thun könnte, es aber nicht eben so leicht ist, sich zu helfen, wenn man merkt, daß die Aufgabe zu groß sei.

\quellseite{184}
Was die biblische Geschichte anlangt, so ist man, wie gesagt, weniger berathen als bei dem Katechismus. Daß man erst Geschichten aus der Geschichte vortragen soll, um später die Geschichte, d. i. den Zusammenhang der Geschichten zu lehren, ist eine gute, aber für die Praxis zu unbestimmte Lehre, denn es müssen doch die Geschichten so ausgewählt werden, daß es später leicht wird, an ihnen den geschichtlichen Zusammenhang zu zeigen, den eigentlichen geschichtlichen Unterricht anzuknüpfen. Ja es ist vorne herein aller Vortrag der einzelnen Geschichten ganz von dem eigentlichen geschichtlichen Unterricht abhängig, und zwar in so hohem Maße, daß erst noch die Frage ist, ob nicht schon bei dem Unterricht der jüngeren Kinder eine ganz kurze Uebersicht des geschichtlichen Verlaufs sammt einer kleinen Summe von Zahlen neben dem Vortrage der Geschichten hergehen soll, d. i. ob nicht doch dem jungen Alter im Grunde schon Geschichte in Geschichten vorzutragen sei\footnote{Was oben als noch in Frage stehend hingestellt ist, hat für den Schreiber dieses aufgehört in Frage zu stehen. Er hat die nöthige Erfahrung indes (während der 8 Jahre, welche zwischen Aufl. I. und II. dieser Schrift inne liegen) gemacht, daß die kurze Uebersicht des Reiches Gottes für das junge Kind faßlich, anziehend und ergreifend ist, wie die Geschichte Josephs oder Davids. Es läßt sich ganz wohl Stufen aufwärts von der Uebersicht zur Einsicht gehen. Auch ist eine kleine Summe von Merkzahlen selbst für den kindlichen Geist ein Ordnungs- und Klarheitsprincip, welches ganz nothwendig vorhanden sein muß. Und wie leicht ist hier das Nöthige erreicht! Wie lebendig wird dem Kinde jede einzelne Geschichte durch Einordnung in die Uebersicht der Reichsgeschichte Gottes und unter die herrschenden Merkzahlen. Wer es versucht, wird schnell zur Einsicht und Zuversicht gelangen und seinen Schülern „Geschichten aus der Geschichte,“ d. i. im Zusammenhang des großen Ganzen geben. — Geschichten aus der Geschichte, oder Geschichte in Geschichten!}. Der vorhandenen Hilfsmittel sind unzählige. Wer
\quellseite{185}
kennt nicht die Calwer, die Baseler, die Zahn’sche, Buchdrucker’sche biblische Geschichte? Aber ein Lehrbuch, das einfach, kurz und leicht mehr als Regel und Wegweiser durch den für Lehrer und Schüler in der heiligen Schrift oder einem biblischen Geschichtsbuch vorliegenden Stoff Curse und Pensa anzeigte, und gewisse Tritte zum Ziele thun lehre, scheint noch zu mangeln, wenigstens hat sich keines die allgemeine Anerkennung verschafft.

Des Beispiels und allgemeinen Maßes wegen dürfte einige Empfehlung verdienen: „Kirchengeschichte des Alten und Neuen Testaments, zum Gebrauch für Schulen und Confirmandenstunden, von Fr. Dümichen. Breslau bei Trewendt.“

\section{29. Vorbereitung zur Katechese}

Bei einem jeden Unterricht ist es von der höchsten Wichtigkeit, daß der Lehrer ganz Meister seines Stoffes sei. Dies aber ist ohne Vorbereitung eine unmögliche Sache, und zwar geziemt sich Vorbereitung auf jede Lehrstunde, jedes Lehrpensum. Der Lehrer muß wissen, was er in jeder Stunde zu geben, was er zu leisten hat. Darauf, und daß ers leisten könne, gründet sich alle Ruhe, Sicherheit, Leichtigkeit und Lebhaftigkeit seines Benehmens beim Unterricht.

\quellseite{186}
Insgemein ist die Vorbereitung durch bestehende und eingeführte Katechismen sehr erleichtert. Man bedarf keiner Invention des Inhalts. Wenn aber, wie es zuweilen der Fall ist, Katechesen über gewisse Themata oder freie Katechesen zu halten sind, so gilt rücksichtlich der Vorbereitung alles, was wir oben von der Vorbereitung auf die Predigt zum Lobe der Dialektik sagten. Es muß dann auch ein Thema, ein Eingang, eine Partition und eine Peroratio gefunden und vorbereitet werden, und die Vorbereitung auf die Katechese unterscheidet sich alsdann von der auf die Predigt nur dadurch, daß man das rhetorische Element nicht zu berücksichtigen hat und überhaupt außer genauer Disposition keiner Ausführung bedarf. Denn es ist eine rein unmögliche Sache, der Katechese, welche so sehr von den zu katechisierenden Kindern, ihrer Gabe und ihrem guten Willen abhängt, eine Ausführung bis in die einzelnen Fragen schon voraus zu bestimmen.

Man wird kaum sagen können, daß durch die Katechese mehr als durch die Predigt gewirkt werde, oder daß es eine höhere Gabe sei zu katechisieren als zu predigen. Im Gegentheil, der erste Segen und die höhere Begabung wird der Predigt zuzueignen sein. Es bedarf auch zur letzteren einer reicheren und vielseitigeren Gabe. Aber allerdings gibt es eine Gabe, welche, obschon beiden nöthig, von dem Katecheten weniger als von dem Prediger scheint entbehrt werden zu können, nemlich die dialektische. Der Prediger kann seinen Mangel in diesem Stücke zudecken und verhüllen oder durch großen Fleiß der Vorbereitung ersetzen, während
\quellseite{187}
der Katechet bei der Richtung, die er auf seine Schüler in den oft durch ihre Fragen und Antworten bedingten Gedankengang zu nehmen hat, eine große Leichtigkeit und Behendigkeit in Behandlung der Begriffe und Sätze bedarf, wenn er seine Disposition und den allgemeinen Gang der Katechese festhalten und zum bestimmten Ziele leiten will. Und hierfür gibt es keine Vorbereitung als den fortwährenden Fleiß in dialektischer Ausbildung und die dauernde Freude an dialektischer Uebung im Allgemeinen. Werden die Katechesen nach einem Leitfaden gehalten, so wird dennoch eine Vorbereitung immer nöthig bleiben, nemlich die, welche in genauer Kenntnißnahme und lebendiger Durchdringung des ganzen Leitfadens und jedes treffenden Pensums nach seinem Gedankengang und Zusammenhang besteht. Daß mit der Zeit die Vorbereitung leichter wird, sich oft auf ein bloßes Orientieren im Leitfaden und Unterrichtsgang beschränken, in wenigen Minuten geschehen sein kann, ist gewiß. Es muß aber dennoch immer eins festgehalten werden, nemlich daß man nicht die nöthige Vorbereitung unterlasse, so lang oder kurz sie werde.

Eine Gefahr hat die katechetische, wie die homiletische Vorbereitung. Nemlich die Vorbereitung wird oft nicht im Interesse des Zweckes der Predigt oder Katechese, sondern in einer Art von wissenschaftlichem Interesse vorgenommen. Man kommt dadurch zu eigenem reicheren Wissen, tieferer Erkenntniß und völligerer Befriedigung und kramt dann beim katechisiren gerne selbstgefällig viele, zum Zwecke nicht dienliche Weisheit aus, oder muthet den Schülern zu, was man ihnen
\quellseite{188}
nie zumuthen sollte. Es geht hier den Katecheten in ihrem Maße wie den Universitätsprofessoren. Wie diese, in eigenes fortgehendes Lernen vertieft, oftmals vergessen, daß sie nicht Theologen ihres Gleichen, sondern Pfarrer bilden sollen; so thut zuweilen ein Katechet, als sollte er nicht Christen, sondern Pfarrer bilden. Man vergesse also bei der Vorbereitung nie den Zweck der Vorbereitung, nemlich sich tüchtig zu machen, daß man dem Kinde das Nöthige in der faßlichsten, nachhaltigsten und segensreichsten Weise geben könne.

Haben wir nun diese Gefahr hervorgehoben, so dürfen wir auch eine andere nicht vergessen, die nemlich, welche im Extrem eines getreuen Haltens am vorgeschriebenen Pensum liegt, die des Schleudrians und eines handwerksmäßigen, todten Treibens. Wer weiß und will, was er soll, den wird Gott auch behüten, daß er in dieser Gefahr nicht hängen bleibe.

Die hauptsächlichste Gabe und die Bildung des Katecheten ist also die dialektische, nemlich im Unterschiede vom Homileten und Liturgen. Nun ist es bekannt, wie wenig, ja gar nicht formale Logik und Dialektik auf unseren Schulen und Universitäten getrieben zu werden pflegt, wie wenig dialektische Ausbildung die Candidaten des Predigtamts in das Amt mitbringen. Es ist aber auch bekannt, daß die Bildungsweise und der Bildungsgang unserer Schullehrer noch ungleich weniger eine dialektische Ausbildung gewähren kann. Die formale Bildung zum Katecheten fehlt deshalb dem
\quellseite{189}
jungen Schullehrer nicht minder, sondern eben so sehr und weit mehr als dem Candidaten des Predigtamts, welcher doch durch die Art seines Lernens und durch seine grammatische und sprachliche Bildung wie unbewußt zu größerer dialektischer Thätigkeit geführt wird. Es läßt sich drum schließen, wie viel an der oft wiederholten Behauptung sei, daß Schullehrer insgesamt bessere Katecheten seien, als die Pfarrer. Es könnte mit dieser Behauptung im Grunde nur gesagt werden wollen, daß Noth schwimmen lehre, und daß die Schullehrer vermöge ihres Berufs in dieser Noth öfter seien, als die Pfarrer, und deshalb sich kräftiger getrieben fühlen müssen, sich irgend eine katechetische Form anzueignen.

\section{30. Die Form der Katechese}

Schon § 24 ist auf die Vereinigung des akroamatischen und dialogischen Elements bei der Katechese hingewiesen worden. Diese Vereinigung würde vielleicht in der folgenden Weise am einfachsten und natürlichsten geschehen.

Wie die Predigt einen Eingang hat, so hat auch die Katechese einen längeren oder kürzeren homiletischen Eingang der zum Thema überführt. Dieses muß genannt werden, auch wenn man nicht über eigentliche Themata, sondern über Pensa eines Leitfadens oder Katechismus zu katechisiren hat.

Nach dem Eingang wird es gut sein, den ganzen Inhalt des Pensums oder der Katechese in kurzer, klarer, nicht homiletisch, sondern dialektisch gehaltener Uebersicht vorzutragen.

\quellseite{190}
Die Schüler werden den kurzen Vortrag beachten, weil sie ja wissen, daß sich die Katechese an ihn anschließen und über seinen Inhalt verbreiten wird.

Hat man den Inhalt des zu lernenden Pensums vorgetragen; so dürfte es auch dann noch sehr zweckmäßig sein, aus dem Leitfaden, wenn man nemlich einem solchen zu folgen hat, das Pensum lesen zu laßen.

Darauf laße man die erklärenden, verständigenden, examinierenden Fragen folgen und behalte dabei immer das Ziel im Auge, nemlich daß das Pensum gelernt werden soll.

Ist dies Ziel erreicht, so kann man noch einmal den Inhalt der Katechese zusammenfassen und in der Peroratio den Werth des Gelernten hervorheben, zur Ausführung und Uebung des Gelernten ermuntern, was gut ist, auch wenn man derartiges schon während des Katechisierens vorgebracht hätte und es als Wiederholung erschiene.

Die angegebene Weise kann vielleicht manchem wie eine Erschwerung für katechetische Kunststücke vorkommen. Allein was liegt am Ende an denen? Man katechisiret nicht, um zu katechisiren, und kann nicht die letzte Absicht haben, bei der Katechese seine Meisterschaft in der katechetischen Kunst und Form zu beweisen, sondern es handelt sich davon, daß man durch Katechesieren den Zweck heilsamer Erkenntniß erreiche, daß seliglich gelernt werde. Dazu dient die oben angegebene Form, die man aber deshalb nicht pedantisch festgehalten wissen will, sondern selbst wiederum unter die Regel: Summa utilitas omnis regula, stellt. — Hiebei kann man nicht
\quellseite{191}
umhin, auch hier wieder auf die vortrefflichen Katechismuspredigten\footnote{Aber auch sie, wie andere Predigten jener Zeit, begehren in ihrer Einfalt keine homiletischen Meisterstücke zu sein, sondern verdienen unsern Preis nur als katechetische Vorträge und in Anbetracht des in ihnen festgehaltenen geringen Pensums. Titel: „Katechismus oder Kinderpredig. Gedruckt durch Johann Petreium, anno 1533.“ Vorwort: „Dieser Katechismus oder Kinderpredig ist nicht um der Pfarrherrn oder Prediger willen geschrieben; denn man fast wohl weiß, daß der meiste Theil aus ihnen so viel christliches Verstandes wohl hat, Gott hat Lob, daß sie solche Lehr selbst ziemlich und nützlich könnten handeln. Sondern um der jungen Kinder willen ist er aus großer Noth also verfaßt und zusammen gebracht. Sintemal das die täglich Erfahrung gibt, daß, wer Kinder lehren will, der muß ihnen nicht zu viel auf Ein Mal vorgeben, und dasjenige, das er ihnen vorgibt, oft und dick, immerdar in einerlei Weis und Worten wieder vortragen. Denn wenn man ihnen jetzt dies, jetzt jenes, jetzt mit diesen, jetzt mit andern Worten vorsagt, so behalten sie nicht allein nichts davon, sondern werden auch nur ungeschickter und unachtsamer zu lernen, denn sie vor waren.“ „Dieweil denn wohl zu besorgen gewesen, es würde nicht ein jeder Pfarrherr oder Prediger die Mühe auf sich nehmen, den Katechismus also zu fassen, daß er ihn einmal wie das andremal lehrete, dazu nicht vonnöthen, einen jeden mit solcher Arbeit zu beladen, — auch damit es an einem Orte wie am andern gehalten würde, welches bei dem gemeinen Mann viel Aergernis verhütet; so ist ihnen hie mit dieser Arbeit gedienet, auf daß die Kinder den Anfang christlicher Lehr von uns desto leichter mögen begreifen und behalten. Denn wenn das geschieht, werden sie nicht allein alle andre Lehr desto besser verstehen, sondern auch feine, fromme, christliche, geschickte und weise Leute zu allerlei christlichen Ständen und Aemtern werden, dazu billig Vater und Mutter, Pfarrherr und Prediger und alle christliche Obrigkeit fleißig fördern und helfen sollen. Da gebe Gott Seine Gnade zu.“ Vgl. Löhe’s Erinnerungen aus der Reformationsgeschichte von Franken. Nürnberg 1847. p. 131.} der brandenburgischen Kirchenordnung hinzuweisen, an denen man sich für katechetische Einleitungsvorträge ein Muster nehmen kann, wenn gleich sie bei aller ihrer Kürze doch noch
\quellseite{192}
für unsere Zwecke zu lang sein würden. Die brandenburgischen Predigten setzen aber nicht einen gedruckten Katechismus voraus, wie unsere katechetischen Einleitungsreden, sondern sie standen selbst an der Stelle des Katechismus.

\section{31. Dialogische Form}

Die dialogische Form dient, sie herrscht nicht, wie wir sahen. Fragt man nun, was das Urbild und Ziel sei, wonach man in der dialogischen Katechese zu ringen hat, so ist es das gegenseitige freie Gespräch. Der Katechet hat es am weitesten gebracht, welcher seine Katechumenen nicht bloß zur Auffassung und Einhaltung eines Dialogs, und zu Antworten, sondern zu eigenem Fragen, zum Lehrgespräch brachte. Christus im Tempel als zwölfjähriger Knabe zeigt die vollkommenste Gestalt des rechten Katechumenen und der rechten Katechese. Ist ein solches Interesse gewonnen, daß die Schüler selbst mit in die Sache eingehen und mitreden und mitfragen, dann wird die Katechese leicht, fruchtbar und seliges Vergnügen. Die Kunst verschwindet und verklärt sich zur geheiligten Natur gottseliger Unterhaltung. Aber freilich die wenigsten Lehrer und die wenigsten Schüler kommen zu diesem Ziel, ja mancher Lehrer ist thöricht und abgeschmackt genug, die zarte Pflanze
\quellseite{193}
sich regenden Lebens und eigener freierer Theilnahme von Seiten des Schülers für ungebührlich anzusehen und den Schüler auf die schulmäßige Frage und Antwort zu verweisen.

Dies Ziel im Auge, aber auch die traurige Beschaffenheit der meisten Confirmanden und Schüler, von welcher man nothwendig ausgehen muß, müssen wir uns denn vor allem fragen und klar zu machen suchen, wie man die dialogische Form der Katechese zu gestalten habe.

Wenn man nach einem gehaltenen katechetischen Vortrage und über denselben katechisiren sollte, so würden die Fragen zuerst Examinationsfragen sein, an die sich Verständigungsfragen anschließen würden, worauf am Ende wieder Examinationsfragen kämen. Hat man aber nach einem Katechismus oder Leitfaden zu unterrichten, so werden nach Vorlesung und Aufsagen des treffenden Pensums entsprechend Constructions- und exegetische Fragen beginnen, nach Befund der Sache Verständigungsfragen folgen, Examinationsfragen schließen. Man könnte allerdings auch die Constructions- und exegetischen Fragen Examinationsfragen nennen, doch sind sie nicht immer und eröffnen vielmals dem Schüler erst das Verständnis. — Hiebei kann man wohl die Bemerkung nicht verhalten, daß am Ende alle katechetischen wie überhaupt alle Fragen in die Constructionsfrage eingeschlossen sind, und daß die Lehre vom Satzbau und vom Styl dem Katecheten für seine katechetische Ausbildung sehr förderlich sein kann. Bei kleineren Kindern wird man ohnehin alle Katechese geradezu mit den
\quellseite{194}
einfachen Constructionsfragen eröffnen müssen. Man bahnt so sichere Auffassung und Verständnis der Sache an und arbeitet allen andern Fragen vor. Neigt man sich dabei zum Kinde und läßt man immer die Absicht blicken, es ins Verständnis des Gelesenen oder Auswendiggelernten einzuführen, so wird es vielleicht zutraulich, fragt am Ende selbst und es entsteht das rechte katechetische Verhältnis, die gemeinsame Bestrebung des Lehrers und Schülers, dem letzteren die Wahrheit zur Erkenntnis zu bringen. Bei Katechesen über biblische Geschichte ist außer dem Examinieren noch eine Uebung sehr ersprießlich, nemlich das Erzählenlaßen des Gelernten. Nicht bloß wird dadurch formale Bildung gefördert, der mündliche und eben damit auch der schriftliche Ausdruck der Gedanken geregelt und geübt, sondern es prägt sich auch das, was man selbst erzählt, der Seele desto tiefer ein. Doch wird man darauf verzichten müssen, alle Kinder zum Erzählen zu bringen, weil nach der Erfahrung manche Individualität, die vortrefflich faßt und bewahrt, den Mangel ihrer Begabung gerade darin hat, daß sie nicht zusammenhängend wieder zu geben vermag, was sie innerlich erfaßt hat und bewegt. Damit ist freilich auch gesagt, daß es manche Natur, die sonst trefflich begabt sein kann, nicht zum geordneten schriftlichen Ausdruck der Gedanken bringt.

Man hat oftmals mit Luther hervorgehoben, daß man eine gewisse stehende Form der Fragen und Antworten anstreben und ganze Geschlechter in denselben heranwachsen lassen müsse. Es ist dabei natürlich keine sinnlose Dressur gemeint,
\quellseite{195}
vermöge welcher das Kind eilig auf den bekannten Klang der Fragen die bekannte und doch unerkannte Antwort gibt. Es ist auch nicht gemeint, daß alle und jede Fragen immer in gleicher Weise wiederkehren sollten. Das wäre, wie es an sich unmöglich ist, das geisttödtendste Beginnen, das sich denken ließe. Aber die kirchliche Terminologie muß in allen Fällen und das stehende Wort eines und desselben Leitfadens oder Katechismus desgleichen festgehalten werden. Es fördert die Einheit der Erkenntniß bei dem Geschlechte, das man zu erziehen hat. Ueber den Einen Inhalt und bei Festhaltung der Einen Terminologie muß sich jedoch das katechetische Gespräch in großer Verschiedenheit und Freiheit verbreiten können. Denn es ist Ziel und Vollendung der katechetischen Leistung, über das Feststehende die eingehendste, verständigendste Unterhaltung zu eröffnen.

Will man freilich das erreichen, so wird man wenigstens dem Landkind erlauben, ja es ermuntern müssen, zu reden, wie es kann, sei es hochdeutsch oder im Dialekt. Es kann sogar zuweilen dienlich und nöthig sein, daß der Lehrer selbst im Gespräch zum Dialekte heruntersteige, wie man es hie und da im Württembergischen und in der Schweiz hören kann. Ist doch der Dialekt an und für sich nicht unedel, sondern wird es nur dadurch, daß man ihn zur Bauernsprache macht, statt daß ihn die gebildeteren Stände pflegen und durch ihren Gebrauch vor den Auswüchsen bewahren sollten, die er im Munde des rohen und unberathenen Volkes so leicht bekommt. Er wird gemein durch gemeinen Sinn und Mund. Im Munde
\quellseite{196}
des Gebildeten wird er herzlich annahend und schön, Gründe genug, weshalb man sich des Dialekts nicht zu schämen braucht, ihn gestatten und selbst üben kann, zumal, wenn er für Einführung der Jugend ins Verständnis der heiligen Lehre dienlich ist, und das ist er gewiß, wenn das Kind, das nun einmal hochdeutsch nicht reden kann, oder aus falscher Scham sich weigert es zu thun, auf diese Weise zum Reden gebracht ist.

Zum Schluß des Paragraphen noch die wohl überflüssige Bemerkung, daß man Fragen, auf welche nur mit Ja oder Nein zu antworten ist, nicht sehr oft bringen soll. Zuweilen ist zwar ein Ja und Nein auch keine kleine Sache, und es mag daher dem Katecheten überlassen bleiben, in solchem, d. i. im rechten Fall, auch ein Ja oder Nein hervorzurufen. Aber es ist dennoch jedermann bekannt oder doch leicht zu erkennen, was man mißbilligt, wenn man gegen die immerwährende Wiederkehr von ja und nein angeht.

\section[32. Ist die Katechese von der Kanzel?]{32. Ist die Katechese von der Kanzel oder ambulando zu halten?}

Ist das Ideal einer Katechese das lebensvolle Gespräch Zweier oder die anständige Unterhaltung einer Gemeinschaft von Heiligen über herrliche Dinge aus dem Reiche Gottes, so scheint es auch gar keinem Zweifel zu unterliegen, daß man bei derselben unter der Gemeinde der Kinder und Katechumenenschar wandeln soll. Die Wechselbeziehung gibt sich da
\quellseite{197}
viel leichter und natürlicher. Es läßt sich auch Ruhe und Lebhaftigkeit des Katecheten, wie es nöthig ist, auf diese Weise wohl vereinigen. Eine Unterhaltung, ein Gespräch von der Kanzel, wie wunderlich würde man das finden, wenn man nicht von Jugend auf gewöhnt wäre zu hören, daß es sich hie und da finde! Die Kanzel ist nicht zur Unterhaltung, sondern zur Predigt, die Katechese aber fordert die Nähe des Lehrers.

— Freilich bei den Kindern, so wie sie sind, gewährt die Kanzel einen Standpunkt, von wo aus man alle übersehen und leichter in Ruhe und Ordnung halten kann, während beim Umhergehen des Katecheten der Muthwille hinter seinem Rücken sich gehen lassen kann: ein zufälliger Vor- und Nachtheil, welcher die naturgemäße Form, das Ambuliren, zwar nicht verdächtigen kann, aber dennoch oft genug, zumal in kleinen Kirchen, den Sieg über die ambulirende Lehrweise davon tragen wird.

\section{33. Confirmandenunterricht}

In der älteren lutherischen Kirche hatte man größtentheils eine sehr einfache Confirmation. Wenn der Superintendent seine Diöcese zum Fastenexamen bereiste, wurden ihm auch die Kinder vorgeführt, welche das erstemal zum heiligen Sacramente gehen sollten. Er examinierte sie im geringsten Maße katechetischer Erkenntniß, und wenn sie das Examen bestanden, betete er unter Handauflegung über ihnen und sprach sie reif zum ersten Abendmahlsgang. Bei dieser Sitte war
\quellseite{198}
insofern die antike Weisheit beibehalten, daß man nicht dem jedesmaligen Pfarrer selbst über die Reife seiner Katechumenen das letzte Urtheil zugestand, sondern Examen und Confirmation der Katechumenen der kirchlichen Aufsichtsbehörde in die Hände legte. Dagegen aber trat eins völlig zurück, nemlich die Erneuerung und Bestätigung des Taufbunds, die eigentliche Confirmation. Die amtliche Handlung des Superintendenten bezog sich bloß auf die Abendmahlsfähigkeit. Durch diesen Mangel kam die Confirmation zu keinen Ehren, zu keinem Werthe und konnte es auch nicht. In der neuen Zeit freilich geschieht zu großem Schaden an vielen Orten das Gegentheil, das Sacrament tritt hinter die Confirmationsfeier weit zurück. Es scheint hie und da ein viel höherer Festtag zu sein, wenn das Kind confirmiert wird, als wenn es seinen ersten Gang zur Gnadentafel Jesu antritt. Das Richtige wird sein, einer jeden Handlung ihre Würde, also dem heiligen Mahle den unbegrenzten Vorrang zu laßen.

Von diesem Standpunkt aus wird man auch den Confirmandenunterricht zu halten haben. Dieser Unterricht ist nichts anderes, als ein Unterricht zur Vorbereitung auf die Confirmation und den erstmaligen Genuß des heiligen Abendmahls. Weil er das ist, so ist es eine verkehrte Auffassung, ihn als Religionsunterricht im allgemeinen zu betrachten. Der allgemeine Religionsunterricht sollte vielmehr bei dem Confirmandenunterricht vorausgesetzt werden können, wenn dieser seinen eigenthümlichen Charakter haben und behalten soll. Es ist aber freilich in den wenigsten Gemeinden so, daß der
\quellseite{199}
vorauszusetzende Unterricht und die entsprechende Erkenntnißstufe auch wirklich vorausgesetzt werden darf. Doch können wir hier, wo wir von dem Confirmandenunterricht reden, wie er sein soll, keine Rücksicht auf diesen Mangel nehmen. Zwar wird es auch dann, wenn bei vorhandener genügender Vorbildung der Kinder ein wirklicher Confirmandenunterricht gegeben werden kann, nicht an einer gewissen Vollständigkeit und am Eingehen auf die ganze christliche Lehre fehlen, es wird bei dem Confirmandenunterricht die gesammte Lehre großen Theils wieder vorkommen, aber in einer eigenthümlichen Weise, zu besonderem Zweck und in besonderer Beziehung. Die eigentlichen Themata des Confirmandenunterrichts, bedingt von der doppelten Bedeutung desselbigen, sind die Sacramente, Taufe und Abendmahl, und was von dem einen zum andern überleitet, Confirmation und Absolution. Der Confirmandenunterricht ist daher seiner Natur nach ein sacramentlicher, und der Unterrichtsgang könnte etwa folgender sein:

1) Was ist die Taufe?\\
2) Von der Kindertaufe.\\
3) Von Pathenpflichten und Pathendank.\\
4) Ordnung der Taufhandlung.\\
5) Jachtaufe und Anleitung dazu.\\
6) Was ist die Confirmation im Allgemeinen?\\
7) Was ist der Taufbund, der erneuert und bestätigt werden soll?

\quellseite{200}
8) Die Abrenuntiation der Confirmation insonderheit.

9) Das Credo der Confirmation.

10) Die wahre Kirche mit Bezug auf die 7te Confirmationsfrage, welche nach der Treue gegen dieselbe fragt.

11) Zulässigkeit und wahrer Werth der Confirmation.

12) Confirmationsordnung.

13) Absolution. Warum Absolution vor Confirmation?

14) Von der Beichte. Anleitung zu den verschiedenen Arten derselben.

15) Von der heiligen Zucht der Bruderliebe. Nach Matth. 18.

16) Stellung der Beichthandlung vor der Confirmation, sammt Ordnung und Form der Beichte und Absolution.

17) Stiftung des heiligen Abendmahls.

18) Das heilige Abendmahl. Confessionelle Unterschiede.

19) Würdigkeit und Selbstprüfung.

20) Abendmahlsgemeinschaft.

21) Abendmahlsordnung.

22) Die Lehre vom Sacrament, sammt Würdigung der sieben Sacramente anderer Kirchen\footnote{Nr. 22 könnte am Anfang des ganzen sacramentlichen Unterrichts stehen.}.

23) Vom Gebet und dessen seliger Uebung.

\quellseite{201}
Nimmt man nun den Confirmandenunterricht in dieser eigenthümlichen Form und Begrenzung\footnote{Geistliche, welche die Sache fassen, sollten ihren Amtsbrüdern vorgehen, und ihnen durch schriftliche Ausarbeitung des Confirmandenunterrichts zeigen, was und, sofern es möglich ist, das schriftlich darzulegen, wie man es geben soll. — Der Verfasser dieser Blätter hat es 1860 versucht, „den sacramentlichen Theil des Confirmandenunterrichts“ nach obiger Reihenfolge auszuführen. S. Wilh. Löhe’s Tractate für die Seelsorge. VI. Der sacramentliche Theil des Confirmandenunterrichts. Zur Repetition für Confirmierte. Nürnberg. u. E. Sebald’sche Buchdruckerei und Verlagshandlung 1860. Der höchst unvollkommene Versuch erlebt so eben (1868) eine 2. Auflage. Er hat von mancher Seite Anerkennung gefunden, die er bei seiner Beschaffenheit nicht verdiente. Er ist und bleibt auch in der neuen Auflage ein armer Versuch, einen richtigen Gedanken auszuführen, den nemlich, daß der Confirmandenunterricht von dem übrigen Religionsunterricht zu scheiden und als Einführung in das kirchliche und sacramentliche Leben zu behandeln sei.}, und setzt man einen vorausgegangenen allgemeinen Religionsunterricht, so begreift es sich, weshalb man in der älteren Zeit dem Confirmandenunterricht nur eine kleinere Summe von Lehrstunden widmete. Ein Sechs-Wochenunterricht, wie ihn Markgraf Georg Friederich von Ansbach (1557—1603) in Franken einführte, ist zumal bei täglich gegebener Stunde vollkommen ausreichend, um das oben angegebene Ziel zu erreichen. Ja, es möchte nicht einmal gut sein, einen so bestimmt gefaßten Unterricht mehr in die Länge zu ziehen. Der Confirmandenunterricht ist feierlicher Natur, auf der ahnungsreichen Schwelle der Abendmahlsgemeinschaft gegeben, und eben deswegen erwartet auch die Gemeinde von ihm etwas besonders Hebendes und Anregendes.

\quellseite{202}
Wie wäre es aber möglich, ein ganzes oder halbes Jahr hindurch die feierliche Beziehung auf die kommenden heiligen Handlungen gleich fest und frisch zu halten? Maß ist hier allerdings nöthig, damit nicht der Zweck verfehlt werde.

Es ist übrigens nicht die Meinung, daß der Geistliche, weil er einen kurzen Confirmandenunterricht geben soll, damit auch des vorausgehenden allgemeinen Unterrichts entbunden und die Last seiner Arbeit gering gemacht werden soll. Er kann und soll die Kinder von Jugend auf unterrichten, so aber, daß er zwischen dem Religionsunterricht im allgemeinen und dem Confirmandenunterricht den rechten Unterschied aufrecht erhalte, beide nicht vermenge, sondern dem letzteren seine eigenthümliche Weihe und Würde laße.

Aus der Natur des Confirmandenunterrichts ergiebt sich endlich, daß die Präparanden noch in den allgemeinen Unterricht gehören, beim Confirmandenunterricht nur Hörer sein können, so wie, daß die Repetenten hörend wiederholen. Es ergiebt sich aber nicht, daß das Institut der Präparanden und Repetenten aufhören solle. Vorfeier und Nachfeier ist hier ganz gut. Und wohl denen, welche sie haben können und benützen.

\quellseite{203}
\chapter{Liturgisches}

\section{34. Die heiligen Zeiten}

Der Mensch lebt in Zeit und Raum; auch das gottesdienstliche Leben, wovon die Liturgik handelt, muß deshalb nach Zeit und Raum und nach der Art und Weise, die es selber besitzt, betrachtet werden. Reden wir also zuerst von der gottesdienstlichen Zeit, dann vom gottesdienstlichen Raum, zuletzt vom gottesdienstlichen Leben selbst.

Bei der gottesdienstlichen Zeit haben wir vom Tage, von der Woche, und vom Jahre zu reden. Herr unserer Zeit, unserer Tage, Wochen und unseres Jahres ist der gekreuzigte Christus, der auferstanden ist und zur Rechten des Vaters sitzt. Am Charfreitag starb der Herr. Was natürlicher, als daß man den Tag darauf ganz in Erinnerungen des Tags vorher verlebte. Das blieb der Christenheit, ein jeder Tag vergeht im Andenken des Leidens Jesu. Der Verlauf seines Leidens und Sterbens regelt den gottesdienstlichen Tag. Und wie nun der Todestag des Herrn in allen unseren Tagen herrscht, so gibt seine große Woche unseren Wochen Bedeutung, Maß und Ziel. Der Sonntag, da er sich zum Leiden einstellte, wird zwar von dem Auferstehungssonntag überstrahlt; aber jeder unserer Sonntage ist nun ein Bruder des Ostertags geworden. Der Mittwoch mit dem offen erklärten
\quellseite{204}
Entschluß des Verraths, der Donnerstag mit der Einsetzung des heiligen Mahls, der Freitag mit der Geschichte des Leidens und Sterbens, der Sonnabend mit der Grabesruhe Jesu drücken allen unseren Mittwochen, Donnerstagen, Freitagen und Samstagen den liturgischen Charakter auf. Und wie jeder Tag für sich im Gedächtnis des Leidens Jesu vergeht, so jede Woche im Gedächtnis seiner letzten Woche bis zum Tag der Auferstehung. Und wie ein jeder Tag den Charfreitag, jede Woche Jesu letzte Woche, so wiederholt das Jahr feiernd das ganze Leben Jesu. Die Wartezeit begehen wir im Advent. Von Weihnachten bis Septuagesima sehen wir ihn geboren werden, wachsen, zum Manne reifen, ins Amt treten, im Amte wirken. Von Septuagesima bis zum Osterabend feiern wir sein Leiden, dann folgt die Feier seiner Auferstehung, seiner vierzig österlichen Tage, seiner Auffahrt, die Ausgießung seines heiligen Geistes und im Verlauf des Ueberrestes vom Jahre das Gedächtnis des Baues und Lebens seiner heiligen Kirche. So leben wir alle Tage, Wochen und das ganze Jahr unserem Herrn Jesu.

Und zwar begehen wir immer das Gedächtnis der Leiden und der Auferstehung in derselben Weise. Um elf Uhr Mittags läuten alle Glocken im Lande die heiße Mittagsstunde Jesu ein, wo er am Kreuze hieng, Nacht um ihn her\footnote{So müssen wir das Elfuhrläuten faßen, weil bei uns die Ansicht durchgegriffen hat und allgemein geworden ist, daß JEsus Christus nicht um 11 Uhr, wie die Alten annahmen, sondern um 9 Uhr gekreuzigt worden sei.};
\quellseite{205}
ernste Feier seiner Leiden in Buße für die Sünden und Lobpreisung seines allerheiligsten Verdienstes füllt die Seele des Christen. Um drei Uhr denkt man im Stundengebete seines Todes und betet um einen guten Tod. — In jeder Woche begeht man am Mittwoch und Freitag Jesu Leiden mit Buße, jeden Sonntag seine Auferstehung mit Freuden. — Und im Jahre ist die vierzig-, ja siebenzigtägige Gedächtniszeit der Leiden Jesu Bußzeit, Fastenzeit, Betzeit; die fünfzig Tage vom Ostertage an hohe Freudenzeit. — Wir sehen, die Geschichte Jesu, Jesus selbst ist unser Leben, — das Kirchenjahr, die Woche, der Tag eitel Gedächtnis Jesu, und wer sich ins Kirchenjahr hineinlebt, lebt sich in das Leiden Jesu hinein und in sein ganzes Leben.

Jedoch dürfen wir uns nicht verhehlen, daß das Kirchenjahr noch einen gedoppelten Sinn hat, oder wenn es ihn bei uns Protestanten verloren hat, doch gehabt hat und in der That auch noch haben sollte. Denn einmal zeigt jeder Kalender mit seinen vielen, jedem Tage beigeschriebenen Namen, die noch überdies gehäuft und gemehrt werden könnten, daß das ganze Jahr des Gedächtnisses der Heiligen voll ist\footnote{S. das „Calendarium sanctorum oder die Zeugenwolke des Neuen Testamentes,“ im Haus-, Schul- und Kirchenbuch von W. Löhe. Thl. II. 1859. p. 115 ff.}. Wie die Sterne den Sonnenlauf begleiten, so begleitet das vielfache Gedächtnis frommer Ueberwinder das hohe Gedächtnis des Lebens, Leidens, Sterbens Jesu. Und jeder Heilige spiegelt in seinem Leben, Leiden und Sterben das Leben und
\quellseite{206}
die Vollendung seines Herrn wieder. Wie gar nicht die lutherische Kirche eine Feindin dieses Gedankens gewesen ist, kann aus Luthers und anderer Schriften und aus der reichen Kalenderliteratur der lutherischen Kirche erwiesen werden. Von Luther an, auf dessen Mahnen und mit dessen Vorrede Georg Major die gereinigten Vitae patrum herausgab\footnote{Vitae patrum, in usum ministrorum verbi, quoad ejus fieri potuit repurgatae. Per D. Georgium Majorem. Cum praelatione D. Doctoris Mart. Luth. Witebergae 1562.}, bis in die neue Zeit herauf gehen zahlreiche Schriften, welche die Geschichten und den Tod der Heiligen Gottes feiern und dem Herrn nachahmen, von welchem geschrieben steht, daß vor ihm das Gedächtnis und der Tod seiner Heiligen werth gehalten ist\footnote{Die Literatur der protestantischen Kirchen über das Leiden und Sterben der Märtyrer hat in den wenigen Jahren, seitdem die erste Auflage dieses Buches erschienen ist, sehr zugenommen. Eine Hinweisung auf Piper’s Kalender und Fliedner’s Märtyrerbuch reicht hin, auf die Spur einer ganzen Reihe von ähnlichen literarischen Erscheinungen zu helfen. Es gibt nun zwar Darstellungen von Märtyrergeschichten, welche den Beispielen der Märtyrer nicht zu Frucht und Segen helfen können. Aber gewiß ist es doch, daß neben der heiligen Schrift, die außer Vergleich steht, keine Lectüre so ergreifend und segensreich ist, als die Passionsgeschichten der Märtyrer. Die gesammte Fluth der christlichen Unterhaltungs- und Romanlectüre rauscht vor uns vorbei. Aber die Geschichten der Helden Jesu werden bleiben und bis an’s Ende wirken. Auch die Beispiele des heroischen Glaubens aus dem Mittelalter werden auf diejenigen, welche Zeiten und Zustände zu unterscheiden und zu würdigen wissen, immer erfrischend und stärkend wirken.}.

Ferner feierte die Kirche vor Zeiten aber auch den natürlichen
\quellseite{207}
Verlauf des Jahres, seinen Fortschritt von Samen zur Ernte, von Sommer zum Winter, von Frost zur Hitze und wieder zum Froste mit ihrem Quatemberfasten und feiert ihn auch in der neuern Zeit mit ihren Erntefesten, Sylvesterabenden \&c.

Jeder Kalender schattet diese dreifache Jahresfeier nach dem Lauf der Natur, des Lebens Jesu und der Geschichte seiner Heiligen ab. Sonne, Mond und Sterne zeigen darin ihren Lauf, wandeln aber alle ehrerbietig um Jesum und seine Heiligen, deren Gedächtniß und Feiern dem natürlichen Jahre und seinen Zeiten und Tagen heiligen Sinn und Bedeutung, Werth und Leben verleihen. Es gibt drum auch nicht leicht ein Buch, welches so den Totaleindruck machen könnte, wie die gesammte Zeitengeschichte und die große Zeitenuhr, der gestirnte Himmel, dem HErrn und seiner Kirche unterthänig sei, als eben der Kalender, versteht sich mit Hinweglaßung alles Schnörkels und Aberglaubens, davon auch ältere und neuere Kalender dem abergläubig verwöhnten Volke zu Gefallen voll sind.

\section[35. Entstehung der Hauptfestzeiten]{35. Entstehung der Hauptfestzeiten des Jahrs}

Das Kirchenjahr ist ein Ganzes, aber kein gesuchtes, kein gemachtes, kein ängstlich menschliches System. Es fehlt demselben hie und da an Rundung und gleichmäßiger, harmonischer Ausbildung, wie ein einfacher Fingerzeig auf die lange Reihe der Trinitatissonntage nach dem Festhalbjahre
\quellseite{208}
zeigen kann. Was aber an Rundung des Jahres fehlt, wird durch die schöne Feier des gewöhnlichen Tages und der gewöhnlichen Woche ersetzt, deren stille Herrlichkeit auch ihr großes Recht und ihre süße Wohlthat haben. Denn wenn allezeit Fest wäre und es ein Festsystem gäbe, welches jeden Tag festlich machte, sei’s auch durch die mannigfaltigste und verschiedenste Feier, so wäre doch am Ende gar kein Tag festlich. Es würde das Außerordentliche zum Alltäglichen, Gewöhnlichen gemacht und ebendamit seiner Besonderheit und Eigenthümlichkeit beraubt. Ist nun aber nicht das ganze Jahr abgerundet, so sind es doch zwei kleinere Ganze innerhalb des Jahres, welche wie Sonne und Mond am Kirchenhimmel, stehen und das Jahr regieren. Wir meinen die eigentlich große Festzeit des Jahres von Septuagesima bis Trinitatis und die Weihnachtsfeier von Advent bis zum Epiphanienfeste mit ihren Nachklängen in den Epiphaniensonntagen.

Es war ganz natürlich, daß sich von Anfang an vor allem, was der Herr gelebt und gethan hatte, das Gedächtnis der Leiden und seiner Verherrlichung hervorhob. Lag doch am Leiden, Sterben, Auferstehen und der Verherrlichung des Menschgewordenen all unser Heil. Hier blüht und trägt für uns der Lebensbaum. Die zweite festliche Zeit, die der Weihnachten, trat bedeutungsvoller und ernster erst mit und nach jenen großen Streitigkeiten hervor, welche über Gottheit und Menschheit, über die Naturen und über die Person des Herrn geführt wurden. Als sich die göttliche Wahrheit in Betreff dieser hochwichtigen Punkte durch gewaltige Kämpfe
\quellseite{209}
hindurch Bahn gemacht hatte, konnte die Kirche ihren neugewonnenen und tiefer erkannten Glauben und die große Gewißheit von Christo, dem Immanuel, in kirchlicher Feier nicht besser ausdrücken, als durch das Fest seiner Geburt. Sie kniete vor aller Welt an der Krippe des Kindes Jesu nieder, denn es ist Gott und Mensch in Einer Person.

Anfangs feierte man nur die große Woche unter den Wochen, wie den Freitag unter den Tagen, bald aber dehnte man die Feier der Leiden weiter aus — auf 40 Tage. Waren es doch ungefähr 40 Stunden, die man auf die Zeit vom Tode bis zur Auferstehung Christi rechnen mußte. Jeder Stunde einen Tag vom Jahre zum Gedächtnis! Hatte doch der Herr selbst 40 Tage, 40 Nächte gefastet; ebenso Moses, ebenso Elias. 40 war dadurch eine geheiligte Zahl. Also 40 Tage lang dauert von alter Zeit her die Zeit des Gedächtnisses der Leiden Jesu. Da nun aber jeder Sonntag Bruder des Ostertags war, dem Gedächtnis der Auferstehung gewidmet, voller Freuden, so konnte er doch in den 40 ernsten Tagen des Gedächtnisses von Jesu Leiden nicht zählen. Es mußten also die Sonntage vom Gedächtnis der Leiden ausgenommen werden, Freudenpunkte mitten in den Trauertagen sein (wenn schon die Freude gemäßigt erschien durch die Nachbarschaft der Trauer) und die Tage, die 40, die Quadragesima, mußte bis auf den Mittwoch vor Invocavit, bis auf Aschermittwoch zurückgezählt werden. Mit diesem Tage begann die Quadragesima. Im Morgenlande, wo man an Sonnabenden nicht, wie im Abendland fastete, sondern den alttestamentlichen
\quellseite{210}
Sabbath wie einen Bruder des neutestamentlichen Sabbathstags eher mit Freuden begieng, wo man auch den Donnerstag als Einsetzungstag des heiligen Mahles von den Fasten ausnahm, mußte man die 40 Tage der Fasten vor Ostern rückwärts noch weiter erstrecken, als im Abendland. Aus der alten Quadragesima entwickelte sich so eine Quinquagesima, Sexagesima, Septuagesima, so daß die Gedächtniszeit der Leiden zehn Wochen vor Ostern begann. Auch im Abendlande ließ man sichs gefallen, nicht bloß die Quadragesima zu feiern, eine gewisse Feier schon mit dem Sonntag in der Septuagesima zu beginnen. Man dachte dabei gerne an die 70 Jahre des babylonischen Exils und nahm die ernste, traurige Feier auch wie ein Exil der Freuden, wie eine Prüfungs- und Kasteiungszeit, aus der man sich nichts desto weniger heraussehnte, dem schönen Ostertag entgegen. — Ganz dieser eben dargelegten Ansicht von der Fastenzeit entspricht der Charakter des Sonntags Lätare, des mittleren in der Quadragesima, der einen freudigen Charakter hat, weil er nach halbvollbrachter Kasteiung wie ein fröhlicher Ruhepunkt erscheint, der auf das nahende Ende der Trauer sieht, während die andern Sonntage der Fastenzeit doch immer noch nach dem Wermuthgarten duften, in dessen Mitte ihre Freuden blühen.

Der Osterabend, obwohl noch ein Fasttag, doch auch Taufstag der Katechumenen, war schon freudenvoll. Man war ja schon nahe am offenen Freudenthor, beim freien Freudengipfel. Da vergaß man die Kasteiung, die noch währte, im Hinblick auf die Freude, welche nach Mitternacht mit dem
\quellseite{211}
Halleluja eintrat. Es folgte eine Quinquagesima der Freuden. Jeder Tag war ein Sonntag. Von Freude zu Freude ging man, bis mit Pfingsten, später mit Trinitatis, das Jahr zum gewöhnlichen Laufe sich wendete. Die 40 Tage vor, die 50 Tage nach Ostern geben zusammen ein wundervolles System der Feier, welches in seiner hellen Pracht vor aller Welt so anerkannt stand, daß die Kirchenväter in ihren Schriften die Heiden aufforderten, in irgend einem heidnischen Cultus etwas von der Art aufzuzeigen.

Wie vor der großen Freudenzeit der Ostern, so ging auch vor Weihnachten eine Bereitungszeit, ein Abbild der 400 Jahre vor Christi Erscheinung vorher, gefeiert wie die Bereitungszeit vor Ostern mit ernster Einkehr in das eigene Herz, mit Fasten, Beten und Almosengeben. Und wie auf Ostern die Pentecoste, so folgt auf die Geburt Jesu die Epiphanienzeit, eine Freudenzeit nicht von gleichem Waßer mit der Pfingsten, aber dennoch eine Freudenzeit.

Schön und ernst schlang sich in die Zeit der größten Feste die Quatemberfeier ein. Im Dezember feierte man am Quatember des sogenannten zehnten Monats das Fest der vollendeten Ernte, und zwar mit Buße, Fasten, Beten, Almosen. Da kam also zum Ernste der Wartezeit auf Christum auch noch der Ernst des Erntefasten. Beim Beginn des Quadragesima sah man im Quatember, der am Mittwoch vor Reminiscere eintrat, hinaus auf das werdende Jahr, begleitete die Hoffnung einer Ernte mit Buße, Fasten, Beten, Almosen. Und wenn Pfingsten vorüber war, wenn das
\quellseite{212}
Jahr in seiner Pracht stand, die geschlossene Freudenzeit aber zu ernster Prüfung des Verhaltens während der Freude einlud, dann kam wieder, wie auch später im September, ein Quatemberfasten. Vier große Faste, Buß- und Bettage kehrten im Jahre an bezeichnenden Stellen wieder und luden zur Einkehr ein.

Es könnte auffallend scheinen, daß man, wie die Leidenszeit, so die Erntezeit und den Fortschritt der Natur und ihrer jährlichen Entwickelung mit dem Ernste der Buße bezeichnete. Allein, wenn wir in den Leiden Christi mit Recht unserer Freude Grund sehen, so ist doch der alte Gedanke, die Ursache unserer Freuden mit Buße zu begehen, zart, treu und innig: man litt in der Leidenszeit mit, indem man Buße that für alle Sünde, die Christo so großes Leid gebracht hat, und sich dafür an Ostern und Pfingsten von Herzensgrunde freute. Man schloß sich feiernd innig an die Worte an: „Er ist um unserer Sünde willen dahingegeben und um unserer Gerechtigkeit willen auferstanden,“ und fügte zur Sündenstrafe gebührendes Leid, zur Gerechtigkeit der Auferstehung die Herzensfreude. Und wenn wir die Ernte mit Lobgesang und Dank feiern, so wußte das Alterthum gewiß nicht weniger tief und geziemend die göttliche Wohlthat nicht besser zu empfangen, als mit Enthaltung und Buße und nicht besser für sie zu danken, als mit Almosen und Collecten\footnote{Man vergleiche die herrlichen Sermonen Leo’s des Großen de Collectis et Eleemosynis, de Jejunio decimi mensis ac eleemosynis etc. etc. Opera D. Leonis Magni, Romani Pontificis, ejus nominis primi etc. etc. Antverpiae, apud Philippum Nutium. 1583. p. 11 ff. 18 ff.}.

\quellseite{213}
\section{36. Von den Lectionarien der Alten}

Je und je wurde in den Versammlungen der Christen die heilige Schrift gelesen und zwar gestaltete sich, wenigstens im Abendlande, die Sache so, daß man die heiligen Bücher nicht fortlaufend und Wort für Wort las, sondern ausgewählte Lectionen gebrauchte. Diese Sitte läßt sich bis in das graue Alterthum zurück verfolgen, und spricht sich namentlich in vielen uns übrig gebliebenen Lectionarien aus. Diese Lectionarien sind nicht allein deshalb merkwürdig, weil man aus ihnen ersieht, was man in der Kirche zu lesen pflegte, sondern auch deswegen, weil aus ihnen erhellt, welche Feste man feierte, und weil aus den gewählten Lectionen auf den Sinn geschlossen werden kann, in welchem man sie wählte, also auf die den Festfeiern zu Grunde liegenden Festgedanken.

Die Lectionen der Alten sind nun aber verschieden, es gibt Horenlectionen und Meßlectionen. Jene erstrecken sich über das ganze Jahr und gaben Gelegenheit, den größten Theil der heiligen Schrift in Jahresfrist nach einer dem Kirchenjahr angepaßten schönen Ordnung durchzulesen. Sie sind ohnehin von den Meßlectionen dadurch unterschieden, daß sie nicht in unbeschränkter Willkür je nach Bedarf und Ermessen bald aus diesem, bald aus jenem biblischen Buche genommen sind, sondern daß eine bestimmte Folge der Bücher bei der Auswahl eingehalten ist. Sie vertreten im Abendlande das Prinzip der Lectüre ganzer Bücher, das cursorische Bibellesen, jedoch excerptorisch oder auszugsweise. Die Meßlectionen sind
\quellseite{214}
nicht, wie die Horenlectionen auf alle Tage vertheilt, sondern es gibt ihrer nur für die gefeierten Tage. Wir haben deshalb die Tage, für welche es Meßlectionen gibt, auch für gefeierte Tage zu erkennen. Da ist es nun besonders hervortretend, wie sorgfältig die Feier der Septuagesima, die Adventszeit und die Quatembertage mit Lectionen begabt ist. Man gewinnt deshalb aus dem Reichthum der Lectionen die im vorigen Paragraphen ausgesprochene Gewißheit von der hohen Feier der genannten Zeiten, und überhaupt kann man vom Kirchenjahr der Alten, von dessen Ueberresten wir am Ende jetzt noch kirchlich leben, keinen bessern Begriff bekommen, als in den Lectionarien.

Aus ihnen ist namentlich auch das Urtheil geschöpft, welches oben von der minderen Rundung des Kirchenjahres als eines Ganzen, dagegen aber auch von der Quadragesima und Pfingstquinquagesima als Schmuck und Krone des Jahres gefällt ist. Von Septuagesima bis Pfingsten greifen nicht allein die Horen, sondern auch die Meßlectionen herrlich zusammen, Ein Zug und Gedanke geht durch alle. Eben aus ihnen aber zeigt es sich auch, daß die Festgedanken der Alten nicht immer die sind, welche wir bei den Festzeiten einzuhalten pflegen. Die eben angedeuteten Gedanken der Fest- und Quatemberzeiten, die so wenig die unsrigen sind, gründen sich z. B. ganz auf die Lectionarien, wenngleich sie auch durch andere geschichtliche Umstände bestätigt und bewährt sind.

Für die neue Zeit, welche kein Lectionarium von Werth hat schaffen können, eben weil sie es durch Reflexion schaffen
\quellseite{215}
wollte, kann das Studium der alten Lectionarien und des alten Kirchenjahrs nur empfohlen werden. An jenen Erzeugnissen des Alterthums selbst müßte man sich klar werden, wo allenfalls das alte Kirchenjahr eine andere Auffassung zuließe, als die altgewohnte. An ihnen würde man aber auch Verstand und Weisheit der Alten bewundern lernen und sich um so lieber an sie anschließen, als man selbst keine Kraft und Weisheit besitzt, es ihren Schöpfungen voraus oder gleich zu thun\footnote{Einen Versuch der Wiederherstellung s. W. Löhe’s Haus-, Schul- und Kirchenbuch für Christen des lutherischen Bekenntnisses. 2. Thl. Stuttgart. Bei C. Bertelsmann. 1859. p. 141—158.}.

\section{37. Von dem heiligen Raum}

So gewiß die Ordnung der Zeit für die Liturgie von der größten Bedeutung ist, so gewiß ist auch eine angemessene Einrichtung und Ordnung des Raumes von einer großen Wichtigkeit für denselben Zweck. Der Christ kann überall beten und ein geistliches Leben führen, es läßt sich auch kein Ort von einer solchen Häßlichkeit denken, daß er nicht dennoch durch Ueberschwang des Geistes zu einem Bether werden könnte. Die Höhlengottesdienste der Alten beweisen das reichlich. Andererseits läßt es sich aber auch nicht leugnen, daß die Oertlichkeit mit ihrer Beschaffenheit der menschlichen Andacht Hindernisse entgegensetzen und umgekehrt sie in einem gewissen Maße auch fördern kann. Wir leben ein geistliches Leben im
\quellseite{216}
Räume wie in der Zeit, und zwar hat der Leib an demselben als unzertrennlicher Genosse der Seele seinen Antheil, so daß Leiblichkeit und Räumlichkeit mit nichten im Gegensatz zum Gottesdienst im Geist und in der Wahrheit zu stellen sind. Im Gegenteil, wo irgend der Geist es vermag, wird er auch das leibliche und die räumliche Umgebung und Einrichtung dem geistigen Zwecke der Anbetung entsprechend gestalten. Eben das zeigt sich auch in der ältesten Zeit, welche von denen nicht recht erkannt ist, die Spiritualismus und Abstraction mit Einfalt verwechseln. Es ist gerade die Einfalt eines kirchlichen Totallebens, was die Alten lehrte, ihre kirchlichen Räume des Zweckes würdig auszustatten. Wo überall deshalb in der ersten Zeit die Verfolgung ruhte, schmückte sich der Altar und das Gotteshaus und zwar in einem Maße, die man heutzutage hie und da ebenso für Luxus ausgeben würde, wie Judas die Narde des salbenden Weibes für Unrath und Verschwendung erklärt.

Wir haben aber dieses Orts nicht vor, vom Kirchenbau im Allgemeinen und von der Außenseite des kirchlichen Gebäudes zu reden, sondern es ist rein der liturgische Raum- das Innere der kirchlichen Gebäude, welches wir ins Auge fassen. Wie die Alten die Außenwand ungeschmückt ließen, aber den inwendigen Raum zum Paradiese umschufen, so bescheiden auch wir uns hier des Aeußeren, ohne es zu verachten oder gleichgültig von ihm zu denken.

Die Hauptfrage ist nun: Was in einer lutherischen Kirche die erste und wichtigste Stelle sei, ob Altar oder Kanzel?

\quellseite{217}
Allein diese Frage zu entscheiden ist leicht. Auch im lutherischen Gottesdienst ist nicht die Predigt, sondern das heilige Abendmahl die größte Feier. Die Predigt führt selbst zum heiligen Mahle, als zur innigsten geheimnißvollen Vereinigung der Christen mit ihrem Christus. Was aber das Ziel alles Gottesdienstes ist, heiligt auch die Stelle, an welcher es vollzogen wird, zur ersten in einem christlichen Gotteshause. Darum ist und bleibt auch für den lutherischen Kirchenbau und für die Einrichtung unserer gottesdienstlichen Gebäude der Altar die wichtigste Stelle, auf welche sich der gesammte heilige Raum zu beziehen hat. Eben daher ist die Anordnung der alten Gotteshäuser im Ganzen auch der lutherischen Kirche zuzueignen.

Unbewußt, aber kenntlich liegt der alten Einrichtung gottesdienstlicher Locale die Einrichtung der Stiftshütte oder des Salomonischen Tempels zu Grunde. Vor dem Heiligtum stand der Brandaltar, Golgathas Vorbild, welcher von dem ganzen Tempeldienst vorausgesetzt wird. Vom Brandaltar führt den Priester der Weg zum Waschfaß, das Hände und Füße reinigt, und hiermit ist die Taufe mit ihrem Nachklang in der Absolution vorgebildet, in der uns die Kraft des Opfers Jesu zugeeignet wird. Nur der Gewaschene tritt aus dem Vorhofe zum Heiligtum, nur der Getaufte und in der Gnade der Vergebung Stehende kann zum Gottesdienste des neuen Testamentes eingehen. Im Heiligtum steht der Leuchter, der uns Gottes Wort sinnbildet, die Schaubrode, Vorbilder unserer Dank- und Liebesopfer, sowie des Brotes, das wir dem Herrn
\quellseite{218}
zu seinem Sacramente bringen, — und der Rauchaltar — auf der Schwelle des Allerheiligsten, — der uns an das große Opfer des Gebets, des Dankes und des Lobgesanges erinnert. Im Allerheiligsten ist die Lade mit dem ganz goldenen Sühnedeckel, der das Gesetz bedeckt, der von Cherubim beschaut und geehrt wird, mit der Gelte und Aarons grünendem Stabe: alles Vorbild, wie des ewigen Lebens, so des heiligen Mahles, in welchem uns nahet, der in der Wolke wohnte, Er, der unsre Versöhnung ist, und dessen wunderbare Vereinigung mit der Menschheit Cherubim und alle Engel zu schauen gelüstet. Ganz dem entsprechend haben die Kirchen im neuen Testamente einen Vorhof mit dem Urbild des Waschfasses, dem Taufstein, der auf der Schwelle des Vorhofes zum Heiligthum steht. Im Heiligen wandelt die priesterliche Gemeinde, und wird von Gottes Wort, des Leuchters Urbild, erleuchtet. Da übt sie auch ihre gottesdienstlichen Gebete, da dankt und lobsingt sie und bringt Schaubrodt und Opfer ihrem Herrn. Und im Allerheiligsten steht ihr Altar, an dem sie vielmehr empfängt, als gibt, an welchem sie ihre Opfermahlzeiten hält und ihrem Christus zur innigsten Gemeinschaft begegnet. Offen steht vor ihr der Ort der höchsten Offenbarung und Vereinigung, und das ist der Unterschied des alten und neuen Testamentes. So ergiebt sich die dreifache Eintheilung einer jeden christlichen Kirche. Der Altar deutet das Allerheiligste der streitenden Kirche. Die Kanzel aber steht entweder an einem schicklichen Ort im Schiff der Kirche, oder an der Schwelle des Chors.

\quellseite{219}
Am Altare und in dessen Nähe ist der natürliche Platz derer, die am Altare dienen. Ob der Hirte und Bischof der Gemeinde bei seinen amtlichen Handlungen nach Osten gewendet steht, oder hinter dem Altare von Osten nach Westen gekehrt, ist gleich und hat beides eine schöne Deutung. Nicht gleich aber ist es, sondern liturgisch erfordert, daß die ganze Gemeinde nach Osten zum Altare schaue, also, daß alle Sitze und Bänke dem Altare zugewendet seien; daraus geht hervor, daß auch die Kanzel am Westen im Osten des Schiffes angebracht wird.

Die Chorsänger sammt dem Cantor stehen wohl am besten an der Spitze der singenden Gemeinde, auf der Schwelle des Chores oder Schiffes, etwa auch in der sogenannten Vierung. Die Orgel aber fände wohl hier ihrer Beschaffenheit nach keinen Platz. Sie gehört an's entgegengesetzte Ende des Chors, über den Vorhof. Man hat es auch zweckmäßig gefunden, daß der Chor sich bei der Orgel sammle, und alles, was zur Leitung des gemeindlichen Gesanges dienen soll, aus der Höhe schalle.

Die ganze Einrichtung des Gottesdienstes weist darauf hin, daß der Ort des Altars über das Schiff erhaben sei, damit die ganze Gemeinde Augenzeuge dessen werden könne, was am Altar geschieht. Doch werden die Protestanten kein Interesse haben, den Altar allzu hoch zu stellen, weil sie kein Meßopfer, keine Anbetung der Hostie gestatten und auch nicht darauf ausgehen können, eine Scheidung des sacramentlichen Ortes von der priesterlichen Gemeinde anzudeuten.

\quellseite{220}
Die Bänke sollten bequem zum Knien eingerichtet sein, weil das Knien eine kirchliche Gebärde ist, und allerdings die der Andacht förderlichste Stellung, wenn es nicht durch Schmerz und Unbehagen der Kniee stört. Eine Unterstützung der Arme der Knienden ist eben deshalb nöthig.

Emporen sind eine alte Sitte, doch fand es das Alterthum mit Recht schicklicher, die Frauen statt der Männer auf die Emporen zu verweisen.

Gut ist es, wenn der gottesdienstliche Raum nicht zu klein ist. Engigkeit widerspricht dem liturgischen Gedanken, der seine Differenzen und Intervalle auch durch eine gewisse räumliche Scheidung ausdrückt. Eine zu kleine Kirche fordert auch nicht zur Andacht auf, zumal wenn sie in protestantischer Weise mit Bank und Bank bis auf die Altarstufen hin bedeckt ist. Insbesondere aber ist es dem Liturgen selbst angenehm, die Gemeinde nicht zu nahe zu haben: es ist eigen, aber wahr, daß man sich inniger mit ihr vereinigt fühlt, wenn sie leiblich etwas ferner, als wenn sie näher steht. Es gibt hierfür, versteht sich, keine dogmatischen Gründe, sondern nur psychische, die sich aber doch geltend machen, obwohl es keinem Zweifel unterliegt, daß das Mißbehagen überwunden werden muß, wenn die Ursache davon nicht zu beseitigen ist.

\quellseite{221}
\section{38. Heiliges Geräthe}

Unter diesem Namen wünschte ich, weil wir keinen anderen zusammenfassenden haben, ebensowohl die unbeweglichen als die beweglichen Geräthschaften unserer Kirchen zusammenzufassen zu dürfen, welche zu unseren gottesdienstlichen Zwecken dienen, also ebensowohl den Taufstein, die Kanzel und den Altar, als Taufbecken und Kannen, Kelch und Ciborium, sammt all den Hüllen und Decken, welche für beides dienen. Ich fasse sie nicht bloß wegen ihres gemeinschaftlichen Zweckes dem christlichen Gottesdienste zu dienen zusammen, sondern auch deshalb, weil ich über sie, soweit sie nemlich aus der neuern Zeit stammen, und über das Verständniß, welches die von ihnen haben, die sie brauchen, im allgemeinen ein und daßselbige Urtheil glaube äußern zu müssen: ich finde nemlich alles beides sehr gering und unter aller Erwartung. Was uns Menschen sonst theuer und lieb ist, darauf verwenden wir Sorgfalt, lassen uns Kosten und Unkosten nicht gereuen; auch ein Armer ist zum Beispiel erfinderisch, um die Gräber der Seinigen zu schmücken und zu ehren und dergleichen. Nun ist ohne Zweifel das, was uns in unseren Kirchen gegeben wird, das Größte und Beste, was uns in der ganzen Welt gegeben werden kann, nemlich in der Taufe und im Worte der Geist mit seinen Kräften, und im Sacramente des Altares der Leib und das Blut des Herrn. Darum sollten wir die Orte, von denen aus uns die himmlischen Gaben gereicht werden, ehren
\quellseite{222}
und schmücken, und ebenso die Geräthe, in denen sie uns mitgetheilt werden. Wir schmücken die Gräber, in denen die todten Leiber unserer Väter ruhen, während wir die Orte und Geräthe, von und aus welchen wir lebendige Himmelsgüter nehmen, gerade so behandeln, als malte unter uns in Anbetracht ihrer eine ganz andere, umgekehrte Schluß- und Gedankenfolge. Wer sich davon überzeugen will, der mache nur eine Wanderung durch die Kirchen seiner Umgegend, schaue Taufsteine, Kanzeln und Altäre, lasse sich die heiligen Gefäße und die Paramente zeigen, und er wird erstaunen, wie gering alles ist, und welch ein Widerspruch zwischen dem Zwecke und der Erscheinung des heiligen Geräthes hervortritt. Sollte aber auch hier und da sich das Gegentheil finden, und an manchem Orte die angestammte Liebe der Gemeinden zu ihren Kirchen sich in einem gewissen Aufwande oder gar in einer Art von Verschwendung zeigen, so wird einem vielleicht greller und störender als hier und da die Armuth und Aermlichkeit, die Geschmacklosigkeit und Häßlichkeit des Geräthes bei allem Reichthum ins Auge fallen. Jeder Protestant traut sich ein Urtheil über die Einrichtung römischer Kirchen zu, und hält es für sein angestammtes Recht, seine Meinung sagen zu dürfen, während er es ganz erträglich findet, wenn seine eigene Pfarrkirche an Mängeln leidet, die für einen vorurtheilslosen Menschen vielleicht ebenso widerwärtig und noch widerwärtiger sind, als die geschmacklose Ueberladung einer römischen Kirche. Woher kommt nun dies alles und was man dieser Art mehr sagen könnte, wenn nicht von einer grenzenlosen Unwissenheit,
\quellseite{223}
und fast möchte man sagen Rohheit der Pfarrer und der ganzen protestantischen Geistlichkeit. Wie lange her ist es, daß, ich will nicht sagen die Gemeinden, denen, mit ehrenwerthen Ausnahmen, in ihren Kirchen alles gut genug zu sein pflegt, sondern die Pfarrer selbst es nur in den Bereich ihrer Fürsorge und ihres Eifers gezogen haben, wie die heiligen Geräthe ihrer Kirchen beschaffen seien. Man kann es gewiß eine ganz neue Lebensregung in unserer Kirche nennen, daß sich hier und da nur Unzufriedenheit und einiges Streben kund gibt. Wo sich aber auch dieses Streben zeigt, hat es denn Licht vor sich, und wissen denn wirklich die Pfarrer, wie ihr heiliges Geräthe beschaffen sein soll? Ganz abgesehen von der Schönheit, nötigt doch schon die Praxis, nicht alles sein und gehen zu lassen, wie es mag und kann, und wer einigermaßen darauf ausgeht, bequem und zweckmäßig amtieren zu können, dem möchten sich doch gewisse Anforderungen an die Beschaffenheit der Geräthe aufdringen. Ist nun aber ein Diener der Kirche nicht bloß darauf bedacht, daß alles praktisch, sondern auch, daß es würdig und schön sei, und wird einer bei wirklich vorhandener Liebe zum Amt und den heiligen Sacramenten durch das Abgeschmackte, Unwürdige und Geringe der Geräthe und der Ausstattung seiner Kirche gestört, so muß doch auch der Drang und das Verlangen nach Abhilfe und Besserung erzeugt werden; nicht bloß der Schönheitssinn, sondern die Frömmigkeit selber muß gegen alles allzu geringe, unwürdige oder gar abgeschmackte Geräthe protestiren. Kommen nun solche Regungen im Herzen eines protestantischen
\quellseite{224}
Geistlichen empor, was wird er thun? Da er von alle dem nichts gelernt hat und deshalb nichts versteht, so wird er bei seinem Eigendünkel und seiner Willkür in die Schule gehen; vielleicht traut er auch seiner Frau, die selbstverständlich noch weniger weiß als er, Takt und Geschmack zu, und läßt sich von ihr berathen, oder er fragt seinen Meßner, oder tappt und greift sonst im Nebel herum, da er nicht einmal eine Erkenntnißquelle zu finden weiß, aus der er schöpfen könnte. Die Früchte einer solchen Unwissenheit und Unberathenheit kann man allenthalben finden und, wenn man will, aufzeigen.

Die römische Kirche, welche durchaus nicht unberathen ist, sondern in ihren kirchlichen Büchern Anweisung und Unterricht genug aufgespeichert besitzt, hat in diesem Stücke nicht weniger als die protestantische eine Zeit des Verfalls gehabt, und auch sie hat sich im allgemeinen noch nicht wieder zurechtgefunden, aber die Reaction zum Besseren entstand doch zuerst in ihr, und sie hat es bei ihren reichen Quellen allerorten leicht, sich zu erholen, während die protestantische Kirche fast gar keine Quellen hat, und die reformatorischen Kirchenordnungen so arm sind, daß sich für Anordnung der Kirchen und ihrer Geräthe aus ihnen und der gleichzeitigen Literatur schier nichts schöpfen läßt. Auch die wenigen Studien, die in den ersten Jahrhunderten nach der Reformation auf diese Seite hin gemacht wurden, sind verschollen und vergessen, ganz abgesehen von ihrem Werthe, wie groß oder klein derselbe gewesen sein möge. Später trat dann die jämmerliche Zeit ein, in welcher man an dergleichen auch nicht mehr dachte. Uns muß es daher
\quellseite{225}
in diesem Stücke gerade so gehen, wie in Betreff des Kirchenbaues. Noch vor kurzer Zeit plagte man sich mit Beantwortung der Frage, ob nicht der protestantische Kirchenbau ein ganz anderer als der römische, ob nicht bei uns die Kanzel ebenso das Maßgebende sei, wie bei den Römischen der Altar, und es ist schier ein Wunder, daß man so schnell erkannte, es sei bei uns, den Lutheranern, wie in der kirchlichen Führung das Sacrament, so im Kirchenbau der Altar das Höchste und Maßgebende nicht weniger als bei den Römischen. So werden wir uns auch in Gottes Namen bequemen müssen, zu den Schriften der alten Kirche und ihren Traditionen rücksichtlich der Kirchengeräthe zurückzukehren, und was praktisch und schön sei, aus der Erfahrung der Jahrhunderte zu lernen. Es muß nur alles mit reformatorischem Sinn studiert, und liegen gelassen werden, was dem Dogma der reformatorischen Kirchen widerspricht.

Wir wollen ein Beispiel geben. Der römische Taufstein muß ein wirklicher Stein sein, weil die römische Kirche das Taufwasser im Taufstein jährlich nur einmal consecrirt und zum Gebrauch für’s ganze Jahr im Taufstein aufbewahrt. Aus diesem Grunde hat der römische Taufstein einen verschließbaren Unterdeckel, der auch wirklich von Taufe zu Taufe sorgfältig verschlossen und der Schlüssel von dem Pfarrer selbst aufbewahrt wird. Ganz anders ist der Taufstein der lutherischen Kirche. Da in ihr der Grundsatz herrscht, daß eine Consecration der Elemente nur für die Dauer jeder einzelnen sacramentlichen Handlung wirksam sei, so hat sie von ihrem
\quellseite{226}
Anfang her die jährliche oder doch für längere Zeiten geltende Consecration des Taufwassers verworfen, ja sie hat in grellem Widerspruch gegen ihr eigenes Verfahren beim heiligen Abendmahl die Consecration des Taufwassers ganz und gar aufgegeben und consecrirt nicht mehr, ohne daß den meisten Pfarrern nur einfällt, daß das Wasser so gut als beim Sacrament des Altars Brod und Wein consecrirt werden sollte. Bekanntlich haben erst Männer wie Hößling in der neuen Zeit die Kirche auf die rechte Bahn gewiesen. Dieser Grundsatz der lutherischen Kirche und ihre daraus folgende Praxis hat den protestantischen Taufstein zu etwas ganz anderem gemacht. Da man immer nur soviel Wasser braucht, als nötig ist ein einzelnes Kind zu beträufeln, so braucht man gar keinen steinernen Stein, es kann ein hölzerner sein oder gar keiner, man kann auf jedem Tisch taufen; daher es auch wirklich Tauftische gibt, die dann bequemer zum Ein- und Auswickeln der Kinder sind (wo man nemlich die Kinder zur Taufe auswickelt). Man braucht aber auch keine Tische, und es ist ganz genug, wenn man nur ein Taufbecken hat, das versteht den Stein und den Tisch, und wenn der Meßner es beim Taufen unterhielte, der Pfarrer selber die Taufkanne handhaben würde, so könnte man desto bequemer nach der Taufe das Wasser gleich weggschütten (!). Denn das Taufbecken hat doch seinen alten Brauch verloren und wäre so protestantisch sachdienlich angewendet. Bei einem römischen Taufstein wendet man das Taufbecken allerdings auch an, das Wasser, womit getauft wurde, aufzufangen, da es aus dem Taufstein geschöpft wird
\quellseite{227}
und nach dem Gebrauch nicht wieder in das übrige consecrierte Wasser kommen soll. Bei uns jedoch ist der Gebrauch des Taufbeckens kein verhütender, sondern es steht bloß im Dienste der Schicklichkeit, da man doch das Wasser nicht auf die Erde will fallen lassen. Wiewohl auch das am Ende fast ebenso gut geschehen könnte, als wenn der Meßner die paar Tropfen Taufwasser hernach auf den Kirchenboden oder vor die Kirchenthür hinschüttet: er hat eben nicht soviel Takt und Gefühl, das Element nach dem Sacrament mit Ehrerbietung zu behandeln, weil es doch zum Besten gedient hatte, wozu es dienen konnte. — Hieraus kann man sehen, wie sehr verschieden ein protestantischer Taufstein von einem römischen ist und sein muß. Dennoch aber wird der reine kirchliche Geschmack und die rechte Praxis der heiligen Taufe, sowie die Pädagogik der Kirche sich niemals mit dem armen Taufbeckensystem vertragen, sondern für jede Kirche ein Baptisterium fordern, einen Ort, der symbolisch am Eingang der Kirche gewölbt, umschlossen und mit einem erhöhten oder vertieften Taufstein versehen ist. Sobald man aber soweit gekommen ist, diese Forderung zu machen, wird man von dem ganzen Bau des römischen Taufsteines nach weggelaßenem Unterdeckel und Schloß das Beste und Schönste lernen, was es des Falls geben kann. Man wird den elenden Widerspruch eines hölzernen Taufsteines fallen lassen, ihn mit der alten Piscina unter seiner Basis versehen, protestantisch über dem Taufstein taufen, das abfließende Wasser durch die Röhre im Taufstein am besten in die Piscina fallen lassen (damit ist aller Miß%
\quellseite{228}
brauch, aller abergläubische Gebrauch des Taufwaßers, alle Unehrerbietigkeit roher Meßner abgeschnitten) und man wird dann den Deckel des Taufsteins nach alter Weise in voller oder geminderter Zier, mit oder ohne Flaschenzug anbringen können, zur Freude der Gemeinde und zur Verherrlichung des Sacramentes.

Mit diesem Beispiele meinte ich zu zeigen, was ich unter reformatorischem Studium des alten liturgischen Geräthes verstehe. Es ist schier mit allen Geräthen und mit allen protestantischen Paramenten dasselbe, woran mit dem Altar, mit Ciborium und Kelch, und man muß es in der That, namentlich in dieser Zeit eingerissener Unwissenheit und Rohheit, einem jeden jungen Geistlichen zur Pflicht machen, Studien über die heiligen Geräthe zu machen und das Seine beizutragen, daß die uralte Weisheit, Schicklichkeit und Schönheit gereinigt und gesäubert auch in unsere Kirchen einziehe und Platz gewinne. — Die Schmach, daß alte Geistliche nicht einmal wissen, wie sie einen Altar-Anzug fertigen lassen sollen, ja nicht einmal wie ein Kelch beschaffen sein soll, und wo man ihn anfaßt etc., dürfte von unserer Kirche und ihren Dienern genommen werden.

Man könnte sagen, die Kirchenoberen sollten Bestimmungen geben. Aber die Kirchenoberen sind in der Regel altersreife Männer, die selbst aus den Zeiten der Unwissenheit stammen, und es könnte die Sache nicht gründlicher verderbt werden, als wenn deren subjective Meinungen, ehe es noch in der Kirche für diese Dinge Tag geworden ist, den Gemeinden
\quellseite{229}
in Form von Anordnungen aufgedrängt würden. Gott Lob sind hie und da schöne Anfänge vorhanden, welche aber durch die Hand der Gewaltigen schnell unterdrückt werden könnten. Man hat in der neuesten Zeit Erfahrungen genug gemacht, wohin es kommt, wenn in der Form der Anordnung und des Befehls das Gute und Nöthige eingeführt werden will. Da kann man sich die Erfahrung, die bittere, ersparen, die kommen würde, wenn man das Schickliche und Schöne, auch wenn man es selbst wüßte und verstünde, befehlen würde. Auf dem kirchlichen Gebiete ist die Freiwilligkeit zu ehren, zu wecken, zu erziehen, nicht aber der Gehorsam gegen Befehl zu fordern. Nachdem es einmal so weit gekommen ist, daß der heilige Geist das Gefühl für das kirchlich Schickliche und Schöne in Betreff der Geräthe und Paramente geweckt hat, wird er auch weiter helfen, und alles, was man von Regiments wegen thun sollte, wäre nur bescheidene Abwehr des offenbar Verkehrten.

Vielleicht ist es hier am Ort, junge Geistliche auf die Altarstudien von Laib\footnote{Studien über die Geschichte des christlichen Altars von Fr. Laib und Dr. Schwarz. Stuttgart. Rümelius Wttb. 1857.} und Schwarz hinzuweisen, die fast selbst reformatorisch sind und von allen aus der Zeit des Heiligen- und Reliquiendienstes stammenden Altarbildungen zu der schönen Einfalt der heiligen Mensa, des Tisches Jesu führen.

Hieher gehörte übrigens statt eines einzelnen Paragraphen, der zur Anregung geschrieben ist, ein ganzes Buch, eine Instructio pastoralis liturgica.

\quellseite{230}
\section[39. Bestandtheile des Gottesdienstes]{39. Von den Bestandtheilen des Gottesdienstes}

Von den ersten Gemeinden heißt es: „Sie blieben allezeit in der Apostel Lehre, in der Gemeinschaft, im Brodtbrechen und im Gebet\footnote{Ἦσαν δὲ προσκαρτεροῦντες τῇ διδαχῇ τῶν ἀποστόλων καὶ τῇ κοινωνίᾳ καὶ τῇ κλάσει τοῦ ἄρτου καὶ ταῖς προσευχαῖς. act. 2. 42.}.“ Etliche haben es angefochten, daß unter Gemeinschaft die Mittheilung des kirchlichen Almosens, mit Brodtbrechen das Sacrament des Altars gemeint sein solle, ohne es doch besser erklären zu können, und ohne den Sprachgebrauch der Schriften des Alterthums mehr auf ihrer Seite zu haben. Allein wenn auch die Auffassung des Spruches, zu der wir uns bekennen, nicht richtig wäre, wenn also auch nicht zugestanden werden könnte, daß in dem Spruche die Bestandtheile des gottesdienstlichen Lebens angegeben seien, so könnte man doch nicht leugnen, daß Apostellehre, kirchliches Almosengeben, Sacrament und Gebete die Hauptbestandtheile des gottesdienstlichen Lebens seien, die namentlich auch den Sonntag zum Sonntag machen, und die große Versammlung der Gemeinde zu dem, was sie ist.

Es lassen sich nun aber außerhalb der großen Versammlung, der Communio, ganz wohl Gottesdienste denken, bei welchen einer von den vier Bestandtheilen über die andern vorherrscht. Es gibt Predigtgottesdienste, bei welchen Gebet
\quellseite{231}
und Gesang Nebensache sind; Gebetsversammlungen, bei denen die Predigt und Apostellehre, — Abendmahlsfeiern, wo anderes zurücktritt; ja, es ließen sich auch Zusammenkünfte zum Almosen denken, bei denen die übrigen Stücke des gottesdienstlichen Lebens im Hintergrunde ständen. Die Trennung und Hervorhebung einzelner Theile ist auch gar nicht zu beanstanden. Doch ist sie keine vollständige, und es liegt am Tage, daß ganz kein Bestandtheil den andern entbehren kann. Dies gilt insonderheit von dreien. Wer kann zu irgend einem gottesdienstlichen Geschäfte zusammenkommen, ohne z. B. zu beten? Welche gottesdienstliche Versammlung wäre ganz ohne das Wort Gottes? Andacht und Gottes Wort sind von allen gottesdienstlichen Versammlungen unabtrennbar. So weit es nun natürlich ist und möglich, die genannten Bestandtheile zu trennen, mag man es thun, wie bereits gesagt, wenn man nemlich Grund dazu hat. Aber sie sollen auch vereinigt werden, und das geschieht, wie gleichfalls schon angedeutet, bei der Communio. Da erscheint das gottesdienstliche Leben in seiner Vollständigkeit. Die Gemeinde empfängt durch Wort und That des Herrn (Sacrament) und gibt durch Wort und That (κοινωνία), und die Begegnung Gottes, des Herrn, und seiner Gemeinde im feierlichsten Ernst ist Lebenshöhe. In der Communio ist nicht bloß ein Wortgetöne, nicht ein Feiern von allen Werken, da im Gegentheil die höchsten und besten Werke im tiefsten Frieden ausgeübt werden. Schier ist die Sache zu groß und hoch, um den Namen zu gebrauchen, aber wir wollen doch wagen, es so auszudrücken: Hier ist die höchste
\quellseite{232}
Poesie, und was Epik, Lyrik und Dramatik aus dem Stoff des gewöhnlichen menschlichen Lebens machen und bilden können, es reicht doch nicht an die Vereinigung aller der heiligsten Lebensblüthen und Früchte in der Communio. Kein Wunder, wenn sich nun in der Anordnung der vier gottesdienstlichen Bestandtheile zur Communio das höchste Kunstwerk erweist, das je aus Menschenhänden kam. Hier ist Natur eines Lebens, das übernatürlich ist, und bei und neben der höchsten Einfalt findet sich eine wunderbar reiche Pracht, welche werth ist, aus ihren verborgenen Kammern, namentlich der alten Liturgien des Morgenlandes hervorgeholt, beschaut und genossen zu werden.

\section[40. Die alte Communio]{40. Die alte Communio und überhaupt die Liturgie der Alten im Verhältniß zu unserer Zeit}

Die lutherische Kirche hat, was die Liturgie anlangt, insofern echt reformatorisch verfahren, als sie nicht vor großem Ueberdruß an der geistlosen Weise, in welcher man im 16. Jahrhundert bei den Römischen die Liturgie zu bloßen Formen und Formelndienst herunterwürdigte, das Schöne übersah und wegwarf, was auch noch von den Römischen so vielfach beibehalten war. Wenn auch Luther durch seine Schrift von 1526 („deutsche Messe“) über die Grenze des reformatorischen Verfahrens, betreffend die Liturgie, hinüber schritt und einen andern Weg betrat, als den er echt reformatorisch 1523 betre%
\quellseite{233}
ten hatte; so wirkte sein Beispiel doch nicht so entscheidend, daß man ihm allewege nachgefolgt wäre. Wenigstens in sehr vielen Landeskirchen enthielt man sich, irgend etwas wegzuwerfen, worinnen Segen sein konnte. Allein es fehlte denn doch das liturgische Studium und diejenige Einsicht in den Gang des großen liturgischen Ganzen der alten Zeit, auch bei den alten Lutheranern, zu sehr, als daß nicht schon damals manches untadelige und herrliche Erbstück der alten Zeit beseitigt worden wäre. Ich wenigstens kann nirgends finden, daß bei den lutherischen Agenden ein lutherisches Prinzip angewendet worden wäre, außer dem negativ dogmatischen. — Späterhin traten Umstände ein, durch welche die lutherische Kirche in liturgischen Dingen nothwendig ärmer werden mußte. Die lateinischen Gesänge verschwanden, eben damit der gregorianische Gesang, da man deutsche Texte zu den alten Noten schon wegen Versunkenheit der deutschen Sprache nicht schaffen konnte; an die Stelle des gregorianischen Gesangs trat das geistliche Volkslied, der sogenannte Choral, und in den zog sich die liturgische Kraft der Kirche zurück. Je länger, je mehr mußte man das Verständniß der alten Gottesdienstordnung, welche die Hauptsache bei den alten Liturgien war, verlieren und deshalb froh sein, von all dem Nichtverstandenen allmählich ein Stück nach dem andern zu verlieren; endlich mußte es so weit kommen, daß jeder Ueberrest der alten Zeit wie störend wirkte, und daß man nur die einfachste Vereinigung der obengenannten vier Hauptstücke des Gottesdienstes ertrug und schön fand. Man kann wohl nicht anders sagen, als daß noch vor
\quellseite{234}
kurzer Zeit reformierte Leerheit und Armuth das Ziel zu sein schien, nach welchem sich die lutherische Liturgie ausstreckte.

Dagegen konnte eine Reaction nicht ausbleiben. Hier und da fing einer an, die lutherischen Agenden einzusehen, und für den ersten Anfang erschien dann schon die in ihnen sich findende Fülle und Gliederung paradiesisch schön und hochpoetisch, im Vergleich mit der herrschend gewordenen Armuth und ihrem Schlendrian, der unerträglicher wurde, als es der Schlendrian der alten Fülle gewesen war. Aeußerte man sich nun warm für’s Gute der alten Zeit, so fanden die Zeitgenoßen schon die Liebe zu altlutherischen Agenden wie romanisierend, und man glaubte kaum, daß Fülle und Gliederung des Gottesdienstes lutherisch sein könnte. Die Unwissenheit und Unbekanntheit mit der eignen Armuth macht dreist und unverschämt genug, diejenigen als römisch auszuschreien, welche vom lutherischen Alterthum mehr wußten und für das, was sie wußten, warme Liebe bezeugten. Wurde nun schon die Liebe zum eigentlich Lutherischen so arg mißverstanden, wie konnte man böser Nachrede entgehen, wenn man bei aller offen dargelegten Mißbilligung römischen Verderbens der alten Liturgie doch auch noch bei den Römischen aus dem grauen Alterthum liturgische Anordnungen oder Erzeugnisse fand und pries, welche mit Unrecht antiquirt waren, und heute noch der Andacht ausgezeichnet dienen könnten, wenn man sie nur erst wieder mit Lust und Andacht brauchen würde. Nicht der geringste, sondern weitaus der bedeutendste Grund oder vielmehr Ungrund, warum so manche treue Lutheraner, welche mit allem
\quellseite{235}
Ernste die Resultate der lutherischen Reformation festhielten, für römisch ausgeschrieen wurden, ist ihre offen dargelegte Liebe und Bewunderung der alten Liturgien. Das übersah man natürlich, daß diese verschrieenen Männer die uralten morgenländischen Liturgien noch weit mehr, als die römischen bewunderten. Es fiel den Leuten nicht ein, daß man dieselben Männer mit dem gleichen Grunde auch der Hinneigung zur morgenländischen Kirche und zwar gerade zu ihren Secten (denn die ältesten Liturgien haben die Secten aufbehalten) zeihen konnte und mußte\footnote{S. Liturgiarum orientalium collectio opera et studio Eusebii Renaudotii Parisini. Editio secunda correctior. 2 tomi. Frankfurt a/M. 1847.}.

Gegenwärtig hat sich die Zeit gewendet. Wenigstens für die alte lutherische Liturgie hat man in weiteren Kreisen Partei genommen und sogar angefangen gegliederte Liturgien, zum Theil auch neu aus alten Stücken geordnete, den unbereiteten Gemeinden zuzumuthen. Dieselben Männer, welche zuerst wieder den Sinn für Liturgie geweckt hatten, mußten nun umgekehrt warnend am Wege des neuen liturgischen Fortschritts stehen. Es ist wahr, daß Volk gewöhnt sich endlich an liturgische Formeln, welche man ihm alle Sonntage beharrlich vorträgt und läßt sich dieselben endlich wieder als etwas Unvermeidliches gefallen, aber hiermit ist eben doch nichts gewonnen. Was hilfts, daß man sich an Unverstandenes und Mißliebiges gewöhnt? Wenn die amerikanischen Gemeinden von Frankenmuth, Frankenhilf, Frankenlust u. s. w. die gegliederte
\quellseite{236}
Liturgie der Löhe’schen Agende gern und freudig annahmen und übten, so ist das aus ihrer Zusammensetzung und der kirchlichen Heranbildung ihrer Glieder leicht zu begreifen. Es ist bei ihnen ein gewisses Maß von Liebe und Eingehen in die Liturgie, wodurch sie vor Ekel und Schlendrian bewahrt bleiben. Bei andern Gemeinden, die nicht dieselben liturgischen Voraussetzungen haben, muß man langsam gehen, nicht durch Zwang, sondern durch Verständigung und Ueberzeugung wirken, wie man aus Art. 28 der Augsburger Confession solche Weisheit wohl lernen könnte. Man muß dabei aber allerdings ein sicheres Ziel und Ganzes vor Augen behalten, und einen Weg einschlagen, auf dem man zuerst die Hauptsache in gute Ordnung bringt, und einen einfachen Gottesdienst anfangen kann, aus welchem sich hernach die völligere Zier und Schöne mühelos entwickeln läßt.

Bei weitem am wichtigsten ist es, das gemeine Gebet in seiner schönen Mannichfaltigkeit die Gemeinden zu lehren, und es mit ihnen zu üben. Eine Gemeinde, welche Lust hat am Bitt-Gebet, Lob und Danksagung für alle Menschen, ist eine wahrhaft liturgische Gemeinde, auch wenn ihre Gottesdienste von der antiken Herrlichkeit der Liturgie nichts haben noch üben. Wer seine Gemeinde zu dieser Stufe geführt hätte, hätte auch den wesentlichen Mangel der gegenwärtigen Gemeinden, welche nur hören und innerlich angeregt werden, nicht selbst thätig sein und drum nicht selbst beten wollen, überwunden. An das Gebet schließt sich von selbst der Psalm, der Hymnus, das Credo u. s. w. an, und es ist der geistigere
\quellseite{237}
und geistlichere Fortschritt vom Beten zum Singen, nicht vom Singen zum Beten, weil das Singen durch die Schönheit des Klangs mehr vom Sinn des Gesangs abzuziehen, als in denselben hineinzuziehen pflegt. Nach dem Gebete dürfte die κοινωνία in Angriff kommen, es ist gewiß wohl gethan, die Gemeinden aus den Episteln Pauli über das Opfer des kirchlichen Almosens zu belehren, und dafür zu gewinnen, sie zu einer lauteren, von Werktheiligkeit entfernten Uebung desselben anzuleiten.

Was die Feier des heiligen Abendmahls anbelangt, so bringen die Gemeinden auch jetzt noch so viel Sinn und Ehrerbietung zu derselben mit, daß erfahrungsmäßig auch die, welche im übrigen Gottesdienst schnell Romanismus wittern, bei dieser heiligen Handlung eine reichere Entfaltung heiliger Formen und heiligen Inhalts ertragen.

\section{41. Vom täglichen Gottesdienste}

Dieffenbach läßt in seiner Hausagende den täglichen Hausgottesdienst nach Formen vor sich gehen, die der Communio entlehnt sind. Hätte er die lutherischen Gottesdienste am Morgen und Abend und ihren Zusammenhang mit den alten Horen gewürdigt, so würde er ohne Zweifel der Communio ihre Formen gebührend gelassen, und für den Hausgottesdienst lieber aus der lutherischen Messe und Vesper geschöpft haben, wiewohl sich der Hausgottesdienst auch nicht
\quellseite{238}
nothwendig an diese anschließen muß, weil er keiner kirchlichen Form bedarf, und sein eigenes Recht hat.

Anders ist es natürlich, wenn man tägliche Hausgottesdienste in der Gemeinde herstellen will. Für diese Absicht können die Metten und Vespern der alten lutherischen Kirche Winke genug geben. — Mit der Lehre von der Verdienstlichkeit der Werke, also auch des Gebets, sank auch, und zwar sogleich mit Eintritt der Reformationszeit, die regere Theilnahme am Gottesdienst. In vielen Kirchenordnungen erscheint Mette und Vesper rein als Gottesdienst der Schule, weshalb auch der Fall besonders vorgesehen ist, daß Erwachsene bei der Mette und Vesper erscheinen. Da sollten die Lectionen auch deutsch gelesen werden, welche sonst nur lateinisch recitiert wurden. Ob man nun bei Einrichtung täglicher Gottesdienste heut zu Tage mehr Glück haben werde, ist zweifelhaft. Hie und da wird sich ausnahmsweise eine größere Theilnahme finden; insgemein wird man schwerlich von großem Glücke sagen können. Wäre der Gedanke des Opfers in der lutherischen Kirche anerkannt und mehr in Uebung, so würde ein größerer und kräftigerer Trieb da sein, zusammen zu beten, wenngleich dieser Gedanke eine solche Höhe hat, daß man von der Menge nicht hoffen kann, sie werde sich auf demselben halten können. Sie wird aus angeborner Trägheit immer auf eine Werkheiligkeit gerathen und auch das Beste verkehren, denn was wahrhaft geistlich ist, ist eben nichts für Massen.

Eine ganz andere Frage ist es aber, ob man deshalb die täglichen Gottesdienste, nemlich wo man im Falle ist,
\quellseite{239}
Einrichtung zu machen, ganz unterlassen soll. Finden sie auch keine allgemeine Theilnahme, so finden sie doch ein gewisses Maß von Theilnahme und eröffnen vielen die Möglichkeit der Erbauung. Auch wird ein eifriger Pfarrer vielleicht liturgische Vereine bilden können, welche nicht bloß zu liturgischer Gesangübung, sondern ganz ernstlich zum Gebete zusammenkämen, und durch welche in der Gemeinde Theilnahme und Geschick für Liturgie geweckt und erhalten werden könnte. Es ist gewiß kein Grund vorhanden, etwas Gutes bloß deshalb zu unterlassen, weil es insgemein unterlassen wird; im Gegentheil, es muß das Gute das Recht der freien Theilnahme in der Gemeinde behalten und genießen und darum müssen auch freie Vereine ihre Berechtigung innerhalb der Gemeinde haben. Ist das im allgemeinen wahr, so sind doch auch liturgische Vereine zulässig und anzuerkennen. Für sie könnten die täglichen Gottesdienste Sammelpunkte sein, wenn es auch nicht Satzung werden dürfte, daß ein Vereinsglied jeden täglichen Gottesdienst besuchen müßte. Es soll damit ohne alle sanguinischen Hoffnungen nur ein Mittel mehr angegeben werden, dem täglichen Gottesdienste Theilnahme zu erwecken, und zwar eine freie und ernste, ohne daß deshalb hart darauf gedrungen werden sollte. Außerdem wird die Schule wie in der Reformationszeit, so auch jetzt noch am leichtesten zum täglichen Gottesdienste zu bringen sein, wiewohl einer Schule so wenig als einem andern Menschenhaufen zuzutrauen ist, daß sie niemals aus dem Gebete und Gesange ein Geplärr und einen Singfang machen werde. Was den Mönchen bei den Choren,
\quellseite{240}
das ist nicht minder den Schülern bei dem täglichen Gottesdienst zuzutrauen und nachzureden.

Was die Anordnung des täglichen Gottesdienstes anlangt, so hat man Reichthums halber eben so vorsichtig zu gehen, als bei der Communio\footnote{Vgl. Löhe’s Agende für christliche Gemeinden des lutherischen Bekenntnisses. Thl. I. (Nördlingen 1853) p. 69, die Unterweisung über Matutin und Vesper und die verschiedenen Anordnungen der 4 Theile derselben (Psalm, Lection, Hymnus Oration). — Sehr spürt man bei Ueberlegung über die alte Mette und Vesper unsere große Versündigung, daß wir die Psalmen haben völlig fallen lassen und die große, schwere Unterlassungssünde, daß wir bei bereits wieder gewordenem besseren Lichte, so gar nichts thun, die Psalmen wieder in öffentlichen Gebrauch zu bringen. Wir könnten eher alle unsre Gesangbücher missen als den Psalter.}; es gibt auch hier eine vorbereitende einfache Form, die man um so unbedenklicher erwählen kann, als auch bei den alten täglichen Gottesdiensten ohnehin kein so zusammenhängender Gedanke herrscht, als bei der Communio.

\section[42. Mangel des Symbolischen]{42. Mangel des Symbolischen überhaupt und des symbolischen Handels insonderheit}

Ganz entblößt von allem symbolischen Handeln ist die Liturgie der lutherischen Kirche nicht; doch ist sie im allgemeinen, namentlich bei den Benedictionen so ganz auf den alleinigen Gebrauch der Rede verwiesen, daß dadurch die Benedictionen wenig selbstsprechend und ansprechend werden. Es kann dies
\quellseite{241}
am deutlichsten an dem Beispiel einer protestantischen Gottesäckerweihe oder Kirchenweihe gezeigt werden. Der Consecrator steht an seinem Ort, redet und betet von da aus, und alles, was er thut, geschieht durchs Wort. Wenn nur z. B. die singende Procession, wie sie doch auch bei Leichen in Gebrauch ist, angewendet werden könnte, und noch besser eine singende den Ort, das Gebäude von außen oder innen umziehende Procession, deren Gesang an bedeutsamen Stellen vom Gebet des Consecrators unterbrochen würde; so könnte eine solche Weihung bei weitem entsprechender und in ihrer eigentlichen Bedeutung kenntlicher werden. Es ist nicht Absicht, diesen Umstand hier des weiteren auszuführen, aber angeregt dürfte der Gedanke wohl werden, zumal die lutherische Kirche früherhin gar manches Symbolische hatte, was immer mehr verschwindet. Wir erinnern z. B. an das Westerhemd bei der Taufe. Aber freilich, an die unschuldigste Ceremonie hängt sich Aberglaube nur desto leichter an, je entsprechender die Absicht und je sprechender sie ist.

\section[43. Sogenannte Principien der Liturgie]{43. Sogenannte Principien der Liturgie in der lutherischen Kirche}

Freiheit und Liebe nennt man die Principien der Liturgie in der lutherischen Kirche. Die Freiheit wirkt Mannigfaltigkeit, die Liebe verhütet, daß die Mannigfaltigkeit nicht allzubunt werde. Das Princip der Freiheit ist leicht
\quellseite{242}
sehr mißverständlich auszudeuten. Wenn Luther die Freiheit rühmt, so meint er damit nur dem Zwang der Menschenlehren entgegenzutreten und dem Gewissen seine Rechte zu wahren, welches Christus befreit hat, damit wir nicht mehr der Menschen Knechte würden. Seine Meinung ist aber keineswegs die, daß unter dem Namen der christlichen Freiheit jede liturgische Willkür Lob und Ehre haben solle. Auch in erlaubten und bloß menschlichen Dingen gibt es Regel und Maß, die man nicht ungestraft übertreten, verachten und mißhandeln darf, — und jedes Gebiet des menschlichen Lebens und Wissens kann recht und falsch bebaut und bepflanzt werden. Nicht das ist Luthers Meinung, daß es einerlei und gleichgültig sei, in welchen Formen man Gott diene, sondern nur der Mangel eines göttlichen Gebotes, also eines positiven gesetzlichen Zwanges soll behauptet werden. Man kann daher die Freiheit der liturgischen Formen festhalten und doch nach den besten, schönsten, Gottes und Seiner Kirche würdigsten ringen, ja um die Gedanken und Ceremonien bitten, um Formen, welche der Andacht am förderlichsten sind.

\section{44. Liturgischer Vortrag}

Es gibt einen homiletischen, einen katechetischen und einen liturgischen Vortrag. Der leichteste, das ist am leichtesten zu fördernde und einzuhaltende ist der katechetische, weil er, wenn er dialogisch ist, dem Leben am nächsten steht, wenn er aber akroamatisch ist, durch die Absicht der Belehrung auf der
\quellseite{243}
bestimmtesten Bahn geführt wird. Zwischen inne steht der homiletische, weil er, obwohl der gewöhnlichen Rede verwandt, doch seinem Inhalte nach sich so sehr im genus dicendi grande bewegt, daß sich ein Kanzelton so leicht und sicher natürlich einschleicht. Die heilige Rede könnte und sollte ganz in natürlicher Individualität des Redners vorwärts schreiten, aber freilich in geheiligter Natürlichkeit, und eben das macht den homiletischen Vortrag schon schwerer als den katechetischen. Am schwersten aber ist ohne Zweifel der liturgische Vortrag. Der liturgische Vortrag ist schwierig, man singe oder spreche. Ein Alter sagt, man singe die Collecten und Lectionen, weil für die Liturgie kein Vorwalten der Individualität, keine subjective Betonung gestattet werden könne; so müsse gesungen werden, daß ein jeder seinen eigenen Ton hineinlegen könne. Um nun das zu erreichen, macht man den liturgischen Gesang heutzutage zu einem puren Recitativ und schult die jungen Geistlichen für denselben ein. Da gibt es dann junge Leute die geistlos an den Altären tönen, womit weder der Gemeinde gedient ist, noch Gott geehrt wird. Der liturgische Gesang ist trotzdem daß die Subjectivität zurücktreten muß, dennoch einer Betonung fähig, denn es gibt eben eine Betonung wie sie aus dem Verhältnisse der Gemeinde Gott dem Herrn gegenüber hervorgeht, und der Liturg am Altare muß sich in seiner geistlichen Vorbereitung gewissermaßen von seiner eigenen Stimmung und seinem Gefühlsleben loswickeln und den Ton suchen, der allen gemein ist und gemein sein sollte; diesen kirchlichen Sing- und Redeton zu finden, ist keinem ungeistlichen Menschen möglich,
\quellseite{244}
und wie ich glaube auch keinem, der die Aufgabe und das Bewußtsein nicht in sich trägt, die Stimme der Gemeinde von Gott zu sein. Zu der Zeit, wo auch die Protestanten an den Altären nach Tradition sangen und redeten, wo jeder jüngere Geistliche alle Tage hören konnte, was traditionell liturgischer Ton sei, hatte man noch weit mehr echt liturgischen Gesang und Vortrag als jetzt, wo man sich erst durch Studium und Reflection hineinfinden muß. Es mag vor den Ohren der jungen Geistlichkeit fast beleidigend klingen, aber am Ende ist doch etwas wahres daran, wenn man sagt: wie man liturgisch singen solle, könne man weitaus am besten kirchlich gebildeten römischen Priestern ablernen. — Indeß haben wir es hier mit dem Gesang am Altar nur deshalb zu thun, weil er mit dem redenden Vortrag am Altare den einen Hauptgrundsatz gemein hat, daß die Subjectivität, subjectives Fühlen und Betonen vermieden werden müssen. Um am Altare richtig lesen zu können, muß man im Falle sein, am Altare lesen zu sollen. Die Schule der liturgischen Sprache geht erst an, wenn die Studentenjahre hinter einem liegen, und das amtliche Leben eintritt. Ein neu ordinierter Diener Jesu und seiner Gemeinde muß dann in Demuth und Andacht faßen, was er soll, nemlich nicht zu eigener Andacht, sondern im Namen der Gemeinde zu Gott zu sprechen, und Gottes und seiner Apostel ipsissima verba dem Volke so vorzutragen, daß es die Stimme des guten Hirten erkennt und sie faßen lerne. Beherrscht dieser Beruf den Liturgen, so wird ihm ein Geistlicher, der mit und neben ihm dient, am leichtesten die
\quellseite{245}
Fehler und Mängel sagen können; das Gespräch mit einem solchen über das Gelingen oder Mißlingen seiner Leistung wird seine Fühlhörner erziehen, daß er je länger je mehr inne wird, was kirchlich liturgischer Vortrag ist. Gelinge es aber auch so viel es wolle, so wird doch Fähigkeit und Geschick schnell wieder abnehmen und verschwinden, wenn nicht die geistliche Vorbereitung auf den liturgischen Dienst und die eigene Andacht vor und bei demselben dem Liturgen treu verbleibt; wer im Namen der Gemeinde in Geist und Wahrheit beten und lesen will, der bedarf vor allem subjectiver Frömmigkeit und geistlicher Uebung; wem die eigene Andacht und Vorbereitung zum liturgischen Dienste kein Ernst ist, der lernt am allerwenigsten die große Kunst heiliger Diener Gottes, am Altar richtig zu lesen und zu betonen, und wer es unterläßt, sich vorzubereiten, der wird bald am Altare ein Handwerker werden, dem der Geist eine Erinnerung in sein Herz bringen wird: „Thu nur weg das Geplärr deiner Lieder, denn Ich mag dein Psalterspiel nicht hören.“ —

\chapter{Zur Seelsorge}

\section{45. Was verstehen wir unter Seelsorge?}

Kein Wort scheint leichter, als das Wort Seelsorge: Seelsorge ist die Sorge eines Pfarrers für die Seelen der ihm anvertrauten Heerde. Doch versteht man unter dieser
\quellseite{246}
Sorge nicht bloß das innerliche Denken und Sinnen eines seelsorgenden Mannes zum Besten seiner Gemeindeglieder, sondern auch die daraus hervorgehende Thätigkeit. In diese Thätigkeit kann man am Ende alles einschließen, was ein Pfarrer zu thun hat: jede Predigt und Rede, jede Katechese, jeder Gebetsgottesdienst, die sacramentliche Versehung ist im allgemeinen zur Seelsorge zu rechnen. Und doch versteht man unter Seelsorge im speciellen Sinn noch etwas anderes. Wer gepredigt, katechisirt, Kommunion gehalten, getauft etc. hat, pflegt das alles noch keine Seelsorge zu nennen. Unter Seelsorge im besonderen Sinn versteht man die auf die pastoralen Zustände der Gemeinde eingehende Bemühung und Theilung des göttlichen Wortes von Seiten des Pfarrers. Pastoral aber nennt man nicht jeden Zustand einer Seele, sondern nur diejenigen, in welche der Mensch auf dem Wege der Heilsordnung gegenüber der Führung des Wortes und Sacramentes gelangt. Ein Jeder, der von dem heiligen Geiste berufen wird, verhält sich gegen den göttlichen Ruf und die göttliche Führung in Folge seiner Eigenthümlichkeit auf eine eigne Weise, bringt seine besondern Hindernisse mit, und es gilt dann, ihn so zu führen, daß die Hindernisse gehoben und die Seele so einfach und sicher als möglich dem allgemeinen Ziele göttlicher Berufung entgegengeführt werde. Was nun der Pfarrer in diesem Sinne und zu dieser Absicht thut, das heißt Seelsorge im besonderen Sinn. —

Es gibt eine seelsorgerische Führung der ganzen Gemeinde und der einzelnen Gemeindeglieder, denn
\quellseite{247}
es gibt Eigenthümlichkeiten und Hindernisse, durch welche sich ganze Gemeinden von einander unterscheiden, und ebenso gibt es Eigenthümlichkeiten und Hindernisse, durch welche sich die einzelnen Gemeindeglieder von einander unterscheiden, weshalb man auch von pastoralen Zuständen ganzer Gemeinden ebenso reden kann, wie von den pastoralen Zuständen einzelner Christen. Z. B. gibt es Gemeinden, in denen das Fluchen, andere, in welchen das Trinken, oder das Stehlen u. s. w. herrschend geworden sind, und diese allgemein herrschenden Sünden geben alsdann der ganzen Gemeinde ihre besondere Physiognomie; denn eine herrschende Sünde ist wie der Gipfel eines Berges, durch welchen der ganze Berg erst die ihm eigenthümliche Gestalt gewinnt. Wenn ein Pfarrer die Führung einer Gemeinde übernimmt, so kann er höchstens durch die Mittheilung Anderer den pastoralen Zustand derselben wissen, aber diese Mittheilung kann ihm höchstens den Blick schärfen, er darf aber die Gemeinde nicht gleich vorn herein nach derselben behandeln; seine Pflicht ist's, sich vorurtheilslos zu erhalten, und eine besondere Behandlung erst dann eintreten zu lassen, wenn er aus eigener Anschauung und Erfahrung die geistliche Physiognomie der Gemeinde erkannt hat. Erst aus der eignen Bekanntschaft mit dem pastoralen Zustand der Gemeinde kann ihm eine Erkenntniß seiner pastoralen Ziele entstehen. Diese Ziele sind für die verschiedenen Gemeinden verschieden, und es kann daher kein Pfarrer in der Führung seines Volkes einfach einem andern folgen. Doch soll damit nicht gesagt sein, daß die verschiedenen Ziele einander völlig aus%
\quellseite{248}
schließen, oder daß nicht über die nächsten Ziele der Gemeinden hinaus andere gemeinschaftliche Ziele und zuletzt ein höchstes, allen gemeinschaftliches da wäre. Wenn z. B. der pastorale Zustand einer Gemeinde insbesondere durch Geiz und irdischen Sinn bedingt wird, so ist es gewiß das nächste Ziel für die pastorale Führung derselben, sie zur Barmherzigkeit und Mildigkeit zu erziehen. Ist nun aber auch dies nächste Ziel erreicht, so zeigt sich hinter demselben gewiß ein ferneres und höheres, denn die Gemeinde könnte barmherzig werden und die Genußsucht beibehalten wollen, die ihr vielleicht schon früher bei allem Sparen und Scharren und Geizen eigen war. So arbeite dann ein Pfarrer einer höheren Vollendung der Gemeinde dadurch entgegen, daß er sie mit aller Geduld zur Entsagung und Entbehrung und zu einem Leben anleitet, welches mäßig und nüchtern ist zum Gebet. Man könnte sagen, ein Pfarrer solle, weil er sich in der Schätzung des pastoralen Zustands der Gemeinde so leicht irren und in Folge dessen sich falsche Ziele stecken kann, ganz einfach, ohne Ziel zu setzen, thun was ihm unter die Hand kommt. Es mag auch Pfarrer geben, denen ein anderer Weg als der, ziellos zu thun was ihnen unter die Hand kommt, gar nicht übrig bleibt; im allgemeinen aber wird man doch diesen Weg nicht rathen können.

Sei langsam in deinem Urtheil über den pastoralen Zustand deiner Gemeinde und in der Aufstellung deines pastoralen Zieles: es sei dir eine hohe Angelegenheit, Zustand und Ziel recht zu erkennen. Bis du die sichere Erkenntniß hast,
\quellseite{249}
kannst du thun, was dir unter die Hände kommt, nach bestem Wissen und Gewissen; hast du aber einmal Gewißheit über den Zustand der Gemeinde und über dein nächstes pastorales Ziel, so wird sich daraus eine starke Kraft entwickeln, und es wird dir deine Erkenntniß wie ein Gurt deiner Lenden werden, gewisse Tritte zu thun.

Ganz ähnlich ist es mit der pastoralen Führung der Einzelnen; auch rücksichtlich ihrer wirst du dein Urtheil nicht beschleunigen dürfen; fixiere es nicht zu schnell; es ist weder der beste, noch der weiseste, der bald fertig ist mit Wort und Urtheil; der pastorale Zustand der Einzelnen ist nicht leichter, sondern schwerer zu erkennen, als der einer ganzen Gemeinde; du wirst nie dahin kommen, über alle und jede ein festes Urtheil zu haben, da sich ja auch die Zustände der Menschen nicht allezeit gleich bleiben und die Erkenntniß derselben ebenso wechseln und sich ändern muß wie sie selbst. Da der Einzelnen viele sind und es eine gewaltige Aufgabe für einen Menschen ist, so viele Einzelne richtig zu erkennen, und einem jeden das rechte pastorale Ziel zu stecken, so wäre es bei der Seelsorge der Einzelnen verzeihlicher als bei der einer ganzen Gemeinde, wenn man an der, so gar großen und vielfachen Aufgabe verzagte und sich mit dem wenigsten und geringsten genügen ließe. Dennoch aber gewinnt die seelsorgerische Thätigkeit an den Einzelnen erst dann den rechten Werth, wenn sie auf die Erkenntniß des obschon vielleicht wechselnden pastoralen Zustandes und Zieles gegründet ist; und wenn auch ein jeglicher Seelsorger bescheidenlich darauf verzichten muß, alle
\quellseite{250}
einzelnen so aus dem sicheren Mittelpunkte heraus zu leiten und zu führen, so muß es ihm doch eine bleibende Anforderung an sich selbst werden, so vielen als möglich so zu dienen, wie er es wünschen müßte, allen dienen zu können.

Steckt er sich damit ein hohes Ziel, so wird es ihm nicht bloß zur Aneiferung, sondern auch zur immer neuen Demüthigung dienen, und grade das ist ja für einen Mann, der amtshalber so gar viel Recht haben und behalten muß, weniger eine Strafe als ein Lohn seines Strebens.

\section{46. Pastorale und psychische Zustände}

Ein Seelsorger wurde zu einem tödlich kranken Bauersmann in einem seiner Parochialorte gerufen, um ihm das Sacrament zu reichen; er hielt im Beisein seines aufmerksamen Mesners mit dem Kranken das pastorale Gespräch desto vorsichtiger, weil er die leibliche Krankheit und ihre Wirkung nicht sicher wußte. Mit Ausnahme einer einzigen Rede, welche man als unpassend und ungeziemend gegen den Seelsorger hätte deuten können, sprach sich der Kranke so aus, wie man es für einen Sterbenden wünschen könnte; sein pastoraler Zustand schien der einer gnadenhungrigen Bußfertigkeit. Pfarrer und Mesner stimmten überein, daß man das Sacrament reichen dürfe, und so wurde es auch gereicht. Am andern Tage kam der Bruder des Kranken und bat im Namen des letzteren um einen wiederholten Besuch des Pfarrers: der Kranke wisse kein Wort davon, daß er gestern das Sacrament
\quellseite{251}
empfangen habe; aber er sei nun besonders darüber betrübt, daß er, wie ihm die Seinen gesagt hätten, dem Pfarrer eine ungeziemende Antwort gegeben habe, und möchte nun denselben um Verzeihung bitten. Während also der Pfarrer, seiner Meinung nach, den Kranken vorsichtig behandelt und den pastoralen Zustand desselben richtig erkannt hatte, zeigte sich’s nun, daß der Kranke gar nicht zurechnungsfähig gewesen und daß ein psychischer Zustand für einen pastoralen gehalten worden war.

Aehnliches kann man im 1. Theil von de Valenti’s Medicina Clerica pag. 273 ff. lesen. Eine Wöchnerin, bei der eine Milchversetzung vorgegangen war, ordnete auf eine unbegreiflich eitle Weise ihre Leichenangelegenheit, ihren Sarganzug, disponierte über die, von ihr nachzulassenden Kleider u. s. f. Auf die mißbilligende Einrede eines christlichen Chirurgen fing sie dann an, sich zu bekehren, und es erfolgte nun eine der Heilsordnung gemäße, völlige Geschichte einer anscheinend gründlichen Buße und gläubigen Ergreifung Christi des Herrn. Was aber geschah? Als die Gefahr vorüber war, wußte sie von beidem gar nichts, weder von ihrer Kleidereitelkeit, noch von ihrer Bekehrung, und sie war im gesunden Zustande, wie sie vor der Krankheit gewesen war. Was sich wie die wechselnden Zustände eines Menschen, der von der Eitelkeit zu Gott heimkehrt, gezeigt hatte, war eitel Spiel einer delirierenden Seele gewesen. Die vorhandenen psychischen Zustände waren pastoralen Zuständen täuschend ähnlich. — Wie vergeblich war in beiden Fällen die Seelsorge, und wie
\quellseite{252}
viel besser wäre es gewesen, wenn der sterbende Bauer und die Wöchnerin ganz einfach in's Gebet genommen worden wären; aber freilich konnte man in beiden Fällen, wenigstens bei dem Bauern ganz gewiß, nichts anderes thun, als geschehen war.

Psychische Zustände sehen also, wie gesagt, den pastoralen sehr oft täuschend ähnlich, und hat man das einmal erkannt, so steht man besonders an Krankenbetten desto zurückhaltender und ärmer und dringt desto mehr darauf, daß man die gesunde Zeit des Lebens benutze, einen festen Grund des Heils zu legen.

Wenn nun die Verwechselung der psychischen Zustände mit pastoralen nur bei Kranken vorkäme, so wäre doch nur die Krankenseelsorge damit zu einem bescheidenen Werthe heruntergesetzt, während doch jedermann geneigt ist, auch die meisten Seelsorger selbst, ihren Werth recht hoch anzuschlagen. Schon das ist sehr empfindlich. Allein die psychischen Zustände durchkreuzen die pastoralen auch bei den Gesunden mehr als man denkt. Dadurch ist die Aufgabe des Seelsorgers, den pastoralen Zustand richtig zu erkennen und ihm gemäß das pastorale Ziel zu stecken, im allgemeinen vielfach erschwert. Die psychischen Zustände sind vielfach vom Geiste des Menschen, an den sich die Seelsorge allein zu richten hat, unabhängig, entspringen sehr häufig aus rein leiblichen Ursachen. Sie sind wie Staub um den Wagen des Geistes her, und äffen den Menschen, der sich selbst beurtheilen will, geschweige einen andern, der ein vorurtheilsfreies Urtheil über seinen
\quellseite{253}
Nächsten braucht. Schon diese sich immer und immer aufdringende Wahrnehmung nötigt den Seelsorger, das apostolische Wort: „Seid langsam zum Reden (zum Urtheil) und langsam zum Zorn“ sich besonders anzueignen. Die Schwierigkeit der Privatseelsorge tritt damit in’s Auge und rechtfertigt den Wunsch, daß sich die Seelsorger auch mit dem Studium der psychischen Zustände beschäftigen möchten.

\section{47. Wechselwirkung Leibes und der Seele}

Das Höchste, was der Mensch erreichen kann, ist Unabhängigkeit seines Geistes von den Einflüssen des Leibes und der Seele, das Höchste, aber selbstverständlich auch das Schwerste, nur in den seltensten Fällen Erreichte. Unabhängigkeit des Geistes von diesen Einflüssen und freie, willige, getreue Abhängigkeit des Geistes von dem göttlichen Worte und Willen muß das endliche Ziel aller Selbsterziehung, so wie aller Erziehung und Seelsorge durch andere sein. Wenn das Ziel von niemand völlig erreicht werden kann, so muß es doch immer angestrebt werden, und es kann auch von Stufe zu Stufe, wie man davon so viele Beweise im Leben heiliger Menschen hat, mit Erfolg angestrebt werden. Etwas ganz anders als die Freiheit des Geistes von den leiblichen und psychischen Einwirkungen ist Freiheit von der Wechselwirkung Leibes und der Seele auf einander. Diese ist unter keinen Umständen zu erreichen. Die Wechselwirkung besteht durch Gott. Der Jammer des Menschen besteht auch keineswegs darin, daß diese Wechselwirkung vorhanden, sondern daß der
\quellseite{254}
Geist nicht unbefangen über ihr steht, sondern ohne Ende von allem das Psychische, wo nicht gar das Leibliche mit dem Geistigen und Geistlichen vermengt wird. Hätte der Geist seine volle Freiheit von dem Psychischen und Leiblichen, so hätte er auch die volle Macht über beides, er könnte dann auch in den Regionen des Gemüths- und Seelenlebens so wirken, wie es sein sollte; denn der Geist mit seinem Denken und Wollen sollte den entscheidendsten Einfluß auf das Befinden der Seele, ja des Leibes haben. Wäre der Geist von Gottes Wort beeinflußt, so würde gleichermaßen Seele und Leib vom Geiste und dem göttlichen Worte bestimmt werden, und der ganze Mensch auf der Bahn der Vollendung vorwärts gehen. Der Mensch würde dann, wenigstens dem Anfang nach, seinem herrlichen Zustande entsprechen, der ihm in der Auferstehung der Gerechten als ewiges Erbe geschenkt werden soll. Denn in jenen seligen Zeiten wird Geist und Seel und Leib in vollkommener Harmonie leben, und der Geist wird König über Leib und Seele sein.

Um aber diese Freiheit und Herrschaft des Geistes anzubahnen und um die Vermengselung des Psychischen mit dem Geistigen und Geistlichen abzuscheiden, sollte man sich bestreben, diese Wechselwirkung kennen zu lernen. So weit sie erkannt wird, wird man versuchen können, sie dem Geiste unterthänig zu machen und den Geist von ihrem Einfluß zu befreien. In dieser Meinung hat sich der Verfasser dieser Zeilen schon vor Jahrzehnten an eine Auctorität gewendet, deren Urtheil ihm ebensowohl vom Standpunkt des Arztes, wie von
\quellseite{255}
dem des practischen Theologen Anerkennung zu verdienen schien; er legte die Frage zur Beantwortung vor: ob es nicht gut sein würde, wenn sich zukünftige Seelsorger wenigstens bis zu dem Grade medicinisch ausbildeten, daß sie die vorhandenen leiblichen Krankheiten erkennen und die psychischen Wirkungen derselben voraus wissen und erwarten können.

Die Antwort fiel jedoch durchaus nicht nach dem Sinne der Frage aus: der Pfarrer solle ganz einfach seines Amtes warten und nach bestem Wissen und Gewissen handeln, ohne sich um die Krankheit seines Pfarrkindes zu kümmern. Als Muster wurde dem Fragsteller Demler in seinen bekannten Schriften hingestellt. Demler war dem Frager längst bekannt, und was an der Behandlung und Seelsorge des Mannes zu rühmen war, das war von ihm längst anerkannt und auch benützt; er wußte nur nicht, wie ferne die Hinweisung auf Demler seine Frage beantworten sollte. Die Antwort war für ihn gar keine, obwohl er jungen Geistlichen, an denen er eine Neigung zu medicinieren wahrgenommen hatte, ganz in derselben Weise abwehrend und auf Einfalt im Amte dringend selbst gerathen hatte. Er hatte zur Zeit seiner Frage alles medicinieren längst aufgegeben, obwohl er als Vicar in einer großen Pfarrei längere Zeit unter Beiziehung eines bekannten Arztes mit dem allerbesten Erfolg homöopathischen Rath und Medicin gegeben hatte. Alles, was er wollte, war, sich und andere Amtsträger vor den Täuschungen psychischer Zustände an den Krankenbetten sicher zu stellen. Er wendete sich daher an andere medicinische Auctoritäten, merkte aber bald, daß
\quellseite{256}
bei den Aerzten selbst insgemein überaus wenig Kenntniß der psychischen Wirkungen leiblicher Krankheiten vorhanden war, daher er sich nicht wundern konnte, den gewünschten Rath nicht zu bekommen. Im Verlauf der Zeit lernte er Albert Mathias Vering’s Buch „Ueber die Wechselwirkung zwischen Seele und Körper im Menschen,“ Leipzig 1817, kennen. Das endlich gab ihm, was er wollte. Und obwohl er von einem Arzte jener Zeit nicht erwarten konnte, was man ja auch 50 Jahre nachher nicht erwarten darf, nämlich Anerkennung und Gebrauch der größten aller psychischen Mittel, des göttlichen Wortes und Sacramentes; so erquickte ihn doch der Ernst des wohlmeinenden Mannes, und daß er bei aller Willkür und Subjectivität seiner psychischen Rathschläge so viel Bescheidenheit und gelegentlich auch Respect vor der Religion hatte und äußerte. So groß erschien ihm die Bedeutung der Wechselwirkung Leibes und der Seele für die Seelsorge namentlich der Kranken, daß er auf Grund dessen, was er aus Vering gelernt hatte, den Diaconissenschülerinnen zu Neuendettelsau einen eingehenden Unterricht darüber nicht bloß gab, sondern auch dictierte. Bei diesem Unterrichte suchte er Vering’s Mangel zu erstatten und an die Stelle der rationalistisch-subjectiven Rathschläge und psychischen Mittel die pastorale Arznei aus Gottes Wort zu setzen. So entstand das Dictat, welches der ersten Auflage dieses Buches unter dem folgenden Titel als Anhang beigedruckt wurde (siehe erste Auflage p. 195 ff.): „Von dem Einfluß der leiblichen Krankheiten auf das psychische Befinden des Kranken, sowie von der Anwendung geistlicher Mittel zur Hebung der daraus hervorgebenden
\quellseite{257}
Gefahr der Seele.“ Ihm selbst wurde die Beschäftigung mit diesem Dictate bei der Seelsorge zu großem Segen, und er wurde dadurch insbesondere zur Bescheidenheit in der Anwendung des göttlichen Wortes geleitet. Seitdem jenes Dictat gedruckt wurde, sind acht Jahre vorübergegangen und da nun eine neue Auflage des „Evangelischen Geistlichen II“ zu besorgen war, mußte die Frage erwogen werden, ob, und wenn ja, in welcher Gestalt das genannte Dictat der neuen Auflage wieder beizudrucken wäre. Er bat daher einen ihm vertrauten wissenschaftlichen Arzt, das Dictat von seinem Standpunkt zu prüfen, und bekam dann die Antwort, daß Vering’s medicinischer Standpunkt dem der gegenwärtigen Aerzte nicht mehr entspreche. Dieses Urtheil wurde ihm durch Aufzeigung von Beispielen glaubwürdig gemacht, und die Erwägung desselben wirkte den Entschluß, um so weniger einen neuen Abdruck des Dictates zu veranstalten, als dasselbe in einem noch nicht vergriffenen Sonderdruck vorhanden ist, welcher, wie das ganze Buch vom „evangelischen Geistlichen“ anno 1858 bei G. G. Liesching in Stuttgart (jetzt C. Bertelsmann in Gütersloh) erschienen ist. Zugleich aber entschloß sich der Verfasser, in der neuen Auflage des zweiten Theils vom „evangelischen Geistlichen“ sich so ernstlich und nachdrücklich, als er’s immer vermöchte, für das Studium der Wechselwirkung Leibes und der Seele von Seiten der Seelsorger zu erklären, und gewissermaßen beispielsweise einiges aus dem mehrerwähnten Dictat an einem schicklichen Orte der zweiten Auflage einzuverleiben\footnote{S. § 63.}. Dadurch
\quellseite{258}
soll der Meinung die Thür versperrt werden, als hätte der Verfasser des „Evangelischen Geistlichen“ seine Hochschätzung der seelsorgerischen Bedeutung des gegenseitigen Einflusses von Leib und Seele fallen lassen. Im Gegenteil, die Erfahrung hat ihm diese Wechselwirkung nur desto größer und wichtiger gemacht und die ihm nachgewiesenen Mängel des Veringischen medicinischen Standpunkts haben auch das nicht vermocht, seine Freude an dem Veringischen Buche zu vermindern, dessen Studium er im Gegenteil jungen Amtsbrüdern wie früherhin empfehlen und sie ermuntern möchte, die von Vering angegebenen psychischen Mittel mit den in dem erwähnten Dictate aus Gottes Wort vorgelegten zu vergleichen.

Der ärztliche Freund, welcher wegen des Neudrucks des Anhangs zum II. Theil dieses Buches consultiert wurde, empfahl dem Verfasser ein neueres Buch: „Die psychischen Zustände, ihre organische Vermittelung und ihre Wirkung in Erzeugung körperlicher Krankheiten“ von Doctor Ottomar Domrich, Professor der Medicin, bei Friedrich Mauke in Jena, 1849. Das höchst interessante Buch, welches gewiß auch werth ist, von Seelsorgern, welche die Wichtigkeit des Themas erkennen, gelesen zu werden, trägt ganz den Stempel jetziger Wissenschaftlichkeit und liest sich daher auch nicht, wie das einfach nüchterne Buch des älteren Arztes Vering, aber es behandelt doch in neuerer Weise ganz eingehend dieselbe Sache. Man könnte sagen, Domrich’s Schrift handle ja von dem Einfluß der psychischen Zustände auf den Leib, während doch dem Seelsorger hauptsächlich an der Erkenntniß des Einflusses
\quellseite{259}
leiblicher Zustände auf die Seele liegen müsse, damit er die psychischen Zustände von den pastoralen scheide. Allein wenn ich einen Stab habe, kann ich ihn an zweien verschiedenen Enden fassen, der Stab bleibt doch derselbe. Und gerade so ist es mit dem Thema von der Wechselwirkung Leibes und der Seele, und man kann Domrich’s Buch im allgemeinen leicht in das verwandte Thema vom Einfluß des Leibes auf die Seele übersetzen. Wie es aber dem Manne der Wissenschaft so leicht und oft geschieht, daß er statt der getreuen Beobachtung Schlüsse aus Grundsätzen setzt, so mag es auch dem trefflichen Gelehrten ergangen sein, so daß man eben sein Werk für den seelsorgerischen Zweck mit Verstand und Unterschied benutzen muß; von dem wissenschaftlichen Werthe desselben zu urtheilen, ist der Verfasser dieses Buches, nur ein Mann der Praxis, weniger als andre competent.

O möchte doch das heranreifende Geschlecht von Seelsorgern erkennen, wie wichtig für die Praxis die Erkenntniß der Wechselwirkung Leibes und der Seele sei, und Gott seiner Kirche in baldem etliche gewaltige Seelsorger geben, die auch durch medicinische Studien befähigt seien, ihren Berufsgenossen durch öffentliche Unterweisungen zu dienen!

Welcher einfache Mensch wird den Anspruch der heutigen Aerzte, auf dem Gebiete des psychischen Lebens allein das Wort zu führen, ich will nicht sagen für bescheiden, sondern nur für gerecht erkennen; der pure Rationalismus der heutigen Psychiatrie kann sich so wenig als jeder andre Rationalismus, sei es auch ein wissenschaftlich geformter, Leuten empfehlen,
\quellseite{260}
deren Erkenntnißquelle auch für die psychischen Zustände des Menschen das göttliche Wort ist, das allewege gewissere Tritte thun lehrt, als die Ergebnisse einer Wissenschaft, der Psychiatrie, die noch so sehr mit Anfangszuständen ringt. Aber umgekehrt, wie unleidlich ist auch die Arroganz oft ganz junger Seelsorger, die ihre eignen pastoralen Zustände nicht kennen, geschweige die andrer Leute, und nun gar von psychischen Zuständen und psychischen Krankheiten so reden, als wüßten sie etwas! Woher sollen sie denn etwas wissen? Die Vorbildung unserer jungen Geistlichen ist doch gar nicht im Stande, auf diesem Felde etwas zu leisten; so kommen dann die jungen Männer in’s Amt hinein und reden, wenn sie von psychischen Zuständen reden, wie der Blinde von der Farbe und der Taube vom Schall. Wovon man nichts versteht, davon muß man schweigen, so lange nemlich, bis man etwas gelernt und erfahren hat. Zum Lernen und Erfahren aber sollten sich auf dem Gebiete, von welchem wir reden, alle jungen Seelsorger ermuntern lassen. Niemand, auch kein Arzt, hat so viel Beruf, Psychiatrie zu treiben, als die Seelsorger, welche an dem göttlichen Worte den reichsten Schatz der besten psychischen Mittel besitzen, die weit über den ungewissen rationalistischen Rathschlägen unchristlicher, wohl gar materialistisch gesinnter Aerzte stehen. Aber Studium bedarf es eben, Erfahrung braucht man, und damit man auf dem noch unbebauten Acker christlicher und kirchlicher Psychiatrie etwas leiste, muß man vor allen Dingen sein nichts erkennen und sich ernsten Fleißes und Vorwärtsstrebens nicht schämen.

\quellseite{261}
\section[48. Ärztliche Kenntnisse des Seelsorgers]{48. Ärztliche Kenntnisse des Seelsorgers überhaupt}

In alten Lehrbüchern der Pastoraltheologie, ja sogar in Kirchenordnungen wird den Seelsorgern nicht bloß medicinisches Studium empfohlen, sondern sogar ein gewisses Maß von medicinischer Praxis gestattet. Früher, wo es der Aerzte noch nicht so viele gab als jetzt, war namentlich der Landmann in Krankheiten sehr hülfslos und man mußte froh sein, wenn er seinen Pfarrer um Rath fragte statt aller möglichen Quacksalber und Betrüger. Daher sagt z. B. Andreas Hartmann, nicht der auch jetzt noch gerühmte Pastorallehrer, sondern der Verfasser einer kleinen aber trefflichen Pastoral-Anweisung für Dorfpfarrer\footnote{Unvorgreiflich-einfältig und wohl gemeinter Entwurf, wie ein Dorfpfarrer seiner anvertrauten Gemeinde erbaulich vorstehen möge Christo Jesu, dem Erzhirten unsrer Seele zu Lob und Preiß und redlichen Unterhirten zu Nutz und Dienst. In der Furcht des HErrn ehemals verfaßt, anjetzo aber an’s Licht gegeben von M. Andr. Hartmann, Pfarrer in Truchtelfingen. Ulm 1710.}:

\begin{quote}
„Zu dem Ende (nemlich um aus den Leiden die Schuld zu entwickeln) mag einem Minister einige medicinische Wissenschaft und ein öfterer Umgang mit Medicis und Chirurgis, so er’s haben kann, nicht wenig dienen. Denn ich halte dafür, daß ein Minister einen trefflichen Vortheil gewinne, einen Kranken zur Erkenntniß seiner Sünde zu bringen, wenn er die causas morborum ex principiis tum naturalibus, tum supernaturalibus wohl versteht.“
\end{quote}

Ebenso sagt die Cynosura würtembergica von 1658 von den medicinischen Studien der Pfarrer: „Sie mögen neben Verrichtung ihres Berufes auch
\quellseite{262}
medicinam lesen, und ihnen, auch guten Freunden und Nachbarn zu nutz gebrauchen, aber ohne Erlaubniß der geschworenen Landärzte nicht öffentlich profitieren, noch der Practik nachziehen.“ Trotzdem daß sich unsere Zustände so sehr geändert haben und man an allen Orten, wenn man will, leicht ärztlichen Rath haben kann, wird doch ein jeder, der den Wirkungskreis der Landpfarrer einigermaßen kennt, auch jetzt noch dem guten Andreas Hartmann und der württembergischen Cynosura Beifall geben, und zwar ganz aus den alten Gründen. Niemand wäre geeigneter, wir wagen viel, indem wir das sagen, für die ärztliche Praxis auf dem Lande, als der Pfarrer, vorausgesetzt nemlich, daß er für dieselbe vorgebildet und zugleich geprüfter Arzt wäre. Er wäre der wohlfeilste, der parateste Arzt seiner Gemeinde, und obendrein der treueste, da ihm sein göttlicher Hauptberuf die heiligsten Gründe für die Berathung seiner armen Leute gäbe. Allein davon kann ja unter unsern Verhältnissen auch nicht einmal die Rede sein: die Berufe sind geschieden, und da und so lange es nun einmal so ist, wird einem Seelsorger, auch wenn er Medicin studiert hätte und zur Praxis ganz befähigt wäre, doch nicht besser gerathen werden können, als sich alles Medicinirens und medicinischen Berathens ganz und gar zu enthalten. Und zwar gilt das nicht bloß vom Medicinieren auf dem Gebiete des Allopathen, sondern auch auf dem des Homöopathen. Es ziemt einem Pfarrer, in keinem andern Beruf weder geschickt einzugreifen, noch zu pfuschen, und mit den Aerzten vollkommen unverdorben zu bleiben. Wer den Beruf hat, der hat die Verantwortung, und dem gehört das Wort, und was deines
\quellseite{263}
Amts nicht ist, da laß deinen Fürwitz. So lange der Weg nicht gefunden ist, den Seelsorger und Arzt in einer Person unter öffentlicher Anerkennung zu vereinigen, soll man die Trennung der beiden Berufe desto schärfer einhalten, Liebe und Freundschaft der Ärzte auf dem Wege des Verzichts auf alles medicinische Reden und Wirken suchen.

\section{49. Die seelsorgerischen Mittel}

Alles kann auf den Menschen wirken, gut oder böse, recht oder verkehrt. Wer nicht zieht, verzieht — ist ein Erfahrungssatz der Pädagogen. Ebenso ist es mit der Seelsorge. Wer den Beruf der Seelsorge hat, von dem wird eine Wirkung ausgehen, gleich viel, ob sie nun die rechte sei, oder eine falsche. Es kommt dabei vor allem auf die rechten Mittel an. Die rechten seelsorgerischen Mittel sind aber schon genannt und abgehandelt: Predigt, Katechese, Liturgie, kurz Gottes Wort und die heiligen Sacramente. Es ist allerdings eine zu enge Begrenzung der Seelsorge, wenn man bloß bei dem öffentlichen und gewöhnlichen Gebrauch dieser Mittel stehen bleibt und fertig ist, so wie man gepredigt, katechisiert, liturgisiert und das Sacrament verwaltet hat. Es muß zum allgemeinen und öffentlichen Gebrauch des Wortes auch der besondere und außerordentliche kommen, den man eben unter den Namen Privatseelsorge zusammenfaßt. Andererseits ist es aber auch große Thorheit, wenn man das Außerordentliche zum Ordentlichen machen will, wenn man verkennt, daß Predigt, Katechese und Liturgie, Gottes Wort und Sacrament das Beste in der
\quellseite{264}
Seelsorge thun. Die Privatseelsorge ist etwas Außerordentliches und gehört mit ihrem ganzen Segen erst dem, an welchem die großen Mittel der allgemeinen Seelsorge ihre Wirkung gethan haben. Für unbekehrte Leute gibt es keine andere Seelsorge, als Predigt und Katechese, das empfindet man so oft an Krankenbetten der Gottlosen. Es spürt sich, man müßte eigentlich predigen, wenn gleich nur Einem, katechisieren und lehren, wenn gleich nur einen einzigen Alten. Da ist kein pastoraler Zustand, als der des Unglaubens, kein pastorales Mittel, als das allgemeine Wort Gottes, und sowie man mehr und anderes thun und sagen will, als was man überhaupt dem Gottlosen zu thun und zu sagen hat, fühlt man sich arm und elend und gedrückt. Die großen allgemeinen Mittel der Seelsorge müssen gewirkt haben, bevor man auf besondere Seelenzustände eingehen soll.

Sehr häufig setzt man die Privatseelsorge in die Kraft menschlicher Beredung. Ein junger Geistlicher, der mit seinen Gründen nicht auslangte, um seine Bauern von groben Sünden abzubringen, kam auf den Gedanken, es fehle ihm die populäre, gemeinfaßliche, seelsorgerische Sprache, und nahm daher zu seinen seelsorgerischen Besuchen den Schullehrer mit, welcher die Sprache des Volkes trefflich kannte und konnte, und siehe, mit dessen Hülfe brachte er die Leute allerdings dahin, daß sie thaten, was er wollte. Allein er täuschte sich doch. Die Leute waren mit menschlichen Gründen des zeitlichen Vortheils, der Ehre, der Ruhe und des Gemachs etc. überredet, aber bekehrt und gebessert waren sie nicht. Der alte Mensch wurde durch den Schullehrer vom Wege grober
\quellseite{265}
Sünden abgelenkt, aber ein neuer Mensch war nicht vorhanden, geschweige vorwärts gebracht im Guten. Eben weil sie blind und todt waren, weil sie noch zu keinem geistlichen Leben und Verständnis gekommen waren, verstanden sie den jungen Geistlichen nicht, auch wenn er klar und einfach redete, und konnten ihn auch nicht verstehen. Es fehlte ihnen der Segen des gepredigten Wortes, darum gab es für sie keine Seelsorge. An diesem Beispiel zeigt es sich klar, was die rechten Mittel der Seelsorge seien, was die falschen. Auch Demosthenes, auch Cicero haben Großes auf den Haufen und im Haufen gewirkt, aber die Leute blieben, wer sie waren. Auch Heiden konnten also wirken, aber nicht sittlich, nicht für die Ewigkeit. Seelsorge kann man ihre Wirkung nicht nennen. Menschengründe und Menschenweisheit helfen nicht, nicht bei der Predigt, nicht bei der Privatseelsorge. Gottes Wort ist das allein seelsorgerische Mittel, alleine Gottes Wort dringt durch. — Es versteht sich, daß es deshalb keinem Seelsorger zum Verbrechen gemacht werden soll, auch menschliche Gründe unter dem göttlichen Worte auf seine Pfarrkinder wirken zu lassen; er kann auch gar nicht anders. Aber er muß wissen, daß ihm die Seelsorge nur dann wahrhaft gelingen kann, wenn Gottes Wort Anklang im Menschenherzen fand, und nur dann gelungen ist, wenn es seine ihm eigene göttliche παιδεία auf die Menschenseele ausübt. Es muß darum Grundsatz des Seelsorgers sein: „Wer da redet, daß er es rede als Gottes Wort. Loquatur eloquia Dei.“

\quellseite{266}
\section{50. Methodismus der Seelsorge}

Es kann niemand andere Mittel der Wirksamkeit für einen Seelsorger ausfindig machen, als die gegebenen und oben erwähnten. Aber man kann am Ende auf den Gedanken kommen, die Mittel bringen nicht durch, wenn man sie nicht auf eine besondere Weise anwendet. Das heißt, man kann zur Einseitigkeit seine Zuflucht nehmen, um im Ganzen und Allgemeinen etwas Erkleckliches auszurichten. Ein Verfahren, wie wenn man ein Messer mit der Spitze wirken lassen wollte, weil die Schneide nicht hilft. — Spener sah die todte Art, Predigt zu hören, welche allenthalben eingerissen war. Da gebrauchte er besondere Versammlungen, die Predigt zu repetiren\footnote{Spener’s Predigt am 6. Sonntag nach Trinitatis 1669 hatte einen Einfluß, wie vielleicht seit der ersten Predigt Petri in Jerusalem keiner. Sie leitete das große Werk ein, das Spenern von Gott vertraut war.}, und siehe — es half — eine Weile. Die menschliche Theilnahme wurde durch die besondere Maßregel für das Wort Gottes besonders angeregt. Da fand das Wort Gottes mehr Aufmerksamkeit, als bei der öffentlichen Predigt, drang deshalb ein und wirkte ungewöhnlich. Das Mittel verlor aber seine Wirksamkeit, als es aufhörte, außerordentlich zu sein. Ja es schlossen sich an den methodistischen Gebrauch desselben eigenthümliche Uebel an. Ebenso kann einmal eine methodistische Versammlung auf den Menschen, der so etwas nicht gewohnt ist, einen Eindruck machen, daß er dem Worte weniger Widerstand leistet, und es ihm
\quellseite{267}
deshalb mehr hilft. Es kann einmal kommen, daß ein Camp-meeting sogar vielen Leuten nützt, ja es kann auf dem Wege eines solchen einseitigen Methodismus wohl auch eine ziemliche Zeit lang gewirkt werden, — menschlich nämlich, und unter dem Tumult und Aufruhr menschlicher Kräfte kann wohl auch Christus, der Wunderbare, eine Zeitlang Fische fangen. Allein wenn dann später die Sache nicht mehr unbefangen geschieht, wenn sie betrieben wird, wenn das Außerordentliche zum Ordentlichen werden muß, dann hört der Segen auf. Es ist bei all diesen Dingen ein dem Herrn bekanntes, inneres, sittliches Maß, das wir zwar in unserer Thorheit tausendmal überschreiten, dessen Ueberschreitung Er aber, so zu sagen, nicht mehr vergibt, wenn man nach seinem Sinn zur Einsicht gekommen sein sollte, daß es genug ist.

Gerade so ist es mit derjenigen Anwendung des göttlichen Wortes, welche wir Privatseelsorge nennen. Liest man eine alte lutherische Pastoraltheologie, so kann das Rubrum der Privatseelsorge gänzlich mangeln. Karg scheinen überhaupt die Rubren, und gering die Pflichten eines Predigers und Seelsorgers nach diesen Autoren. Liest man dagegen einen Gottfried Arnold, einen Baxter: Ja, was für Pflichten, welche Verantwortung, welch eine unerträgliche Noth, daß man tausendmal mit Chrysostomus rufen möchte: Mirum, si sacerdos salvetur! Wie aber? Kennen Arnold, Baxter \&c. neue Mittel der Seelsorge? Haben sich ihnen vormals verborgene Schätze geöffnet? Nein, sondern was die Alten haben, haben sie, nur auf besondere Weise, nur in besonderer specialisirender Anwendung.

\quellseite{268}
Was sich von selbst versteht, wird hervorgehoben; was zuweilen vorkommt, wird über das Maß betont. Während die alten Pastorallehrer in ihrer Einfalt die geistlichen Waffen zu sicherer Führung in die Hand des Anfängers niederzulegen wissen und Muth dazu machen, sie zu brauchen, und dem guten Geiste zutrauen, daß Er Kraft, Ausdauer und allmählich Erfahrung gebe: ist bei jenen eine in’s einzelne gehende Darstellung der Arbeit und Verantwortung eines Seelsorgers, vor der man erschrecken kann. Diese gesetzliche Weise demüthigt tief, nimmt aber Muth und Freudigkeit. Der Methodismus, auf den in bester Absicht gedrungen wird, hat die umgekehrte Wirkung. Er lähmt, weil die Leistungen unmöglich scheinen.

Die einfache Regel ist: Gebrauche die alten Mittel in alter Weise und bleibe im Lehren, Lernen und Erfahren, in Anfechtung und Gebet, auf daß du zum Seelsorger reifest. Du wirst öffentlich und sonderlich, vielleicht in hundert und tausend Weisen deinen Pfarrkindern nahe kommen können, aber übertreib es auf keine Weise, mit keinem Mittel, mit keiner Gabe. Thue in Einfalt das Deine. Brauche betend die uralten Mittel auf jede Weise, die sich indicirt, und laß Gott sorgen, wie es gerathen werde. Du kannst Hausbesuche machen, und wieder keine, gerufen und zuweilen ungerufen zu den Kranken gehen, dies und das thun, und nicht thun, wie du es nach ruhiger Ueberlegung aller Umstände vor Gott für das Beste hältst. Aber mach dir nichts zur Fessel und Gewissenslast. Meide jeden Methodismus.

\quellseite{269}
\section{51. Vom Theilen des göttlichen Wortes}

Schon der Prediger hat wohl Acht zu geben, daß er das Wort Gottes richtig theile. Das Wort ist Gesetz und Evangelium. Es ist nütze zur Lehre, zur Strafe, zur Besserung und zur Züchtigung. Da gilt es nun aufmerksam, daß man nicht Gesetz predige, wo man Evangelium predigen sollte, nicht lehre, wo man strafen, nicht strafe, wo man lehren sollte. Mehr noch gilt es aber, in der Privatseelsorge Gottes Wort zu theilen. Wie ein Arzt zuerst den Kranken und seine Krankheit kennen lernt, ehe er Arznei verordnet, und genau nach der Krankheit die Arznei einrichtet, so erforscht ein Seelsorger den Zustand des Menschen und seine Stufe in der Ordnung des Heils, und genau nach dem Befund des pastoralen Zustandes reicht er die göttliche Arznei. Denn als Privatseelsorger ist der Pfarrer ganz Arzt. Es handelt sich hier nicht um Speiseaustheilen und Kleideranziehen, sondern um Arznei, um die geistliche Diät des Kranken u. dgl. Es muß daher ein Pfarrer die verschiedenen Stufen des christlichen Lebens und seine Krankheiten und Schwachheiten kennen lernen, und für alles die Arznei aus Gottes Wort. Wie ein Apotheker in viele Gläser und Büchsen vertheilt hat, was er braucht, und mit gewandter Hand bald dieses, bald jenes hervorlangt, so der Seelsorger, der Seelenarzt. Hiezu nützt das Studium der Bücher solcher Autoren, welche pastorale Weisheit besaßen, aber es reicht nicht aus, und ein guter Seelsorger wird man nur durch Erfahrung. Das vergehe
\quellseite{270}
man nicht. Insbesondere aber vergesse man ja nicht, daß die Privatseelsorge Seelenkur, der Privatseelsorger der Seelenarzt ist, daß man also die Privatseelsorge ebensowenig für etwas Ständiges und bei jedem Pfarrkind Fortdauerndes halten darf, als man die Arznei mit der Speise, den Arzt mit dem Speisemeister, und den Apotheker mit dem Koch vermengen darf.

Hat man den Zustand seines Kranken klar erkannt, so suche man ihn selbst zur Erkenntniß desselben zu führen; das ist bei der Seelenkur nöthig, nicht ebenso bei der Leibeskur. Erkennt der Kranke sein Uebel, so zeige man ihm die rechte Arznei, denjenigen Theil des göttlichen Wortes, welcher seinen Mangel erstattet und seine Krankheit heilt, und helfe ihm belehrend, züchtigend, ermunternd, erziehend zur Anwendung und Ergreifung des Mittels, zur Geduld und Ausdauer bei der Anwendung, und bleibe mit ihm und für ihn im Gebet, bis Hülfe geworden ist.

Immerhin wird man aber mit klaren Worten und deutlichen Gedanken fahren müssen, denn der Mensch ist so beschaffen, daß zwar Anfangs, vor dem Fall das geschenkte Leben auch sein Licht bei sich hatte und mit sich brachte, nun aber alles Leben durch das Licht und die Wärme gegeben wird, die vom Lichte ausstrahlt. Jeder Gedanke aus Gottes Wort hat seinen Segen, und übt ihn um so mächtiger, je klarer er gegeben wird. Nur auf Grund bestimmter Sätze und Wahrheiten kann Strafe, Trost, Ermunterung u. s. w. gegeben werden, und es lasse sich deswegen niemand das Auge durch ein entgegengesetztes Verfahren blenden.

\quellseite{271}
Sollten der verordneten Arznei Hindernisse im Wege stehen, so muß auch hier wieder zu deutlicher Erkenntniß derselben geführt, und durch Kraft des göttlichen Lichtes das weggethan werden, was die Genesung der Seele hindert. Auch hier vertreibt Licht die Nacht, und aus dem Lichte strahlt die Wärme, welche das Eis der Seele schmilzt. Jeder Fortschritt im Guten, jeder Rückschritt vom Bösen ist durch Licht, durch Wahrheit, durch Gedanken, durch Worte Gottes bedingt.

\section{52. Casuistik}

Manche haben in der neuern Zeit von Casuistik nichts wissen wollen, weil von dem allgemeinen Wissen und Erkennen der Wahrheit der Schritt zum Besonderen leicht sei, und die Salbung allerlei lehre. Allein ist es denn nicht am Tage, daß einerseits eben in der Casuistik sich die Salbung bewährt, die allerlei lehrt, andererseits aber trotz der Salbung, die wir haben, gar mancher schwache, auch gar mancher sich reich und groß dünkende Geist im vorkommenden Fall sich doch nicht zu rathen und zu helfen weiß? Mögen Jesuiten die Casuistik in Verruf gebracht haben, sie konnten eigentlich doch nur ihre Casuistik in Verruf bringen, während die Casuistik selbst bei allen denen im Werthe bleiben muß, die da wissen, wie schwach wir sind, und wie nöthig uns deshalb Rath und Weisung derer ist, welche durch Gottes Salbung klar und heiter sahen. Weit entfernt also, die Casuistik für überflüssig zu erkennen, ist sie es gerade, welche dem jungen Geistlichen am meisten helfen, die er auch lernen kann, und
\quellseite{272}
die ihn bei rechtem Studium bescheiden und demütig macht, eben weil der Fälle ein unendliches Meer ist, und oft zwei Fälle, die einander ähnlich scheinen, doch verschieden ausgedeutet werden müssen. In früheren Zeiten der lutherischen Kirche hat auch niemand von der Casuistik so geringschätzig gedacht, wie heut zu Tage viele.

Dr. Johann Franz Buddeus in seiner Isagoge historico-theologica ad theologiam universam singulasque ejus partes, novis supplementis aucta, Leipzig 1730, pag. 638 beklagt sich über die ungerechten Vorwürfe der Römischen, daß die Protestanten nicht wüßten, was erlaubt sei oder nicht; Didacus de Bezal sage: In England, Holland und Deutschland gebe es die reichsten Kaufleute lutherischer Confession, die nicht bloß keinen Rath wüßten, sondern auch keinen Menschen hätten, den sie zu Rath ziehen könnten, wenn es sich darum handle in Sachen öffentlicher Verträge oder der Wiedererstattung den rechten Weg zu betreten; überhaupt sei das Studium der Moral ganz und gar vernachlässigt. Solche Vorwürfe beruhten aber entweder auf Ignoranz, oder sie kämen aus der größten Unverschämtheit. Johann Joachim Zentgraf diene aber in der Vorrede zu seiner Summa juris divini nach Verdienst. Buddeus gibt zu, daß man bei uns so gar dickleibige Bände über Moral und Casuistik nicht habe wie bei den Römischen, deshalb aber fehle es keineswegs an einschlägigen Schriften; man könne im Gegentheil eine gar nicht unbedeutende Anzahl anführen. Er theilt nun die hieher gehörigen Schriften in eigentlich casuistische, in solche, welche den Decalog auslegten, und endlich
\quellseite{273}
lich in theologische Bedenken, und führt alsdann zwar nicht die ganze hieher gehörige Litteratur, aber doch die hauptsächlichsten Schriften bis auf seine Zeit herab an. Da man nun voraussetzen muß, daß manchem jungen Geistlichen damit ein Dienst geschehe, so erlauben wir uns die hervorragendsten Werke, wie sie Buddeus und andere anführen, zu verzeichnen. Da diese Bücher in der Regel nicht gesucht werden, so stehen sie auch nicht hoch im Geldwerth, und wer da will, kann sich dieselben ohne bedeutende Kosten anschaffen, wenn er nur beim Durchgehen antiquarischer Cataloge ein Auge darauf hat. Wir geben sofort das Verzeichnis:

1. Der berühmte Wittenberger Theologe Friedrich Balduin las von 1622 bis 1626 über Gewissensfälle, und nach seinem Tode erschienen im Jahre 1628, und hernach öfter, von ihm 4 Bücher de casibus conscientiae. Das Werk hat die letzte Feile von ihm selbst nicht erfahren, so daß es den übrigen Schriften Balduin’s nicht würdig zur Seite steht; die Wittenberger Theologen aber erkannten die Nothwendigkeit eines solchen Werkes an.

2. Auf Balduin folgte der Gießener, später Coburger Theologe Caspar Finck. In seinem Todesjahre 1631 ließ er erscheinen: Sylloge quaestionum illustrium theologicarum, theoreticarum, maximam vero partem practicarum et quotidianarum ex assidua bibliorum, orthodoxorum patrum, sinceriorum theologorum, jure consultorum etc. lectione, experientia propria etc. rotunde decisarum. Das Werk hat drei Sectionen, deren letzte 1632 nach dem Tode des Autors an’s Licht trat.

\quellseite{274}
3. Im Jahr 1632 ließ der Reval'sche Pastor Ludwig Dunte zu Lübeck erscheinen: Decisiones mille et sex casuum conscientiae e diversis theologorum scriptis collectae, contractae et in brevitatem redactae, et juxta ordinem locorum theologicorum positae. Das Werk fand wegen seiner compendiarischen Natur Beifall und wurde öfters gedruckt.

4. Der Coburger Theologe Andreas Keßler schrieb in deutscher Sprache eine theologia casuum conscientiae, hodierno cumprimis tempori accomodatorum, welche sein Amtsnachfolger Johann Christoph Selde im Jahre 1658 zu Wittenberg drucken ließ. Das Buch ist gewissermaßen ein Supplement der vorigen und enthält eine Aehrenlese von Fällen, welche in anderen ähnlichen Schriften übergangen worden waren.

5. Georg König gab 1654 heraus: Casus conscientiae, qui in sex capitibus doctrinae catecheticae, una cum tabula oeconomica subinde solent occurrere, erudite et fideliter decisae. Sein Sohn Georg Matthäus König veranstaltete 1676 eine neue, aus dem Manuscript seines Vaters vervollständigte Ausgabe.

6. Von demselben Autor erschien nach seinem Tode 1665: Heptas casuum conscientiae miscellorum.

7. Arnold Mengering ließ mehrere Schriften aus dem Bereiche der Gewissens-theologie in deutscher Sprache erscheinen: „Erneuerter evangelischer Gewissenswecker 1645, Erneuerte evangelische Gewissensruhe 1647. Geistesblicke und Sündenregister nach dem Katechismus Luthers 1652. Evangelisches Gewissensrecht 1661.“

\quellseite{275}
8. Seine bisher genannten Vorgänger übertrifft Johann Conrad Dannhauer, dessen Collegium decalogicum zuerst 1639 erschien, dann aber im Jahre 1669 von Balthasar Bebel unter dem Titel Deuteronium Dannhauerianum sehr verbessert und vermehrt erschien. Seine Theologia casualis, die im Jahre 1706 von Johann Friedrich Mayer zu Greifswalde herausgegeben wurde, ist nichts als der erste Entwurf von Dannhauer’s Liber conscientiae apertus oder Theologiae conscientiariae tomi duo.

9. Schon vor Dannhauer ließ in Straßburg Johann Schmid sein Collegium casuum conscientiae erscheinen (1634 und 1635), und schrieb 1645 auch conciones conscientiariae.

10. Insbesondere zu empfehlen dürfte unseren jüngeren Geistlichen ein kurzes und dennoch sehr reichhaltiges Buch in Duodez sein. Der Verfasser ist Andreas Bruckner von Schweinfurt; der Titel: Manuale mille quaestionum illustrium theologicarum, praecipue practicarum, decem centuriis inclusarum et succincte quidem, nervose tamen ex praecipuis theologis decisarum. Es erschien 1679 zu Nürnberg mit Noten vom Jenenischen Theologen Philipp Müller. Wenn es nur gegenwärtig einen theoretisch und praktisch genugsam gebildeten Geistlichen gäbe, der einen Prucknerus redivivus nach den Bedürfnissen der gegenwärtigen Zeit schreiben könnte. Das wäre ein köstliches Geschenk für das zukünftige Geschlecht junger Geistlichen.

11. Dem Prucknerischen Buche gegenüber steht die zu Tübingen von 1660 bis 1662 in sechs Theilen erschienene Theologia casualis Johann Adam Osianders, in qua
\quellseite{276}
quaestiones, dubia et casus conscientiae circa credenda et agenda enucleantur. Derselbe Mann hat auch noch andere Schriften geschrieben, in welchen vorgelegt wird, was Juris Divini sei.

12. 1692 erschien zu Frankfurt und Leipzig in deutscher Sprache Samuel Schelwigs Cynosura conscientiae sive perspicua et scripturae sacrae congrua multorum, maximam partem singularium nec e trivio petitorum item quorundam nondum satis enucleatorum casuum conscientiae decisio. (Leitstern des Gewißens. Stettin. 1692.)

13. In Buddeus Aufzählung macht den Schluß der hierher gehörigen Bücher Friedemann Bechmanns Theologia conscientiaria, seu tractatus de casibus conscientiae. Jena. 1705.

14. Dr. Johann Nicolaus Mieslers Opus novum quaestionum practico-theologicarum sive casuum conscientiae. Frankfurt. 1676. Dieses Buch dürfte wohl am wenigsten zu übergehen sein, da es neben dem Prucknerischen Duodezbande, obwohl selbst ein Band von 668 Seiten in klein Folio, doch sehr bequem zu gebrauchen ist.

15. Zu besonderem Gebrauche in unserer Abendmahlsmengerischen Zeit dürfte wohl auch zu empfehlen sein August Pfeiffers Informatorium; conscientiae Eucharisticum, complectens triginta quaestiones in administratione sacrae coenae piis mystis adprime utiles. 1687. Derselbe Pfeiffer hat aber auch 1717 eine Gewißensschule erscheinen laßen.

\quellseite{277}
16. Buddeus selbst hat 1717 das bekannte Buch: Gottholds Manuale casuisticum bevorwortet.

— Auch die reformirte Kirche hat ihre Casuisten, unter denen die Werke von Johann Heinrich Altstedt, Wilhelm Amesius, Hall, Perkins, Taylor, Hoorenbek nicht mit den Erzeugnissen der lutherischen Kirche verwechselt werden dürfen.

Dabei ist jedoch zu merken, daß in allen, auch den lutherischen casuistischen Büchern die Antworten auf die gestellten Fragen nicht immer garantiert werden; es kann kommen, daß in einem gut confessionellen Buche eine Antwort eines Theologen von ganz anderer Confession zu finden ist, weil eben gerade die Antwort eines orthodoxen Casuisten nicht zu finden war. Solche Antworten dienen dann selbstverständlich zu nichts anderem als zur Schärfung des Nachdenkens, wie man denn überhaupt auch die gesamte Casuistik der lutherischen Kirche weniger rühmen als zum Studium empfehlen sollte. Es geht wie z. B. mit den Definitionen des großen Gelehrten Melanchthon, die man auch nicht mehr einfach wiederholen oder unterschreiben kann, während sie doch die Ausgänge des Nachdenkens und des Bemühens, Besseres herzustellen, sein und werden könnten. Die ganze Casuistik der lutherischen Kirche bedarf einer neuen unbefangenen und schriftgetreuen Bearbeitung, sie wäre es aber schon um des großen Segens willen, den sie stiften könnte, werth, noch einmal frisch in die Esse gestoßen, ja von den Theologen jeder folgenden Zeit immer neu zu ihrer Zeit Nutz und Frommen bearbeitet zu werden.

Warum Buddeus den eigentlichen Casuistiken insonderheit die Erklärungen des Decalogs einreiht, wie oben erwähnt
\quellseite{278}
ist, springt in die Augen: Ad jurisprudentiam divinam vel maxime spectant, qui Decalogum seu Summam juris Divini commentationibus suis illustrarunt. Auch sie wie die Casuistiken haben es mit der göttlichen Jurisprudenz zu thun. Dannhauer und Dorsch sind auf diesem Arbeitsfelde gefeierte Namen.

Ebenso ist es unverkennbar, weshalb Buddeus die Sammlungen theologischer Bedenken und Rathschläge hieher zieht; sie könnten ja ganz einfach in die casuistischen Bücher an ihrem Orte eingestellt sein. Hieher gehörig sind folgende bekannte Bücher:

1. Des Württemberger Theologen Felix Bidembach in Gemeinschaft mit seinem Bruder Johann Moriz im Jahr 1608 herausgegebenen Decades novem consiliorum theologicorum; die neunte gab Moriz nach dem Tode seines Bruders heraus und ließ 1614 zu Frankfurt noch eine zehnte folgen.

2. Abraham Calov, Johann Meißner, Johann Andreas Quenstedt und Johann Deutschmann ließen 1664 zu Frankfurt am Main die Consilia theologica Wittenbergensia erscheinen.

3. Der Hamburger Theolog Johann Dedeken unternahm es, einen Thesaurus consiliorum et decisionum in 3 Foliobänden heraus zu geben. Johann Ernst Gerhard von Jena bereicherte und vervollständigte das Werk, starb aber vor vollendetem Druck, so daß es Christian Grübel 1671, obendrein mit einem Anhang vermehrt, in Jena erscheinen ließ.

Wenn man nun zum Schluß noch den edlen-frommen Philipp Jacob Spener um seiner kirchenbekannten
\quellseite{279}
zahlreicher Bedenken wegen in die Zahl der casuistischen Theologen eingezeichnet hat, wird man wohl glauben dürfen, die Meinung vollständig niedergelegt zu haben, wie wenn die lutherische Kirche je und je die Casuistik gescheut und vielleicht gar gehaßt hätte, etwa deshalb weil abgeschmackte Leute auf dieses Gebiet der hellen Praxis oftmals ganz unpractische Fragen bloß der dialectischen Uebung wegen hinübergezogen haben. Die Seelsorge besteht freilich nicht allein in der Lösung von casuistischen Fragen, sie hat weit mehr zu thun, wirkt vornehmlich auf den Willen ein und regiert das Herz. Aber das bleibt doch auch wahr, daß man die Casuistik nicht entbehren kann. Ist sie aber nicht entbehrlich, so wird sie wohl auch nützlich und erneuerten Fleißes werth sein.

\section{53. Vertrauen}

Gottes Wort kann ich aus jedem Munde, Gottes Sacrament aus jeder Hand empfangen, wenn es mir unverfälscht ist, ich mag ein Vertrauen zu dem Pfarrer, dem Prediger oder Diener haben oder nicht. Aber zum Seelenarzt kann ich ebensowenig und noch weniger einen Seelsorger wählen, zu welchem ich kein Vertrauen habe, als ich mir einen leiblichen Arzt wähle, dem ich kein Vertrauen schenke. Man kann sogar sagen, daß für die Pflege geistlicher Kranken das Vertrauen weit nöthiger sei, als zur leiblichen Krankenkur. Wie könnte ich einem Unweisen, Unredlichen, Untreuen meine Seelennoth anvertrauen! Und ob einer weise, redlich und treu wäre, aber ich wüßte oder erkennete es nicht, hülfe es mir denn?

\quellseite{280}
Könnte ich glauben und vertrauen? Es ist das so einfach und natürlich, daß man es nur nicht begreifen kann, wie ein Seelsorger auf Vertrauen Anspruch machen kann und mag, wenn er sieht, er besitze keines. Wo ist Bescheidenheit leichter, schon natürlich leichter, schon aus Vernunftgründen leichter, als bei Vertrauenssachen? Aber freilich, auch Seelsorger sind Menschen, es kann ihnen Menschliches begegnen. Sie sollten nur mit Demuth ihre Schwachheit und Sünde erkennen, und sich selbst heilen lassen, statt um Vertrauen zu buhlen, und es sich durch Mittel zu erwerben, die bei den besseren Pfarrkindern Mißtrauen mehr erwecken müssen, und auch erwecken, als sie Vertrauen schaffen können.

Hat ein Seelsorger Vertrauen, so wird es durch Ernst und heilige Strenge nur gesteigert. Man wünscht vom Seelsorger, daß er es genau nehme, und es tröstet selbst kein Trost aus dem Munde eines Mannes, der alles leicht nimmt. Aber eben, weil das Vertrauen sogar mit mit dem versöhnen kann, was eigentlich zurückschrecken könnte, erwächst dem, der es besitzt, eine so hohe Verantwortung. Vermöge des Vertrauens kannst du recht führen und irre führen! Ein jeder muß Rechenschaft geben für die, welche seinem seelsorgerischen Rath vertrauen.

Da helfe uns der barmherzige Gott durch Jesum Christum. Amen.

\section{54. Beichten}

Vor nunmehr gerade dreißig Jahren schrieb der Verfasser dieses Buches eine kleine Schrift unter dem Titel:
\quellseite{281}
„Einfältiger Beichtunterricht für Christen evangelisch-lutherischen Bekenntnisses“, Nürnberg bei J. Ph. Ram · 1836, in der Absicht, in seinen Lebenskreisen, wenn es möglich wäre, eine neue Liebe und Lust zur Beichte, insonderheit zur Privatbeichte, zu wecken. Die Schrift ist indessen in Vergessenheit geraten, damals aber that sie ihren Dienst in dem Maße, daß, als er ein Jahr darauf Pfarrer in Neuenbettelsau wurde, die Leute nicht anders dachten, als sie würden nun alle wieder privatim beichten müssen. Der erste Mensch, der ihn auf der Gasse ansprach, war ein alter Schwachsinniger, und die ersten Worte, die er sprach: „Ich kann meine Beichte.“ Der neue Pfarrer war jedoch gar nicht gewillt, mit Sturmeseile die Privatbeichte wieder aufzurichten, und es dauerte nicht weniger als sechs Jahre, während welcher er allerdings jede Gelegenheit benützte, die Privatbeichte und ihre Segnungen anzupreisen. Im Frühling des Jahres 1843 gab ihm eine Stelle eines kirchenregimentlichen Rescripts, welches er abzukündigen hatte, und durch welche der Gebrauch der Privatbeichte da, wo man sie wünschte, freigegeben wurde, den Anlaß, diejenigen, welche Lust hätten, privatim zu beichten, für den nächsten Sonnabend einzuladen. Er vermuthete nun, es würden die wenigen einzelnen Leute, die sich schon früherhin die Privatabsolution ausgebeten hatten, zur Privatbeichte erscheinen. Das ging aber ganz anders. Die Kirche war am nächsten Sonnabend von lauter Privatbeichtenden angefüllt, und der Pfarrer hatte gleich zum ersten Male ganz unvermuthet eine übermäßige Beichtarbeit. Der erste, welcher in die Sacristei trat, um zu beichten, war der damals älteste Mann der
\quellseite{282}
Gemeinde, der zweite gleichfalls einer der ältesten, Kirchenpfleger eines Filials, ein Mann, mit welchem nicht zu spaßen war, einer der schwierigsten Charactere in der Gemeinde. Beim Anblick dieser Männer kam dem Pfarrer die Vermuthung, es möchte seine Abkündigung so verstanden worden sein, als sollte die Privatbeichte zwangsweise eingeführt werden, weshalb er sich nun auch gegen diese Auffassung wehrte. Darauf antwortete der ersteingetretene: er habe sich schon lange nach der Privatbeichte gesehnt, denn er habe ja in seiner Jugend schon privatim gebeichtet; der zweiteingetretene, ein langer und starker Mann, holte mit seiner Rechten aus und schlug auf seine Brust, und sagte: „Ich hab Sie schon recht verstanden, aber ich will, ich will.“ Das klang wie ein recht tief aus der Brust ertönendes mea culpa, mea culpa.

Von jenem Tage an, also ganz mit einem Male war die Privatbeichte die herrschende, so daß Fälle vorkamen, daß ungefähr 100 privatim beichteten und einer die allgemeine Beichte verlangte, welche dann selbstverständlich auch nur wieder eine Art von Privatbeichte wurde. Es dauerte lange, bis der widermärtige Theil der Gemeinde, der an Zahl nicht groß war, sich von dem allgemeinen Eindruck dieser Wendung der ganzen Gemeinde zur Privatbeichte so weit erholte, daß sich die ihm Zugehörigen verabredeten, immer an denselbigen Tagen miteinander zur allgemeinen Beichte zu gehen. Da gingen dann diese Widermärtigen zusammen zur allgemeinen Beichte und versicherten je länger je lauter, wieviel schöner und ergreifender schon durch die Ansprache die allgemeine Beichte sei als die Privatbeichte. So kam nun aber auf die allgemeine
\quellseite{283}
Beichte eine Art von Schmach, als wäre sie eben die Versammlung der Widerwärtigen und der pastoralen Führung der Gemeinde Widerstrebenden, — und das durfte doch auch nicht geduldet werden.

Zum Anfang des Jahres 1848 wurde das Filial R. nach Neuendettelsau eingepfarrt, womit den neuen Parochianen, die früher eine weite Strecke an ihren Pfarrort zu gehen hatten, eine große Wohlthat erzeigt wurde; denn nach Neuendettelsau hatten sie kaum eine halbe Stunde auf ganz ebenem Weg zu gehen. Als nun das erste Mal in diesem Filial Beichte gehalten wurde, siehe, da kam nun auch gleich die ganze Gemeinde, Mann für Mann, zur Privatbeichte; auch wer keine Beichte wußte und unfähig war, aus dem Herzen zu beichten, stellte sich doch zur Privatbeichte ein und ließ sich von dem Pfarrer, so gut es ging, helfen. Diese allgemeine Betheiligung an der Privatbeichte sollte eine Art von Dank für den Pfarrer von Dettelsau sein, der zur Einpfarrung treulich geholfen hatte; denn daß ihm mit dem allgemeinen Anschluß an die Privatbeichte etwas recht Wohlgefälliges geschehen würde, davon hatten sie alle die Ueberzeugung.

Was war nun aber erreicht? Die Privatbeichte war durchgeführt und es schien, daß sich auch die Widerstrebenden durch irgend eine Bemühung geschickter Art zu dem allgemeinen Wege würden führen lassen. Allein das alles war gar nicht beabsichtigt, und der Pfarrer hatte kein Wohlgefallen an dieser ausschließenden Herrschaft der Privatbeichte, weil sein pastorales Ziel ein ganz anderes war. Er wollte eine reichlichere Betheiligung der Gemeinde am Sacrament des Altares
\quellseite{284}
herzuführen und that, was er konnte, um den Schlendrian des jährlich zwei- oder viermaligen Abendmahlgehens zu zerstören. Schon lange schwebte ihm die sonntägliche Communion und die Betheiligung an derselben nicht nach Gewohnheit und traditionell festgehaltenen Abendmahlstagen, sondern je nach Bedürfniß als ein pastorales Ziel der Gemeindeführung vor Augen. Wenn nun bei der Mehrung der Abendmahlsgänge und Tage die jedesmalige Privatbeichte aller einzelnen Communicanten festgehalten werden sollte, so wäre damit ein unüberwindliches Hinderniß für das eigentliche pastorale Ziel erwachsen.

Wenn so ein Landmann anstatt zwei Mal im Jahre zwanzig Mal oder auch nur zwölf Mal zum Sacramente gehen will, was soll er dann beichten? Es fehlt ihm ja freilich nicht an Sünden, aber wie an der nöthigen Bildung sich darüber auszusprechen, so auch an der Bildung, bei der Einförmigkeit seines Lebensganges den Wechsel seiner sündigen Zustände, geschweige seiner sündigen Thaten und Worte zu bemerken. Schwatzen und lügen soll er doch am allerwenigsten in der Beichte, — wo soll er also Stoff hernehmen? Da liegt dann die Gefahr, sich durch Formelbeichten aus der Noth zu helfen ganz nahe; der Schlendrian, den man aufheben wollte, bricht mit einer unwiderstehlichen Macht herein, und obendrein gibt es für den Pfarrer eine so unerträgliche Last, alle diese Formelbeichten anzuhören und überdies zu bescheiden, daß nichts anderes in sicherer Aussicht steht, als das Hinsterben des ganzen Instituts der Privatbeichte unter dem nächsten Nachfolger, wie ja früher schon in der
\quellseite{285}
protestantischen Kirche aus gleichen Ursachen die Privatbeichte hingestorben ist, und die allgemeine Beichte sich Pfarrern und Gemeinden als Rettungsanker aufdrängte. Es mußte dahin gewirkt werden, daß das Ziel einer reichlicheren Betheiligung am Sacramente möglich wurde. Die Privatbeichte ruht ja in der protestantischen Kirche auf ganz andern Ansichten und Gründen als bei den Römischen. Der Protestant ist nicht gehalten, nach einer Art von vollkommener Beichte zu ringen; das Beichtinstitut kann daher auch niemals in unserer Kirche zu dem allgemeinen und gewaltigen Erziehungsmittel werden, zu welchem die Ohrenbeichte im Mittelalter heranwuchs. Ist doch auch für die römische Kirche eine ganz andre Zeit gekommen! Unser Katechismus verlangt nicht, daß alle Sünden gebeichtet werden, die man weiß, sondern die, die man weiß und fühlt im Herzen, und auch die nicht um der Absolution willen, welche nicht bloß die dem Beichtvater unbekannten, sondern auch die vom bußfertigen Beichtkinde selbst nicht erkannten Sünden umfaßt. Nicht die Beichte, von so großem Segen sie ist und sein kann, sondern die Buße oder vielmehr Bußfertigkeit der Seele ist das von der protestantischen Kirche aufgestellte unerläßliche Erforderniß eines rechten Communicanten.

Es ist eine, auf die Augsburgische Confession gegründete kirchliche Ordnung, daß unverhört, ungebeichtet, unabsolviert niemand zum Sacrament gehen soll, und obwohl diese Ordnung weiter gar nichts als Menschenwerk ist, das also nach der Lehre des Protestantismus selber dahin fallen kann, so darf man doch nur froh sein, daß die Ordnung noch besteht
\quellseite{286}
und sich nicht alles ohne kirchliche Vorbereitung zum Abendmahl stürzt. Man muß diese Ordnung wenigstens für den großen Haufen unsrer Gemeinden so lange als möglich festhalten, wenn es auch weiter nichts als Thorheit und Gewaltthat ist, sie von allen und für alle ohne Unterschied festgehalten haben zu wollen. Will man sie aber festhalten, so muß man die Festhaltung möglich machen und nichts Unzweckmäßiges fordern. Die ersten Gemeinden gingen täglich zum Sacrament; da wurde vorausgesetzt, daß sie täglich und immer bußfertig seien. Hätten sie aber, Mann für Mann, alle Tage privatbeichten sollen und können? Man sieht, die Privatbeichte setzt, wenn sie vor jedem Abendmahlsgang geschehen soll, entweder ein selteneres Abendmahlgehen voraus, oder sie muß selber seltener werden und also der Grundsatz dahinsfallen. Wer demnach ein möglichst reichliches Abendmahlsgehen erreichen will, der muß die Privatbeichte auch von dem Haufen der Gemeinde nicht als unerläßliche Voraussetzung fordern, sondern sie mit andern Formen der Beichte wechseln lassen und seine Gemeinde dahin führen, daß sie auf die mannichfaltigste Weise beichte und nur auch die besonders wichtige und segensreiche Privatbeichte, mit der Augsburgischen Confession zu reden, „nicht fallen lasse.“ Wenn das erreicht wird, ist der Schlendrian so viel als möglich entkräftet, die Bußfertigkeit so viel als möglich gefördert, die Gemeinde so viel als es eben möglich ist, zum würdigen Genuß des Sacramentes vorbereitet.

(Es versteht sich von selbst, daß der Pfarrer mit dem, was er nun in seiner Gemeinde erreicht hatte, durchaus nicht
\quellseite{287}
zufrieden sein konnte: die Privatbeichte war ziemlich durchgeführt, aber es fehlte der nöthige Geist und das rechte Maß, und das nächste pastorale Ziel war so erreicht, daß es für das fernerliegende und höhere zum Hinderniß wurde. Es entstand also die Aufgabe, Maß zu setzen und die Gemeinde durch Unterweisung dahin zu bringen, daß sie die Privatbeichte in ihrem Verhältnis zur allgemeinen erkennen möchte.

Diese Aufgabe war eine überaus schwere. Hauptsächlich wurden die Beichtansprachen, die, wenn sie gleich nur die allgemeine Beichte bevorworteten, auch von denen angehört wurden, die bereits privatim absolviert waren, benützt, um Licht zu schaffen. Was gefaßt werden und in’s Leben der Gemeinde kommen sollte, war folgendes:

\begin{quote}
„Die Privatbeichte ist recht und heilsam, sie hätte nie abkommen sollen, sie muß auch allenthalben wieder emporkommen und ist ein hohes Bedürfniß der Kirche. Aber auch die allgemeine Beichte hat ihr Recht und ist auch je und je, schon in den Zeiten, in welchen die Privatbeichte im Schwange ging, in irgend einer Form gebraucht worden, wenn auch nicht in der jetzt gewöhnlichen Form der Abendmahlsvorbereitung. — Wenn die Privatbeichte allein und immer gebraucht wird, wird sie erfahrungsmäßig zum unerträglichen Schlendrian. Wenn die allgemeine Beichte allein gebraucht wird, wird auch sie zu einem noch unerträglicheren Schlendrian. — Man muß beide gebrauchen, je nach der Seelen wechselndem Bedürfniß, damit eine der anderen zur Abwehr der puren Gewöhnung und des Schlendrians diene. Ein jedes Beichtkind muß sich vor der
\quellseite{288}
Anmeldung besinnen, welche Art der Beichte ihm für diesmal die dienlichste sein wird.“
\end{quote}

Man sollte denken, die Sache sei nicht so schwer zu fassen und doch war es in der Wirklichkeit, wie bereits gesagt, ein äußerst schwer zu erreichendes und thatsächlich auch wirklich nicht erreichtes Ziel. Die Leute vom Filial, bei denen die Gewöhnung am kürzesten war, von dem oben die Sprache war, wurden zuerst, so zu sagen, in die Kur genommen. Aber weil sie nicht lebendig und willig genug waren zu der von Gott und der Kirche erforderlichen und also pflichtmäßigen Selbstprüfung vor jedem Abendmahlsgang noch eine Verpflichtung auf sich zu nehmen, nämlich rücksichtlich der Wahl der Beichtweise mit sich selber gewissenhaft einig zu werden; weil sie zu träge waren, so viel religiöse Bewegung in sich aufzunehmen; suchten sie das gedoppelte Ziel des Seelsorgers mit sich selbst in Widerspruch zu setzen. Vorher hatte alles nur allgemein gebeichtet; nun beichtete man der Mehrzahl nach privatim; jenes war nicht recht, dieses ist auch nicht recht: was soll man thun? Da neigt man sich einfach wieder zur allgemeinen Beichte, die ja doch noch viel bequemer ist als die Privatbeichte und obendrein unter Umständen, wenn nemlich der Pfarrer die Gabe der Rede hat, anregender und genußreicher werden kann. So wenig gelang es dem Pfarrer, die Seinen zum besseren Ziel zu führen.

Während die Sache in diesem Stadium schwebte, fiel der Pfarrer in eine Krankheit, die ihn auf lange Zeit hin untüchtig machte, die frühere Thätigkeit als Beichtvater fortzusetzen. Es wurden eine gute Weile nur allgemeine Beichten
\quellseite{289}
gehalten. Als diese Periode vorüber war, mußte ganz neu gebaut werden. Die Filiale, die eigene Beicht- und Abendmahlstage haben, waren während der kranken Zeit des Pfarrers nach der größeren Mehrzahl ihrer Glieder zum alleinigen Gebrauch der allgemeinen Beichte zurückgesunken. Die Parochialorte theilen sich in solche Glieder, welche, wie vorher schon, die allgemeine Beichte, und in solche, die vorherrschend die Privatbeichte gebrauchen. Im Pfarrort selbst ist größeres Verständnis: es gibt zwar solche, die wie früher, so auch jetzt nur zur allgemeinen Beichte gehen, aber es wird auch sehr viel privatim gebeichtet, und es fehlt gar nicht an solchen, welche eine löbliche Abwechslung zwischen den beiden Arten der Beichte beobachten, so daß man zwar nicht sagen kann, das pastorale Ziel sei erreicht, wohl aber dennoch annehmen darf, es sei eine Annäherung zu dem Ziele erfolgt; das heißt, derjenigen Gemeindeglieder, welche die verschiedenen Arten, zu beichten, unterscheiden und mit Verstand benützen, seien mehr geworden. Eine Gemeinde, deren große Mehrzahl über diese Unterschiede klar wäre und sich recht zu verhalten bemühte, wäre freilich nicht bloß eine christliche angeregte, sondern sie hätte bereits eine hohe Stufe des kirchlichen Lebens erreicht. Ist es nun schon eine große Seltenheit, daß eine Gemeinde christliche Anregung empfängt, so ist es jedenfalls noch weit seltener, eine rechte Beicht- und Abendmahlsgemeinde zu finden. Ob irgend ein Seelsorger dies große Erdenglück hat, eine solche zu bedienen, ist die Frage; das wäre ja ganz etwas anderes und eine weit höhere Stufe, als eine Gemeinde zu haben, die sich, wie römisch-katholische Gemeinden, irgendwie
\quellseite{290}
an die Ohrenbeichte gewöhnt und sich mit ihr vertraut gemacht hätte. Man wird wohl darauf verzichten müssen, jemals ein solches pastorales Ziel zu erreichen; dennoch aber bleibt es alles Strebens werth, und ist es in Wahrheit schon viel, auch nur irgend eine Stufe auf diesem Wege zu erreichen. Ein fleißiger Arbeiter kann sich in Anbetracht des Mißlingens immerhin damit trösten, daß Gottes Weg allenthalben schmal ist, und derer wenige sind, die auf ihm wandeln. Der Pfarrer, von dem die Rede gewesen, hat durch lange und viele Erfahrung die Ueberzeugung gewonnen, daß das, was als allgemeines Ziel gemeindlicher Führung unerreichbar ist und nur mit der Resignation auf völliges Gelingen angestrebt werden kann, dennoch den besseren einzelnen Gemeindegliedern als pastorales Ziel und als ein Ziel der Selbstverziehung hingestellt werden kann, daß alles Ringens werth ist und von Segen trieft. Der freie Gebrauch aller möglichen Arten der Beichte je nach Bedürfniß der Seele ist eine herrliche Stufe christlicher und kirchlicher Bildung.

\section{55. Verschiedene Arten der Privatbeichte}

Die meisten unter den jetzt lebenden Menschen protestantischen Bekenntnisses, so wird man wohl sagen dürfen, sind schlechte Beichtkinder, wenn man nemlich ein gutes Beichtkind nur denjenigen nennen darf, der einen rechten Gebrauch von der Privatbeichte machen kann und macht. Der Gebrauch der Privatbeichte setzt eine Erziehung zu derselben und in ihr voraus, eine Erziehung in der Aufrichtigkeit im
\quellseite{291}
allgemeinen und in der Aufrichtigkeit gegen den Beichtvater insonderheit. Eine solche Erziehung aber gibt es unter uns gar nicht, und wir werden wohl sagen dürfen, daß wir alle zumal weder selbst so erzogen worden sind, noch unsre Kinder so erzogen haben, daß Aufrichtigkeit überhaupt und insonderheit Aufrichtigkeit gegen den Beichtvater auch nur ein Ziel der Erziehung gewesen wäre; wir sind aufgewachsen, ohne daß es uns auch nur einmal eingefallen wäre, aus dem Schweigen über unsre Sünden uns einen Vorwurf zu machen. Daher kommt es auch, daß es so vielen unter uns überaus schwer wird, das Herz zu eröffnen und privatim zu beichten, und daß man in die größte Verlegenheit kommt, wenn man es auch einmal probieren will, zur Privatbeichte zu gehen. Mißlingt es nun etwa gleich das erste Mal, so ist man für sein Leben mit der Privatbeichte fertig, ja man wird vielleicht von Stund an ein Uebelredner der Privatbeichte, und obwohl man eine Erfahrung gemacht hat, aus der man sich abnehmen könnte, daß man gar nichts von der Sache versteht und völlig incompetent ist, sie zu beurtheilen, so urteilt man doch und verdirbt nicht bloß sich selbst, sondern auch andern Lust und Geschmack an dem vielleicht edelsten Gewächs, das aus dem Boden der Kirche Gottes frei hervorgewachsen ist. Daß es nun also geht und in Folge dessen die Privatbeichte so überaus schwer in der Kirche wieder empor kommen kann, ist hoch zu beklagen, aber nicht schwer zu verstehen. Man wird ja doch auch sagen müssen, daß es gegenwärtig unter den Pfarrern selbst nur sehr wenige gibt, die etwas von der Privatbeichte verstehen, selbst privatim
\quellseite{292}
beichten, ja nur einmal in ihrem Leben sich in der Lage befunden haben, privatim zu beichten, gar nicht davon zu reden, daß sie andre in die Privatbeichte einführen, in derselben und für sie erziehen könnten. So gilt es denn ganz natürlich: qualis rex, talis grex. Und doch ist die Privatbeichte etwas so Vortreffliches und ein solches Erforderniß des innern Lebens, daß sich aller Orten und bei allen Confessionen auf der Stelle wenigstens Surrogate erzeugen, so wie christliches Leben entsteht. Man kann eben auf dem Wege des Heils die Führung nicht entbehren, ohne daß man den Mangel theuer büßen muß, und die Frage: „Ihr Männer, lieben Brüder, was sollen wir thun?“ läuft seit den Zeiten der Apostel bei allen denen von Mund zu Mund, die vom heiligen Geiste ergriffen worden sind. Daher es ohne allen Zweifel ganz nöthig ist, daß man sich in der lutherischen Kirche besinne, ob man nicht lieber anstatt der Surrogate die Privatbeichte selber mehr in’s Auge fassen sollte und ihr den lohnenden Fleiß zuwenden, welchen sie fordern kann.

Die geringste Form der Privatbeichte ist allerdings die Formel, und sie ist es ja, welche auch da, wo man besseres erstrebt, bei weitem am meisten gebraucht wird. Dennoch aber kann auch sie ganz wohl benützt werden; sie verhält sich zu der rechten Privatbeichte wie eine Fibel zum Lesebuch, wer nur etwas aus ihr zu machen versteht. Bei der niedrigen Bildungsstufe, welche die meisten Landleute einnehmen, kann es kommen, daß die oder jene arme Seele nicht einmal eine Formel kann, sondern anstatt derselben z. B. einen Pathendank aufsagt. Das könnte man alsdann für eine große
\quellseite{293}
Schmach der Privatbeichte nehmen. Und doch gibt so etwas Gelegenheit zur Belehrung; der den Pathendank statt der Beichtformel spricht, kann bei dieser Gelegenheit den Pathendank von der Beichtformel unterscheiden lernen, und geht in Folge des weiser weg, als er gekommen ist. Auch er rechtfertigt das Urtheil eines Beichtvaters, daß am Ende die schlechteste Privatbeichte noch nützlicher sei als eine allgemeine Beichte gewöhnlicher Art, bei welcher so oft die Absolution sammt der Beichtrede verschlafen wird. Die geringste Beichtformel ist eine Tafel der Buße und der Bekehrung, so daß ein Beichtkind aus ihr auch mit ganz geringem Aufwand von Zeit und Kraft über das unterrichtet werden kann, was ihm ewig noth thut. So summarisch, und doch so klar und deutlich wie in der Beichtformel, kann einer armen Seele ihre hohe Nothdurft in anderer Weise gar nicht vorgebracht werden, und dabei hat man erst den großen Vorzug, daß man dem Beichtkinde aus seinem eigenen Munde, denn es spricht ja seine Formel, also gewissermaßen aus eigner Ueberzeugung zu Herzen sprechen kann. Wie viel man an die Beichtformel anhängen, in welch’ weitem Kreise man sie benützen kann, beweisen die sogenannten Absolutionsbücher der lutherischen Kirche, so sehr sie doch nur ein Zeichen der untergehenden, seelsorgerischen Thätigkeit der lutherischen Kirche im Beichtstuhl genannt werden müssen. Man darf also mit den Formelbeichten zwar allerdings nicht zufrieden sein, aber sie auch nicht so verachten und solchen Ekel vor ihnen fassen, daß man sie lieber ganz wegwürfe, was doch im Grunde nichts anders hieße, als mit dem Bade das Kind ausschütten. Ein Pastor
\quellseite{294}
lege sich nur drauf, seine Elementarchristen geduldig in das Abcbuch, in die Fibel der Beichtformel einzuweisen, und lasse sich die Mühe nicht verdrießen, so kann er den Segen finden. Er wird wenigstens aufmerksamere Schüler und Zuhörer nirgends finden, als im Beichtstuhl, auch nirgends bereitere Herzen und willigere Aufnahme dessen, was er zu sagen hat.

Aber allerdings ist die Beichtformel nicht die höchste Stufe der Beichte, obwohl sie allerdings oft auch dem gefördertesten Christen dienen kann, und viele von denen, die von der freien Uebung der Privatbeichte den größten Segen bekommen haben, endlich wieder mit beruhigter Seele zum Gebrauch der Formel zurückkehren. Zwischen den beiden Zuständen des Elementarchristenthums und der regelmäßig und im Frieden fortschreitenden geheiligten Seele liegen Zustände mitten inne, für welche die Formel durchaus unzulänglich ist. Da muß nothwendig die freie Aussprache der Seele gefordert werden, und sie erzeugt sich auch wie von selbst. Wenn die Seele lebendig ergriffen wird, wie die Hörer der ersten Predigt in Jerusalem, dann wird sie auch reden; die Beichten erweckter Seelen geschehen mit den eigensten Worten derselben, Scheu und gêne fallen weg und die Noth lehrt beichten und beten, im Dialecte oder wie es eben gehen will. Tritt irgend einmal so eine Zeit ein, so wird ein einigermaßen tüchtiger Beichtvater sie auch zu benutzen wissen, der armen Seele entgegenkommen, ihr Verlangen klären und sie dahin führen, daß sie endlich in dem feierlichen Ausspruch ihres Bedürfnisses und in der Absolution die größte Beruhigung und Stärkung findet, und daß sie sich fortan mit Lust und
\quellseite{295}
Eifer der besten Schule übergibt, die unter dem Himmel ist, dazu der besten Erziehung und Selbsterziehung, nämlich der Schule und Erziehung der Privatbeichte.

\section[56. Arbeit eines Beichtvaters]{56. Arbeit eines Beichtvaters bei der Privatbeichte}

Wenn ich an meinem Orte stehe und auf die Beichtkinder warte, und dann eines um das andere herzutritt, um zu beichten und nicht bloß Absolution, sondern auch irgend einen Rath oder sonst ein Wort zu empfangen, das Licht, Kraft oder Trost bringen kann; so weiß und fühle ich, daß meine ernsteste Stunde in der ganzen Woche gekommen ist, und daß ich vor meiner schwierigsten Aufgabe stehe. Mein Auge hängt an der Pforte und wartet auf den Eintretenden. So wie ich ihn sehe, tritt mir die Gestalt seiner Seele vor Augen; ehe er noch redet, frage ich mich schon, was ich ihm zu sagen habe; ich forsche und studiere mehr, als wenn ich Bücher aufgeschlagen habe und darin nach Weisheit suche. Fängt nun das Beichtkind an zu reden und seine Seelenangelegenheiten zu offenbaren, so werde ich ganz Ohr; — nichts verloren gehen, nichts unbeachtet zu lassen, ist mein ganzer Sinn. Dabei kann ich nicht anders, ich muß forschen, was ich zu sagen habe. Ich weiß, daß ich nun bin wie eine Blume, die eben aus der Knospe brechen, ihre Gestalt und ihren Ruch geben soll. Wie ernst ist mir's bei meiner Seligkeit, geben zu dürfen und zu sollen, nicht aber zu nehmen! Und diese meine innere Lage wiederholt sich, so oft ein neues Beichtkind erscheint. Ich muß in so viel besondere
\quellseite{296}
Lebensverhältnisse eindringen, als Beichtkinder kommen. Nichts ist in meinen Augen schwerer als das: belastet von meinem eigenen pastoralen Zustand, muß ich doch anderer Leute Zustände zu den meinigen machen, so daß ich drin lebe, aus ihnen heraus und für ihre Heilung rede. Ich stehe still, kein Fuß, keine Hand regt sich; aber wenn es so einige Stunden fort gegangen ist, — dann wills nicht mehr gehen; ich fühle, daß es mir zu viel wird. Müde geworden, möchte ich ausruhen von diesen Reisen meiner Seele in fremdes Leben hinein; ich möchte wieder eine Weile mit gar nichts außer mir selbst zu thun haben, ja auch mit mir selbst nicht: pur ruhen möchte, absehen von allem und jedem. Ich höre und rede oft mit steigender Freude, ich fühle gewisser Maßen, daß Lebenswasser von mir fließt, — ich merkte Segen von oben, und daß ich im Hause meines Herrn diene; aber dennoch gehts nicht auf die Länge. Ich kenne keine Arbeit, wie die der Privatbeichte, bei der man, wenn nicht vor Jammer, so doch vor Freude müde wird. Leib und Seele, Nerven und Gedanken fühlen ihre Grenzen. Darum würde ich auch einem jeden, bei dem es mit Beichtgehören Ernst wird, rathen, sich nicht zu übernehmen, sondern es einzutheilen und einzurichten. Es ist für die Beichtkinder selbst vom größten Nutzen, wenn der Beichtvater munter und kräftig ist, zu hören und zu reden. Das kann er aber nicht, wenn er zu viele nach einander hören muß und bescheiden: die letzten spüren es, sie bekommen nichts. Darum trau Dir und Deiner Kraft nicht zu viel. — Freilich, wenn Du nur Formelbeichten hast und auf alle die Formeln mit aufgelegten Händen nur wieder eine
\quellseite{297}
Formel, die Absolution, spricht, das kann Dir am Ende auch widerwärtig werden; aber die Arbeit, von der ich rede, und das Ermüden, das ich kenne, lernst Du nicht kennen. Man wird des Hand- und Mundwerks gewohnt.

\section{57. Nöthige Weisheit eines Beichtvaters}

Schon die Anstrengung der Seele, der Nerven, ja des Leibes ist groß, wenn man ernstlich Beichte hören soll. Schon das Wollen, Suchen, sich anstrengen ist schwer. Aber wenn man nur kann, was man will und soll, — findet, was man sucht, und die Anstrengung ihr Ziel erreicht! Allein da hilft bei dem, der nichts weiß und kann, kein Besinnen. Da helfen keine Phrasen, und das bißchen Pulver, das einer aus seiner armen Erfahrung mitbringt, ist bald verschossen. Man wird nach Rath und Unterweisung hungrig. Da helfen auch keine Collegienhefte; die Schule hat wenig Rath für die Geschäfte des höchsten Amtslebens. Zwar ja doch: der Schreiber dieses hat keine Collegienhefte übrig, aber doch hat er in seiner Jugend etwas gehört, das er schwarz auf Weiß nach Hause getragen hat und nun im Alter doch oft brauchen kann. Es ist ein Collegium über Pastorale und Katechetik, vom seligen Hofprediger Strauß in Berlin gelesen. Die Commilitonen mochten's nicht; aber er hats nachgeschrieben, weil es ihm Segen und Nutzen fürs Amt verhieß. Immerhin ists ihm aber doch auch oft für das Amtsleben nicht genießbar gewesen. Es gibt eben keine protestantischen Bücher, von denen ein Seelsorger den erwünschten nöthigen Segen nehmen könnte, —
\quellseite{298}
zumal wenn es sich von der Seelsorge im höhern Chor, von der Seelsorge im Beichtstuhl handelt, es sei denn, daß man zu der theologia conscientiaria, zu der Casuistik, Zuflucht nehme. Ein Beichtvater braucht kaum etwas nöthiger, als casuistische Studien. Aber ich will nicht wiederholen. Dagegen will ich zum ehrlichsten Diebstahl ermahnen, den es gibt, — nämlich zur Beraubung der römischen Pastoralschriften und Beichtspiegel. Hier sind leicht zu hebende Schätze in Menge. Der große Thomas von Aquino hat ein kleines Büchlein hinterlassen, aus dem ein jeder sich überzeugen kann, was alles noch vom Mittelalter, namentlich in Sachen der Seelsorge und Beichte zu lernen ist. Es gibt ja solcher Bücher viele: was ist z. B. von Charlier Gerson zu lernen und aus seinen reichen Schriften zu nehmen. Aber ich will einmal zunächst nur auf Thomas von Aquin verweisen. Die kleine Schrift, von welcher ich rede, ist betitelt: Confessionale seu libellus peroptimus beati Thomae de Aquino de modo confitendi et de puritate conscientiae, cuilibet confessori et confiteri volenti perutilis et necessarius. Möchte es bald auch unter uns wieder Leute geben, welche solche Schriften anzuwenden und für andere zurecht zu machen verstehen! — Aehnlich verhält es sich aber auch mit andern der römisch-katholischen Kirche gehörigen Schriften. So viel man in ihnen nicht brauchen kann, so viel kann man brauchen. Die Regula Aurea Confessariorum oder Instructiones saluberrimae Caroli Borromaei wiegen für den Zweck der Beichtväter schwerer, als manche pastorale Bibliothek. — Wer so etwas sagt, kann verdammt, und seine Worte können mißdeutet werden.

\quellseite{299}
Auf die Länge aber wird man keine Verketzerung zu fürchten haben, weil wahr bleibt, was wahr ist. — Kurz, wir brauchen Studium und Unterweisung zum beichtväterlichen Amte.

\section{58. Gefahr für Beichtväter}

Wer als Beichtvater Vertrauen findet und die Macht der Privatbeichte kennen lernt, kommt allerdings auch in Gefahr, — nemlich der Unbescheidenheit. Es gibt keine Macht auf Erden, welche der eines Beichtvaters über willige Seelen gleichkommt. Vom Beichtstuhl aus können bei den Protestanten nicht weniger als bei den Römischen alle natürlichen Verhältnisse gesegnet werden. Aber vom Beichtstuhl aus können auch alle natürlichen Verhältnisse unterminiert und gesprengt werden. Und zwar kann das Leuten kommen, die sich so etwas gar nicht vorgenommen haben und es sich nicht nachsagen lassen, auch wenn es noch so wahr ist. Der Beichtvater gleicht einem Landmann, der seine Wasser in die fernsten Winkel seines Gartens leitet, — und der, froh seiner Kraft und Kunst, alles aufbietet, um mit seinem Ziele fertig zu werden. So leitet der Beichtvater die Wasser des göttlichen Wortes bis in die speciellsten Lebensverhältnisse hinein. Gelingensfroh kann er dann recht mit vollem Bewußtsein darnach streben, Herr aller Verhältnisse zu werden. Es kann ihm gelingen und ist ja auch oft genug gelungen, Vater und Mutter, Mann und Frau aus dem Sattel zu heben, wider einander zu hetzen und in Aufruhr zu bringen, den Frieden der Familien zu vertilgen und ein Feuer anzuzünden, von dem
\quellseite{300}
Christus nicht gemünzt hat, es brennete schon. Junge Seelsorger, die anfangen, ihre Macht zu spüren, lassen sich oft verleiten, in alles hineinzuschauen, alles zu besprechen und so über das Volk zu herrschen. Davor behüte aber Christus seine armen Knechte, die zu allen ihren Mühseligkeiten und Sünden nicht auch noch die brauchen, die Herren und Regenten auf allen Gebieten des Lebens werden zu wollen oder gar zu werden. Heilige Bescheidenheit und Zurückhaltung ist für Beichtväter nötig. Du mußt auch als Beichtvater in kein Verhältnis eingreifen, für das zu sorgen für dich kein Beruf vorhanden ist. Ein rechter Beichtvater will nicht alles wissen, alles bereden, ringt darnach, sein beichtväterliches Amt so zu führen, daß er niemand beschwert, daß er jedem Lebensverhältnis ferne bleibt, unwissend, und gerade dadurch voll unschuldigen Einflusses auf alle seine Beichtkinder und deren Verhältnisse wird. Er wartet ab, was an ihn kommt, — und thut dann sein Amt, als thäte ers nicht. Es gibt allerdings ein solches Verhalten, bei welchem man der beste Freund einer Familie und doch ferne von ihr bleiben kann. Erlöse Gott die Seinen von dem Wahn, als müsse und solle ein Beichtvater Hausfreund und Kamerad aller Beichtkinder werden. Das Finden — und Fernhalten, lehre uns Gott!

\section{59. Beichtgeheimniß}

Wo wird der Beichtvater sein, dem nicht zuweilen auch von seinen besten Beichtkindern der Vorwurf gemacht wird, er habe das Beichtgeheimniß verletzt? Natürlich! Jede Beichte
\quellseite{301}
gehört in irgend ein Bereich allgemeiner Sünden. Wenn nun der Seelsorger allgemeine Sünden in oder außer der Beichte vorbringt, so bilden sich manche ein, er bespreche ihre Verhältnisse, die er doch nur beichtweise kenne. Fromme Männer lernen auch das ertragen sammt allem, was daraus an Argwohn und Haß gegen sie entspringt. — Viele Beichtkinder haben ferner die Unart, ihre Umstände dem Seelsorger beichtweise, andern aber im guten Vertrauen der Freundschaft mitzutheilen. Kommt nun etwas an’s Licht, so mißtrauen sie nicht ihren Freunden, bei denen sie Ursache hätten, sondern ihrem Beichtvater, der vielleicht sogar die Gabe und Gnade hat, die Beichten zu vergeben, und vielleicht schon deshalb keine Silbe verräth. Wenn die Beichtväter unglücklich genug sind, in Verdacht des Schwatzens zu kommen, so hilft manchmal keine Versicherung noch Beweis, sondern Abneigung und Widerwille bleibt unüberwindlich. Diese und dergleichen Erfahrungen entschuldigen, ja rechtfertigen zuweilen die Entschließung, ganz wie es bei den Römischen geschieht, absolutes Schweigen und Verhüllen auch derjenigen Worte und Auslassungen der Beichtkinder zu beobachten, die gar nicht zur Beichte gehören und gehören können. — O die Beichte! Wie viele mißbrauchen sie zu Mittheilungen, welche mit dem Sündenbekenntniß gar nichts zu thun haben! Wie viele mißbrauchen sie zu Lügen, Verläumdungen und jeglicher Bosheit, — wie viele lügen ganz absichtlich, um entweder den Beichtvater zu hintergehen oder böse Absichten anderer Art zu erreichen. Die Beichte ist oft ein Pfuhl der Sünde und Selbstsucht und ein Tummelplatz aller Leidenschaften. Da muß sich ein Beicht%
\quellseite{302}
vater sicher stellen und das Schweigen und Vergeben studieren, ja nicht bloß schweigen, sondern auch sein lernen. — So einzig ist die Stellung eines Beichtvaters, daß er nicht einmal von seinem in der Beichte genommenen Urtheil über ein Beichtkind Gebrauch machen darf. Er muß nicht wissen, auch was er weiß; nicht Falschheit, sondern hohe Redlichkeit ist es, über dieselben Leute ein Beichturtheil und ein Urtheil außerdem zu haben und von einem auf das andere den Einfluß abzuwehren.

— — Wie unklar sind auf diesem Lebensgebiete unsre Gedanken, wie verkehrt unser Verhalten! Wir können nur erst unklar sein; wir haben ja blutwenig Beichterfahrung. Wie lange wird es dauern, bis wir nur wieder wissen, was Beichtgeheimniß ist und was nicht, und bis Beichtväter und Beichtkinder lernen, was für Pflichten sie gegenseitig übernehmen und haben! — Und wie lange wird es nur dauern, bis man sich die Mühe gibt, sich zu besinnen und dem edelsten Institut der Kirche, der Privatbeichte, die Lebensbedingungen zu schaffen, welche nöthig sind, welche es nicht entbehren kann.

\chapter{Krankenseelsorge}

\section{60. Deren Werth überhaupt}

Leute, welche ihre Urtheile über den Werth der Krankenseelsorge nicht aus der täglichen Erfahrung schöpfen oder schöpfen können, sondern nach Vorurtheilen oder Eindrücken zu sprechen gewohnt sind, pflegen den Werth der Krankenseelsorge eben so zu über%
\quellseite{303}
schätzen, wie gerade die, welche täglich an Krankenbetten stehen und das heilige Amt ausüben, in der entgegengesetzten Versuchung sind, nemlich gar zu gering vom Werthe derselben zu halten und zu sprechen. Die Wahrheit ist, daß die Krankenseelsorge selten augenfällige Erfolge hat. Die meisten unserer Gemeindeglieder sind Elementarchristen, die mit einem Minimum geistlichen Lebens ganz zufrieden sind, sich, ihre innere Noth und ihre geistlichen Bedürfnisse nie erforscht haben und auch zu roh und ungebildet, zu stumpf und zu gleichgültig sind, als daß sie sich zu einer Selbstprüfung treiben ließen. Schon in ihrem gesunden Zustand findet man keinen Eingang zu ihren Seelen, und wenn sie krank werden, will es noch härter gehen. Zuweilen fördert allerdings die Krankheit und lehrt, wie andere Anfechtungen, aufs Wort merken; aber gewöhnlich bringt sie statt Fördernisse der Bekehrung oder der Erweckung aus dem trägen, dumpfen Seelenschlafe nur Hindernisse. Der Seelsorger hat daher alle Ursache, zu predigen: „Spare deine Buße nicht, bis du krank wirst.“ Das Wort und die Weisheit des Seelsorgers findet, ich wiederhole zur Bestätigung, den Weg zur Seele im gesunden Leibe schwer, nun erst zur Seele im kranken Leibe! Ganz besonders aber sind die Seelen kranker Landleute schwer zu finden; der Mangel an allgemeiner Bildung erschwert alle seelsorgerische Arbeit, wie ein schlechter Weg das Vorwärtsgehen hindert. O ist die Seelsorge der Kranken so erfolgarm, erfolglos, — und, versteht sich, so schwer, so schwer! Da gibt es sogenannte Demüthigungen und Enttäuschungen! Es ist so trostlos, unter den Kranken herumzugehen und keine weitere Genüge
\quellseite{304}
zu haben, als daß man doch recht gethan und Liebe geübt hat, wenn man die edle Zeit im Krankendienst verbrachte. Leg daher vornherein Deine sanguinischen Hoffnungen nieder und hoffe weder für Dich, noch für die Kranken von Deinen Besuchen sehr viel; Du gewöhnst die Besuche sammt dem Sterbensangesicht, — und die Leute selbst sind langsam, ihre Herzenspforte zu öffnen. Es ist wahr, was Sailer (Pastorale III. p. 6) sagt: „Dem Sterben eines Menschen beiwohnen, ist noch weit lehrreicher, als dem Aufgang der Sonne beiwohnen, so sehr auch dieses empfehlenswerth ist, oder vielmehr keiner Empfehlung bedarf.“ Aber nimm’s nur nicht gar zu groß und hoch. Nimm lieber die Krankenpflege für recht schwer und wundere Dich nicht, wenn Du nirgends Deine Ohnmacht und Kraftlosigkeit so sehr erkennst, als gerade an Kranken- und Sterbebetten. Der ist am glücklichsten, der wenig erwartet; er ist dann desto dankbarer für jeden, wenn auch kleinen Erfolg. Und werden Dir nun die Krankenbesuche geringer, weil Du so geringes leistest; so folge dem weisen und guten Sailer desto lieber, wenn er am angeführten Orte p. 14 sagt: „Gehe nie zu einem Kranken, ohne den Vater des Lebens mit kühnem, starkem Vertrauen anzuflehen um die seltene Gabe, überall den Weg in das Herz eines trostbedürftigen Mitmenschen zu finden und den Sünder von jeder Straße zu seinem Gott zu führen.“ — Ohne Gottes ganz besondern Segen wirst du wenig wirken, und Du wirst an Deinen Krankenbesuchen große Last haben, wenn Du nicht etwa aus purem Mitleid und
\quellseite{305}
Verlangen, Liebe zu erweisen, aushältst und Dich deshalb immer und immer wieder bei den Krankenbetten einfindest.

\section{61. Besonderer Zweck der Krankenpflege}

Die Seelsorge bezieht sich auf alle Christen, also auf leiblich Kranke eben sowohl wie auf Gesunde. Da die Kranken keinen andern Weg zur Seligkeit haben, als die Gesunden, alle den christlichen Heilsweg gehen müssen, auch für die Kranken keine andern Gnadenmittel verordnet sind, als für die Gesunden; so ist die Seelsorge der Kranken und der Gesunden wesentlich eine und dieselbe. Der Seelsorger muß eben die Person des Kranken und seinen pastoralen Zustand kennen lernen, aus Gottes Wort das rechte Heil- oder Förderungsmittel auf dem Weg zum Leben wählen, es weislich und treulich anwenden. Und doch ist die geistliche Krankenpflege von der besondern Seelsorge der Gesunden verschieden. Der Seelsorger hat nemlich nicht bloß den Kranken auf den Heilsweg zu führen und auf demselben zu fördern und zu erhalten bis an’s Ende; sondern man verlangt von ihm auch, daß er die eigenthümlichen Hindernisse kennen und heben solle, welche jede Krankheit dem Menschen auf dem Weg zum ewigen Leben entgegenstellt. So hat z. B. die Bauch- oder Hautwassersucht das Eigene, daß sie den Kranken nicht bloß stille und gelassen, sondern auch zuweilen stumpf macht; scrophulöse Leiden machen Kinder, die mit ihnen behaftet sind, oft eigensinnig, zänkisch, launisch, verschmitzt, herrschsüchtig; Lungenschwindsüchtige hängen sehr oft am
\quellseite{306}
Leben wie an dem höchsten Gute \&c. Da verlangt man nun von dem Seelsorger nicht bloß, daß er den allgemeinen pastoralen Zustand des Menschen kennen und bessern solle, er soll nicht bloß den Unbußfertigen zur Buße, den Bußfertigen zum Frieden Gottes, den Gläubigen zum Fleiß in der Heiligung \&c. führen; sondern er soll die psychischen und moralischen Einwirkungen der Krankheiten in’s Auge fassen und durch die Wunderwirkung des göttlichen Wortes die Verstimmungen der Seele des Kranken und die ihr drohenden Versuchungen und Sünden überwinden. Der Wassersüchtige soll aus der Gefahr der Stumpfheit zu einem heitern, regen Interesse an dem Wohlergehen der eigenen und anderer Seelen gefördert, — das scrophulöse Kind sanftmüthig und voll Energie gegen seine Versuchungen, der Lungenschwindsüchtige todesmuthig und todeslustig werden. Dieser besondere Theil der Krankenpflege setzt aber Kenntniße, Beobachtungsgabe, eingehendes Erfahren und großes Geschick voraus, wenn etwas geleistet werden soll, und wer deswegen ein wenig in diesen Beruf hineingestiegen ist, wird bald merken, was für eine Einladung zur Bescheidenheit und Demuth ihm entgegenkommt.

\section{62. Ueber Anfechtung und Besessenheit}

Man theilt die Anfechtungen in drei verschiedene Klassen:\\
1) in leibliche,\\
2) in geistliche,\\
3) in gemischte.

\quellseite{307}
Anfechtungen der ersten und zweiten Gattung sind selten, die häufigsten sind die der dritten Gattung. Wenn man von leiblichen Anfechtungen redet, so versteht man darunter Anfechtungen dämonischer Art, Einwirkungen des Satans auf das leibliche Befinden, sei es durch ein unnatürliches Wohl- oder Uebelbefinden. Insbesondere versteht man unter leiblichen Anfechtungen auch die dämonischen Krankheiten.

Zu den geistlichen Anfechtungen gehören die übermäßigen und unordentlichen Einwirkungen gewisser Gedanken, Begierden und Gefühle auf den gesammten Zustand des Menschen, mögen nun diese Gedanken, Begierden und Gefühle von Menschen oder von Dämonen stammen.

Gemischt nennt man die Anfechtung, wenn sie von dem einen Theil des Menschen, also dem Leibe oder der Seele, ausgehen und seinen zweiten Theil in ihr Bereich und in ihre Macht mit hineinziehen. Es kann sich dabei eine Anfechtung dieser Art nicht bloß an eine Anfechtung der ersten Art anschließen, sondern es kann auch irgend ein von Gott gewollter, leiblicher Zustand des Menschen zur Ursache der Anfechtung für die Seele dienen.

Die Kennzeichen der leiblichen Anfechtung sind nicht leicht. Bei unseren Vätern hielt man dafür, daß dämonisch-leibliche Anfechtung vorhanden sei, wenn kein arzneiliches Mittel seine Dienste that. Doch ist dieses Zeichen an und für sich selbst ein unzuverlässiges, weil möglicher Weise das richtige arzneiliche Mittel nicht getroffen sein kann. Es wird daher auf diesem Wege selten mehr als eine bloße Wahrscheinlichkeit erreicht werden können. Die Römischen gebrauchen
\quellseite{308}
nicht bloß gegen die wirkliche Besessenheit, sondern auch gegen den Zauber, den man zu den dämonischen Anfechtungen rechnen kann, ihre Exorcismen und nehmen an, daß dann Zauber vorhanden sei, wenn sich das Uebel durch die Beschwörungen mächtiger regt.

Allein da nicht bloß der Teufel ein Lügner ist, sondern auch die menschliche Natur voll Truges, Heuchelei und Gleisnerei zu sein pflegt, so gibt auch die Anwendung der Exorcismen, auch wenn man gegen dieselben nichts sonst einzuwenden hätte, kein sicheres Zeugniß vorhandenen Zaubers oder dämonischer Anfechtung. Man wird daher auf gewisse Zeichen verzichten und sich mit der größeren oder geringeren Wahrscheinlichkeit begnügen müssen.

Leichter erkennbar sind die rein geistlichen Anfechtungen, wofern man nemlich nur herausbringen will, nicht von wem die Anfechtung stammt, sondern nur, daß Anfechtung da ist. Das letztere herauszubringen ist die Hauptsache; am ersteren liegt weniger. — Eine geistliche Anfechtung beruht entweder auf einem falschen Gedanken oder auf einer falschen Begierde oder auf einem falschen Gefühl: ich bin entweder durch Gedanken oder Begierden oder Gefühle angefochten. Ist ein Mensch sonst verständig, er hegt aber trotz vorhandener guter Begabung und sonstiger Nüchternheit einen kenntlich falschen Gedanken und ist geneigt, ihm Folge in seinem Leben zu geben, so ist er angefochten, woher auch seine Anfechtung komme. Hat einer im allgemeinen einen guten Vorsatz und eine lautere Begier, er kann aber im einzelnen einer unstatthaften Begierde nicht los werden und ist in der Versuchung, ihr nach%
\quellseite{309}
zugeben, so steht er in der Anfechtung, woher sie auch immer stamme. Hat jemand irgend ein Gefühl, zu welchem kein Grund vorhanden ist, dem vielleicht geradezu vorhandene Gründe widerstreben, er fühlt aber in sich eine Neigung, dem Gefühle Statt zu geben, so ist er in der Anfechtung, woher sie auch stamme.

Hat eine geistliche Anfechtung, von welcher Art sie auch sei, Einwirkung auf das leibliche Befinden, so kann ich vielleicht eine Weile zweifeln, ob dies der Fall sei, oder ob vielmehr vorhandenes leibliches Befinden auf die Seele den Einfluß übe; es wird aber die Mischung des Geistlichen und Leiblichen die Erkenntniß der vorhandenen Anfechtung nicht erschweren, sondern vielmehr erleichtern.

Eine Hoffnung auf Freiheit von der Last der Anfechtung kann man einem jeden Angefochtenen insofern geben, als schon das Wort Anfechtung nicht den höchsten Grad des Leidens andeutet, sondern im Grunde nur einen Anfang der Leiden bedeutet. Doch wird man wohl thun, die Hoffnung selber in’s Ungewisse zu stellen und theils von dem Willen Gottes, theils von dem Verhalten des Angefochtenen abhängig zu machen, da es offenbar ist, daß Gott nicht einen jeden Menschen, nicht einmal seinen Apostel Paulus, von Anfechtung befreien will, die er zugelassen, und ebenso, daß ein widerstrebendes und ungeschicktes Verhalten des Kranken die Wirksamkeit aller Mittel abstumpfen kann. Was insbesondere die dämonischen Anfechtungen betrifft, so wird man wohl thun, auf sie mit aller Bestimmtheit die Belehrung zu beziehen, welche man dem Angefochtenen überhaupt über die Behandlung seiner Anfechtung
\quellseite{310}
und deren Wichtigkeit für sein ewiges Heil zu geben hat. Dämonische Anfechtungen, und zwar ebensowohl die leiblichen als die geistlichen, sind zufällige, wenn auch noch so schwere, Lasten, auf welche ganz und gar der Spruch anzuwenden ist:

\begin{quote}
Denen, die Gott lieben, müssen alle Dinge zum Besten dienen (Römer 8.).
\end{quote}

Durch Darlegung dieser Anschauung muß dem Kranken vor allen Dingen die Ueberschätzung seiner Last genommen werden, welches die Basis für jede Heilung und Genesung von solchen Uebeln ist. Es wird also dem Angefochtenen ungefähr folgendes gesagt und eingeprägt werden müssen:

\begin{quote}
„Deine Krankheit ist an und für sich kein unheilbares Uebel; insofern hast du Hoffnung auf Besserung. Doch ist es möglich, daß dein andauernder Leidenszustand dir größeren Segen und Gott größere Ehre bringt; in diesem Falle ist die Heilung nicht wünschenswerth, auch kannst du durch Mißtrauen gegen deine Berather und Seelsorger, sowie durch Ungehorsam gegen ihre Anordnungen die Heilung selbst hindern oder unmöglich machen; in diesem Falle trägst du die gerechte Strafe deines Verhaltens. Du bedarfst daher vor allen Dingen Ergebung in die Hand des Herrn, Geduld, wenn dein Leiden länger währt, und wenn du dahin kommen könntest, wohin St. Paulus kam, Gott in deiner Schwachheit mächtig sein zu lassen, so wärst du ein Schauspiel der Engel und eine Freude Gottes, sintemal Dieser durch freudige Ertragung der Uebel von Seiten schwacher Christen mehr gepriesen wird als durch Aufhebung des Uebels selbst. Auch dir wäre im ersten Fall so viel
\quellseite{311}
geholfen als im letzten, weil einem Hochbeladenen gleiche Hilfe geschieht durch Vermehrung der Geduld wie durch Verringerung der Last.“ —
\end{quote}

Auf diese Weise begrenzt man die Hoffnung auf Genesung richtig, und hat die Hilfe schon vorbereitet, ja begonnen.

Was die Behandlung der Angefochtenen anlangt, so wird immer das erste bleiben, den Kranken über seinen Zustand aufzuklären und ihn mit der Wahrheit zu bedienen. Sowie der leibliche Arzt das Vertrauen der Patienten verliert, der sie mit Lügen bedient, so verdient der Seelenarzt kein Vertrauen, der seine Angefochtenen oder Geisteskranken über sich selbst in Zweifel läßt. Es kann und wird Ausnahmen geben, aber im allgemeinen wird man wohl sagen dürfen, daß man sich bei Angefochtenen und Geisteskranken an den noch vorhandenen, wenn auch vielleicht geknechteten, gesunden Sinn mit Darlegung der vollen Wahrheit wenden müsse; der vollen Wahrheit also auch rücksichtlich der Hoffnung oder Hoffnungslosigkeit des Uebels, wie schon im vorigen Paragraphen gesagt ist. Die Wahrheit muß aber dem Kranken mit Ruhe und Zuversicht einerseits und mit Sanftmuth andererseits vorgetragen werden. Diese Ruhe und Zuversicht aber darf man nicht eher heucheln als haben; man muß dem Kranken so lange nicht verhehlen, daß man keine Aussicht von seinem Leiden habe, als man keine hat. Der Erfahrene wird es für keinen Nachtheil achten, unwissend und unerfahren zu erscheinen, er wird desto seltner getäuscht. Nur der Eitle will gern in allen Fällen den Wissenden spielen.

\quellseite{312}
Ist erkannt, daß eine leiblich dämonische Anfechtung vorhanden ist, so ist das rechte Mittel gegen dieselbe:

1) die schon für die Vorbereitung angegebene Belehrung über die Wichtigkeit und Bedeutung dämonischer Anfechtungen, welche zur Geringschätzung und Verachtung desselbigen führt;

2) das oft wiederholte und anhaltende Gebet gegen die Anfechtung;

3) ob Exorcismus — wird in der lutherischen Kirche die Frage sein.

Auch psychische Einwirkung auf anderem Wege ist möglich. Wenn es nämlich gelingt, den Kranken für eine große Sache oder Wahrheit zu begeistern, so wird er dadurch über sein Leid hinübergeführt und die Macht des Guten kann dadurch in ihm so groß werden, daß der Dämon sein vergebliches Bemühen, ein Kind Gottes zu fällen, aufgibt.

In Betreff der geistlichen Anfechtungen haben wir es zuerst mit der Klasse derjenigen zu thun, deren Anfechtungen entweder in falschen Gedanken, oder in einer falschen Anwendung richtiger Gedanken bestehen. In beiden Fällen ist es das erste und nöthigste, dem Angefochtenen, so wie man selbst Bescheid weiß, die volle Wahrheit über seinen Zustand zu sagen, den vollen Widerspruch gegen seine falschen Gedanken eintreten zu lassen, und denselben mit den nöthigen Beweismitteln zu versehen. Zu diesen Beweismitteln selbst darf man nicht immer neue setzen, sondern es handelt sich darum, recht klar und einfach diejenigen, welche man einmal aufgestellt hat, dem Kranken zu wiederholen, bis er sie gefaßt hat. Da
\quellseite{313}
aber ein Angefochtener ein angehender Geisteskranker ist, so hat man es selten allein mit seinem Verstande und dessen falschen Vorstellungen zu thun, sondern sehr oft auch mit seinem Willen oder vielmehr Unwillen. Der Kranke will sehr häufig nicht geheilt sein, sondern er folgt dem Zuge, der in ihm ist, und klammert sich an die Gedanken an, die ihn plagen. Da hat man es dann auch nicht bloß mit Lehre, sondern auch mit Strafe, Ermunterung und Erziehung zu thun. Diese verschiedene Anwendung einer und derselben Wahrheit muß fortgesetzt werden, bis der Sieg errungen ist. Zur Anwendung jeder einzelnen Stufe der Behandlung gehört Weisheit und Ueberlegung, so wie auch große und andauernde Geduld. Je länger die Anfechtung schon dauert, desto mehr Zeit bedarf es, sie zu überwinden; man kann daher keine schnellen Fortschritte erwarten. Möglich ist es, daß da, wo es in der Macht des Seelenarztes steht, es so zu ordnen, eine völlig entsprechende Umgebung sehr förderlich wird. Ist die Anfechtung in der Begier zu suchen, so hat man nicht zu denken, als ob es nur an dem einer jeden Begier zu Grunde liegenden Gedanken läge; liegt es an dem, so gehört die Anfechtung in die erste Klasse. Man versteht unter der Anfechtung der Begier nichts anderes, als ein verkehrtes Maß derselben, wie z. B. ein Mensch durch Mangel an Liebe oder durch Uebermaß der Liebe zu einem andern in große Seelennoth gebracht werden kann. Bei der Behandlung solcher Angefochtenen ist, wie immer, vollkommene Wahrheit zu gebrauchen. Ist z. B. zu wenig Liebe da, so raube man sich das Vertrauen des Angefochtenen nicht durch den falschen Beweis, als hätte er Liebe
\quellseite{314}
genug, und erkenne es nur nicht. Hat aber der Kranke allerdings Liebe genug und erkennt es nur nicht, so lasse man sich auch nicht hindern, dem Kranken diese Wahrnehmung mitzutheilen, auch wenn er das Vertrauen verwirkt, denn er kann es wiederfinden. Auf vollkommene Wahrheit muß die ganze Behandlung gegründet werden. Ist ein Ueberschwang der Begierde vorhanden, so ist die Behandlung leichter, als wo ein Mangel hervortritt, weil der Widerspruch leichter ist, dagegen aber oft gar lange gewartet werden muß, bis das angefochtene Herz seinen Mangel durch die Kräfte des göttlichen Worts erstatten läßt. In allen Fällen bleibt das göttliche Wort das Hauptmittel, die Aenderung der Umstände und Verhältnisse aber ein zweckdienliches Nebenmittel.

Unter den gemischten Anfechtungen versteht man hauptsächlich solche, die ihren Ursprung zunächst entweder im Leibe oder im Geiste haben, die aber vermöge der Wechselwirkung Leibes und der Seele auch den andern Theil des ganzen Wesens ergreifen und mit in die Leiden hineinziehen. Dabei ist es gleichgültig, ob es dämonische Anfechtungen sind, oder ob dieselben auf natürlichem Wege ihre Erklärung finden. So wie nun bei Erscheinung eines Leidens immer der Arzt oder Seelsorger oder Pfleger zuerst mit sich in’s Reine zu kommen hat, was für ein Uebel er vor sich habe, so muß er, das versteht sich von selbst, auch dahinter zu kommen suchen, ob die ihm entgegentretende Anfechtung seines Pfleglings einfach oder gemischt ist, und wenn gemischt, ob sie den Anfang vom Leibe oder von der Seele genommen habe. Bei der unverkennbaren und unaufhaltsamen Wechselwirkung Leibes und der
\quellseite{315}
Seele wird es allerdings keineswegs erfolglos sein, wenn die nächste Wirkung des Seelenarztes oder des Arztes überhaupt nicht auf denjenigen Theil geht, in welchem das Uebel seinen Ursprung nahm; man wird auch überhaupt nicht einseitig wirken wollen, etwa zuerst auf den Leib und dann auf die Seele oder umgekehrt; dennoch aber wird es von Segen sein, wenn man vornherein denjenigen Theil des menschlichen Wesens vorzugsweise angreifen kann, von welchem die ganze Anfechtung ausgegangen ist. Aus dem bereits Gesagten geht hervor, daß im Falle der Unklarheit über den Entstehungsort des Uebels auch der Angriff auf denjenigen Theil gesegnet sein kann, der von dem andern nur in’s Schlepptau genommen ist, und es wird daher jedenfalls für eine Pflegerin ganz passend sein, wenn sie die gemischte Anfechtung für’s erste als eine einfache behandelt und der Seele diejenige Arznei zuwendet, welche sie bedarf. Man behandelt also die unklare gemischte Anfechtung als einfache. Ist man aber bis zu diesem Punkte gekommen, so ist in dem zuvor gesagten die weitere Weisung schon gegeben.

Um der Verwandtschaft des Gegenstandes willen reden wir hier gleich von dem Zustande der Besessenheit. Bereits aus dem vorigen ist es klar, daß man dämonische Krankheiten und Anfechtungen nicht mit Besessenheit verwechseln darf. Viele Dämonische, welche der Herr in dem Evangelium geheilt hat, waren nicht besessen, sondern sie hatten nur dämonische Krankheiten. Hier reden wir auch nicht von denen, von welchen man sagt, sie seien der Seele nach besessen, sondern ganz einfach von den leiblich Besessenen, in deren Leibe ein oder mehrere
\quellseite{316}
Dämonen die Wohnung aufgeschlagen haben. Leiblich besessen ist also der, dessen Leib eine Behausung des Teufels ist, oder eines Dämons. Was die Wirkung der Besessenheit anlangt, so kann auch diese sehr verschieden sein; es läßt sich denken, daß ein Mensch besessen ist, ohne daß der in ihm wohnende Dämon auch nur irgend ein Zeichen seiner Anwesenheit gibt. Man könnte sich also Besessene denken, von denen andere nicht wissen, daß sie es sind, und die es auch selbst nicht wissen. Andern Falles gibt es aber auch Wirkungen und Zeichen der Besessenheit genug; sie sind sehr mannigfaltig und vor lauter Mannigfaltigkeit sehr unsicher. Es können Zeichen sein, welche sich an der Erkenntnis des besessenen Menschen ereignen, so daß also z. B. ein Besessener weiß, was kein Mensch auf natürlichem Wege wissen kann. Dahin gehören auch die fremden Sprachen, die Besessene sprechen.

Die Besessenheit kann sich aber auch in einer mächtigen Erregung der Begier oder des Willens zeigen, namentlich einer Hingerissenheit und gewaltigen Neigung zum Bösen, oder zu Widerwärtigem, Unnatürlichem, so daß der Mensch vollführt, was er unter anderen Umständen jedenfalls lassen würde. Ebenso erweist sich die Besessenheit in einer starken und unnatürlichen Erregung von Gefühlen; namentlich sind es Gefühle der Furcht, welche den Anfang des bösen Zustandes zu bezeichnen pflegen. Die Zeichen der Besessenheit sind übrigens nicht bloß geistlich, sondern auch leiblich, und erweisen sich z. B. in unnatürlichen Bewegungen der innern oder äußeren Glieder des Leibes, in einer unnatürlichen Kraft, Behendigkeit und Schnelligkeit des Körpers und seiner Glieder. Doch sind
\quellseite{317}
alle diese Zeichen an und für sich selber, besonders wenn sie vereinzelt vorkommen, trügerisch, weil einerseits die menschliche Fähigkeit des Geistes und Gemüthes ein unerforschtes Meer ist, und anderntheils auch Geisteskrankheiten, die mit der Besessenheit nichts zu schaffen haben, oftmals eine ganz unerklärliche und auffallende Wirkung auf den Leib haben. Man wird daher jedenfalls langsam thun müssen mit der Erklärung, daß sich bei dem oder jenem Menschen die Zeichen der Besessenheit finden. Zu den leiblichen Zeichen der Besessenheit gehört auch der in manchen neueren Beispielen vorliegende Umstand, daß Dinge und Stoffe, welche sich in dem menschlichen Leibe natürlicher Weise unmöglich erzeugen können, z. B. Glasscherben, Nägel, Nadeln und dergl. entweder durch den Mund oder durch andere Glieder haufenweise abgeben, und alle, die es sehen, in die größte Verwunderung und Verlegenheit setzen.

Was die Anfänge der Besessenheit betrifft, so wird der Satan oftmals durch leibliche Zustände gefördert, und macht sich hauptsächlich an diejenigen, welche die Lehre vom Kreuz nicht gelernt haben, noch sie anzuwenden verstehen. Wer sein Kreuz tragen kann, wie Christus der Herr nicht bloß lehrt, sondern auch gibt, der ist ein Wunder Gottes in der Welt und die Anfechtung des Teufels prallt von ihm ab. — Die alten Lehrer sagen, die Besessenheit entstehe durch ein persönliches Eingehen oder Einfahren des Teufels in den Leib des Menschen, sei es nun durch den Mund oder welch anderes Glied des Leibes. Nach ihrem Zeugniß und allerdings vielen Beispielen geht dem Acte des Einfahrens allerlei voraus, was die Aerzte mit dem Namen „Hallucinationen“ bezeichnen.

\quellseite{318}
Der Angefochtene sieht besonders nächtlicher Weile Gestalten, wechselnde, thierische, auch menschliche, die sich ihm nahen, ihn schrecken, ihn locken, und oft in fürchterliche Angst versetzen. Er hört Dinge, die andere zuweilen hören, zuweilen auch nicht. Er fühlt unbegreifliche Schmerzen an seinem Leibe und macht Erfahrungen, die sich andere nicht zu enträthseln wissen. Endlich kommt es dahin, daß er entweder der Sache müde wird, oder im Gebete und Glauben nicht mehr widerstrebt, und so erzählt er dann einmal vielleicht nach einer schweren Nacht, es sei nun alles anders, der Teufel sei in ihn gefahren. Ob die dem Einfahren vorausgehenden ungeheuren sinnlichen Erfahrungen wirklich Hallucinationen sind, oder von wesenhafterer Art, ob allemal ein solches Einfahren stattfindet, wenn ein Angefochtener davon spricht, das ist eine andere Sache. Jedenfalls aber geht der Bericht der Erfahrenen aller Zeiten dahin, daß die Besessenheit damit beginne.

Die Behandlung der Besessenen anlangend, gibt es einen doppelten Weg, den ordentlichen der Belehrung und Seelsorge und den außerordentlichen des Exorcismus. Die gründlichste Hülfe geschieht wohl auf dem ersten Wege. Der Besessene, oder der es scheint, muß eine eingehende Belehrung über seinen Zustand bekommen und man darf bei der Belehrung die Anwendung auf den eigenen Fall, in welchem der Besessene ist, durchaus nicht scheiden. Durch diese praktische Belehrung kommt Licht in die niedergedrückte Seele und sie lernt sich gegenüber der Last ihres Drängers verstehen und stellen.

Es muß aber nicht allein eine Belehrung gegeben werden über die Natur des Uebels, sondern auch über die Wichtig%
\quellseite{319}
keit und Unwichtigkeit für die Leitung der Seele zum ewigen Leben. Jeder Besessene überschätzt die Last und Schwierigkeit seines Leidens und geht ungern auf den wahren Satz ein, daß die Besessenheit ein Leiden sei wie ein anderes, ein zeitliches Uebel, welches die Entwickelung der Seele für die Ewigkeit nicht hindert, sondern wohl gar fördert. Lernt der Besessene diese Wahrheit, und nimmt in Folge deß sein allerdings schweres Kreuz muthig und geduldig auf sich, so trägt sein Zustand zur Ehre Gottes mehr bei, als wenn er schnell und plötzlich von demselben befreit würde. Ist der Besessene ein unbelehrter Mensch, so ist er in Kraft seines Leidens desto mehr den Weg der Buße und des Glaubens zu führen. Wird Christus seine Gerechtigkeit, so wird er auch seine Freiheit, seine Gesundheit sein können, und es wird mitten im Leid die Freude wachsen. Dieser seelsorgerische Weg, der allerdings mit der Belehrung nicht beschlossen ist, sondern auch eine Kräftigung des Willens bei sich führen muß, ist der edelste und gesegnetste. Wird das Haus voll Licht und Leben, das zuvor voll Dunkelheit oder Zwielicht gewesen ist, so wird es dem höllischen Bewohner in demselben nicht mehr wohl sein und er wird es vielleicht verlassen, ohne alle die leidigen und beschwerlichen Demonstrationen, welche der zweite Weg mit sich bringt. Schon mit dem ersten Wege soll das Fasten und Beten verbunden sein, welches der eigene Mund des Herrn für die Heilung des Leidenden anordnet. Insbesondere aber soll es bei dem zweiten Weg eintreten, sowohl von Seiten des Leidenden als des Seelsorgers oder Exorcisten. Sollte jemand die Rechtmäßigkeit des zweiten Weges in Zweifel ziehen, so
\quellseite{320}
wird ihm aus den Worten und dem Verhalten Jesu leicht die Gewißheit zu verschaffen sein, daß auch dieser Weg im Reiche Gottes seine volle Berechtigung hat. Wird er betreten, so kann es allerdings Gaben und Menschen geben, die einen heroischen, die Sache schnell zu Ende führenden Glauben und eine wunderbare Kraft des Gebets besitzen, allein es können bei dieser Behandlung auch Monden und Jahre vergehen, scheinbare Heilungen und wirkliche Rückfälle, wirkliche Heilungen und scheinbare Rückfälle erfolgen, und es wird überhaupt kein Uebel so voll Täuschungen und Enttäuschungen, Betrug und Entmuthigungen sein, wie gerade dieses. Eben daher muß man mit Hingebung und großer Langmuth sich der Sache widmen und neben der Belehrung und Seelsorge alles Ernstes das Gebet oder auch den Exorcismus fortsetzen.

Die alten Kirchenordnungen, z. B. die pommerische, enthalten Formeln, bei deren Gebrauch sich’s herausstellen soll, ob der Mensch wirklich besessen ist oder nicht. Die pommerische Formel ist mehr werth, als es scheint, aber eine völlige Befriedigung bietet sie nicht. Diejenigen, welche in der rechtgläubigen Kirche sich mit diesem Uebel befassen, können durch eigene Erfahrung dem nachfolgenden Geschlechte noch gar viel Weisheit überliefern. Die pommerische Formel enthält auch keine Anleitung zum Exorcismus, auch nicht zu dem in der lutherischen Kirche üblichen Kirchengebet für Besessene. Die Zeit muß erst an Einsicht reifen und die Geister sich über die Sache klären, ehe die Lücken der lutherischen Kirche in diesem Stücke gebüßt werden können.

\quellseite{321}
Kann man weder den ersten noch den zweiten Weg gehen, so verlege man sich einfach auf das Gebet für den Besessenen und mit ihm und übe darin die nöthige Treue\footnote{Der vorhergehende § ist dem Aufsatz des Verf. über Geisteskrankheiten einverleibt, der sich im I. Bande von Vilmar’s Pastoral-theologischen Blättern p. 172 ff. findet.}.

\section{63. Hypochondrie und Hysterie}

\begin{center}
(Aus dem Anhang zur ersten Auflage dieses Buches p. 229 ff.)
\end{center}

Obwohl man sagt, daß die erstere Krankheit dem männlichen, die zweite aber dem weiblichen Geschlechte zugeschrieben werden müsse, so sind doch beide unleugbar verwandt, und sie verdienen ohne Zweifel die größte Berücksichtigung, weil der Patienten, die an ihnen leiden, unzählige sind. Der Ausgangspunkt des Leidens ist das vegetative System des Organismus. Hier aber hat es keinen fest und eng begrenzten Wohnsitz; es verändert seine Stätte leichter und öfter als irgend ein anderes körperliche Leiden. Zuerst an den Vorgang des innern Athmens\footnote{Ueber das innere Athmen f. Dr. G. H. v. Schubert’s „die Krankheiten und Störungen der menschlichen Seele.“ Stuttgart und Tübingen 1845. p. 11, 12, §. 5. p. 132, §. 18. Daselbst f. auch über „Hypochondrie und Hysterie.“} gebunden, kann es überall im Leibe, wo dieses stattfindet, sich äußern. Es kann deshalb, jetzt in der Lunge verweilend, die Form eines Lungenleidens annehmen, kann von der Luftröhre in die Muskelsubstanz des Schlundes, kann in den Darmkanal, kann in die obere Region der Nieren und in’s Gehirn, kann in die Re%
\quellseite{322}
gion der willkürlich bewegenden Muskeln auswandern und sich bald als krankhafte Affection des Darmkanals, bald in der Form der Wahnbilder oder Convulsionen und Zuckungen erscheinen, und erst dann, wann sich in irgend einem dieser Organe ein wirklich bleibendes Leiden, eine Veränderung der wirklich normalen Bildung der Theile und Mischung durchgebildet hat, ist seiner wandernden Wirksamkeit ein Zielpunkt gesetzt und alle eigenthümlichen Erscheinungen dieser Krankheit verschwinden.

Bei dem Leiden, welches wir hier zu betrachten haben, ist vor allem der Wechselverkehr des besondern organischen Verkehrs mit der Außenwelt, die selbstthätige Bewegung nach dieser hin gehemmt und unterbrochen, und hiermit verliert die organische Substanz jene Spannung, durch welche sie zur Verbindung mit der Lebensluft befähigt wird. Jede gesunde, kräftige Thätigkeit des Geistes wie des Lebens, jede freudige Stimmung des Gemüths befördert das innere Athmen; jede anhaltende Unthätigkeit, in welcher kein eigentliches Ausruhen ist, weil keine wahrhafte Arbeit ihr vorangeht, der geschäftige Müßiggang mancher Lebensarten und Stände, eine freudenlose Stimmung des Gemüthes wirken hemmend auf seinen Fortgang. Es darf dabei auch auf die Beschaffenheit der Nahrungsmittel hingewiesen werden, denn einige von diesen gewähren leichter und reichlicher, andere schwerer und spärlicher das brennbare Element, welches von der Lebenskraft durchwirkt, das Athmen unterhält.

Diejenigen Stände, welche bei einer anhaltend sitzenden Lebensweise ihren täglichen Unterhalt nur durch eine mechani%
\quellseite{323}
sche Bewegung, etwa der Hände, erwerben müssen, wobei nicht immer die Aufmerksamkeit durch einen Wechsel in der Arbeit beschäftigt ist, machen insonderheit zur Hypochondrie geneigt; also z. B. der Stand des Webers, Schneiders, Copisten und vieler Beamten. Dasselbe gilt auch von solchen Menschen, welche im Ueberfluß der äußern Lebensgüter nur genießen gelernt haben, nicht aber wirken und arbeiten; bei diesen hört oft alle Selbstthätigkeit völlig auf.

Da das innere Athmen durch die freudige Stimmung des Gemüths befördert wird, diese freudige Stimmung aber sehr häufig auf der gewohnten Befriedigung und Sättigung beruht, welche die Seele in den Genüssen und Einwirkungen der gewohnten Umgebung findet; so entsteht das hypochondrische Leiden häufig aus Heimweh, aus einem ungestillten Sehnen nach einer lieb gewordenen, jetzt unterbrochenen Gewohnheit, nach einem verlornen Besitz.

Zuweilen entsteht hypochondrisches Leiden auch in Folge der Abspannung, welche auf eine große und freudige Anspannung der organischen Thätigkeit zu folgen pflegt. Der zuvor durch innere Selbstthätigkeit wie bis zum Himmel erhoben war, geräth in einen Mangel des innern Athmens, in große Schwermüthigkeit und Betrübniß bloß deswegen, weil nunmehr die Erregung vorüber ist.

Alle solche Lebensverhältnisse und Einflüsse, welche bei Männern die Hypochondrie begründen, geben bei den Frauen noch mehr Veranlassung zur Hysterie. Die Lebensweise, sowie die Beschäftigung eines großen Theils der Frauen ist der oben geschilderten der Männer sehr ähnlich; die Hilflosigkeit
\quellseite{324}
und einsame Stellung der entweder gar nicht oder unglücklich vermählten Frauen, die natürliche Schwäche des Geschlechts gibt viel öfter und leichter zu Hemmungen der innern selbstkräftigen Bewegungen des Lebens, zur Unterdrückung der freudigen Stimmung Anlaß, als das gewöhnliche Loos der Männer; die Wirksamkeit der Lebenskraft des Weibes ist aber nach der natürlichen Bestimmung desselben mehr auf ein Werden und Bewegen in der vegetativen als in der psychischen Region des Organismus gerichtet. Daher nehmen beim weiblichen Geschlecht die Erscheinungen des hypochondrischen Leidens vorherrschend ihren Ausgang und Verlauf in der vegetativen Region des Organismus, vornehmlich in jener Partie desselben, welche für seine natürliche Bestimmung die wesentlichste ist. Und so wird die Hypochondrie des Weibes dem Namen wie dem Wesen nach zur Hysterie, welche mehr zu Krämpfen und Convulsionen sich hinneigt, während die Hypochondrie des stärkeren Geschlechtes zwar von diesem Leiden frei ist, aber viel öfter sich in wunderlichen Einbildungen und vorübergehenden Verwirrungen in der obern Sphäre der sinnlichen Wahrnehmungen und Vorstellungen kund gibt.

Die selbständige Lebensthätigkeit, von welcher der gesunde Fortgang des innern Athmens abhängt, wird bei dem Manne vorherrschend durch psychische Elemente angeregt und erhalten, beim Weibe vorherrschend durch leibliche, z. B. durch die Abspannung des Organismus von öfterer Schwangerschaft und langem Stillen der Kinder oder durch den Uebergang aus einer viel bewegten Lebensweise zu einer stillsitzenden. Daher geht denn auch die Hysterie häufiger in tief
\quellseite{325}
eingehende Zerstückungen einzelner Organe der vegetativen Region über, und die Lungensucht macht dann dem hysterischen Leben und Leiden ein Ende, während beim Mann der Uebergang zu Melancholie und Wahnsinn gewöhnlicher ist.

Die einzelnen Erscheinungen der Hysterie deuten darauf hin, daß der Grund derselben Hemmung des inneren Athmens ist. Da wird oftmals plötzlich das äußere Athmen, und damit der Quell des innern gehemmt; Angst und Beklemmung steigen mit oder ohne Bewußtsein bis zu Zuckungen; die Glieder werden starr und kalt; überhaupt klagen hysterische Frauen fast beständig über Kälte, haben eisig kalte Füße, trockene Haut, bis dann plötzlich dazwischen eine fliegende Hitze kommt, und damit oft Hemmung des Athmens und Beklemmung der Brust. Mit diesen unvermittelten Uebergängen hängt dann auch jener seltsame Wechsel der Launen zusammen: stundenlang ohne alle Veranlassung ist die Patientin lustig gestimmt, lacht ohne Aufhören, bis sie dann ohne alle Veranlassung weint und sich ängstigt. Mit der Hemmung des inneren Athmens hängen dann auch alle übrigen Erscheinungen der Krankheit zusammen. Hysterische haben wenig Bedürfniß zu trinken; oft entsteht Appetit wie ein Heißhunger, der dann durch einige Bissen völlig gestillt ist; schnell schwillt der Unterleib gewaltig an, schnell senkt er sich wieder ein; den Darmkanal und namentlich den Schlund befallen Zusammenschnürungen, welche sammt dem Gefühle, als ob ein runder Ballen im Schlunde stäcke, zu den charakteristischen Kennzeichen der Krankheit gehören. Zuweilen wird die Hysterische von einem schnell eintretenden heftigen Kopfweh
\quellseite{326}
befallen; sie weiß sich nicht zu fassen; da kommt ein Freund oder ein interessanter Gegenstand, ein merkwürdiges Ereignis zeigt sich; da ist auf einmal das Kopfweh verschwunden.

Anders ist es bei dem hypochondrischen Mann. Wenn seine Arbeitszeit, sein Geschäftstag vorüber ist, sinkt er zurück in die Betrachtung seines leiblichen Befindens; seine Welt ist sein Leib; mit krankhafter Spannung und ängstlicher Sorge, fast beobachtet er alle Veränderungen in seinem Körper, in einem jeden von ihnen ahnt er Gefahren eines nahenden Uebels, ja des Todes, welchen er fürchtet, weil ihm zum Leben, wie zum Tode der Muth fehlt. Dabei leidet er an Verdauungsbeschwerden, an Aufstößen, an Unregelmäßigkeit der Ausleerungen, an Eßlust, sowie an merkwürdigen Wahnbildern, die er zu sehen glaubt, von denen er aber doch in der Regel weiß, daß es Wahnbilder sind (Tasso). Den eben erwähnten Hallucinationen des hypochondrisch Schwermüthigen wohnt bei schwächeren Kranken die ernste Gefahr bei, daß sie sich des Willens und der Thatkraft des Kranken bemächtigen und zu einem Mittelpunkt werden, um welchen sich das ganze Leben dreht.

Die Hallucinationen der Frauen scheinen von mehr plastischer Art zu sein und beziehen sich häufig auf eine eigenthümliche, seltsame Bewegung der Hohlmuskel des Darmkanals. Eine Maus ist durch den offenen Mund in den Magen gekrochen, ein lebendiger Frosch, eine Kröte, eine Unke, ein Wassersalamander, eine Schlange ist in ihren Eingeweiden, sie spüren’s ja, der kalte Körper des Thieres streift an ihren Magen an, sie fühlen die Füße; jetzt spüren sie die
\quellseite{327}
Bisse, Laute vernimmt man wie eines Frosches, einer Unke; wo das Thier kriecht, hebt sich ja der Leib, und das alles ist so täuschend, daß auch andere bedenklich werden können, und doch sind die Thiere nur Gase.

Das Leiden der Hypochondrie und Hysterie äußert, wenn es sich höher steigert, und wenn der selbstkräftige Wille sich nicht zu einem ernsten Widerstand gegen dasselbe ermannt, einen krankhaft verändernden Einfluß, nicht allein auf die Empfindungen und Gefühle, sondern auch auf die Gesinnung des Menschen. Die Seele des Hypochondristen bekommt allmählig eine Hinneigung zu einer Selbstsucht, welche sein Verhältnis zur Außenwelt in hohem Grade stört. Seine Theilnahme an all’ dem, was nicht er selbst ist, wird geschwächt, die Theilnahme anderer an ihm selbst scheint ihm in jedem Falle zu klein. Das kräftige, freudige Bewegen der Gesunden regt ihn zum Unwillen, ja zur Bitterkeit auf; er wird mißtrauisch, ja wohl auch gehässig gegen andere, oder er wird wenigstens in Liebe und Haß einseitig, liebt den einen leidenschaftlich, haßt den andern eben so leidenschaftlich, und weil Liebe und Haß weder auf dem Grunde eines vernünftigen Erkennens, noch auf dem eines vernünftigen Willens ruhen, so hat weder Liebe noch Haß Bestand, sondern ist wechselnd und wandelnd wie die Zustände der Krankheit selber. Schlimm, wenn in solchen Fällen auch diejenige Zucht des Geistes mangelt, die bei einer höhern Bildung wie von selber kommt. — Zuweilen führt das Leiden den Kranken zu einer Verzweiflung an ihm selber, welche in vielen Fällen nicht die Form der demüthigen Selbsterkenntniß an
\quellseite{328}
sich trägt, sondern die mit der hochmüthigsten Selbsterhebung bestehen kann, wenn der Kranke den Grund seines Mißlingens und die Vereitelung seiner selbstsüchtigen Wünsche nicht in sich selber, sondern außer sich sucht. Sehr nahe stehen die Zustände der Hypochondrie und Hysterie an den Grenzen der Melancholie. Der Kranke kommt zuweilen zu einem solchen Lebensüberdruß, daß ihn der Tod mit einer Gewalt anzieht, wie nur immer den Lüstling eine Sinnenlust. Es folgt zuweilen auch wirklich Selbstmord. Ein anderer Höhenpunkt der Krankheit, nämlich der Uebergang in Somnambulismus darf hiebei nicht verschwiegen werden. Somnambule Frauen tragen in ihren Leiden sehr häufig die Bestrafung ihrer Trägheit und ihres Unwillens gegen die ihnen angebotene Seelenkur.

Die Aerzte haben bei Anordnung ihrer leiblichen Mittel nach dem oben angegebenen hauptsächlich darauf zu merken, daß innere Athmen wieder herzustellen und zu heben. Die Aerzte rathen deshalb einerseits Seeluft, Seebäder, Seereisen, weil die Luft an und auf der See viel comprimirter ist, als die auf den Höhen und eine weit größere Menge von Lebensluft in sich enthält. Dasselbige ist der Fall bei der kalten Luft, weshalb es den Patienten dieser Art so oft in kalter Witterung besser geht als in warmer; andererseits rathen aber die Aerzte gerade die Luft der Hochgebirge, welche, wenn auch nicht so substanzreich als die comprimirte Luft zur See, doch vermöge einer besondern Art von Spannung auf den Vorgang des Athmens einwirkt.

\quellseite{329}
Nöthiger aber als die medicinische Hülfe, ist den Kranken die psychische. Die Selbstthätigkeit muß wieder erweckt werden, und dazu bedarf es einer Anregung des Willens, die zwar zuweilen auch auf anderem Wege gelingt, am sichersten aber durch vollkommene Aufrichtigkeit und Wahrhaftigkeit im Betreff seines Leidens und durch den Ruf des Evangeliums zu einer völligen Bekehrung erreicht wird. Man belehre den Kranken über die Ursachen und die Art seiner Krankheit und zeige ihm, daß die wunderlichen Gefühle, die ihn plagen, zwar ihren Grund in leiblichen Uebeln haben, nicht aber auf Uebel weisen, die sehr bedeutend und gefährlich wären. Man zeige ihm, daß die qualvollen Seelenleiden, die er zu tragen hat, zur Krankheit selbst gehören. Ist man so glücklich, den Kranken so weit zu bringen, daß er über seine Lage ruhig nachdenkt, so ist schon viel gewonnen. Man zeige dem Kranken namentlich auch die große Kraft des Willens und die Nothwendigkeit einer Reaction gegen die sich aufdrängenden Gefühle und suche seinen Willen zum Guten und zum Gehorsam zu erziehen. Dabei aber wird es immerhin ganz an der Stelle sein, auch die einzelnen Ausgeburten der Krankheit als das auffassen zu lehren was sie sind. Man höre die Klagen des Kranken geduldig an, und lasse sich ja nicht die Einsicht in die augenblickliche Gefahrlosigkeit des Uebels zu einer leichtfertigen, nachlässigen, oder gar spöttischen Behandlung desselben reizen. Mit Theilnahme zeige man in jedem einzelnen Fall dem Kranken die wahre Ursache des Leidens. Dabei hüte man sich ja, ihn in Gesellschaft anderer Hypochondrischer oder Hysterischer zu bringen, da es ja bekannt ist, daß sogar Gesunde
\quellseite{330}
von dem hysterischen Uebel angesteckt werden. Ebenso bewahre man ihn vor dem Lesen medicinischer Bücher, aus denen sich seine krankhafte Phantasie nur alle Tage immer neue Angst und eine andere Krankheit nehmen würde. Dagegen wird es ganz richtig sein, wenn man seine Gedanken von seinen eingebildeten Leiden und Schreckensbildern abzuziehen sucht, und ihn allmählich zur Thätigkeit und Arbeit gewöhnt.

Sehr oft hilft dem Kranken eine Veränderung seiner Lage, zumal wenn der Same der Krankheit in der bisherigen Lage war. Das hysterische Mädchen genest sehr oft durch die Ehe, der durch ökonomisches Unglück hypochondrisch Gewordene durch Versetzung in eine sorgenfreie Lage. Ueberhaupt ist gründliche Veränderung der Lage ein treffliches Gegenmittel dieser Krankheit, ohne Zweifel aber das trefflichste, wenn der Weltmensch in ein anderes Verhältnis zu Gott, der im Wahn der Buße Steckende in den Frieden der Versöhnung, der wegen Mangel an Fortschritt in seinem Glauben irre Gewordene zu einem trefflichen Fortschritt in der Heiligung gebracht wird. Unaufhaltsames Fortschreiten ist das beste Gegengift gegen alle Geisteskrankheiten.

\quellseite{331}
\section[64. Eine alte Regel für Seelsorger]{64. Eine alte Regel für Seelsorger beim Krankenbesuch}

(Aus Felix Bibenbach’s Manuale Ministrorum Ecclesiae. Tübingen 1603. pag. 642.)

In actione Ministri cum aegroto tria sunt consideranda, videlicet:

Initium, quod fieri potest hoc modo, ut ostendatur infirmo\\
\quad 1. Occasio nostri accessus.\\
\quad\quad Quia adest Minister, vel\\
\quad\quad\quad Non vocatus, tunc indicabit: Συμπαθείαν suam Christianam; Officii rationes, quibus tenetur accedere, monere, consolari etc. aegrum.\\
\quad\quad\quad Vocatus ab aegroto: Audiat aegrotum facientem initium. Ipse ostendat suum propositum. Nota. Initium colloquii commodissime fieri potest ex voluntate Dei. Matth. 10. v. 30.\\
\quad 2. Causa omnium morborum: scilicet, peccatum. Idque vel in\\
\quad\quad Genere: si nihil aliud vel aeger fateatur, vel ipsi Ministro constet.\\
\quad\quad Specie: si forsitan enorme aliquod facinus, admissum ab infirmo, innotuerit.\\
\quad 3. Remedium contra peccata: poenitentia scilicet, quae constat\\
\quad\quad Agnitione, Contritione, Confessione.\\
\quad\quad FIDE in Christum.\\
\quad\quad Proposito vitae emendandae, si longior hujus lucis usura concedatur.

Medium, ad quod pertinent haec:\\
\quad Absolutio. Facta confessione.\\
\quad Communicatio Coenae Domini. Gratiarum actio.\\
\quad Declaratio usus Coenae, et consolatio etc.

Finis, qui est varius, pro ratione personae aegrotantis, vel etiam nostri accessus. Huc pertinent\\
\quad 1. Exhortationes ad:\\
\quad\quad Patientiam, praesertim, si morbus sit gravis, diuturnus.\\
\quad\quad Mortem alacriter obeundam, si Deo ita videatur.\\
\quad 2. Consolationes pro ratione tentationum, vel de\\
\quad\quad Peccato admisso etc.\\
\quad\quad Praedestinatione etc.\\
\quad\quad Perseverantia finali.\\
\quad\quad Tribus his capitibus variae etiam aliae tentationes, quae ex his oriuntur, comprehensae sunt.\\
\quad 3. Extrema adminicula, quibus subveniendum jam jam morituris, ut sunt potissimum\\
\quad\quad παράκλησις ad agentes animam.\\
\quad\quad Pracationes, vel: Ipsius mortui. Aliorum pro morituro.
\quellseite{332}
\chapter{Anhang}

\begin{titelseite}
(Zu pag. 103 * des „Evangelischen Geistlichen“ II.)

\begin{center}
{\large D. ANDR. AD. HOCHSTETTERI,}\\[0.6em]
{\large [THEOLOGI TUBINGENSIS]}\\[1.2em]
DEI\\
RECTA\\
CONCIONANDI\\
TEXTUMQUE SACRUM\\
CUM\\
EXPONENDI\\
TUM\\
APPLICANDI\\
RATIONE\\
COMMENTARIOLUS.\\[1em]
TERTIUM EDITUS\\
CURANTE\\
CHRISTOPH. FRIDRICO SARTORIO,\\
S. THEOL. DOCT. ET. PROF. PUBL. ORD.
\end{center}
\end{titelseite}

\quellseite{333}
\section{Gedruckt}

\begin{center}
1. 1701.\\
2. Göttingen. 1741.\\
3. Tübingen. 1767.\\
4. Gütersloh. 1876.
\end{center}

\quellseite{334}
\section{PRAEFATIO EDITIONIS PRIMAE}

Habes hic, amice Lector, de recta ratione formandi conciones commentatiunculam: quam ab Ipso Auctore editam, ab eoque praesente exornatam ego quidem maluissem. Sed nescio qua is verecundia prohibitus, quod jam octennio abhinc absolverat, hactenus intra sua servavit scrinia, ut tandem rogantibus ita nonnullis, huic docendi perspicuitati adsuetis, ipso fere invito in Ecclesiae conspectum extrudatur. Sincero omnia animo, et in Lectorum aedificationem scripta sunt: ut de fructu, quem qualiscumque haec opella habitura est, nulli dubitemus. Faxit DEUS, ut eam de recta catechizandi ratione Libellus propediem sequatur. Vale in Domino JESU, et precare DEUM, qui Ecclesiae suae afflictissimae, et in partes distractae suppetias ferat ipse, perpetuoque adsistat opitulator.

\quellseite{335}
\section{PRAEFATIO NOVA}

Commentariolus hic A. MDCCI sine Auctoris nomine primum editus, et quod parvis libellis, etiam optimae notae, communiter solet accidere, oblivioni fere datus, b. D. JAC. GUIL. FEUERLINO, Theologo Goettingensi, dignus visus est, quem Goettingae A. MDCCXLI recudi curaret, et concionandi praecepta ad illius ductum Studiosis explicaret. Eum denuo prelo subjicere visum, ut habeant Nostri, quem in re homiletica ducem sequantur: qui si uberiora desiderant, cum libello hoc monitorem aeque probatum, b. D. JOH. CHRISTIANI KLEMMII nostratis Centuriam paragraphorum regulas, cautelas, et monita homiletica complectentium, utiliter conjungent, quae Centuria A. MDCCXLIV. primum edita, ac diu desiderata, commodum, et prout plane existimamus, bono omine, hoc ipso mense recuditur. Praesentem quod attinet libellum, eum prout primum fuit editus, nihil mutatum reddimus, nisi quod menda sublata, tituli librorum, qui citantur, plene expressi, paragraphi commodius distincti numeris, et sub finem, quo expeditior libelli usus sit, Conspectus additus est. Qui cognoverint Commentariolum, ii procul dubio alterum de recta catechizandi ratione libellum enixe nobiscum desiderabunt: sed eum unquam editum esse, nobis quidem non constat. Quo magis optandum, ut Nostri hunc insignem libellum assidue volvant; et cum ex reliquis monitis atque consiliis, tum ex sale spirituali, praeceptis passim, cum primis, ultimis duobus paragraphis adsperso, quamplurimum salutaris fructus percipiant. Tubingae, a. d. IV Maji, A. MDCCLXVII.

\quellseite{336}
\begin{center}
JESU JUVA.
\end{center}

\section{DE RATIONE CONCIONANDI}

\subsection{§ 1}

Quanquam non veterum modo et in Domino jam defunctorum, sed et hodienum de Ecclesia orthodoxa praeclare merentium Theologorum studio deducta eo est res omnis homiletica, ut novo ad eandem adparatu opus non esse videatur: assurgere tamen ad Amicorum desideria visum modo est, enitique aliquid in timore Domini de materia tam utili tamque necessaria, quod promovere illorum commoda, animosque ad sacrum in Ecclesia munus parare facilius posset. Faxit divinum Numen feliciter! in cujus devota invocatione ad rem ipsam accingimur, brevi aliquot paragraphorum delineatione finiendam.

\subsection{§ 2}

MATERIA Concionis non nisi biblicus esto Textus, isque adaequatus Concioni, ne vel excedat verborum quantitas vel deficiat. Dividi commode solet in statum seu definitum, et electum liberoque Oratoris arbitrio positum.

\subsection{§ 3}

Duo sunt, eaque generalia, quae maxime observari a concionatura ante omnia debent.

Alterum, naturalis, ut sic dicam, textus notitia (§ 4).

\quellseite{337}
Alterum vero artificium est (§ 19), circa textum ipsum in proponendi eloquendique ratione occupatum.

\subsection{§ 4}

TEXTUS sacri NATURA ipsa ex his quinque dignoscitur: Investigandum primo Objectum, progrediendumque deinceps ad Genus et Speciem, inquirendum in Argumenta, ipsumque denique Thema ex istis judicandum. (§ 5. 6. 9. 10. 12.)

\subsection{§ 5}

OBJECTUM in textu sunt:

1. Articuli fidei christianae seu dogmata ad credendum in verbo DEI proposita.

2. Christiani hominis mores, sive virtuosi sive vitiosi.

3. Afflictiones in via ad coelum nobis obvenientes.

Neque quidquam in Scripturis uspiam leges, quod ad aliquam e tribus hisce classibus referri non possit, imo debeat.

\subsection{§ 6}

Quod ipsum liquido patescit, si GENERA dicendi, ex 2 Tim. III. 16. et Rom. XV. 4. ab Oratoribus sacris petita, paulo altius repetamus. Quanquam enim quinque illorum numerentur, non tamen circa plura versantur objecta, quam modo recensuimus.

1) Fidei enim sive historicae Jac. II. 19. sive fiducialis et salvificae articulos definit et inculcat Genus didascalicum; quos elenchticum defendit et adversa dogmata refellit.

2) In Moribus, si boni sint, commendandis ac insinuandis; si mali sint, taxandis et dissuadendis distinguitur Paedeuticum: cum hos admissos in Auditoribus vehementius reprehendat Epanarthoticum.

\quellseite{338}
3) Afflictionibus autem et malis solatia opponit Paracleticum.

\subsection{§ 7}

Atque haec Genera dicendi eo lubentius amplectimur, quo certius, in sacris praecipue, Spiritus S. ductum sequi praestat, quam hominum inventis dictamina DEI adstringere. Servemus nempe vetustum illud tritumque S. Irenaei, οὐδὲν ἄτερ γραφῆς. Tametsi enim sint, qui in his usus potius et applicationes, non ipsa genera quaerenda esse existiment, facile tamen iis satisfit, si distinguamus formalem doctrinae, elenchi, paediae, epanorthoseos et paracleseos in textu situm, a fluente e textu, doctrinali, elenchtica etc. conclusione; de priori nunc nobis sermo est, de usibus deinceps dicemus.

\subsection{§ 8}

Caeterum non ignoramus, profanorum Rhetorum dicendi genera observari a nonnullis in formanda sacra Oratione solere; quod ipsum non a B. HUNNIO saltem\footnote{in methodo concionandi, p. 8.}, sed et a B. L. OSIANDRO factum est in elegantissimo, quem de Ratione concionandi An. 1582 evulgavit libello, qui tamen ipse non semel monet, accommodari conciones ad ista genera posse, sed sine superstitione. Cumque quartum Causarum genus, Judiciale puta, difficilius sit et obscurum, Studiosis suasor est, ut illas, quas in sacris deprehendunt, Controversias ad genus didascalicum referant. Neque enim semper feliciter et commode istas praeceptiones ad theologicas res accommodari posse.

\quellseite{339}
\subsection{§ 9}

SPECIES nil aliud sunt quam modi varii, quibus Genera sive exprimuntur, sive determinantur.

In Didascalico, verbi causa, Species sunt: narrationes, testificationes, asseverationes, demonstrationes, descriptiones omnes etc.

In Elenchtico: accusationes, defensiones expostulationes, disputationes, negationes.

In Paedeutico: laudes, suasiones, dissuasiones, adhortationes, dehortationes.

In Epanorthotico: imprecationes, comminationes, dehortationes etc.

In Consolatorio: sympathiae, querelae, erectiones, divini auxilii promissiones, afflictionis recensiones; unde et ad hoc genus preces referre possis.

\subsection{§ 10}

Sequuntur ARGUMENTA\footnote{de Argumentis prolixe disserit supra laudatus Io. Christ. KLEMMIUS in Centuria Regg. Homilett. § 38 ss.} per quae textus deducitur. Hic accurate dispiciendum oratori, quid in textu explicet, quid probet, quid illustret seu amplificet. Huic enim numero Doctores Argumentorum genera incluserunt, ut

a. alia dicantur Explicantia, circumstantiae nimirum vulgari versiculo comprehensae: Quis, quid, ubi, quibus auxiliis, cur, quomodo, quando? inter quas eminet:

1) Quis? seu subjectum (aut subjecti potius, quibus id in textu proxime videtur vel vestitum vel vestiendum, conditiones).

2) Quid? seu objectum.

\quellseite{340}
3) Quomodo? scil. actionis forma, quae ex antecedentibus et concomitantibus dijudicari solet.

β. Alia Probantia, i. e. medii termini, quibus sententiae in textu assertae stabiliuntur, et contrariae prosternuntur.

γ. Alia denique Illustrantia, (quo et refero Amplificantia) quae vel simili, vel dicto, vel exemplo thema dilucidius reddunt: suntque vel insita in ipso textu, vel dissita seu aliunde assumpta, e quibus NB. insita potius, dein adscititia (et ex his, textui viciniora) attendenda.

\subsection{§ 11}

Et haec quidem si quis diligentius observaverit, formandis Locis communibus et Collectaneis aptissima deprehendet, faciliusque in re tam ardua tamque controversa progredietur.

\subsection{§ 12}

Denique et ad THEMA respiciendum, quod vel simplex est vel compositum. Simplex appellari communiter solet, ad quod, cum nihil quaerat, neque affirmative neque negative responderi potest: Compositum autem, quod quaestionem in se continens vel affirmari potest vel negari.

Paucis, si rem recte judicemus, in omni genere praeter didascalicum sunt themata composita; cum omne aut aliquid refutandum, aut suadendum dissuadendumque aut vehementius taxandum habeat: solum didascalicum partim themata conjuncta, articulos sc. fidei probandos, partim simplicia habet, circa quae rerum credendarum natura explicatur.

\subsection{§ 13}

MEDIA, ad hoc omne quod praefati sumus, feliciter perficiendum, sunt ANALYSES.

\quellseite{341}
\subsection{§ 14}

1) Analysis Grammatica, quae vocum phrasiumque emphases in fonte praecipue enucleat, et hoc fine parallelismos vocum aeque ac rerum scrutatur, neque vel ullum, quantum fieri potest, praetermittit syncategorema, quando quidem copulativa etiam, et, etc. saepius non illustrat tantum, sed plane exponit sensum, si studiose attendatur. Qua in re consulendus imprimis GLASSIUS, Philos. S. L. IV. Tract. 1 seqq. Suoque exemplo luculenter praevit Mart. GEIERUS in eruditis ad Psalmos, Proverbia, Cohelet, Danielemque Commentationibus.

\subsection{§ 15}

2) Analysis Logica, quae singularum vocum atque Periodorum connexionem aut disjunctionem dijudicat, inde argumentatur et concludit, et isthac via tamdiu progreditur, donec sub specie et objecto tanquam constituentibus Propositionem naturalem omnia textus membra accurate conspiciat digesta. Quodsi tamen textus sit prolixior, in aliquot enunciationes beneficio hujus Analyseos resolvi potest. De contextu vero, cujus potissima hic habenda ratio, legatur GLASSII de scripturae sensu eruendo tractatus in Phil. S. P. II. Lib. II. p. m. 351. (Ed. nov. p. 492.) consulaturque insignis Exegeta Seb. SCHMIDIUS in Commentariis passim, praecipue vero in Paraphrasi Evangelii S. Johannis, quam prae aliis ejus Operibus, commendatam sibi habeant Nostri.

\subsection{§ 16}

3) Analysis Rhetorica, quae Tropos Figurasque evolvit in sacro Textu contentas, de quibus consulatur GLASSIUS in Rhetorica sacra.

\quellseite{342}
\subsection{§ 17}

Dispiciendum quoque hic fuerit de Affectibus movendis, quia, quod AUGUSTINUS\footnote{I. 4 de Doctr. Chr. c. 11 et 17.} dixit, „Oratoris „sacri non est docere tantum, sed etiam flectere, et „movere animos.“ V. Io. Bened. CARPZOVII, Patris, Hodegeticum Membr. I. Aphor. 9 ubi praecipua quoque argumenta ad movendum peccati odium, amorem virtutis, spem denique ac fiduciam in DEO annotavit.

\subsection{§ 18}

Facient huc, quae de Affectibus anglico sermone scripsit eruditissimus Eduardus REINOLDUS, et gallico, le Pere SENAULT, in bellissimo Tr. des Passions de l’Ame; et Abbas BRETTEVILLANUS in Eloquentia cathedraria et forensi, Parisiis Anno 1689 evulgata, qui inprimis elegantissime disquirit de remediis Affectuum; quam inprimis materiam olim excoluit in Collegio Theol. moralis privato Germanicorum Moralistarum Phoenix D. Io. Ioach. ZENTGRAVIUS; et in den Lebens-Pflichten singulare aliquid hoc in argumento praestitit D. SPENERUS, quod frustra alibi quaeras. Videatur etiam B. GERHARDUS in Schola Pietatis, et Homiletae passim. Placet B. DANNHAUERI consilium (in Hermeneutica sacra), qui de Affectu ipsius Scriptoris sollicitum Homiletam Interpretem esse, illumque, quantum fieri potest, imitari jubet. Mire heic commendandus, cui me multa debere libens gratusque profiteor, Dn. D. SELIGMANNUS, Lipsiensis.

\quellseite{343}
\subsection{§ 19}

Alterum a Concionaturo probe observandum vocavimus ARTIFICIUM, quod circa Methodum (§ 20) imprimis, Catenam partium, (§ 23) Inventionem (§ 46) Ornatumque (§ 47) occupatur.

\subsection{§ 20}

METHODUS duplex est, (servamus distinctionem a nonnullis hodie impugnatam) Analytica et Synthetica. Analyticae officium vi vocis est ἀναλύειν, h. e. resolvere textum in partes suas, verborumque ac phrasium vim exponere, denique et doctrinas inde fluentes subnectere. Synthetica autem locum quendam communem e textu eruit, eumque argumentis tum insitis (atque his praesertim), tum dissitis confirmat, amplificat.

\subsection{§ 21}

Atque haec generalis communissimaque Methodi divisio est, cui non repugnare quadrimembrem in paraphrasticam, dogmaticam, articulatam et thematicam demonstraverunt alii: Cum paraphrastica quam vocant, et dogmatica nec non articulata ad analyticam referri possit, thematica autem idem esse videatur, quod alii Syntheticam appellant.

\subsection{§ 22}

Caeterum discrimen hoc inter Methodum analyticam et syntheticam a dispositione et propositione peti debet, quippe quae aut eadem est cum textuali, sicque analytica; aut textuale πόρισμα seu synthetica: pertinetque huc illud PLINII: „Arbitror ORDINEM rem unam esse, „e qua toti Orationi plurimum afflari soleat lucis, gratiae, „jucunditatis.“

\quellseite{344}
\subsection{§ 23}

Inter PARTES Concionis eminet primo EXORDIUM (supervacaneum enim est hic agere de praeloquiis, quae non ubique in usu sunt, saepiusque ingratam Concionem, haud raro tamen alacrem excitatumque auditorem reddunt) quod cum quaedam veluti janua et introitus Orationis sit, argumento textus et scopo conveniens esse oportet, commendarique in eo imprimis meretur brevitas et perspicuitas; longa enim exordia auditorem defatigare solent, ut reliqua minus attente audiat. Quanquam eruditae insinuationes rejiciendae non sunt, ubi necessitates tale quid requirunt.

\subsection{§ 24}

Sumuntur autem exordia

1) vel a Re ipsa, sive rei, quam tractamus, commendatione: aut proceditur

2) a thesi generali ad specialiorem et Hypothesin, sive in factis sive in promissis sive in comminationibus; v. g. si agamus in Propositione de captivitate Pauli, in exordio dicatur de incommodis et miseria captivorum in genere;

3) item ab antecedenti ad consequens; v. gr. si de benignitate Redemptoris loquamur, exordium duci potest ab antecedente nostra miseria, dicto aliquo pathetice descripta;

4) vel ab Exemplis aut Oppositis, tum relative, tum contrarie, tum privative talibus; v. gr. si Dominica XXII post. Trin. Propositio est: Homines sunt servi Dei, ad reddendam rationem vocandi, tum exordium peti posset ab Opposito relativo, quod Deus sit Dominus noster;

\quellseite{345}
5) vel etiam a Circumstantia vel loci, v. gr. inaugurationibus templorum, vel temporis, v. g. in festis aut anniversaria celebratione beneficii a DEO collati, Deut. IV. 9. 10, vel personarum: ubi tamen monendum est, non satis commode benevolentiam in sacris concionibus a persona dicentis captari, sicut etiam ab Auditorum persona saepius sine suspicione adulationis captari non potest, nisi forte unum alterumque casum excipias, v. g. si Pastor reconvaluerit, commode ordietur ex Philipp. II. 27.

6) Vel a Cohaerentia textus, et ubi continuus textus biblicus explicatur, vicem exordii supplet apud nonnullos repetitio illorum, quae in antecedente textu praecipua fuerunt, quam ἀνακεφαλαίωσιν brevem esse oportet, ne Concionator ignaviae suspicionem incurrat: in qua quidem re prudentia Oratoris opus est, dabitque is potius operam, ut novos semper nervos veluti flammulas penetralia ipsa afficientes in exordio afferat, quo ardentius ad Deum auditorum animi extollantur. Sunt, qui exordia ex abrupto ducant successu nonnunquam haud contemnendo, vid. Nic. CAUSSINUS de Eloquentia sacra et civili, VI. 7. p. 242. et 251.

\subsection{§ 25}

Caeterum partes Exordii fere tres sunt:

1) Exordialis, ut sic dicam, propositio (πρότασις), in qua eleganter nonnunquam ipsa principalis propositio continetur.

2) Ἔκθεσις, brevis aliqua explicatio, vel a simili, vel a dicto, vel ab exemplo.

3) Est digressio seu transitus ad propositionem. Is non omissa textus enarratione promte quasi fluere
\quellseite{346}
debet ex ipso exordio. Io. OLEARIUS\footnote{in Oratoria ecclesiastica.} VII recensuit partes, sed praeter necessitatem.

\subsection{§ 26}

Secunda Orationis sacrae Pars, atque illud Punctum, ut sic loquar, Archimedaeum, ex quo tota verti dirigique potest Concio, est PROPOSITIO, eaque pro differentia Methodi vel Analytica vel Synthetica.

\subsection{§ 27}

Analytica dicitur, quae scopum textus argumentis dicendorum ponit: estque nihil aliud quam (definiente CARPZOVIO) ipsum textus argumentum brevibus et perspicuis verbis indicatum; v. g. si quis velit proponere ex Evangelio Dominicae XXI post Trinitatis, tum scopus illius, qui est sanatio filii Reguli (die Gesundmachung des Königischen Sohns), loco propositionis quoque adduci posset.

\subsection{§ 28}

Propositio autem Synthetica illa appellatur, quae locum quendam communem e textu erutum proponit. Ita v. g. ex eodem Evangelio synthetice posset proponi „fides in cruce probata,“ (der im Kreuz bewährte Glaube), quo apprime conveniret exordium 1. Petr. I. 7. Observandumque hic obiter, propositionem eleganter nonnunquam phrasibus exordialibus exprimi: quod si fieri commode non possit, cogitandum saltem est, ipsam eo fore gratiorem, quo minus remota fuerit ab exordiali sententia. Atque hinc facile patet, assensum praebere nos non posse Eruditorum nonnemini, qui locum communem in exordio memorari posse negavit, cum id ipsum
\quellseite{347}
syntheticae Propositionis munus sit. Distinguendum enim inter Locum communem, qui proponi potest, et ejus usum, qui non.

\subsection{§ 29}

Sequitur PARTITIO, tertia concionis pars; estque illa nihil aliud quam distributio propositionis in ea, quae proxime continet: atque ideo ab Oratore sacro non insuper habenda, ut distinctius Auditor intelligat, quid expectandum sibi de propositione facta et quodnam istud sit, ad quod in specie omnia audienda illa referre debeat. Pertinetque huc monitum B. HUNNII\footnote{Meth. conc. p. 12.}: „Partitio „sic instituatur, ut ambitu suo universum textum includat, „nec ulla, vel minima particula textus praetereatur, „quae non sub aliqua parte, cujus facta hic mentio est, „contineatur.

De terminis homoeoptotis et homoeoteleutis (in quibus nonnulli gratiam pene omnem quaerunt, sed quam prudenter?) dicemus in discursu; advertetque Homileticae cultor, explicari dividique propositiones per argumenta topica pene omnia v. c. per opposita, per causas, effectus et adjuncta, per circumstantias, per partes, tum essentiales tum integrales, ab utili, necessario etc. per integras denique enunciationes modosque alios, de quibus privatim.

\subsection{§ 30}

Est Propositionis, Partitionisque quaedam appendix VOTUM: quo DEI auxilium ad perficiendum cum fructu sacrum opus imploratur. Quanquam enim effundant preces sub initium, usitatissime ex 2. Cor. XIII.
\quellseite{348}
13, ipsaque deinceps oratione Dominica continuentur, consultum tamen decensque est ut et tertiae addantur, cum domus DEI, domus precationis, opusque quod exsequitur perorans aeque ac audiens, coelicum omnino sit et divinum, insuperque longe gravissimum. Quare improbamus merito illorum morem, qui suspiria haec talia aut vota omittunt (quod a Viro et Concionatore celebri, factum esse memini) inepto sane brevitatis studio: sed poterant utique alia omitti forsan minus utilia, aut contrahi saltem, quam preces, quas ἀδιαλείπτως fundere jussit Salvator noster.

\subsection{§ 31}

Quarta eaque principalissima sacrae Orationis pars est TRACTATIO (confirmationem vocant profani Rhetores): estque nihil aliud, quam argumenti propositi justa enarratio, quae duabus constat partibus EXPLICATIONE nimirum (§ 32) et APPLICATIONE (§ 34).

\subsection{§ 32}

In EXPLICATIONE argumenti rite instituenda triplex observabitur Analysis, de qua § 13 ss. Ea vero nequaquam in verbis tantum occupatur, sed si suscipere cum fructu eam velis, Contenta textus analytice:

a) Primo debeat considerari, ordinari dividique in particulas, quarum cuilibet deinceps suum tribuendum (quae ipsa secunda partitio est CARPZOVIO membr. 2. Aph. 3. § 1. et partes omnes textus rite inter se ordinat), ostendendo nimirum, quis scopus sit, quae occasio, quae circumstantiae loci temporisque, v.g. si attendatur tempus scriptionis Epistolarum, de quo C. H. SANDHAGEN in Harmonia\footnote{Einleitung die Geschichte unseres Herrn Jesu Christi nach der Zeit-Ordnung zu betrachten. Berol. 1702. 8.}
\quellseite{349}
nunquam satis commendata, — quod subjectum, quae vis congruentiae?

β) Deinde ad ipsam descendi phraseologiam debet, explicarique tum genus loquendi tum emphasis vocum phrasiumque, tum denique si dissensus sit interpretum, ea sola in Concione afferri debet interpretatio, quam ex proprietate linguarum, ex phrasium similiumque locutionum collatione, ex contextu, ex praecedentium et sequentium diligenti consideratione (huc etiam refero fidem optimorum manuscriptorum) ad praesentem scopum Spiritus Sancti omnium maxime collimare perspexeris; et cum dubia nonnunquam ex adversariorum ψευδερημηνείαις locisque in speciem pugnantibus oriri in animis auditorum soleant, tollenda ista prudenter sunt, ostendendumque breviter, qua ratione responderi ad illas solide ab homine etiam plebejo, et qui conciliari textus cum hoc illove dubio possit ac debeat. Atque ex his demum vera exsurget Paraphrasis (distinguenda accurate a Metaphrasi, qualis de LUTHERO Doctorum regula est, illum non Metaphrasten, sed Paraphrasten egisse), cujus inter veteres exemplum luculentissimum dedit ERASMUS Roterodamus in Paraphrasi N. T., ex recentioribus autem Sebastianus SCHMIDIUS, ut et in Manuscripto ἀνεκδότῳ Io. Adamus OSIANDER, Cancellarius Tubingensis, cujus Operis summa cum laude mentionem fecit, editionemque, spero, differri non patietur ejus Parentator, HOCHSTETTERUS.

\subsection{§ 33}

Neque praetereunda hic sunt, quae § 10 diximus de Argumentorum triplici genere: haec enim in promptu semper esse debent, ut quae dicta sunt amplificentur et
\quellseite{350}
illustrentur per dicta, per similia, per exempla, Moralium item, ubi commode id fieri potest, inspersione: quo et digressiones referendae, de quibus in discursu. Denique tametsi § 20 in synthetica methodo Locum tractari communem dixerimus, observandum tamen est, explicationem textus in ea quoque pariter atque analytica obtinere. Quanquam enim in hac (analytica) plus dandum Exegesi, quoniam tamen in synthetica Locus communis aut Porisma ex textu propositum fuit, ideo ipsa explicatio ostendere deductionis fundamentum debet, ne quid praeter textum dixisse videamur. Debent Loci communes e textu quasi palmites enasci et textui tanquam fundamento inniti: ne praepostera ratione currus bovem trahere videatur.

\subsection{§ 34}

Altera tractationis pars est debita illius, quod explicatum modo e textu fuit, APPLICATIO ad Locos communes § 35 (cujus situs non definitur, sed liber esto) et consectaria (atque inter haec speciatim [§ 36] ad observationes) tandemque per quintuplicem, aut saltem e textu maxime fluentem Usum (§ 37) NB. auditorio accommodatum.

\subsection{§ 35}

Loci communes e textu efferri non inferri in eundem debent: ita enim intelliget Auditor, nos non nisi divina proponere Dogmata, monente id PETRO: Qui loquitur in Ecclesia, tanquam verbum DEI loquatur. 1. Petr. IV. 11. Ex DEI enim verbo fides est, teste PAULO ad Rom. X. 17. Caeterum non omnia in deductione Loci communis aut consectarii dicenda sunt, quae dici possent (hinc enim prolixae non saltem sed et
\quellseite{351}
impertinentes plerumque atque Auditorio molestissimae conciones ortum habent); sed ea potissimum, quae ex ipsis textus penetralibus sponte consequuntur, nec tam multitudo argumentorum quam illorum pondus considerari debet, quorum unum alterumque adduxisse, idque dicto, vel simili vel exemplo illustrasse, sufficiat.

\subsection{§ 36}

Observatio vocatur, quando praemissa textus expositione e circumstantiis partibusque textus plura eliciuntur porismata, auditorique consideranda proponuntur: quae tractandi ratio elegans quidem et utilis est, sed difficilior tamen et foecundis praesertim ingeniis usitatior. Exempla dabimus in praxi, simulque considerabimus, quae de modo miscendi et contemperandi tum explicationem textus tum ejus applicationem memoravit CARPZOVIUS Memb. 2. Aph. 7. quanquam circa ea, ipsomet fatente, vel nihil certi determinari possit, neque facile (si quid addere licet), tale quid nisi ab exercitatis Theologis tentari debeat.

\subsection{§ 37}

His ergo decenter observatis descendendum ad Usus est, atque id, quod vel formaliter in textu sacro continetur vel elicitum ex eo per locum doctrinalem aut observationes fuit, ad conscientias praesentis auditorii applicandum, attendendumque, quomodo orthodoxa Ecclesiae nostrae doctrina inde confirmari, adversa refutari, virtutes commendari, reprehendique vitia denique et erigi possint afflicti et infirmi.

\subsection{§ 38}

Primus itaque Usus Didascalicus est, quo dogma aliquod fidei sive morum, in loco communi vel obser%
\quellseite{352}
vationibus propositum ad praesens auditorium applicatur, ejusque piae ac devotae meditationi ita commendatur inculcaturque, ut de suo simul officio circa hanc rem faciendo instruatur, quale est (ut exemplo rem declaremus) si de fidei articulo sermo est, grata revelatae atque in nostra Ecclesia usque huc personantis divinae veritatis agnitio: ubi utiliter simul ad Confessionem Augustanam et Symbola, inprimis Catechismum nostrum provocari posset.

\subsection{§ 39}

Secundus Elenchticus est, seu dogmatis oppositi confutatorius, in quo sollicite caveat orator, ne quid adversario tribuat quod ille non statuit: unde non inutiliter ad publicam Ecclesiae Romanae doctrinam in Tridentina Synodo, in Catechismo Romano, vel qui plerisque in manibus est, Canisiano, propositam (breviter tamen) provocatur, si cum Pontificiis res sit. Satius etiam videtur nonnullis plane tacere obsoletas antiquasque haereses: et prudentiae est, non monere, cum monitis opus non est: complura enim mala sunt, quae Ecclesiae a modernis hostibus imminent, quam ut de non necessariis et fere antiquatis simus solliciti. Quam circumspecte in Ecclesia pressa procedi circa hunc usum debeat, docuerunt hactenus reverendi Argentinensium Antistites, interque hos inprimis Isaacus FAUSTIUS, qui clamantibus acerbissime debacchantibusque Jesuitis Ioanni DEZ, Gallo, et P. BAUMGARTNERO, Bavaro, placide illorum objecta discutere, quam paria clamitando referre, maluit.

\subsection{§ 40}

Tertius venit ordine Epanorthoticus, sive pecca%
\quellseite{353}
torum contra textum aut doctrinam inde propositam militantium reprehensorius, ad quem utiliter suscipiendum

1. Repraesentanda inprimis peccata in praesenti auditorio grassantia, provocandumque ad conscientias Auditorum, annon praedicato hactenus verbo inobedientes plerique fuerint, cum de caetero nihil quicquam omiserint ejusdem praecones: conf. Jerem. XXV. 3. Hos. XI. 7. Act. XX. 26.

2. Ostendenda peccatorum istiusmodi gravitas, et quam indignum Christiano sit illis indulgere: argumentis maxime e textu petitis.

3. Subjungenda denique comminatio, qualis e. g. Lev. XXX. 14. seqq.

Inprimis ab affectu omni acerbiori aliena esse debet epanorthosis, ut intelligant Auditores, NB. nos non nisi ex illorum amore studioque promovendae salutis istiusmodi correctiones suscepisse, quo etiam utiliter adhiberi possent Figurae, v. g. Obsecratio, ut roget Orator per amorem DEI, per vulnera IESU Domini et salutis propriae curam; item Detestatio, Praeteritio, Apostrophe etc. de quibus in Praxi.

\subsection{§ 41}

Quartus est Paedeuticus sive admonitorius, quo ex textu aut doctrina inde proposita Auditores sui admonentur officii, ut desistant a malo, bonumque omni animi contentione sectentur. Hebr. XII. 14. Non autem sufficit nuda admonitio: addenda simul sunt

1. Argumenta seu motiva, quibus Auditores ad praestandum obsequium efficaciter excitentur; qualia sunt a necessario, ab honesto, ab utili, a jucundo.

2. Ostendenda etiam media, quorum observatione
\quellseite{354}
laetius faciliusque in studio pietatis progredi possimus: nonnunquam haud infeliciter memoria monitorum certo eoque quotidiano fere symbolo commendatur, quod luculentis exemplis a Theodoro DASSOVIO nobis olim ostensis declarabimus.

3. Solent denique certas dicendi figuras adhibere Oratores in hoc usu, tametsi forsan tale non cogitantes, uti sunt, Apostrophe, Licentia, Epizeuxis, Exclamatio, Expolitio, Epiphonema, Epistrophe, Incrementum, Anaphora, Aetiologia etc.

\subsection{§ 42}

Quintus denique Paracleticus sive Consolatorius est, qui malis corporalibus spiritualibusque, quibus premi in hac vita solent fideles, opponitur. Ad quem recte instituendum

α) Praemitti non incommode potest descriptio doloris, ejusque pro re nata, inprimis si ex spirituali afflictione oriatur, prudens exaggeratio. Parum enim apud afflictos (propria teste experientia) proficiet, qui causam statim dolendi omnem illis negaverit.

β) Petenda e textu argumenta solandi, aut si haberi ex illo non possint, derivandi aliunde e fontibus Israelis rivuli, opponendaque calamitatibus solamina, quorum etiam memoria uno alterove symbolo commendari a nonnullis solet.

γ) Haud raro etiam Figuris utuntur Oratores sacri, quales sunt, Apostrophe, aliaeque etc. sed apud optimos quosque fere praeter opinionem et absque ulla affectatione id contingit. Minime enim ostentari in his cura debet, et „evitare artem summae artis est,“ ait QUINTIL. in Instit. Orat., ubi de celanda arte loquitur.

\quellseite{355}
\begin{quote}
„Suspicio artificiorum derogat Orationi fidem, Oratori „adimit autoritatem,“ CICERO L. I. de Inventione.
\end{quote}

\begin{quote}
„Affectata Oratio gratiam amittit, et luxurians in rebus „sacris Rhetoricatio sic concionatorem decet ut gravem „aliquem Oratorem canis venerandum decet vestitus juvenis, ex variis coloribus consarcinatus.“ Lucas OSIANDER de Ratione Concionandi p. 71.
\end{quote}

\subsection{§ 43}

Restat EPILOGUS, quinta eaque ultima sacrae Orationis portio, cujus duae in Rhetoricis constitui solent partes, Enumeratio scilicet, et Admonitio; quae utraque in Conclusione locum reperit.

\subsection{§ 44}

1. Enumeratio, definiente CICERONE, „dicitur „illa, per quam res disperse et diffuse dictae unum in „locum coguntur, et reminiscendi causa unum sub aspectum subjiciuntur, quo fit, ut simul in memoriam redeat Auditor, putetque nihil esse praeterea, quod debeat desiderare. Commune autem praeceptum (pergit- „CICERO, summus Oratorum) hoc datur ad enumerationem, ut ex unaquaque argumentatione, quoniam „tota iterum dici non potest, id eligatur, quod erit „gravissimum, et unumquodque quam brevissime transeatur, ut memoria, non oratio renovata videatur.“ L. I. de Invent. c. 52.

\subsection{§ 45}

2. Enumerationi severa et gravis subjungitur Admonitio, quae excitet Auditores, ut quae dicta sunt, animis reponant, et ad eorum praescriptum vitae suae rationes instituant.

\quellseite{356}
Nonnunquam tamen omissa partium enumeratione (quae quidem utilissima est) posterior haec Epilogi pars tractationi protinus subjicitur: qua animi Auditorum ostensa rei gravitate et necessitate percellantur, et quidam quasi aculei iisdem infigantur.

\subsection{§ 46}

Dicendum nunc quoque erat de INVENTIONE textumque disponendi ratione vera; sed plus satis jam dictum de eo a CARPZOVIO est. Caeterum ad quid tandem dispositionis supra unum textum variatio, ad centenarium usque producta, prodesse queat, aliorum est judicare. Id saltem moneo, non spernendum Concordantiarum ad inveniendas dispositiones usum esse, in quibus, si quis eminentem prae caeteris in textu vocem evolvere velit, vix unquam deerit, quo Concionem exordiatur, suppeditatamque a naturali forma propositionem vestiat. Cavendum tamen, ne textus videatur accommodandus exordio, cum hoc accommodari textui debeat.

\subsection{§ 47}

De ORNATU non est, quod sollicitus sit studiosus Scripturae S. lector, quippe quae ipsum argumentis omne genus exornantibus abunde instruet. Praesuppono tamen figurarum Rhetoricarum notitiam, triplicisque, de qua superius dictum, resolutionis experientiam.

\section{APPENDIX}

\subsection{§ 48}

Solet denique appendicis loco Disquisitio ab Oratoriae sacrae scriptoribus institui de Elocutione, de
\quellseite{357}
Memoria, de Pronunciatione, ipsaque demum Ratione scribendi Conciones, de quibus nos breviter dicemus.

\subsection{§ 49}

Primo circa GENUS DICENDI seu ELOCUTIONEM vehementer variant Concionatores. Sunt enim (uti studiose observavit B. HUNNIUS\footnote{Meth. conc. p. 31.}), et quotidiana docet experientia) qui grandi et prope aulico dicendi genere gaudent, alii vero ex quadam levitate scurras verius quam Concionatores agunt; quidam odiosa tautologia unum idemque subinde repetunt, verbis paulo mutatis, multisque sine peculiari emphasi congestis Synonymis: qui omnes sane decens in Ecclesia dicendi genus non assequuntur. Etenim qui nimiam in dicendo elegantiam affectant, sine fructu plerumque concionantur, cum intelligi a vulgo haud possint, obliti certe illius Apostolici 1 Cor. II. 1 ὃ καθ’ ὑπεροχὴν λόγου ἢ σοφίας.

\subsection{§ 50}

Parem, imo majorem merentur reprehensionem, qui scurrilem in Ecclesia sermonem usurpantes semetipsos et conciones suas contemtui Auditorum exponunt, ignari prorsus, quomodo conversari rite debeant in domo DEI viventis 1 Tim. III. 15. Plus denique verborum quam rerum habent, qui nimia sermonis copia luxuriant, scopumque adeo concionis primarium, qui est aedificatio, non attingunt.

\subsection{§ 51}

Optimo ergo dicendi genere utuntur, qui sermone perspicuo simplicique, doctis simul et indoctis perceptibili concionantur, quam fere dicendi rationem tenuit
\quellseite{358}
Servator noster, dum parabolis e medio rerum usu petitis familiariter admodum sine pompa et splendore sermonis res maximas, ipsaque regni coelorum mysteria est complexus: neque enim difficiles ejus Parabolae sunt, si, quid Regno DEI intellectum voluerit, didiceris Vide Harmoniam SANDHAGENII et USSERIANAM\footnote{Harmonia Evangelica, wie solche in Iacobi Usserii Bibliothek gefunden, cum Praef. b. FRANCKII, Halae 1699. 8.}. Eum secutus gentium Apostolus de se profitetur, quod sua inter Corinthios praedicatio non fuerit sita in persuasoriis humanae sapientiae verbis sed ostensione spiritus et potentiae, 1 Corinth. II. 13 ne inanis reddatur crux Christi. 1 Cor. I. 17. cap. II. 4.

\subsection{§ 52}

Laudandi ergo imprimis sunt, qui ita sacras literas diligenter legerunt, ut pleraque, quae dicere volunt, verbis et phrasibus scripturae, iisque Auditori notis efferre possint; quanquam haec non ita intelligenda sunt, quasi sententiam suam propriis enunciare verbis non liceat Oratori, cum non ad syllabas, sed ad sententiam Scripturae Oratio formanda sit.

Commendari insuper a gravissimis Theologis solet stylus B. LUTHERI, itemque BALDUINI, (quem post Parentem in Hodegetico Membr. I. Aphor. XI. mirifice auditoribus suis nostro tempore commendavit D. CARPZOVIUS,) ut et B. Io. ARN DII in Postilla et Psalterio plane incomparabili.

\subsection{§ 53}

Secundum de quo tractari in Appendice solet, MEMORIA est. Quia enim Conciones non ex charta legi,
\quellseite{359}
(quod apud Anglos in usu est) sed coram populo memoriter recitari hodie consueverunt, memoriae utique beneficio opus est, quo res fideliter servare, et ad populum convenienter enunciare valeamus.

\subsection{§ 54}

Ea vero nulla re melius juvatur confirmaturque, quam si ipsam Concionis structuram animo recte concipiamus: si ordinata partium dispositio, et concinna, ut sic dicam, textura sit, et quae singulae parti dederimus argumenta, tum explicantia, tum probantia, tum amplificantia, probe teneamus. Ita sane fiet, ut invocato devote aeterno Numine, de memoria non habeat, quod conqueratur Orator noster.

\subsection{§ 55}

Atque hoc ego tutissimum praestantissimumque memoriae praesidium merito dixerim. Imaginum enim conceptiones suadeat qui velit, quod facere video Dan. Georg. MORHOFIUM in Polyhist. L. II. c. 6. p. 364. 365. Nos liberam hic optionem cuivis relinquimus. Locos commendat Auctor Commentariorum Rhetoricorum ad HERENNIUM (L. 3. c. 16. 17.) in iisque collocandas imagines esse existimat: sed tale artificium locum non habere potest, ubi bis aut ter vel saepius etiam in una septimana concionari oportet. Nonnulli in Oratione memoriter recitanda litteris et vocibus certis iisque paucissimis quasi signis et characteribus utuntur, vid. MORHOFIUS l. c. p. 376.

\subsection{§ 56}

Tertium in appendice PRONUNTIATIO est. Etsi enim auditores verbum DEI attente audire et ei obsequi debent, quomodocunque id illis recitetur: ut tamen
\quellseite{360}
libentius ac sine molestia audiant Concionatorem, plurimum in pronunciatione situm est, ut grata ea ac decens sit. Quanquam enim vox fingi non possit, sed ea, quam cuique dedit Dominus, utendum sit, decenter tamen eandem moderari et corrigere ex parte possumus. Id quod facilius fiet, si familiares amice rogaverimus, ut libere de defectibus pronunciationis nos admoneant, nosque ea, quae ingrata sunt, modeste deponamus aut emendemus.

\subsection{§ 57}

Caeterum jam suo quidem tempore questus CICERO, omnes putasse, vix posse de voce et gestu dilucide scribi, cum haec res ad sensus nostros pertineat; et postquam praecepta de pronunciatione dedisset, ita concludit:

\begin{quote}
„Non sum nescius, inquiens, quantum susceperim negotii, qui motus corporis exprimere verbis, et imitari scriptura conatus sim voces. Verum nec hoc confusus sum posse fieri, ut de his rebus satis commode scribi possit, nec si id fieri non posset, hoc quod feci fore inutile putabam, propterea quod hic admonere voluimus, quid oportet, reliqua trademus exercitationi.“
\end{quote}

(l. 3. ad Herenn. c. 11 et 15.) Itaque et nos pauca solum monuisse contenti erimus, reliqua in praxi diligentius observaturi.

\subsection{§ 58}

Primum, quod CICERO (ibid. c. 12.) monet, potissimumque hoc est. „Utile est ad firmitudinem vocis „sedata vox in principio; quid insuavius, quam clamor „in exordio causae? Intervalla vocem confirmant, eadem „sententias concinniores divisione reddunt, et auditori „spatium cogitandi relinquunt.“

\quellseite{361}
In principio igitur concionis sedata esse debet pronunciatio, quae in medio cursu nonnihil pro ipsarum rerum qualitate attollatur, ita tamen ne unquam ultra vires elevetur.

\subsection{§ 59}

Adhibebit etiam Concionator hanc prudentiam, ut cum simpliciter docet, oratione quoque moderata utatur: quando vero increpat aut severius exhortatur, majorem vocis contentionem adhibeat, rebus ipsis accommodam. Breviter: omnis fugiat affectatio: ingratum enim est, quicquid non naturale sed adscititium est; naturalis vero pronunciatio (nisi ipsa sit vitiosa) omnium gratissima est, quo accedere in vultu debet modestia et gravitas Theologo digna, qui se non propriam, sed Dei causam agere meminerit.

\subsection{§ 60}

Non ergo somnolenta sit pronunciatio, nec uno tenore recitetur Concio; variam enim vocis flexionem rerum postulat diversitas; neque una vocabuli syllaba sic elevetur, ut reliquae vix exprimantur aut parum distincte audiantur; nec ita dissimulentur ultimae periodorum syllabae, ut eas deglutire videamur, sed intelligibiles et illae sunto; nec denique ulla a principio impertinens litera apponatur a-an n-a-dass-n’dass, quae vitia probe annotavit Lucas OSIANDER p. 84. 85. Diximus, naturalem esse debere vocem, et in hoc potissimum circa pronunciationem artificium est, si fugiamus in voce artificium. Itaque non clamosa sit vox, neque vel praecipitans vel lenta nimis et tarda sit pronunciatio, neque commata confundenda, ne intempestive contextus rumpatur.

\quellseite{362}
\subsection{§ 61}

De GESTIBUS denique ita olim judicavit OSIANDER. „Uti non omnes gestus simpliciter rejiciuntur, „ita modestia requiritur. Lingua enim, non manibus aut digitis est loquendum, neque ea in re imitandi sunt histriones, vel etiam Oratores veteres et profani, qui causas civiles, non theologicas, coram populo egerunt.“

\subsection{§ 62}

Superest, ut de RATIONE SCRIBENDI CONCIONES paucula attingamus. Dignus sane laude omni atque admiratione illorum zelus laborque est, qui integras conciones ad verbum conscribunt, illasque iis propemodum verbis recitant, quibus consignatae sunt. Ita enim haud facile quidquam excidet per imprudentiam, quod offendere auditores possit.

\subsection{§ 63}

Quoniam autem haec dicendi ratio ingravescente aetate et crescente negotiorum multitudine servari difficulter potest: ideo satius esse existimant Duumviri clarissimi B. Lucas OSIANDER, Parens, et Aegidius HUNNIUS, ut Concionator discat praemeditate ea efferre, quorum Argumenta et Summam in memoria habet: ita ut inter dicendum deligat verba et phrases, quae ad exprimendam animi sententiam requirantur et huic loco conveniant.

\subsection{§ 64}

Meminerint tamen Nostri et qui vegeta adhuc memoria gaudent, consultius esse, suscipere Ecclesiae causa hunc laborem et diligenter mandare memoriae, quae scripto prius expresserint. Quapropter severissime
\quellseite{363}
omnino reprehendendi sunt, qui nescio cui naturali facundiae suae confisi, citra necessitatem, ex tempore, quicquid in buccam venit, effutire, inque re tam necessaria, seria, concernente divini Numinis gloriam et aeternam Auditorum salutem, tam negligenter agere ac loqui in conspectu DEI, Ecclesiae et sanctorum Angelorum non erubescunt: cujus impiae sane et non tolerandae socordiae rationem aliquando Principi Pastorum gravissimam reddere cogentur.

\subsection{§ 65}

E contrario praemeditatus ad dicendum si accedat Concionator, eo majore animo dicere, eoque gravius movere poterit. Sic suam quoque DEO et Ecclesiae fidem approbabit, hacque diligentia excitans donum DEI, quod in ipso est, ulteriora deinceps divinae gratiae sentiet incrementa, consummatoque feliciter cursu, perennantem vitae coronam e manu Domini JESU accipiet: cui sit honor et gloria in secula seculorum.

\trenner

Et haec quidem circa debitam ORATIONIS SACRAE formationem monenda hac vice videbantur. Quae uti in timore Domini intentoque unice in Ecclesiae, quam proprio sanguine redemit DEUS, aedificationem, oculo proposita a me fuerunt, ita non inanem fore hunc laborem in Domino plane confido, Deumque O. M. ardenti denuo prece invoco, qui operarios in messem suam ipse formet emittatque, nostrisque rigantibus coeleste incrementum benedictionemque omnigenam clementer largiatur: Amen.

\quellseite{364}
\section{CONSPECTUS}

\begin{center}
In\\
TRACTATIONE DE RATIONE\\
CONCIONANDI\\
consideranda sunt.
\end{center}

\begin{flushleft}
A. Interna Orationis Structura:\\
\hspace*{1.2em}et quidem\\
\hspace*{1.2em}I. Materia Orationis sacrae. § 2.\\
\hspace*{1.2em}II. Formatio ejusdem debita. § 3.\\
\hspace*{2.4em}ad quam requiruntur\\
\hspace*{2.4em}1) Notitia de Natura Textus S. § 4.\\
\hspace*{3.6em}A. Dignoscitur ex his quinque\\
\hspace*{4.8em}α) Objecto § 5.\\
\hspace*{4.8em}β) Generibus dicendi § 6—8.\\
\hspace*{4.8em}γ) Speciebus § 9.\\
\hspace*{4.8em}δ) Argumentis § 10. 11.\\
\hspace*{4.8em}ε) Themate § 12.\\
\hspace*{3.6em}B. Obtinetur ope Analyseos § 13.\\
\hspace*{4.8em}α) Grammaticae. § 14.\\
\hspace*{4.8em}β) Logicae, § 15.\\
\hspace*{4.8em}γ) Rhetoricae, § 16—18.\\
\hspace*{2.4em}2) Artificium § 19, occupatur circa\\
\hspace*{3.6em}A. Methodum, quae est\\
\hspace*{4.8em}α) Analytica,\\
\hspace*{4.8em}β) Synthetica, § 20—22.\\
\hspace*{3.6em}B. Catenam partium, inter quas eminet\\
\hspace*{4.8em}A) Exordium, cujus recensentur\\
\hspace*{6em}a) Requisita § 23.\\
\hspace*{6em}b) Fontes § 24.\\
\hspace*{6em}c) Partes § 25.
\end{flushleft}

\quellseite{365}
\begin{flushleft}
\hspace*{4.8em}B) Propositio § 26, quae est\\
\hspace*{6em}a) Analytica § 27.\\
\hspace*{6em}b) Synthetica § 28.\\
\hspace*{4.8em}Γ) Partitio, § 29.\\
\hspace*{6em}cujus Appendix Volum § 30.\\
\hspace*{4.8em}Δ) Tractatio ipsa, § 31. quae constat\\
\hspace*{6em}a) Explicatione § 32s.\\
\hspace*{6em}b) Applicatione § 34.\\
\hspace*{7.2em}1. ad Locos communes § 35.\\
\hspace*{7.2em}2. ad Observationes sive Porismata § 36.\\
\hspace*{7.2em}3. ad Usus § 37 qui sunt\\
\hspace*{8.4em}α) Didascalicus § 38.\\
\hspace*{8.4em}β) Elenchticus § 39.\\
\hspace*{8.4em}γ) Epanorthoticus § 40.\\
\hspace*{8.4em}δ) Paedeuticus § 41.\\
\hspace*{8.4em}ε) Paracleticus § 42.\\
\hspace*{4.8em}E) Epilogus § 43, cujus duae sunt partes\\
\hspace*{6em}a) Enumeratio § 44.\\
\hspace*{6em}b) Admonitio § 45.\\
C. Inventionem § 46.\\
D. Ornatum § 47.\\
B. Externus Oratoris Habitus:\\
\hspace*{1.2em}ubi instituitur Disquisitio § 48.\\
\hspace*{1.2em}I. de Elocutione seu Genere dicendi § 49. cujus Requisita\\
\hspace*{2.4em}1) alia sunt remotiva, ne peccet\\
\hspace*{3.6em}α) vel in excessu § 49.\\
\hspace*{3.6em}β) vel in defectu § 50.\\
\hspace*{2.4em}2) alia positiva § 51. 52.\\
\hspace*{1.2em}II. de Memoria, cujus exhibetur\\
\hspace*{2.4em}1) Necessitas § 53.\\
\hspace*{2.4em}2) Adminicula § 54s.\\
\hspace*{1.2em}III. de Pronunciatione, § 56—60.\\
\hspace*{2.4em}et Gestibus, § 61.\\
\hspace*{1.2em}IV. de Ratione scribendi Conciones, § 62.
\end{flushleft}

\quellseite{366}
\section{Ein alphabetisches Sachregister}

\subsubsection{A}

Abkündigungen. 151.

Amtsgaben. 69. 78.

Amtsführung, methodistische. 71.

Amtsmittel. 68.

Anfechtungen. 289.

— deren Behandlung. 295.

Auctorität, persönliche, eines Predigers. 126.

\subsubsection{B}

Beichten. 263.

Beichtgeheimniß. 283.

Beichtvaters Arbeit, große. 278.

— Gefahr. 282.

— Weisheit. 280.

Beispiele und deren Werth. 72. 73.

Beredsamkeit. 82.

Bescheidenheit des Predigers. 127.

Besessenheit. 289. 298.

Bibelübersetzung. 114.

\subsubsection{C}

Casuistik. 254.

Ceremonien. 39. 51.

Concordanz. 14. 117.

Confirmandenunterricht. 133.

— dessen Uebersicht. 182.

\subsubsection{D}

Dialect bei der Katechese. 178.

Dialectik. 91.

Dialogische Form der Lehre. 85. 175.

Domrich's: „Die psychischen Zustände.“ 241.

\subsubsection{E}

Eintheilung der Rede. 97. 98*.

Elocutio. 97. 244.

Ertemporieren. 128. 134*.

Erinnerungen, besondere, an die Gemeinde. 154.

\subsubsection{F}

Form der Katechese. 172.

Formelweisheit. 275.

Formen, verschiedene, der Lehre. 84.

Führung, seelsorgerische, der Gemeinden und der Einzelnen. 229.

\subsubsection{G}

Gaben, deren Erkenntniß. 81.

Gemeines Gebet. 219.

Gemüthsbewegung hervorzurufen. 325.

\quellseite{367}
Genera causarum. 97*.

Geräthe, heilige. 204.

Gesticulation. 150. 345.

Gottesdienst. Bestandtheile. 213.

täglicher. 220.

Gottesdienstlicher Platz für die Abtheilungen der Gemeinde. 202.

\subsubsection{H}

Hauptfestzeiten des Jahres und deren Feier. 190.

Hirten und Lehrer. 67.

Hochstetter’s kleine Homiletik. 103. 315.

Homilie. 103.

Homilie und Rede. Vereinigung beider. 108.

Homiliarium Caroli Magni. 139*.

Horenlectionen. 197.

Hypochondrie und Hysterie. 304.

\subsubsection{I}

Jahreslauf, natürlicher. Dessen Feier. 189.

Inhalt der Katechese. 161.

Invention. 120. 339.

\subsubsection{K}

Katechese ambulando oder von der Kanzel? 179.

Katechesen über biblische Geschichte. 163.

Katechet. Wort. 157.

Katechetische Oden. 160.

Katechetischer Schlendrian. 171.

Katechismus, apostolischer. 162*.

Katechismuspredigten der brandenburgischen Kirchenordnung. 174.

Kenntnisse, ärztliche, des Seelsorgers. 244.

Krankenbesuch, alte Regel für denselben. 314.

Krankenseelsorge. Werth derselben. 285.

\subsubsection{L}

Länge und Kürze der Predigten. 123.

Lectionarien der Alten. 196.

Lehrgabe, Lehrhaftigkeit. 79. 80.

Liturgie, alte, im Verhältnis zu unserer Zeit. 215.

Liturgie, deren Zweck. 88.

Luther’s concionator timidus. 130*.

\subsubsection{M}

Mangel des symbolischen Handelns. 223.

Maß der Gaben. 82.

Maß des katechetischen Inhalts. 164.

Melanchthon’s philosophische Lehrbücher. 91. 93.

Memorieren. 341.

Mißlectionen. 196.

Methodismus der Seelsorge. 249.

Mittel, seelsorgerische. 246.

\subsubsection{N}

Nachahmung Luther’s. 31. 189.

\subsubsection{O}

Oratio. 118.

\quellseite{368}
\subsubsection{P}

Parallelstellen. 117.

Paramente. 204.

Pensa, katechetische. 164.

Postillen, deren Gebrauch. 139. 143.

Predigen, demüthiges. 28.

Predigt, textlose. 101.

Predigtweisen: erste. 98.

— analytische, analytisch-synthetische, synthetische. 110.

— deren Wahl. 110.

Principien der Liturgie. 224.

Privatbeichten, verschiedene Arten. 273.

Privatseelsorge, was sie nicht ist. 248.

Pronunciation. 145. 342.

\subsubsection{Q}

Quatemberfeier. 195.

\subsubsection{R}

Rathen in schweren Händeln. 22.

Raum, heiliger. 198.

Rede, die christliche. 105.

Rede und Homilie, deren Vereinigung. 108.

Reden der Alten. 96.

Rhetorik. 29. 94.

Rhetorik der Apostel. 100.

— Melanchthon’s. 101.

\subsubsection{S}

Satzbau. 135.

Schriften der Propheten und Apostel. 40.

Schriftstudium. 9. 112.

Schriftstudium. cursorisches. 113.

— patristisches. 115.

Seelsorge, was darunter zu verstehen. 228.

Stiftshütte und heiliger Raum des Neuen Testamentes. 200.

Styl. 135.

Symbol und symbolisches Handeln. 223.

\subsubsection{T}

Taufstein in der römischen und in der lutherischen Kirche. 208.

Text. 100.

Textuales theologi. 113.

Theilen des göttlichen Wortes. 252.

II. Tim. 3, 16. 122.

Topik. 121*.

Trösten. 21.

\subsubsection{U}

Unbescheidenheit der Beichtväter. 282.

— der Prediger. 27.

Unverschämtheit derselben. 30.

\subsubsection{V}

Vering „Ueber Wechselwirkung zwischen Seele und Körper im Menschen.“ Leipzig 1817. 239.

Vertrauen zum Seelsorger. 262.

Vorbereitungsstudium auf das heilige Amt. 90.

Vorbereitung zur Katechese. 168.

— Gefahr derselben. 170.

Vortrag, liturgischer. 225.

\quellseite{369}
\subsubsection{W}

Wechselwirkung Leibes und der Seele. 236.

Wirkungsweise des Worts. 75.

\subsubsection{Z}

Zeiten, die heiligen. 186.

Ziele, pastorale. 230. 266.

Zustände, pastorale und psychische. 233.

Zu wem spricht der Prediger? 144.

Zweck, besonderer, der Krankenseelsorge. 288.

\quellseite{370}
\quellseite{371}
\quellseite{372}
\quellseite{373}
\quellseite{374}
\quellseite{375}
